%============================================================
%  Bd. 18 - Halbgruppen
%============================================================

\documentclass[main.tex]{subfiles}

\ifSubfilesClassLoaded{
  \BandSetupStandalone{B18}
}{
}

\title{Bd. 18 - Halbgruppen}
\author{Martin Kunze}
\date{}
\setcounter{file}{18}

\hypersetup{
  pdftitle={Bd. 18 - Halbgruppen},
  pdfauthor={Martin Kunze}
}

\begin{document}

\title{Bd. 18 - Halbgruppen}
\author{Martin Kunze}
\date{}
\hypersetup{pageanchor=false}
\maketitle
\hypersetup{pageanchor=true}
\setcounter{tocdepth}{3}
\tableofcontents
\setcounter{file}{18}
\renewcommand{\FormulaBandID}{18}

\chapter{Einführung}

Dieser Band entwickelt Halbgruppen aus einer assoziativen binären Operation.
Behandelt werden Beispiele, Teilbarkeit, Homomorphismen, Isomorphismen und
die funktorielle Konstruktion der Potenzhalbgruppe. Zusätzliche
Strukturaxiome werden in den folgenden Bänden getrennt untersucht.

\chapter{Halbgruppen}

\section{Definition der Halbgruppe}

\FormulaDefDeltaK[Begriff der Halbgruppe]{%
  \SemigroupAlg{A}{\star}%
}{SemigroupDef}{%
  \DeltaRow{Mengen}{A\dsep x\dsep y\dsep z}
  \DeltaRow{\textbf{Axiome}}{}
  \DeltaRow{Abgeschlossenheit}{%
    x,y\in A\vdash x\star y\in A}
  \DeltaRow{Assoziativität}{%
    x,y,z\in A\vdash (x\star y)\star z=x\star(y\star z)}
  \DeltaRow{\textbf{Neues Symbol}}{}
  \DeltaPrem{Halbgruppe}{\SemigroupAlg{A}{\star}}
}

\FormulaAxiomDeltaK[Zugriff auf die Abgeschlossenheit einer Halbgruppe]{%
  \SemigroupAlg{A}{\star}\dsep x,y\in A
  \vdash x\star y\in A%
}{SemigroupClosureAxiom}{%
  \DeltaRow{Mengen}{A\dsep x\dsep y}
}

\FormulaAxiomDeltaK[Zugriff auf die Assoziativität einer Halbgruppe]{%
  \SemigroupAlg{A}{\star}\dsep x,y,z\in A
  \vdash (x\star y)\star z=x\star(y\star z)%
}{SemigroupAssociativityAxiom}{%
  \DeltaRow{Mengen}{A\dsep x\dsep y\dsep z}
}

\FormulaAxiomDeltaK[Zugriff auf die binäre Operation einer Halbgruppe]{%
  \SemigroupAlg{A}{\star}\vdash
  \mathsf{BinOp}\!\left(A,\star\right)%
}{SemigroupBinaryOperationAxiom}{%
  \DeltaRow{Mengen}{A}
}

\section{Beispiele von Halbgruppen}

\subsection{Die Potenzhalbgruppe}

Aus jeder Halbgruppe entsteht eine weitere Halbgruppe, deren Elemente die
nichtleeren Teilmengen des ursprünglichen Trägers sind. Die Menge dieser
Teilmengen schreiben wir als \(\PowNE{A}\).

Für \(X,Y\in\PowNE{A}\) wird die ursprüngliche Operation mengenweise
fortgesetzt. Der Index \(\mathcal P\) unterscheidet diese Operation auf
Teilmengen von der Operation \(\star\) auf einzelnen Elementen.

\FormulaDefDeltaK[Mengenweise Fortsetzung einer Halbgruppenoperation]{%
  \begin{aligned}[t]
    &\SemigroupAlg{A}{\star}\dsep X,Y\in\PowNE{A}\vdash{}\\[-2pt]
    &X\star_{\mathcal P}Y\coloneqq
      \bigl\{z\in A\bigm|
        \exists x\in X\,\exists y\in Y\,(z=x\star y)\bigr\}
  \end{aligned}
}{PowerSemigroupProductDef}{%
  \DeltaRow{Mengen}{A\dsep X\dsep Y\dsep x\dsep y\dsep z}
  \DeltaRow{Funktionen}{\star\dsep\star_{\mathcal P}}
  \DeltaRow{\textbf{Neues Symbol}}{}
  \DeltaPrem{Mengenweise Operation}{\star_{\mathcal P}}
}

\FormulaThmDeltaK[Elementkriterium für das Mengenprodukt]{%
  \begin{aligned}[t]
    &\SemigroupAlg{A}{\star}\dsep X,Y\in\PowNE{A}\vdash{}\\[-2pt]
    &z\in X\star_{\mathcal P}Y
      \leftrightarrow
      \Bigl(z\in A\land
        \exists x\in X\,\exists y\in Y\,(z=x\star y)\Bigr)
  \end{aligned}
}{PowerSemigroupProductMembership}{%
  \DeltaRow{Mengen}{A\dsep X\dsep Y\dsep x\dsep y\dsep z}
  \DeltaRow{Funktionen}{\star\dsep\star_{\mathcal P}}
}
\begin{tabproofwide}
  \proofstepwidestar[1]{\SemigroupAlg{A}{\star}}{\rA}
  \proofstepwidestar[2]{X\in\PowNE{A}}{\rA}
  \proofstepwidestar[3]{Y\in\PowNE{A}}{\rA}
  \proofstepwidestar[1,2,3]{%
    X\star_{\mathcal P}Y=
    \bigl\{z\in A\bigm|
      \exists x\in X\,\exists y\in Y\,(z=x\star y)\bigr\}}{%
    \FormulaRefAuto{PowerSemigroupProductDef}[def]{1,2,3}}
  \proofstepwidestar[1,2,3]{%
    \begin{aligned}[t]
      z\in X\star_{\mathcal P}Y
      \leftrightarrow
      \Bigl(z\in A\land
        \exists x\in X\,\exists y\in Y\,(z=x\star y)\Bigr)
    \end{aligned}}{%
    \FormulaRefAuto{x\in\{u\in A\mid P(u)\}\eqvdash
      x\in A\land P(x)}[thm]{4}}
\end{tabproofwide}

Im folgenden Beweis fassen \(\exists I^{*}\) und \(\exists E^{*}\) endlich
viele geschachtelte Anwendungen von \(\exists I\) beziehungsweise
\(\exists E\) zusammen. Bei beschränkten Quantoren schließen diese
Kurzschreibweisen die erforderlichen \(\land\)-Schritte ein.

\FormulaThmDeltaK[Die nichtleeren Teilmengen bilden eine Halbgruppe]{%
  \SemigroupAlg{A}{\star}\vdash
  \SemigroupAlg{\PowNE{A}}{\star_{\mathcal P}}
}{NonemptyPowerSemigroup}{%
  \DeltaRow{Mengen}{%
    A\dsep X\dsep Y\dsep Z\dsep u\dsep v\dsep w\dsep x\dsep y\dsep z}
  \DeltaRow{Funktionen}{\star\dsep\star_{\mathcal P}}
}
\begin{tabproofsplitwide}
  \proofpartwideindR[Trägerelement einer nichtleeren Teilmenge]{%
    X\in\PowNE{A}\dsep x\in X\vdash x\in A
  }{%
    X\in\PowNE{A}\dsep x\in X\vdash x\in A
  }[PowerSemigroupCarrierElement]
    \proofstepwidestar[1]{X\in\PowNE{A}}{\rA}
    \proofstepwidestar[2]{x\in X}{\rA}
    \proofstepwidestar[1]{X\subseteq A}{%
      \rAEa{\FormulaRefAuto{NonemptyPowersetMembership}[thm]{1}}}
    \proofstepwidestar[1,2]{x\in A}{%
      \FormulaRefAuto{A\subseteq B,\,x\in A\vdash x\in B}[thm]{3,2}}
  \closeproofpartwideindR

  \proofpartwideindR[Elementare Produkte liegen im Mengenprodukt]{%
    \begin{aligned}[t]
      &\SemigroupAlg{A}{\star}\dsep X,Y\in\PowNE{A}\dsep
        x\in X\dsep y\in Y\vdash{}\\[-2pt]
      &x\star y\in X\star_{\mathcal P}Y
    \end{aligned}
  }{%
    \begin{aligned}[t]
      &\SemigroupAlg{A}{\star}\dsep X,Y\in\PowNE{A}\dsep
        x\in X\dsep y\in Y\vdash{}\\[-2pt]
      &x\star y\in X\star_{\mathcal P}Y
    \end{aligned}
  }[PowerSemigroupProductContains]
    \proofstepwidestar[1]{\SemigroupAlg{A}{\star}}{\rA}
    \proofstepwidestar[2]{X\in\PowNE{A}}{\rA}
    \proofstepwidestar[3]{Y\in\PowNE{A}}{\rA}
    \proofstepwidestar[4]{x\in X}{\rA}
    \proofstepwidestar[5]{y\in Y}{\rA}
    \proofstepwidestar[2,4]{x\in A}{%
      \FormulaRefAuto{PowerSemigroupCarrierElement}[thm]{2,4}}
    \proofstepwidestar[3,5]{y\in A}{%
      \FormulaRefAuto{PowerSemigroupCarrierElement}[thm]{3,5}}
    \proofstepwidestar[1,2,3,4,5]{x\star y\in A}{%
      \FormulaRefAuto{SemigroupClosureAxiom}[ax]{1,6,7}}
    \proofstepwidestar[4,5]{%
      \exists u\in X\,\exists v\in Y\,(x\star y=u\star v)}{%
      \rEI{4,\rEI{5,\rII}}}
    \proofstepwidestar[1,2,3,4,5]{%
      x\star y\in A\land
      \exists u\in X\,\exists v\in Y\,(x\star y=u\star v)}{%
      \rAI{8,9}}
    \proofstepwidestar[1,2,3,4,5]{x\star y\in X\star_{\mathcal P}Y}{%
      \FormulaRefAuto{PowerSemigroupProductMembership}[thm]{1,2,3,10}}
  \closeproofpartwideindR

  \proofpartwideindR[Das Mengenprodukt bleibt im ursprünglichen Träger]{%
    \SemigroupAlg{A}{\star}\dsep X,Y\in\PowNE{A}
    \vdash X\star_{\mathcal P}Y\subseteq A
  }{%
    \SemigroupAlg{A}{\star}\dsep X,Y\in\PowNE{A}
    \vdash X\star_{\mathcal P}Y\subseteq A
  }[PowerSemigroupProductSubset]
    \proofstepwidestar[1]{\SemigroupAlg{A}{\star}}{\rA}
    \proofstepwidestar[2]{X\in\PowNE{A}}{\rA}
    \proofstepwidestar[3]{Y\in\PowNE{A}}{\rA}
    \proofstepwidestar[4]{w\in X\star_{\mathcal P}Y}{\rA}
    \proofstepwidestar[1,2,3,4]{%
      w\in A\land
      \exists x\in X\,\exists y\in Y\,(w=x\star y)}{%
      \FormulaRefAuto{PowerSemigroupProductMembership}[thm]{1,2,3,4}}
    \proofstepwidestar[1,2,3,4]{w\in A}{\rAEa{5}}
    \proofstepwidestar[1,2,3]{X\star_{\mathcal P}Y\subseteq A}{%
      \FormulaRefAuto{A\subseteq B\coloneq
        \forall x\,(x\in A\rightarrow x\in B)}[def]{\rUI{\rRI{4,6}}}}
  \closeproofpartwideindR

  \proofpartwideindR[Das Mengenprodukt ist nichtleer]{%
    \SemigroupAlg{A}{\star}\dsep X,Y\in\PowNE{A}
    \vdash X\star_{\mathcal P}Y\neq\varnothing
  }{%
    \SemigroupAlg{A}{\star}\dsep X,Y\in\PowNE{A}
    \vdash X\star_{\mathcal P}Y\neq\varnothing
  }[PowerSemigroupProductNonempty]
    \proofstepwidestar[1]{\SemigroupAlg{A}{\star}}{\rA}
    \proofstepwidestar[2]{X\in\PowNE{A}}{\rA}
    \proofstepwidestar[3]{Y\in\PowNE{A}}{\rA}
    \proofstepwidestar[2]{X\neq\varnothing}{%
      \rAEb{\FormulaRefAuto{NonemptyPowersetMembership}[thm]{2}}}
    \proofstepwidestar[3]{Y\neq\varnothing}{%
      \rAEb{\FormulaRefAuto{NonemptyPowersetMembership}[thm]{3}}}
    \proofstepwidestar[2]{\exists x\,(x\in X)}{%
      \FormulaRefAuto{A\neq\varnothing\eqvdash
        \exists x\,(x\in A)}[thm]{4}}
    \proofstepwidestar[3]{\exists y\,(y\in Y)}{%
      \FormulaRefAuto{A\neq\varnothing\eqvdash
        \exists x\,(x\in A)}[thm]{5}}
    \proofstepwidestar[6]{x\in X}{\rA}
    \proofstepwidestar[7]{y\in Y}{\rA}
    \proofstepwidestar[1,2,3,8,9]{x\star y\in X\star_{\mathcal P}Y}{%
      \FormulaRefAuto{PowerSemigroupProductContains}[thm]{1,2,3,8,9}}
    \proofstepwidestar[1,2,3,8,9]{X\star_{\mathcal P}Y\neq\varnothing}{%
      \FormulaRefAuto{a\in A\vdash A\neq\varnothing}[thm]{10}}
    \proofstepwidestar[1,2,3,8]{X\star_{\mathcal P}Y\neq\varnothing}{%
      \rEE{7,9,11}}
    \proofstepwidestar[1,2,3]{X\star_{\mathcal P}Y\neq\varnothing}{%
      \rEE{6,8,12}}
  \closeproofpartwideindR

  \proofpartwideindR[Abgeschlossenheit der mengenweisen Operation]{%
    \SemigroupAlg{A}{\star}\dsep X,Y\in\PowNE{A}
    \vdash X\star_{\mathcal P}Y\in\PowNE{A}
  }{%
    \SemigroupAlg{A}{\star}\dsep X,Y\in\PowNE{A}
    \vdash X\star_{\mathcal P}Y\in\PowNE{A}
  }[PowerSemigroupProductClosure]
    \proofstepwidestar[1]{\SemigroupAlg{A}{\star}}{\rA}
    \proofstepwidestar[2]{X\in\PowNE{A}}{\rA}
    \proofstepwidestar[3]{Y\in\PowNE{A}}{\rA}
    \proofstepwidestar[1,2,3]{X\star_{\mathcal P}Y\subseteq A}{%
      \FormulaRefAuto{PowerSemigroupProductSubset}[thm]{1,2,3}}
    \proofstepwidestar[1,2,3]{X\star_{\mathcal P}Y\neq\varnothing}{%
      \FormulaRefAuto{PowerSemigroupProductNonempty}[thm]{1,2,3}}
    \proofstepwidestar[1,2,3]{%
      X\star_{\mathcal P}Y\subseteq A\land
      X\star_{\mathcal P}Y\neq\varnothing}{\rAI{4,5}}
    \proofstepwidestar[1,2,3]{X\star_{\mathcal P}Y\in\PowNE{A}}{%
      \FormulaRefAuto{NonemptyPowersetMembership}[thm]{6}}
  \closeproofpartwideindR

  \proofpartwideindR[Zeugen eines Mengenprodukts]{%
    \begin{aligned}[t]
      &\SemigroupAlg{A}{\star}\dsep X,Y\in\PowNE{A}\dsep
        w\in X\star_{\mathcal P}Y\vdash{}\\[-2pt]
      &\exists x\in X\,\exists y\in Y\,(w=x\star y)
    \end{aligned}
  }{%
    \begin{aligned}[t]
      &\SemigroupAlg{A}{\star}\dsep X,Y\in\PowNE{A}\dsep
        w\in X\star_{\mathcal P}Y\vdash{}\\[-2pt]
      &\exists x\in X\,\exists y\in Y\,(w=x\star y)
    \end{aligned}
  }[PowerSemigroupProductWitnesses]
    \proofstepwidestar[1]{\SemigroupAlg{A}{\star}}{\rA}
    \proofstepwidestar[2]{X\in\PowNE{A}}{\rA}
    \proofstepwidestar[3]{Y\in\PowNE{A}}{\rA}
    \proofstepwidestar[4]{w\in X\star_{\mathcal P}Y}{\rA}
    \proofstepwidestar[1,2,3,4]{%
      w\in A\land
      \exists x\in X\,\exists y\in Y\,(w=x\star y)}{%
      \FormulaRefAuto{PowerSemigroupProductMembership}[thm]{1,2,3,4}}
    \proofstepwidestar[1,2,3,4]{%
      \exists x\in X\,\exists y\in Y\,(w=x\star y)}{\rAEb{5}}
  \closeproofpartwideindR

  \proofpartwideindR[Assoziativität auf Teilmengenelementen]{%
    \begin{aligned}[t]
      &\SemigroupAlg{A}{\star}\dsep X,Y,Z\in\PowNE{A}\dsep
        x\in X\dsep y\in Y\dsep z\in Z\vdash{}\\[-2pt]
      &(x\star y)\star z=x\star(y\star z)
    \end{aligned}
  }{%
    \begin{aligned}[t]
      &\SemigroupAlg{A}{\star}\dsep X,Y,Z\in\PowNE{A}\dsep
        x\in X\dsep y\in Y\dsep z\in Z\vdash{}\\[-2pt]
      &(x\star y)\star z=x\star(y\star z)
    \end{aligned}
  }[PowerSemigroupElementAssociativity]
    \proofstepwidestar[1]{\SemigroupAlg{A}{\star}}{\rA}
    \proofstepwidestar[2]{X\in\PowNE{A}}{\rA}
    \proofstepwidestar[3]{Y\in\PowNE{A}}{\rA}
    \proofstepwidestar[4]{Z\in\PowNE{A}}{\rA}
    \proofstepwidestar[5]{x\in X}{\rA}
    \proofstepwidestar[6]{y\in Y}{\rA}
    \proofstepwidestar[7]{z\in Z}{\rA}
    \proofstepwidestar[2,5]{x\in A}{%
      \FormulaRefAuto{PowerSemigroupCarrierElement}[thm]{2,5}}
    \proofstepwidestar[3,6]{y\in A}{%
      \FormulaRefAuto{PowerSemigroupCarrierElement}[thm]{3,6}}
    \proofstepwidestar[4,7]{z\in A}{%
      \FormulaRefAuto{PowerSemigroupCarrierElement}[thm]{4,7}}
    \proofstepwidestar[1,2,3,4,5,6,7]{%
      (x\star y)\star z=x\star(y\star z)}{%
      \FormulaRefAuto{SemigroupAssociativityAxiom}[ax]{1,8,9,10}}
  \closeproofpartwideindR

  \proofpartwideindR[Dreifachzeugen der links geklammerten Seite]{%
    \begin{aligned}[t]
      &\SemigroupAlg{A}{\star}\dsep X,Y,Z\in\PowNE{A}\dsep
        w\in (X\star_{\mathcal P}Y)\star_{\mathcal P}Z\vdash{}\\[-2pt]
      &\exists x\in X\,\exists y\in Y\,\exists z\in Z\,
        \bigl(w=(x\star y)\star z\bigr)
    \end{aligned}
  }{%
    \begin{aligned}[t]
      &\SemigroupAlg{A}{\star}\dsep X,Y,Z\in\PowNE{A}\dsep
        w\in (X\star_{\mathcal P}Y)\star_{\mathcal P}Z\vdash{}\\[-2pt]
      &\exists x\in X\,\exists y\in Y\,\exists z\in Z\,
        \bigl(w=(x\star y)\star z\bigr)
    \end{aligned}
  }[PowerSemigroupLeftTripleMembershipForward]
    \proofstepwidestar[1]{\SemigroupAlg{A}{\star}}{\rA}
    \proofstepwidestar[2]{X\in\PowNE{A}}{\rA}
    \proofstepwidestar[3]{Y\in\PowNE{A}}{\rA}
    \proofstepwidestar[4]{Z\in\PowNE{A}}{\rA}
    \proofstepwidestar[5]{%
      w\in (X\star_{\mathcal P}Y)\star_{\mathcal P}Z}{\rA}
    \proofstepwidestar[1,2,3]{X\star_{\mathcal P}Y\in\PowNE{A}}{%
      \FormulaRefAuto{PowerSemigroupProductClosure}[thm]{1,2,3}}
    \proofstepwidestar[1,2,3,4,5]{%
      \exists u\in X\star_{\mathcal P}Y\,
      \exists z\in Z\,(w=u\star z)}{%
      \FormulaRefAuto{PowerSemigroupProductWitnesses}[thm]{1,6,4,5}}
    \proofstepwidestar[7]{u\in X\star_{\mathcal P}Y}{\rA}
    \proofstepwidestar[7]{z\in Z}{\rA}
    \proofstepwidestar[7]{w=u\star z}{\rA}
    \proofstepwidestar[1,2,3,8]{%
      \exists x\in X\,\exists y\in Y\,(u=x\star y)}{%
      \FormulaRefAuto{PowerSemigroupProductWitnesses}[thm]{1,2,3,8}}
    \proofstepwidestar[11]{x\in X}{\rA}
    \proofstepwidestar[11]{y\in Y}{\rA}
    \proofstepwidestar[11]{u=x\star y}{\rA}
    \proofstepwidestar[7,11]{w=(x\star y)\star z}{\rIE{14,10}}
    \proofstepwidestar[7,11]{%
      \exists x\in X\,\exists y\in Y\,\exists z\in Z\,
      \bigl(w=(x\star y)\star z\bigr)}{%
      \rEIStar{12,13,9,15}}
    \proofstepwidestar[1,2,3,4,5]{%
      \exists x\in X\,\exists y\in Y\,\exists z\in Z\,
      \bigl(w=(x\star y)\star z\bigr)}{%
      \rEEStar{7\text{--}16}}
  \closeproofpartwideindR

  \proofpartwideindR[Aufbau der links geklammerten Seite]{%
    \begin{aligned}[t]
      &\SemigroupAlg{A}{\star}\dsep X,Y,Z\in\PowNE{A}\dsep{}\\[-2pt]
      &\exists x\in X\,\exists y\in Y\,\exists z\in Z\,
        \bigl(w=(x\star y)\star z\bigr)\vdash{}\\[-2pt]
      &w\in (X\star_{\mathcal P}Y)\star_{\mathcal P}Z
    \end{aligned}
  }{%
    \begin{aligned}[t]
      &\SemigroupAlg{A}{\star}\dsep X,Y,Z\in\PowNE{A}\dsep{}\\[-2pt]
      &\exists x\in X\,\exists y\in Y\,\exists z\in Z\,
        \bigl(w=(x\star y)\star z\bigr)\vdash{}\\[-2pt]
      &w\in (X\star_{\mathcal P}Y)\star_{\mathcal P}Z
    \end{aligned}
  }[PowerSemigroupLeftTripleMembershipBackward]
    \proofstepwidestar[1]{\SemigroupAlg{A}{\star}}{\rA}
    \proofstepwidestar[2]{X\in\PowNE{A}}{\rA}
    \proofstepwidestar[3]{Y\in\PowNE{A}}{\rA}
    \proofstepwidestar[4]{Z\in\PowNE{A}}{\rA}
    \proofstepwidestar[5]{%
      \exists x\in X\,\exists y\in Y\,\exists z\in Z\,
      \bigl(w=(x\star y)\star z\bigr)}{\rA}
    \proofstepwidestar[5]{x\in X}{\rA}
    \proofstepwidestar[5]{y\in Y}{\rA}
    \proofstepwidestar[5]{z\in Z}{\rA}
    \proofstepwidestar[5]{w=(x\star y)\star z}{\rA}
    \proofstepwidestar[1,2,3,6,7]{%
      x\star y\in X\star_{\mathcal P}Y}{%
      \FormulaRefAuto{PowerSemigroupProductContains}[thm]{1,2,3,6,7}}
    \proofstepwidestar[1,2,3]{X\star_{\mathcal P}Y\in\PowNE{A}}{%
      \FormulaRefAuto{PowerSemigroupProductClosure}[thm]{1,2,3}}
    \proofstepwidestar[1,2,3,4,6,7,8]{%
      (x\star y)\star z
      \in (X\star_{\mathcal P}Y)\star_{\mathcal P}Z}{%
      \FormulaRefAuto{PowerSemigroupProductContains}[thm]{1,11,4,10,8}}
    \proofstepwidestar[9]{(x\star y)\star z=w}{%
      \FormulaRefAuto{a=b\vdash b=a}[thm]{9}}
    \proofstepwidestar[1,2,3,4,5]{%
      w\in (X\star_{\mathcal P}Y)\star_{\mathcal P}Z}{\rIE{13,12}}
    \proofstepwidestar[1,2,3,4,5]{%
      w\in (X\star_{\mathcal P}Y)\star_{\mathcal P}Z}{%
      \rEEStar{5\text{--}14}}
  \closeproofpartwideindR

  \proofpartwideindR[Aufbau der rechts geklammerten Seite]{%
    \begin{aligned}[t]
      &\SemigroupAlg{A}{\star}\dsep X,Y,Z\in\PowNE{A}\dsep{}\\[-2pt]
      &\exists x\in X\,\exists y\in Y\,\exists z\in Z\,
        \bigl(w=x\star(y\star z)\bigr)\vdash{}\\[-2pt]
      &w\in X\star_{\mathcal P}(Y\star_{\mathcal P}Z)
    \end{aligned}
  }{%
    \begin{aligned}[t]
      &\SemigroupAlg{A}{\star}\dsep X,Y,Z\in\PowNE{A}\dsep{}\\[-2pt]
      &\exists x\in X\,\exists y\in Y\,\exists z\in Z\,
        \bigl(w=x\star(y\star z)\bigr)\vdash{}\\[-2pt]
      &w\in X\star_{\mathcal P}(Y\star_{\mathcal P}Z)
    \end{aligned}
  }[PowerSemigroupRightTripleMembershipForward]
    \proofstepwidestar[1]{\SemigroupAlg{A}{\star}}{\rA}
    \proofstepwidestar[2]{X\in\PowNE{A}}{\rA}
    \proofstepwidestar[3]{Y\in\PowNE{A}}{\rA}
    \proofstepwidestar[4]{Z\in\PowNE{A}}{\rA}
    \proofstepwidestar[5]{%
      \exists x\in X\,\exists y\in Y\,\exists z\in Z\,
      \bigl(w=x\star(y\star z)\bigr)}{\rA}
    \proofstepwidestar[5]{x\in X}{\rA}
    \proofstepwidestar[5]{y\in Y}{\rA}
    \proofstepwidestar[5]{z\in Z}{\rA}
    \proofstepwidestar[5]{w=x\star(y\star z)}{\rA}
    \proofstepwidestar[1,3,4,7,8]{%
      y\star z\in Y\star_{\mathcal P}Z}{%
      \FormulaRefAuto{PowerSemigroupProductContains}[thm]{1,3,4,7,8}}
    \proofstepwidestar[1,3,4]{Y\star_{\mathcal P}Z\in\PowNE{A}}{%
      \FormulaRefAuto{PowerSemigroupProductClosure}[thm]{1,3,4}}
    \proofstepwidestar[1,2,3,4,6,7,8]{%
      x\star(y\star z)
      \in X\star_{\mathcal P}(Y\star_{\mathcal P}Z)}{%
      \FormulaRefAuto{PowerSemigroupProductContains}[thm]{1,2,11,6,10}}
    \proofstepwidestar[9]{x\star(y\star z)=w}{%
      \FormulaRefAuto{a=b\vdash b=a}[thm]{9}}
    \proofstepwidestar[1,2,3,4,5]{%
      w\in X\star_{\mathcal P}(Y\star_{\mathcal P}Z)}{\rIE{13,12}}
    \proofstepwidestar[1,2,3,4,5]{%
      w\in X\star_{\mathcal P}(Y\star_{\mathcal P}Z)}{%
      \rEEStar{5\text{--}14}}
  \closeproofpartwideindR

  \proofpartwideindR[Dreifachzeugen der rechts geklammerten Seite]{%
    \begin{aligned}[t]
      &\SemigroupAlg{A}{\star}\dsep X,Y,Z\in\PowNE{A}\dsep
        w\in X\star_{\mathcal P}(Y\star_{\mathcal P}Z)\vdash{}\\[-2pt]
      &\exists x\in X\,\exists y\in Y\,\exists z\in Z\,
        \bigl(w=x\star(y\star z)\bigr)
    \end{aligned}
  }{%
    \begin{aligned}[t]
      &\SemigroupAlg{A}{\star}\dsep X,Y,Z\in\PowNE{A}\dsep
        w\in X\star_{\mathcal P}(Y\star_{\mathcal P}Z)\vdash{}\\[-2pt]
      &\exists x\in X\,\exists y\in Y\,\exists z\in Z\,
        \bigl(w=x\star(y\star z)\bigr)
    \end{aligned}
  }[PowerSemigroupRightTripleMembershipBackward]
    \proofstepwidestar[1]{\SemigroupAlg{A}{\star}}{\rA}
    \proofstepwidestar[2]{X\in\PowNE{A}}{\rA}
    \proofstepwidestar[3]{Y\in\PowNE{A}}{\rA}
    \proofstepwidestar[4]{Z\in\PowNE{A}}{\rA}
    \proofstepwidestar[5]{%
      w\in X\star_{\mathcal P}(Y\star_{\mathcal P}Z)}{\rA}
    \proofstepwidestar[1,3,4]{Y\star_{\mathcal P}Z\in\PowNE{A}}{%
      \FormulaRefAuto{PowerSemigroupProductClosure}[thm]{1,3,4}}
    \proofstepwidestar[1,2,3,4,5]{%
      \exists x\in X\,
      \exists v\in Y\star_{\mathcal P}Z\,(w=x\star v)}{%
      \FormulaRefAuto{PowerSemigroupProductWitnesses}[thm]{1,2,6,5}}
    \proofstepwidestar[7]{x\in X}{\rA}
    \proofstepwidestar[7]{v\in Y\star_{\mathcal P}Z}{\rA}
    \proofstepwidestar[7]{w=x\star v}{\rA}
    \proofstepwidestar[1,3,4,9]{%
      \exists y\in Y\,\exists z\in Z\,(v=y\star z)}{%
      \FormulaRefAuto{PowerSemigroupProductWitnesses}[thm]{1,3,4,9}}
    \proofstepwidestar[11]{y\in Y}{\rA}
    \proofstepwidestar[11]{z\in Z}{\rA}
    \proofstepwidestar[11]{v=y\star z}{\rA}
    \proofstepwidestar[7,11]{w=x\star(y\star z)}{\rIE{14,10}}
    \proofstepwidestar[7,11]{%
      \exists x\in X\,\exists y\in Y\,\exists z\in Z\,
      \bigl(w=x\star(y\star z)\bigr)}{%
      \rEIStar{8,12,13,15}}
    \proofstepwidestar[1,2,3,4,5]{%
      \exists x\in X\,\exists y\in Y\,\exists z\in Z\,
      \bigl(w=x\star(y\star z)\bigr)}{%
      \rEEStar{7\text{--}16}}
  \closeproofpartwideindR

  \proofpartwideindR[Umklammern der Dreifachzeugen nach rechts]{%
    \begin{aligned}[t]
      &\SemigroupAlg{A}{\star}\dsep X,Y,Z\in\PowNE{A}\dsep{}\\[-2pt]
      &\exists x\in X\,\exists y\in Y\,\exists z\in Z\,
        \bigl(w=(x\star y)\star z\bigr)\vdash{}\\[-2pt]
      &\exists x\in X\,\exists y\in Y\,\exists z\in Z\,
        \bigl(w=x\star(y\star z)\bigr)
    \end{aligned}
  }{%
    \begin{aligned}[t]
      &\SemigroupAlg{A}{\star}\dsep X,Y,Z\in\PowNE{A}\dsep{}\\[-2pt]
      &\exists x\in X\,\exists y\in Y\,\exists z\in Z\,
        \bigl(w=(x\star y)\star z\bigr)\vdash{}\\[-2pt]
      &\exists x\in X\,\exists y\in Y\,\exists z\in Z\,
        \bigl(w=x\star(y\star z)\bigr)
    \end{aligned}
  }[PowerSemigroupTripleWitnessAssociativityForward]
    \proofstepwidestar[1]{\SemigroupAlg{A}{\star}}{\rA}
    \proofstepwidestar[2]{X\in\PowNE{A}}{\rA}
    \proofstepwidestar[3]{Y\in\PowNE{A}}{\rA}
    \proofstepwidestar[4]{Z\in\PowNE{A}}{\rA}
    \proofstepwidestar[5]{%
      \exists x\in X\,\exists y\in Y\,\exists z\in Z\,
      \bigl(w=(x\star y)\star z\bigr)}{\rA}
    \proofstepwidestar[5]{x\in X}{\rA}
    \proofstepwidestar[5]{y\in Y}{\rA}
    \proofstepwidestar[5]{z\in Z}{\rA}
    \proofstepwide[5]{w}{=}{(x\star y)\star z}{\rA}
    \proofstepwidestarbreak[1,2,3,4,6,7,8]{%
      (x\star y)\star z=x\star(y\star z)}{%
      \FormulaRefAuto{PowerSemigroupElementAssociativity}[thm]{%
        1,2,3,4,6,7,8}}
    \proofstepwidestar[1,2,3,4,5]{%
      w=x\star(y\star z)}{\rChain{9,10}}
    \proofstepwidestar[1,2,3,4,5]{%
      \exists x\in X\,\exists y\in Y\,\exists z\in Z\,
      \bigl(w=x\star(y\star z)\bigr)}{%
      \rEIStar{6,7,8,11}}
    \proofstepwidestar[1,2,3,4,5]{%
      \exists x\in X\,\exists y\in Y\,\exists z\in Z\,
      \bigl(w=x\star(y\star z)\bigr)}{%
      \rEEStar{5\text{--}12}}
  \closeproofpartwideindR

  \proofpartwideindR[Umklammern der Dreifachzeugen nach links]{%
    \begin{aligned}[t]
      &\SemigroupAlg{A}{\star}\dsep X,Y,Z\in\PowNE{A}\dsep{}\\[-2pt]
      &\exists x\in X\,\exists y\in Y\,\exists z\in Z\,
        \bigl(w=x\star(y\star z)\bigr)\vdash{}\\[-2pt]
      &\exists x\in X\,\exists y\in Y\,\exists z\in Z\,
        \bigl(w=(x\star y)\star z\bigr)
    \end{aligned}
  }{%
    \begin{aligned}[t]
      &\SemigroupAlg{A}{\star}\dsep X,Y,Z\in\PowNE{A}\dsep{}\\[-2pt]
      &\exists x\in X\,\exists y\in Y\,\exists z\in Z\,
        \bigl(w=x\star(y\star z)\bigr)\vdash{}\\[-2pt]
      &\exists x\in X\,\exists y\in Y\,\exists z\in Z\,
        \bigl(w=(x\star y)\star z\bigr)
    \end{aligned}
  }[PowerSemigroupTripleWitnessAssociativityBackward]
    \proofstepwidestar[1]{\SemigroupAlg{A}{\star}}{\rA}
    \proofstepwidestar[2]{X\in\PowNE{A}}{\rA}
    \proofstepwidestar[3]{Y\in\PowNE{A}}{\rA}
    \proofstepwidestar[4]{Z\in\PowNE{A}}{\rA}
    \proofstepwidestar[5]{%
      \exists x\in X\,\exists y\in Y\,\exists z\in Z\,
      \bigl(w=x\star(y\star z)\bigr)}{\rA}
    \proofstepwidestar[5]{x\in X}{\rA}
    \proofstepwidestar[5]{y\in Y}{\rA}
    \proofstepwidestar[5]{z\in Z}{\rA}
    \proofstepwidestarbreak[1,2,3,4,6,7,8]{%
      (x\star y)\star z=x\star(y\star z)}{%
      \FormulaRefAuto{PowerSemigroupElementAssociativity}[thm]{%
        1,2,3,4,6,7,8}}
    \proofstepwide[5]{w}{=}{x\star(y\star z)}{\rA}
    \proofstepwide[1,2,3,4,6,7,8]{}{=}{(x\star y)\star z}{%
      \FormulaRefAuto{a=b\vdash b=a}[thm]{9}}
    \proofstepwidestar[1,2,3,4,5]{%
      w=(x\star y)\star z}{\rChain{10,11}}
    \proofstepwidestar[1,2,3,4,5]{%
      \exists x\in X\,\exists y\in Y\,\exists z\in Z\,
      \bigl(w=(x\star y)\star z\bigr)}{%
      \rEIStar{6,7,8,12}}
    \proofstepwidestar[1,2,3,4,5]{%
      \exists x\in X\,\exists y\in Y\,\exists z\in Z\,
      \bigl(w=(x\star y)\star z\bigr)}{%
      \rEEStar{5\text{--}13}}
  \closeproofpartwideindR

  \proofpartwideindR[Zeugenform der links geklammerten Seite]{%
    \begin{aligned}[t]
      &\SemigroupAlg{A}{\star}\dsep X,Y,Z\in\PowNE{A}\vdash{}\\[-2pt]
      &w\in (X\star_{\mathcal P}Y)\star_{\mathcal P}Z
        \leftrightarrow
        \exists x\in X\,\exists y\in Y\,\exists z\in Z\,
        \bigl(w=(x\star y)\star z\bigr)
    \end{aligned}
  }{%
    \begin{aligned}[t]
      &\SemigroupAlg{A}{\star}\dsep X,Y,Z\in\PowNE{A}\vdash{}\\[-2pt]
      &w\in (X\star_{\mathcal P}Y)\star_{\mathcal P}Z
        \leftrightarrow
        \exists x\in X\,\exists y\in Y\,\exists z\in Z\,
        \bigl(w=(x\star y)\star z\bigr)
    \end{aligned}
  }[PowerSemigroupLeftTripleMembershipEquivalence]
    \proofstepwidestar[1]{\SemigroupAlg{A}{\star}}{\rA}
    \proofstepwidestar[2]{X\in\PowNE{A}}{\rA}
    \proofstepwidestar[3]{Y\in\PowNE{A}}{\rA}
    \proofstepwidestar[4]{Z\in\PowNE{A}}{\rA}
    \proofstepwidestar[1,2,3,4]{%
      \begin{aligned}[t]
        &w\in (X\star_{\mathcal P}Y)\star_{\mathcal P}Z
          \leftrightarrow{}\\[-2pt]
        &\exists x\in X\,\exists y\in Y\\[-2pt]
        &\quad\exists z\in Z\,\bigl(w=(x\star y)\star z\bigr)
      \end{aligned}}{%
      \rLRI{%
        \FormulaRefAuto{PowerSemigroupLeftTripleMembershipForward}[thm],
        \FormulaRefAuto{PowerSemigroupLeftTripleMembershipBackward}[thm]}}
  \closeproofpartwideindR

  \proofpartwideindR[Umklammern der Zeugen]{%
    \begin{aligned}[t]
      &\SemigroupAlg{A}{\star}\dsep X,Y,Z\in\PowNE{A}\vdash{}\\[-2pt]
      &\exists x\in X\,\exists y\in Y\,\exists z\in Z\,
        \bigl(w=(x\star y)\star z\bigr)
        \leftrightarrow
        \exists x\in X\,\exists y\in Y\,\exists z\in Z\,
        \bigl(w=x\star(y\star z)\bigr)
    \end{aligned}
  }{%
    \begin{aligned}[t]
      &\SemigroupAlg{A}{\star}\dsep X,Y,Z\in\PowNE{A}\vdash{}\\[-2pt]
      &\exists x\in X\,\exists y\in Y\,\exists z\in Z\,
        \bigl(w=(x\star y)\star z\bigr)
        \leftrightarrow
        \exists x\in X\,\exists y\in Y\,\exists z\in Z\,
        \bigl(w=x\star(y\star z)\bigr)
    \end{aligned}
  }[PowerSemigroupTripleWitnessAssociativityEquivalence]
    \proofstepwidestar[1]{\SemigroupAlg{A}{\star}}{\rA}
    \proofstepwidestar[2]{X\in\PowNE{A}}{\rA}
    \proofstepwidestar[3]{Y\in\PowNE{A}}{\rA}
    \proofstepwidestar[4]{Z\in\PowNE{A}}{\rA}
    \proofstepwidestar[1,2,3,4]{%
      \begin{aligned}[t]
        &\exists x\in X\,\exists y\in Y\\[-2pt]
        &\quad\exists z\in Z\,\bigl(w=(x\star y)\star z\bigr)\\[-2pt]
        &\leftrightarrow \exists x\in X\,\exists y\in Y\\[-2pt]
        &\quad\exists z\in Z\,\bigl(w=x\star(y\star z)\bigr)
      \end{aligned}}{%
      \rLRI{%
        \FormulaRefAuto{PowerSemigroupTripleWitnessAssociativityForward}[thm],
        \FormulaRefAuto{PowerSemigroupTripleWitnessAssociativityBackward}[thm]}}
  \closeproofpartwideindR

  \proofpartwideindR[Zeugenform der rechts geklammerten Seite]{%
    \begin{aligned}[t]
      &\SemigroupAlg{A}{\star}\dsep X,Y,Z\in\PowNE{A}\vdash{}\\[-2pt]
      &\exists x\in X\,\exists y\in Y\,\exists z\in Z\,
        \bigl(w=x\star(y\star z)\bigr)
        \leftrightarrow
        w\in X\star_{\mathcal P}(Y\star_{\mathcal P}Z)
    \end{aligned}
  }{%
    \begin{aligned}[t]
      &\SemigroupAlg{A}{\star}\dsep X,Y,Z\in\PowNE{A}\vdash{}\\[-2pt]
      &\exists x\in X\,\exists y\in Y\,\exists z\in Z\,
        \bigl(w=x\star(y\star z)\bigr)
        \leftrightarrow
        w\in X\star_{\mathcal P}(Y\star_{\mathcal P}Z)
    \end{aligned}
  }[PowerSemigroupRightTripleMembershipEquivalence]
    \proofstepwidestar[1]{\SemigroupAlg{A}{\star}}{\rA}
    \proofstepwidestar[2]{X\in\PowNE{A}}{\rA}
    \proofstepwidestar[3]{Y\in\PowNE{A}}{\rA}
    \proofstepwidestar[4]{Z\in\PowNE{A}}{\rA}
    \proofstepwidestar[1,2,3,4]{%
      \begin{aligned}[t]
        &\exists x\in X\,\exists y\in Y\\[-2pt]
        &\quad\exists z\in Z\,\bigl(w=x\star(y\star z)\bigr)
          \leftrightarrow{}\\[-2pt]
        &w\in X\star_{\mathcal P}(Y\star_{\mathcal P}Z)
      \end{aligned}}{%
      \rLRI{%
        \FormulaRefAuto{PowerSemigroupRightTripleMembershipForward}[thm],
        \FormulaRefAuto{PowerSemigroupRightTripleMembershipBackward}[thm]}}
  \closeproofpartwideindR

  \proofpartwideindR[Assoziativität der mengenweisen Operation]{%
    \begin{aligned}[t]
      &\SemigroupAlg{A}{\star}\dsep X,Y,Z\in\PowNE{A}\vdash{}\\[-2pt]
      &(X\star_{\mathcal P}Y)\star_{\mathcal P}Z
        =X\star_{\mathcal P}(Y\star_{\mathcal P}Z)
    \end{aligned}
  }{%
    \begin{aligned}[t]
      &\SemigroupAlg{A}{\star}\dsep X,Y,Z\in\PowNE{A}\vdash{}\\[-2pt]
      &(X\star_{\mathcal P}Y)\star_{\mathcal P}Z
        =X\star_{\mathcal P}(Y\star_{\mathcal P}Z)
    \end{aligned}
  }[PowerSemigroupAssociativity]
    \proofstepwidestar[1]{\SemigroupAlg{A}{\star}}{\rA}
    \proofstepwidestar[2]{X\in\PowNE{A}}{\rA}
    \proofstepwidestar[3]{Y\in\PowNE{A}}{\rA}
    \proofstepwidestar[4]{Z\in\PowNE{A}}{\rA}
    \proofstepwide[1,2,3,4]{%
      w\in (X\star_{\mathcal P}Y)\star_{\mathcal P}Z}{\leftrightarrow}{%
      \begin{aligned}[t]
        &\exists x\in X\,\exists y\in Y\\[-2pt]
        &\quad\exists z\in Z\,\bigl(w=(x\star y)\star z\bigr)
      \end{aligned}}{%
      \FormulaRefAuto{PowerSemigroupLeftTripleMembershipEquivalence}[thm]}
    \proofstepwide[1,2,3,4]{}{\leftrightarrow}{%
      \begin{aligned}[t]
        &\exists x\in X\,\exists y\in Y\\[-2pt]
        &\quad\exists z\in Z\,\bigl(w=x\star(y\star z)\bigr)
      \end{aligned}}{%
      \FormulaRefAuto{PowerSemigroupTripleWitnessAssociativityEquivalence}[thm]}
    \proofstepwide[1,2,3,4]{}{\leftrightarrow}{%
      w\in X\star_{\mathcal P}(Y\star_{\mathcal P}Z)}{%
      \FormulaRefAuto{PowerSemigroupRightTripleMembershipEquivalence}[thm]}
    \proofstepwide[1,2,3,4]{%
      w\in (X\star_{\mathcal P}Y)\star_{\mathcal P}Z}{\leftrightarrow}{%
      w\in X\star_{\mathcal P}(Y\star_{\mathcal P}Z)}{\rChain{5,7}}
    \proofstepwidestar[1,2,3,4]{%
      \forall w\,\Bigl(
        w\in (X\star_{\mathcal P}Y)\star_{\mathcal P}Z
        \leftrightarrow
        w\in X\star_{\mathcal P}(Y\star_{\mathcal P}Z)
      \Bigr)}{\rUI{8}}
    \proofstepwidestar[1,2,3,4]{%
      (X\star_{\mathcal P}Y)\star_{\mathcal P}Z
      =X\star_{\mathcal P}(Y\star_{\mathcal P}Z)}{%
      \FormulaRefAuto{\forall x\,(x\in A\leftrightarrow x\in B)
        \vdash A=B}[ax]{9}}
  \closeproofpartwideindR

  \proofpartwide{Schluss}
    \proofstepwidestar[1]{\SemigroupAlg{A}{\star}}{\rA}
    \proofstepwidestar[2]{X\in\PowNE{A}}{\rA}
    \proofstepwidestar[3]{Y\in\PowNE{A}}{\rA}
    \proofstepwidestar[1,2,3]{X\star_{\mathcal P}Y\in\PowNE{A}}{%
      \FormulaRefAuto{PowerSemigroupProductClosure}[thm]{1,2,3}}
    \proofstepwidestar[5]{Z\in\PowNE{A}}{\rA}
    \proofstepwidestar[1,2,3,5]{%
      (X\star_{\mathcal P}Y)\star_{\mathcal P}Z
      =X\star_{\mathcal P}(Y\star_{\mathcal P}Z)}{%
      \FormulaRefAuto{PowerSemigroupAssociativity}[thm]{1,2,3,5}}
    \proofstepwidestar[1]{\SemigroupAlg{\PowNE{A}}{\star_{\mathcal P}}}{%
      \FormulaRefAuto{SemigroupDef}[def]{4,6}}
  \closeproofpartwide
\end{tabproofsplitwide}

Die Halbgruppe \((\PowNE{A},\star_{\mathcal P})\) heißt die
\emph{Potenzhalbgruppe} von \((A,\star)\). Ihre Konstruktion setzt keine
Endlichkeit von \(A\) voraus.

\section{Weitere Beispiele}

\FormulaThmDeltaK[Die Addition bildet eine Halbgruppe]{%
  \SemigroupAlg{\mathbb{N}}{+}%
}{NaturalAdditionSemigroup}{%
  \DeltaRow{Mengen}{\mathbb{N}\dsep a\dsep b\dsep c}
}
\begin{tabproof}
  \proofstep{1}{a\in\mathbb{N}}{\rA}
  \proofstep{2}{b\in\mathbb{N}}{\rA}
  \proofstep{1,2}{a+b\in\mathbb{N}}{%
    \FormulaRefAuto{PeanoAddClosure}[thm]{1,2}}
  \proofstep{4}{c\in\mathbb{N}}{\rA}
  \proofstep{1,2,4}{(a+b)+c=a+(b+c)}{%
    \FormulaRefAuto{PeanoAddAssociative}[thm]{1,2,4}}
  \proofstep{}{\SemigroupAlg{\mathbb{N}}{+}}{%
    \FormulaRefAuto{SemigroupDef}[def]{3,5}}
\end{tabproof}

\FormulaThmDeltaK[Die Multiplikation bildet eine Halbgruppe]{%
  \SemigroupAlg{\mathbb{N}}{\cdot}%
}{NaturalMultiplicationSemigroup}{%
  \DeltaRow{Mengen}{\mathbb{N}\dsep a\dsep b\dsep c}
}
\begin{tabproof}
  \proofstep{1}{a\in\mathbb{N}}{\rA}
  \proofstep{2}{b\in\mathbb{N}}{\rA}
  \proofstep{1,2}{a\cdot b\in\mathbb{N}}{%
    \FormulaRefAuto{PeanoMulClosure}[thm]{1,2}}
  \proofstep{4}{c\in\mathbb{N}}{\rA}
  \proofstep{1,2,4}{(a\cdot b)\cdot c=a\cdot(b\cdot c)}{%
    \FormulaRefAuto{PeanoMulAssociative}[thm]{1,2,4}}
  \proofstep{}{\SemigroupAlg{\mathbb{N}}{\cdot}}{%
    \FormulaRefAuto{SemigroupDef}[def]{3,5}}
\end{tabproof}

\section{Potenzen in Halbgruppen}

In einer Halbgruppe gibt es im Allgemeinen kein neutrales Element. Daher
werden zun\"achst nur \emph{positive} Potenzen eingef\"uhrt. Die Konstruktion
wird auf die Rekursionsabbildung aus dem Band \emph{Nat\"urliche Zahlen}
zur\"uckgef\"uhrt. Auf diese Weise ist die Potenzfolge eine wirkliche
Abbildung und nicht lediglich eine informelle Rekursionsvorschrift.

\subsection{Rechtstranslation und Potenzfolge}

\FormulaDefDeltaK[Funktionsvorschrift der Rechtstranslation]{%
  \SemigroupAlg{A}{\star}\dsep a,x\in A\vdash
  \mathsf r_{A,\star,a}(x)\coloneqq x\star a%
}{SemigroupRightTranslationRuleDef}{%
  \DeltaRow{Mengen}{A\dsep a\dsep x}
  \DeltaRow{Funktionen}{\star\dsep\mathsf r_{A,\star,a}}
}

\FormulaThmDeltaK[Typisierung der Rechtstranslationsvorschrift]{%
  \SemigroupAlg{A}{\star}\dsep a,x\in A\vdash
  \mathsf r_{A,\star,a}(x)\in A%
}{SemigroupRightTranslationRuleClosure}{%
  \DeltaRow{Mengen}{A\dsep a\dsep x}
  \DeltaRow{Funktionen}{\star\dsep\mathsf r_{A,\star,a}}
}
\begin{tabproof}
  \proofstep{1}{\SemigroupAlg{A}{\star}}{\rA}
  \proofstep{2}{a\in A}{\rA}
  \proofstep{3}{x\in A}{\rA}
  \proofstep{1,2,3}{\mathsf r_{A,\star,a}(x)=x\star a}{%
    \FormulaRefAuto{SemigroupRightTranslationRuleDef}[def]{1,2,3}}
  \proofstep{1,2,3}{x\star a\in A}{%
    \FormulaRefAuto{SemigroupClosureAxiom}[ax]{1,3,2}}
  \proofstep{1,2,3}{\mathsf r_{A,\star,a}(x)\in A}{\rIE{4,5}}
\end{tabproof}

\FormulaThmDeltaK[Existenz der Rechtstranslation]{%
  \begin{aligned}[t]
    &\SemigroupAlg{A}{\star}\dsep a\in A\vdash{}\\[-2pt]
    &\exists! f\colon A\to A\,
      \forall x\in A\,\bigl(f(x)=\mathsf r_{A,\star,a}(x)\bigr)
  \end{aligned}
}{SemigroupRightTranslationExists}{%
  \DeltaRow{Mengen}{A\dsep a\dsep x}
  \DeltaRow{Funktionen}{\star\dsep f\dsep\mathsf r_{A,\star,a}}
}
\begin{tabproofwide}
  \proofstepwidestar[1]{\SemigroupAlg{A}{\star}}{\rA}
  \proofstepwidestar[2]{a\in A}{\rA}
  \proofstepwidestar[3]{x\in A}{\rA}
  \proofstepwidestar[1,2,3]{\mathsf r_{A,\star,a}(x)\in A}{%
    \FormulaRefAuto{SemigroupRightTranslationRuleClosure}[thm]{1,2,3}}
  \proofstepwidestar[1,2]{%
    \forall x\in A\,\bigl(\mathsf r_{A,\star,a}(x)\in A\bigr)}{%
    \rUI{\rRI{3,4}}}
  \proofstepwidestar[1,2]{%
    \exists! f\colon A\to A\,
    \forall x\in A\,\bigl(f(x)=\mathsf r_{A,\star,a}(x)\bigr)}{%
    \FormulaRefAuto{x\in A\vdash t(x)\coloneq y\dsep
      x\in A\vdash t(x)\in B
      \exists! F\colon A\to B\,\forall x\in A\,F(x)=t(x)}{5}}
\end{tabproofwide}

\FormulaDefDeltaK[Rechtstranslation]{%
  \begin{aligned}[t]
    &\SemigroupAlg{A}{\star}\dsep a\in A\vdash{}\\[-2pt]
    &\rho_{A,\star,a}\coloneqq
      \iota f\Bigl(f\colon A\to A\land
        \forall x\in A\,\bigl(f(x)=\mathsf r_{A,\star,a}(x)\bigr)\Bigr)
  \end{aligned}
}{SemigroupRightTranslationDef}{%
  \DeltaRow{Mengen}{A\dsep a\dsep x}
  \DeltaRow{Funktionen}{\star\dsep f\dsep\rho_{A,\star,a}}
  \DeltaRow{\textbf{Neues Symbol}}{}
  \DeltaPrem{Rechtstranslation}{\rho_{A,\star,a}}
}

\FormulaThmDeltaK[Die Rechtstranslation ist eine Abbildung]{%
  \SemigroupAlg{A}{\star}\dsep a\in A\vdash
  \rho_{A,\star,a}\colon A\to A%
}{SemigroupRightTranslationFunction}{%
  \DeltaRow{Mengen}{A\dsep a\dsep x}
  \DeltaRow{Funktionen}{\star\dsep\rho_{A,\star,a}}
}
\begin{tabproofwide}
  \proofstepwidestar[1]{\SemigroupAlg{A}{\star}}{\rA}
  \proofstepwidestar[2]{a\in A}{\rA}
  \proofstepwidestar[1,2]{%
    \rho_{A,\star,a}\colon A\to A\land
    \forall x\in A\,
      \bigl(\rho_{A,\star,a}(x)=\mathsf r_{A,\star,a}(x)\bigr)}{%
    \FormulaRefAuto{SemigroupRightTranslationDef}[def]{%
      \FormulaRefAuto{SemigroupRightTranslationExists}[thm]{1,2}}}
  \proofstepwidestar[1,2]{\rho_{A,\star,a}\colon A\to A}{\rAEa{3}}
\end{tabproofwide}

\FormulaThmDeltaK[Auswertung der Rechtstranslation]{%
  \SemigroupAlg{A}{\star}\dsep a,x\in A\vdash
  \rho_{A,\star,a}(x)=x\star a%
}{SemigroupRightTranslationEvaluation}{%
  \DeltaRow{Mengen}{A\dsep a\dsep x}
  \DeltaRow{Funktionen}{\star\dsep\rho_{A,\star,a}}
}
\begin{tabproofwide}
  \proofstepwidestar[1]{\SemigroupAlg{A}{\star}}{\rA}
  \proofstepwidestar[2]{a\in A}{\rA}
  \proofstepwidestar[3]{x\in A}{\rA}
  \proofstepwidestar[1,2]{%
    \forall u\in A\,
      \bigl(\rho_{A,\star,a}(u)=\mathsf r_{A,\star,a}(u)\bigr)}{%
    \rAEb{\FormulaRefAuto{SemigroupRightTranslationDef}[def]{%
      \FormulaRefAuto{SemigroupRightTranslationExists}[thm]{1,2}}}}
  \proofstepwidestar[1,2,3]{%
    \rho_{A,\star,a}(x)=\mathsf r_{A,\star,a}(x)}{\rRE{3,\rUE{4}}}
  \proofstepwidestar[1,2,3]{\mathsf r_{A,\star,a}(x)=x\star a}{%
    \FormulaRefAuto{SemigroupRightTranslationRuleDef}[def]{1,2,3}}
  \proofstepwidestar[1,2,3]{\rho_{A,\star,a}(x)=x\star a}{%
    \FormulaRefAuto{a = b,\, b = c \vdash a = c}[thm]{5,6}}
\end{tabproofwide}

\FormulaDefDeltaK[Potenzfolge eines Halbgruppenelements]{%
  \SemigroupAlg{A}{\star}\dsep a\in A\vdash
  \operatorname{pow}_{A,\star}(a)\coloneqq
  \RecFun{a}{\rho_{A,\star,a}}%
}{SemigroupPowerSequenceDef}{%
  \DeltaRow{Mengen}{A\dsep a}
  \DeltaRow{Funktionen}{\star\dsep\rho_{A,\star,a}\dsep
    \operatorname{pow}_{A,\star}}
  \DeltaRow{\textbf{Neues Symbol}}{}
  \DeltaPrem{Potenzfolge}{\operatorname{pow}_{A,\star}(a)}
}

\FormulaThmDeltaK[Die Potenzfolge ist eine Abbildung]{%
  \SemigroupAlg{A}{\star}\dsep a\in A\vdash
  \operatorname{pow}_{A,\star}(a)\colon\mathbb N\to A%
}{SemigroupPowerSequenceFunction}{%
  \DeltaRow{Mengen}{A\dsep a}
  \DeltaRow{Funktionen}{\star\dsep\operatorname{pow}_{A,\star}}
}
\begin{tabproofwide}
  \proofstepwidestar[1]{\SemigroupAlg{A}{\star}}{\rA}
  \proofstepwidestar[2]{a\in A}{\rA}
  \proofstepwidestar[1,2]{\rho_{A,\star,a}\colon A\to A}{%
    \FormulaRefAuto{SemigroupRightTranslationFunction}[thm]{1,2}}
  \proofstepwidestar[1,2]{%
    \RecFun{a}{\rho_{A,\star,a}}\colon\mathbb N\to A}{%
    \FormulaRefAuto{m_0\in M\dsep f\colon M\to M\vdash
      \RecFun{m_0}{f}\colon\mathbb{N}\to M}[thm]{2,3}}
  \proofstepwidestar[1,2]{%
    \operatorname{pow}_{A,\star}(a)=
    \RecFun{a}{\rho_{A,\star,a}}}{%
    \FormulaRefAuto{SemigroupPowerSequenceDef}[def]{1,2}}
  \proofstepwidestar[1,2]{%
    \operatorname{pow}_{A,\star}(a)\colon\mathbb N\to A}{\rIE{5,4}}
\end{tabproofwide}

\subsection{Positive Potenzen}

Ist \(n\in\mathbb N\) und \(n\neq0\), so ist
\(n=\Succ(\Pred(n))\). Der Vorgänger dient daher nur dazu, die bei \(0\)
beginnende Potenzfolge mit den bei \(1\) beginnenden Exponenten zu
synchronisieren.

\FormulaDefDeltaK[Positive Potenz eines Halbgruppenelements]{%
  \begin{aligned}[t]
    &\SemigroupAlg{A}{\star}\dsep a\in A\dsep
      n\in\mathbb N\dsep n\neq0\vdash{}\\[-2pt]
    &a^n\coloneqq
      \operatorname{pow}_{A,\star}(a)(\Pred(n))
  \end{aligned}
}{SemigroupPositivePowerDef}{%
  \DeltaRow{Mengen}{A\dsep a\dsep n}
  \DeltaRow{Funktionen}{\star\dsep\operatorname{pow}_{A,\star}}
  \DeltaRow{\textbf{Neues Symbol}}{}
  \DeltaPrem{Positive Potenz}{a^n}
}

Wegen \(n+1=\Succ(n)\) formulieren wir die Nachfolgeraussagen im Folgenden
mit der Addition auf den natürlichen Zahlen.  Diese Schreibweise passt
unmittelbar zu den anschließenden Potenzgesetzen.

\FormulaThmDeltaK[Nachfolgerindex der Potenzfolge]{%
  \SemigroupAlg{A}{\star}\dsep a\in A\dsep n\in\mathbb N\vdash
  a^{n+1}=\operatorname{pow}_{A,\star}(a)(n)%
}{SemigroupPowerSuccessorIndex}{%
  \DeltaRow{Mengen}{A\dsep a\dsep n}
  \DeltaRow{Funktionen}{\star\dsep\operatorname{pow}_{A,\star}}
}
\begin{tabproofwide}
  \proofstepwidestar[1]{\SemigroupAlg{A}{\star}}{\rA}
  \proofstepwidestar[2]{a\in A}{\rA}
  \proofstepwidestar[3]{n\in\mathbb N}{\rA}
  \proofstepwidestar[3]{n+1=\Succ(n)}{%
    \FormulaRefAuto{PeanoAddOneSucc}[thm]{3}}
  \proofstepwidestar[3]{n+1\in\mathbb N}{%
    \rIE{4,\FormulaRefAuto{PeanoSuccClosure}[ax]{3}}}
  \proofstepwidestar[3]{n+1\neq0}{%
    \rIE{4,\FormulaRefAuto{PeanoSuccNonzero}[ax]{3}}}
  \proofstepwidestar[1,2,3]{%
    a^{n+1}=\operatorname{pow}_{A,\star}(a)(\Pred(n+1))}{%
    \FormulaRefAuto{SemigroupPositivePowerDef}[def]{1,2,5,6}}
  \proofstepwidestar[3]{\Pred(n+1)=n}{%
    \rIE{4,\FormulaRefAuto{PeanoPredSucc}[thm]{3}}}
  \proofstepwidestar[1,2,3]{%
    \operatorname{pow}_{A,\star}(a)(\Pred(n+1))=
    \operatorname{pow}_{A,\star}(a)(n)}{\rIE{8}}
  \proofstepwidestar[1,2,3]{%
    a^{n+1}=\operatorname{pow}_{A,\star}(a)(n)}{%
    \FormulaRefAuto{a = b,\, b = c \vdash a = c}[thm]{7,9}}
\end{tabproofwide}

\FormulaThmDeltaK[Erste positive Potenz]{%
  \SemigroupAlg{A}{\star}\dsep a\in A\vdash a^1=a%
}{SemigroupPowerOne}{%
  \DeltaRow{Mengen}{A\dsep a}
  \DeltaRow{Funktionen}{\star\dsep\operatorname{pow}_{A,\star}}
}
\begin{tabproofwide}
  \proofstepwidestar[1]{\SemigroupAlg{A}{\star}}{\rA}
  \proofstepwidestar[2]{a\in A}{\rA}
  \proofstepwidestar[]{0\in\mathbb N}{\FormulaRefAuto{PeanoZero}[ax]}
  \proofstepwidestar[1,2]{%
    a^{0+1}=\operatorname{pow}_{A,\star}(a)(0)}{%
    \FormulaRefAuto{SemigroupPowerSuccessorIndex}[thm]{1,2,3}}
  \proofstepwidestar[1,2]{%
    \operatorname{pow}_{A,\star}(a)(0)=
    \RecFun{a}{\rho_{A,\star,a}}(0)}{%
    \FormulaRefAuto{SemigroupPowerSequenceDef}[def]{1,2}}
  \proofstepwidestar[1,2]{\rho_{A,\star,a}\colon A\to A}{%
    \FormulaRefAuto{SemigroupRightTranslationFunction}[thm]{1,2}}
  \proofstepwidestar[1,2]{\RecFun{a}{\rho_{A,\star,a}}(0)=a}{%
    \FormulaRefAuto{m_0\in M\dsep f\colon M\to M\vdash
      \RecFun{m_0}{f}(0)=m_0}[thm]{2,6}}
  \proofstepwidestar[1,2]{a^{0+1}=a}{%
    \FormulaRefAuto{a = b,\, b = c \vdash a = c}[thm]{%
      4,\FormulaRefAuto{a = b,\, b = c \vdash a = c}[thm]{5,7}}}
  \proofstepwidestar[]{0+1=1}{%
    \FormulaRefAuto{PeanoAddLeftZero}[thm]{%
      \FormulaRefAuto{PeanoOneInN}[thm]}}
  \proofstepwidestar[1,2]{a^1=a}{\rIE{9,8}}
\end{tabproofwide}

\FormulaThmDeltaK[Nachfolgerregel positiver Potenzen]{%
  \begin{aligned}[t]
    &\SemigroupAlg{A}{\star}\dsep a\in A\dsep
      n\in\mathbb N\dsep n\neq0\vdash{}\\[-2pt]
    &a^{n+1}=a^n\star a
  \end{aligned}
}{SemigroupPowerSuccessor}{%
  \DeltaRow{Mengen}{A\dsep a\dsep n}
  \DeltaRow{Funktionen}{\star\dsep\operatorname{pow}_{A,\star}}
}
\begin{tabproofwide}
  \proofstepwidestar[1]{\SemigroupAlg{A}{\star}}{\rA}
  \proofstepwidestar[2]{a\in A}{\rA}
  \proofstepwidestar[3]{n\in\mathbb N}{\rA}
  \proofstepwidestar[4]{n\neq0}{\rA}
  \proofstepwidestar[1,2,3]{%
    a^{n+1}=\operatorname{pow}_{A,\star}(a)(n)}{%
    \FormulaRefAuto{SemigroupPowerSuccessorIndex}[thm]{1,2,3}}
  \proofstepwidestar[1,2]{\rho_{A,\star,a}\colon A\to A}{%
    \FormulaRefAuto{SemigroupRightTranslationFunction}[thm]{1,2}}
  \proofstepwidestarbreak[1,2,3]{%
    \operatorname{pow}_{A,\star}(a)(n)=
    \RecFun{a}{\rho_{A,\star,a}}(n)}{%
    \FormulaRefAuto{SemigroupPowerSequenceDef}[def]{1,2}}
  \proofstepwidestarbreak[1,2,3,4]{%
    \RecFun{a}{\rho_{A,\star,a}}(n)=
    \rho_{A,\star,a}
      (\RecFun{a}{\rho_{A,\star,a}}(\Pred(n)))}{%
    \FormulaRefAuto{m_0\in M\dsep f\colon M\to M\dsep
      n\in\mathbb{N}\dsep n\neq0\vdash
      \RecFun{m_0}{f}(n)=f(\RecFun{m_0}{f}(\Pred(n)))}[thm]{2,6,3,4}}
  \proofstepwidestarbreak[1,2,3]{%
    \RecFun{a}{\rho_{A,\star,a}}(\Pred(n))\in A}{%
    \FormulaRefAuto{m_0\in M\dsep f\colon M\to M\dsep k\in\mathbb{N}
      \vdash\RecFun{m_0}{f}(k)\in M}[thm]{%
      2,6,\FormulaRefAuto{PeanoPredNatural}[thm]{3}}}
  \proofstepwidestarbreak[1,2,3]{%
    \rho_{A,\star,a}
      (\RecFun{a}{\rho_{A,\star,a}}(\Pred(n)))
    =\RecFun{a}{\rho_{A,\star,a}}(\Pred(n))\star a}{%
    \FormulaRefAuto{SemigroupRightTranslationEvaluation}[thm]{1,2,9}}
  \proofstepwidestarbreak[1,2,3,4]{%
    a^n=\operatorname{pow}_{A,\star}(a)(\Pred(n))}{%
    \FormulaRefAuto{SemigroupPositivePowerDef}[def]{1,2,3,4}}
  \proofstepwidestarbreak[1,2,3,4]{%
    a^n=\RecFun{a}{\rho_{A,\star,a}}(\Pred(n))}{%
    \rIE{\FormulaRefAuto{SemigroupPowerSequenceDef}[def]{1,2},11}}
  \proofstepwidestarbreak[1,2,3,4]{%
    \RecFun{a}{\rho_{A,\star,a}}(\Pred(n))\star a=a^n\star a}{%
    \rIE{\FormulaRefAuto{a=b\vdash b=a}[thm]{12}}}
  \proofstepwidestarbreak[1,2,3,4]{a^{n+1}=a^n\star a}{%
    \FormulaRefAuto{a = b,\, b = c \vdash a = c}[thm]{%
      5,\FormulaRefAuto{a = b,\, b = c \vdash a = c}[thm]{%
        7,\FormulaRefAuto{a = b,\, b = c \vdash a = c}[thm]{%
          8,\FormulaRefAuto{a = b,\, b = c \vdash a = c}[thm]{10,13}}}}}
\end{tabproofwide}

\FormulaThmDeltaK[Abgeschlossenheit positiver Potenzen]{%
  \SemigroupAlg{A}{\star}\dsep a\in A\dsep
  n\in\mathbb N\dsep n\neq0\vdash a^n\in A%
}{SemigroupPowerClosure}{%
  \DeltaRow{Mengen}{A\dsep a\dsep n}
  \DeltaRow{Funktionen}{\star\dsep\operatorname{pow}_{A,\star}}
}
\begin{tabproofwide}
  \proofstepwidestar[1]{\SemigroupAlg{A}{\star}}{\rA}
  \proofstepwidestar[2]{a\in A}{\rA}
  \proofstepwidestar[3]{n\in\mathbb N}{\rA}
  \proofstepwidestar[4]{n\neq0}{\rA}
  \proofstepwidestar[1,2]{%
    \operatorname{pow}_{A,\star}(a)\colon\mathbb N\to A}{%
    \FormulaRefAuto{SemigroupPowerSequenceFunction}[thm]{1,2}}
  \proofstepwidestar[3]{\Pred(n)\in\mathbb N}{%
    \FormulaRefAuto{PeanoPredNatural}[thm]{3}}
  \proofstepwidestar[1,2,3]{%
    \operatorname{pow}_{A,\star}(a)(\Pred(n))\in A}{%
    \FormulaRefAuto{F\colon A\to B\dsep x\in A\vdash F(x)\in B}[thm]{5,6}}
  \proofstepwidestar[1,2,3,4]{%
    a^n=\operatorname{pow}_{A,\star}(a)(\Pred(n))}{%
    \FormulaRefAuto{SemigroupPositivePowerDef}[def]{1,2,3,4}}
  \proofstepwidestar[1,2,3,4]{a^n\in A}{\rIE{8,7}}
\end{tabproofwide}

\subsection{Potenzgesetze}

Die folgenden Gesetze zeigen, dass die Rekursionsklammerung wegen der
Assoziativit\"at keine zus\"atzliche Information tr\"agt.

\FormulaThmDeltaK[Additionsgesetz positiver Potenzen]{%
  \begin{aligned}[t]
    &\SemigroupAlg{A}{\star}\dsep a\in A\dsep
      m,n\in\mathbb N\dsep m\neq0\dsep n\neq0\vdash{}\\[-2pt]
    &a^{m+n}=a^m\star a^n
  \end{aligned}
}{SemigroupPowerAdditionLaw}{%
  \DeltaRow{Mengen}{A\dsep a\dsep m\dsep n\dsep k}
  \DeltaRow{Funktionen}{\star}
}
\begin{tabproofsplitwide}
  \proofpartwideindR[Induktionsanfang]{%
    \begin{aligned}[t]
      &\SemigroupAlg{A}{\star}\dsep a\in A\dsep
        m\in\mathbb N\dsep m\neq0\vdash{}\\[-2pt]
      &a^{m+\Succ(0)}=a^m\star a^{\Succ(0)}
    \end{aligned}
  }{%
    \SemigroupAlg{A}{\star}\dsep a\in A\dsep
    m\in\mathbb N\dsep m\neq0\vdash
    a^{m+\Succ(0)}=a^m\star a^{\Succ(0)}
  }[SemigroupPowerAdditionBase]
    \proofstepwidestar[1]{\SemigroupAlg{A}{\star}}{\rA}
    \proofstepwidestar[2]{a\in A}{\rA}
    \proofstepwidestar[3]{m\in\mathbb N}{\rA}
    \proofstepwidestar[4]{m\neq0}{\rA}
    \proofstepwidestar[]{1=\Succ(0)}{\FormulaRefAuto{PeanoOne}[def]}
    \proofstepwidestar[1,2,3,4]{a^{m+1}=a^m\star a}{%
      \FormulaRefAuto{SemigroupPowerSuccessor}[thm]{1,2,3,4}}
    \proofstepwidestar[1,2]{a^1=a}{%
      \FormulaRefAuto{SemigroupPowerOne}[thm]{1,2}}
    \proofstepwidestar[1,2,3,4]{%
      a^m\star a=a^m\star a^1}{%
      \rIE{\FormulaRefAuto{a=b\vdash b=a}[thm]{7}}}
    \proofstepwidestar[1,2,3,4]{%
      a^{m+\Succ(0)}=a^m\star a^{\Succ(0)}}{%
      \rIE{5,\FormulaRefAuto{a = b,\, b = c \vdash a = c}[thm]{6,8}}}
  \closeproofpartwideindR

  \proofpartwideindR[Induktionsschritt]{%
    \begin{aligned}[t]
      &\SemigroupAlg{A}{\star}\dsep a\in A\dsep
        m,k\in\mathbb N\dsep m\neq0\dsep
        a^{m+\Succ(k)}=a^m\star a^{\Succ(k)}\vdash{}\\[-2pt]
      &a^{m+\Succ(\Succ(k))}=a^m\star a^{\Succ(\Succ(k))}
    \end{aligned}
  }{%
    \begin{aligned}[t]
      &\SemigroupAlg{A}{\star}\dsep a\in A\dsep
        m,k\in\mathbb N\dsep m\neq0\dsep
        a^{m+\Succ(k)}=a^m\star a^{\Succ(k)}\vdash{}\\[-2pt]
      &a^{m+\Succ(\Succ(k))}=a^m\star a^{\Succ(\Succ(k))}
    \end{aligned}
  }[SemigroupPowerAdditionStep]
    \proofstepwidestar[1]{\SemigroupAlg{A}{\star}}{\rA}
    \proofstepwidestar[2]{a\in A}{\rA}
    \proofstepwidestar[3]{m\in\mathbb N}{\rA}
    \proofstepwidestar[4]{k\in\mathbb N}{\rA}
    \proofstepwidestar[5]{m\neq0}{\rA}
    \proofstepwidestar[6]{a^{m+\Succ(k)}=a^m\star a^{\Succ(k)}}{\rA}
    \proofstepwidestar[4]{\Succ(k)\in\mathbb N}{%
      \FormulaRefAuto{PeanoSuccClosure}[ax]{4}}
    \proofstepwidestar[3,4]{m+\Succ(k)\in\mathbb N}{%
      \FormulaRefAuto{PeanoAddClosure}[thm]{3,7}}
    \proofstepwidestar[3,4]{m+\Succ(k)=\Succ(m+k)}{%
      \FormulaRefAuto{PeanoAddSuccRight}[thm]{3,4}}
    \proofstepwidestarbreak[3,4]{m+\Succ(k)\neq0}{%
      \rIE{\FormulaRefAuto{a=b\vdash b=a}[thm]{9},
        \FormulaRefAuto{PeanoSuccNonzero}[ax]{%
          \FormulaRefAuto{PeanoAddClosure}[thm]{3,4}}}}
    \proofstepwidestarbreak[3,4]{%
      m+\Succ(\Succ(k))=(m+\Succ(k))+1}{%
      \FormulaRefAuto{a = b,\, b = c \vdash a = c}[thm]{%
        \FormulaRefAuto{PeanoAddSuccRight}[thm]{3,7},
        \FormulaRefAuto{a=b\vdash b=a}[thm]{%
          \FormulaRefAuto{PeanoAddOneSucc}[thm]{8}}}}
    \proofstepwidestar[1,2,3,4]{%
      a^{(m+\Succ(k))+1}=a^{m+\Succ(k)}\star a}{%
      \FormulaRefAuto{SemigroupPowerSuccessor}[thm]{1,2,8,10}}
    \proofstepwidestar[1,2,3,4,5,6]{%
      a^{m+\Succ(k)}\star a
      =(a^m\star a^{\Succ(k)})\star a}{\rIE{6}}
    \proofstepwidestarbreak[1,2,3,4,5]{%
      (a^m\star a^{\Succ(k)})\star a
      =a^m\star(a^{\Succ(k)}\star a)}{%
      \FormulaRefAuto{SemigroupAssociativityAxiom}[ax]{%
        1,
        \FormulaRefAuto{SemigroupPowerClosure}[thm]{1,2,3,5},
        \FormulaRefAuto{SemigroupPowerClosure}[thm]{%
          1,2,7,\FormulaRefAuto{PeanoSuccNonzero}[ax]{4}},2}}
    \proofstepwidestarbreak[1,2,4]{%
      a^{\Succ(k)+1}=a^{\Succ(k)}\star a}{%
      \FormulaRefAuto{SemigroupPowerSuccessor}[thm]{%
        1,2,7,\FormulaRefAuto{PeanoSuccNonzero}[ax]{4}}}
    \proofstepwidestarbreak[1,2,3,4,5]{%
      a^m\star(a^{\Succ(k)}\star a)
      =a^m\star a^{\Succ(\Succ(k))}}{%
      \rIE{\FormulaRefAuto{PeanoAddOneSucc}[thm]{7},
        \FormulaRefAuto{a=b\vdash b=a}[thm]{15}}}
    \proofstepwidestarbreak[1,2,3,4,5,6]{%
      a^{m+\Succ(\Succ(k))}=a^m\star a^{\Succ(\Succ(k))}}{%
      \rIE{11,\FormulaRefAuto{a = b,\, b = c \vdash a = c}[thm]{%
        12,\FormulaRefAuto{a = b,\, b = c \vdash a = c}[thm]{%
          13,\FormulaRefAuto{a = b,\, b = c \vdash a = c}[thm]{14,16}}}}}
  \closeproofpartwideindR

  \proofpartwide{Induktionsschluss und positiver Exponent}
    \proofstepwidestar[1]{\SemigroupAlg{A}{\star}}{\rA}
    \proofstepwidestar[2]{a\in A}{\rA}
    \proofstepwidestar[3]{m\in\mathbb N}{\rA}
    \proofstepwidestar[4]{n\in\mathbb N}{\rA}
    \proofstepwidestar[5]{m\neq0}{\rA}
    \proofstepwidestar[6]{n\neq0}{\rA}
    \proofstepwidestarbreak[1,2,3,5]{%
      \forall k\in\mathbb N\,
      \bigl(a^{m+\Succ(k)}=a^m\star a^{\Succ(k)}\bigr)}{%
      \FormulaRefAuto{PeanoInductionRule}[ax]{%
        \FormulaRefAuto{SemigroupPowerAdditionBase}[thm]{1,2,3,5},
        \FormulaRefAuto{SemigroupPowerAdditionStep}[thm]{1,2,3,5}}}
    \proofstepwidestar[4,6]{n=\Succ(\Pred(n))}{%
      \FormulaRefAuto{PeanoPredNonzeroValue}[thm]{4,6}}
    \proofstepwidestar[4]{\Pred(n)\in\mathbb N}{%
      \FormulaRefAuto{PeanoPredNatural}[thm]{4}}
    \proofstepwidestar[1,2,3,4,5]{%
      a^{m+\Succ(\Pred(n))}=a^m\star a^{\Succ(\Pred(n))}}{%
      \rRE{9,\rUE{7}}}
    \proofstepwidestar[1,2,3,4,5,6]{a^{m+n}=a^m\star a^n}{\rIE{8,10}}
  \closeproofpartwide
\end{tabproofsplitwide}

\FormulaThmDeltaK[Potenzen desselben Elements kommutieren]{%
  \begin{aligned}[t]
    &\SemigroupAlg{A}{\star}\dsep a\in A\dsep
      m,n\in\mathbb N\dsep m\neq0\dsep n\neq0\vdash{}\\[-2pt]
    &a^m\star a^n=a^n\star a^m
  \end{aligned}
}{SemigroupPowersCommute}{%
  \DeltaRow{Mengen}{A\dsep a\dsep m\dsep n}
  \DeltaRow{Funktionen}{\star}
}
\begin{tabproofwide}
  \proofstepwidestar[1]{\SemigroupAlg{A}{\star}}{\rA}
  \proofstepwidestar[2]{a\in A}{\rA}
  \proofstepwidestar[3]{m\in\mathbb N}{\rA}
  \proofstepwidestar[4]{n\in\mathbb N}{\rA}
  \proofstepwidestar[5]{m\neq0}{\rA}
  \proofstepwidestar[6]{n\neq0}{\rA}
  \proofstepwide[1,2,3,4,5,6]{a^m\star a^n}{=}{a^{m+n}}{%
    \FormulaRefAuto{a=b\vdash b=a}[thm]{%
      \FormulaRefAuto{SemigroupPowerAdditionLaw}[thm]{1,2,3,4,5,6}}}
  \proofstepwide[3,4]{}{=}{a^{n+m}}{%
    \rIE{\FormulaRefAuto{PeanoAddCommutative}[thm]{3,4}}}
  \proofstepwide[1,2,3,4,5,6]{}{=}{a^n\star a^m}{%
    \FormulaRefAuto{SemigroupPowerAdditionLaw}[thm]{1,2,4,3,6,5}}
  \proofstepwide[1,2,3,4,5,6]{a^m\star a^n}{=}{a^n\star a^m}{\rChain{7,8,9}}
\end{tabproofwide}

\section{Idempotente Elemente}

Ein Element ist idempotent, wenn seine zweite Potenz bereits wieder das
Element selbst ist. F\"ur eine feste Halbgruppe sammeln wir alle solchen
Elemente in einer ausgezeichneten Teilmenge des Tr\"agers.

\FormulaDefDeltaK[Idempotentenmenge einer Halbgruppe]{%
  \SemigroupAlg{A}{\star}\vdash
  E_{A,\star}\coloneqq
  \bigl\{e\in A\bigm|e\star e=e\bigr\}%
}{SemigroupIdempotentSetDef}{%
  \DeltaRow{Mengen}{A\dsep e}
  \DeltaRow{Funktionen}{\star}
  \DeltaRow{\textbf{Neues Symbol}}{}
  \DeltaPrem{Idempotentenmenge}{E_{A,\star}}
}

\FormulaThmDeltaK[Elementkriterium f\"ur Idempotente]{%
  \SemigroupAlg{A}{\star}\vdash
  e\in E_{A,\star}\leftrightarrow
  \bigl(e\in A\land e\star e=e\bigr)%
}{SemigroupIdempotentMembership}{%
  \DeltaRow{Mengen}{A\dsep e}
  \DeltaRow{Funktionen}{\star}
}
\begin{tabproofwide}
  \proofstepwidestar[1]{\SemigroupAlg{A}{\star}}{\rA}
  \proofstepwidestar[1]{%
    E_{A,\star}=\{u\in A\mid u\star u=u\}}{%
    \FormulaRefAuto{SemigroupIdempotentSetDef}[def]{1}}
  \proofstepwidestar[1]{%
    e\in E_{A,\star}\leftrightarrow
    \bigl(e\in A\land e\star e=e\bigr)}{%
    \rIE{2,\FormulaRefAuto{x\in\{u\in A\mid P(u)\}\eqvdash
      x\in A\land P(x)}[thm]}}
\end{tabproofwide}

\FormulaThmDeltaK[Ein idempotentes Element besitzt nur eine positive Potenz]{%
  \begin{aligned}[t]
    &\SemigroupAlg{A}{\star}\dsep e\in E_{A,\star}\dsep
      n\in\mathbb N\dsep n\neq0\vdash{}\\[-2pt]
    &e^n=e
  \end{aligned}
}{SemigroupIdempotentPower}{%
  \DeltaRow{Mengen}{A\dsep e\dsep n\dsep k}
  \DeltaRow{Funktionen}{\star}
}
\begin{tabproofsplitwide}
  \proofpartwideindR[Induktionsanfang]{%
    \SemigroupAlg{A}{\star}\dsep e\in E_{A,\star}\vdash e^{\Succ(0)}=e
  }{%
    \SemigroupAlg{A}{\star}\dsep e\in E_{A,\star}\vdash e^{\Succ(0)}=e
  }[SemigroupIdempotentPowerBase]
    \proofstepwidestar[1]{\SemigroupAlg{A}{\star}}{\rA}
    \proofstepwidestar[2]{e\in E_{A,\star}}{\rA}
    \proofstepwidestar[1,2]{e\in A}{%
      \rAEa{\FormulaRefAuto{SemigroupIdempotentMembership}[thm]{1,2}}}
    \proofstepwidestar[1,2]{e^1=e}{%
      \FormulaRefAuto{SemigroupPowerOne}[thm]{1,3}}
    \proofstepwidestar[]{1=\Succ(0)}{\FormulaRefAuto{PeanoOne}[def]}
    \proofstepwidestar[1,2]{e^{\Succ(0)}=e}{\rIE{5,4}}
  \closeproofpartwideindR

  \proofpartwideindR[Induktionsschritt]{%
    \begin{aligned}[t]
      &\SemigroupAlg{A}{\star}\dsep e\in E_{A,\star}\dsep
        k\in\mathbb N\dsep e^{\Succ(k)}=e\vdash{}\\[-2pt]
      &e^{\Succ(\Succ(k))}=e
    \end{aligned}
  }{%
    \begin{aligned}[t]
      &\SemigroupAlg{A}{\star}\dsep e\in E_{A,\star}\dsep
        k\in\mathbb N\dsep e^{\Succ(k)}=e\vdash{}\\[-2pt]
      &e^{\Succ(\Succ(k))}=e
    \end{aligned}
  }[SemigroupIdempotentPowerStep]
    \proofstepwidestar[1]{\SemigroupAlg{A}{\star}}{\rA}
    \proofstepwidestar[2]{e\in E_{A,\star}}{\rA}
    \proofstepwidestar[3]{k\in\mathbb N}{\rA}
    \proofstepwidestar[4]{e^{\Succ(k)}=e}{\rA}
    \proofstepwidestar[1,2]{e\in A\land e\star e=e}{%
      \FormulaRefAuto{SemigroupIdempotentMembership}[thm]{1,2}}
    \proofstepwidestar[1,2]{e\in A}{\rAEa{5}}
    \proofstepwidestar[1,2]{e\star e=e}{\rAEb{5}}
    \proofstepwidestar[3]{\Succ(k)\in\mathbb N}{%
      \FormulaRefAuto{PeanoSuccClosure}[ax]{3}}
    \proofstepwidestar[3]{\Succ(k)\neq0}{%
      \FormulaRefAuto{PeanoSuccNonzero}[ax]{3}}
    \proofstepwidestar[1,2,3]{%
      e^{\Succ(k)+1}=e^{\Succ(k)}\star e}{%
      \FormulaRefAuto{SemigroupPowerSuccessor}[thm]{1,6,8,9}}
    \proofstepwidestar[1,2,3,4]{e^{\Succ(k)}\star e=e\star e}{\rIE{4}}
    \proofstepwidestar[3]{\Succ(k)+1=\Succ(\Succ(k))}{%
      \FormulaRefAuto{PeanoAddOneSucc}[thm]{8}}
    \proofstepwidestar[1,2,3,4]{e^{\Succ(\Succ(k))}=e}{%
      \rIE{12,\FormulaRefAuto{a = b,\, b = c \vdash a = c}[thm]{%
        10,\FormulaRefAuto{a = b,\, b = c \vdash a = c}[thm]{11,7}}}}
  \closeproofpartwideindR

  \proofpartwide{Induktionsschluss}
    \proofstepwidestar[1]{\SemigroupAlg{A}{\star}}{\rA}
    \proofstepwidestar[2]{e\in E_{A,\star}}{\rA}
    \proofstepwidestar[3]{n\in\mathbb N}{\rA}
    \proofstepwidestar[4]{n\neq0}{\rA}
    \proofstepwidestar[1,2]{%
      \forall k\in\mathbb N\,\bigl(e^{\Succ(k)}=e\bigr)}{%
      \FormulaRefAuto{PeanoInductionRule}[ax]{%
        \FormulaRefAuto{SemigroupIdempotentPowerBase}[thm]{1,2},
        \FormulaRefAuto{SemigroupIdempotentPowerStep}[thm]{1,2}}}
    \proofstepwidestar[3,4]{n=\Succ(\Pred(n))}{%
      \FormulaRefAuto{PeanoPredNonzeroValue}[thm]{3,4}}
    \proofstepwidestar[3]{\Pred(n)\in\mathbb N}{%
      \FormulaRefAuto{PeanoPredNatural}[thm]{3}}
    \proofstepwidestar[1,2,3]{e^{\Succ(\Pred(n))}=e}{\rRE{7,\rUE{5}}}
    \proofstepwidestar[1,2,3,4]{e^n=e}{\rIE{6,8}}
  \closeproofpartwide
\end{tabproofsplitwide}

\section{Unterhalbgruppen}

Eine Unterhalbgruppe erbt die Assoziativit\"at von der umgebenden
Halbgruppe. Deshalb gen\"ugen Nichtleerheit und Abgeschlossenheit unter der
eingeschr\"ankten Operation.

Eine \emph{Unterhalbgruppe} ist dabei noch keine Gruppe: Gemeint ist nur
eine nichtleere, produktabgeschlossene Teilmenge. Echte Untergruppen mit
neutralem Element und Inversen werden erst nach Einführung des
Gruppenbegriffs behandelt.

Für ihren Träger verwenden wir die in der Standardliteratur übliche
Bezeichnung (U); das Prädikat schreiben wir als
\(\mathsf{UHGr}(U;A,\star)\).

\FormulaDefDeltaK[Unterhalbgruppe]{%
  \begin{aligned}[t]
    &\SemigroupAlg{A}{\star}\vdash{}\\[-2pt]
    &\mathsf{UHGr}(U;A,\star)\coloneqq
      U\subseteq A\land U\neq\varnothing\land
      \forall x\in U\,\forall y\in U\,(x\star y\in U)
  \end{aligned}
}{SemigroupSubsemigroupDef}{%
  \DeltaRow{Mengen}{A\dsep U\dsep x\dsep y}
  \DeltaRow{Funktionen}{\star}
  \DeltaRow{\textbf{Neues Symbol}}{}
  \DeltaPrem{Unterhalbgruppe}{\mathsf{UHGr}(U;A,\star)}
}

\FormulaThmDeltaK[Tr\"agermenge einer Unterhalbgruppe]{%
  \SemigroupAlg{A}{\star}\dsep\mathsf{UHGr}(U;A,\star)
  \vdash U\subseteq A%
}{SemigroupSubsemigroupSubset}{%
  \DeltaRow{Mengen}{A\dsep U}
  \DeltaRow{Funktionen}{\star}
}
\begin{tabproofwide}
  \proofstepwidestar[1]{\SemigroupAlg{A}{\star}}{\rA}
  \proofstepwidestar[2]{\mathsf{UHGr}(U;A,\star)}{\rA}
  \proofstepwidestar[1,2]{%
    U\subseteq A\land U\neq\varnothing\land
    \forall x\in U\,\forall y\in U\,(x\star y\in U)}{%
    \FormulaRefAuto{SemigroupSubsemigroupDef}[def]{1,2}}
  \proofstepwidestar[1,2]{U\subseteq A}{\rAEa{3}}
\end{tabproofwide}

\FormulaThmDeltaK[Nichtleerheit einer Unterhalbgruppe]{%
  \SemigroupAlg{A}{\star}\dsep\mathsf{UHGr}(U;A,\star)
  \vdash U\neq\varnothing%
}{SemigroupSubsemigroupNonempty}{%
  \DeltaRow{Mengen}{A\dsep U}
  \DeltaRow{Funktionen}{\star}
}
\begin{tabproofwide}
  \proofstepwidestar[1]{\SemigroupAlg{A}{\star}}{\rA}
  \proofstepwidestar[2]{\mathsf{UHGr}(U;A,\star)}{\rA}
  \proofstepwidestar[1,2]{%
    U\subseteq A\land U\neq\varnothing\land
    \forall x\in U\,\forall y\in U\,(x\star y\in U)}{%
    \FormulaRefAuto{SemigroupSubsemigroupDef}[def]{1,2}}
  \proofstepwidestar[1,2]{%
    U\neq\varnothing\land
    \forall x\in U\,\forall y\in U\,(x\star y\in U)}{\rAEb{3}}
  \proofstepwidestar[1,2]{U\neq\varnothing}{\rAEa{4}}
\end{tabproofwide}

\FormulaThmDeltaK[Produktabschluss einer Unterhalbgruppe]{%
  \begin{aligned}[t]
    &\SemigroupAlg{A}{\star}\dsep\mathsf{UHGr}(U;A,\star)\dsep
      x,y\in U\vdash{}\\[-2pt]
    &x\star y\in U
  \end{aligned}
}{SemigroupSubsemigroupClosure}{%
  \DeltaRow{Mengen}{A\dsep U\dsep x\dsep y}
  \DeltaRow{Funktionen}{\star}
}
\begin{tabproofwide}
  \proofstepwidestar[1]{\SemigroupAlg{A}{\star}}{\rA}
  \proofstepwidestar[2]{\mathsf{UHGr}(U;A,\star)}{\rA}
  \proofstepwidestar[3]{x\in U}{\rA}
  \proofstepwidestar[4]{y\in U}{\rA}
  \proofstepwidestar[1,2]{%
    U\subseteq A\land U\neq\varnothing\land
    \forall u\in U\,\forall v\in U\,(u\star v\in U)}{%
    \FormulaRefAuto{SemigroupSubsemigroupDef}[def]{1,2}}
  \proofstepwidestar[1,2]{%
    \forall u\in U\,\forall v\in U\,(u\star v\in U)}{\rAEb{\rAEb{5}}}
  \proofstepwidestar[1,2,3,4]{x\star y\in U}{\rRE{4,\rUE{\rRE{3,\rUE{6}}}}}
\end{tabproofwide}

\FormulaThmDeltaK[Eine Unterhalbgruppe ist eine Halbgruppe]{%
  \SemigroupAlg{A}{\star}\dsep\mathsf{UHGr}(U;A,\star)
  \vdash\SemigroupAlg{U}{\star}%
}{SubsemigroupFormsSemigroup}{%
  \DeltaRow{Mengen}{A\dsep U\dsep x\dsep y\dsep z}
  \DeltaRow{Funktionen}{\star}
}
\begin{tabproofsplitwide}
  \proofpartwideindR[Abgeschlossenheit]{%
    \SemigroupAlg{A}{\star}\dsep\mathsf{UHGr}(U;A,\star)\dsep
    x,y\in U\vdash x\star y\in U
  }{%
    \SemigroupAlg{A}{\star}\dsep\mathsf{UHGr}(U;A,\star)\dsep
    x,y\in U\vdash x\star y\in U
  }[SubsemigroupClosurePart]
    \proofstepwidestar[1]{\SemigroupAlg{A}{\star}}{\rA}
    \proofstepwidestar[2]{\mathsf{UHGr}(U;A,\star)}{\rA}
    \proofstepwidestar[3]{x\in U}{\rA}
    \proofstepwidestar[4]{y\in U}{\rA}
    \proofstepwidestar[1,2,3,4]{x\star y\in U}{%
      \FormulaRefAuto{SemigroupSubsemigroupClosure}[thm]{1,2,3,4}}
  \closeproofpartwideindR

  \proofpartwideindR[Vererbte Assoziativit\"at]{%
    \begin{aligned}[t]
      &\SemigroupAlg{A}{\star}\dsep\mathsf{UHGr}(U;A,\star)\dsep
        x,y,z\in U\vdash{}\\[-2pt]
      &(x\star y)\star z=x\star(y\star z)
    \end{aligned}
  }{%
    \begin{aligned}[t]
      &\SemigroupAlg{A}{\star}\dsep\mathsf{UHGr}(U;A,\star)\dsep
        x,y,z\in U\vdash{}\\[-2pt]
      &(x\star y)\star z=x\star(y\star z)
    \end{aligned}
  }[SubsemigroupAssociativityPart]
    \proofstepwidestar[1]{\SemigroupAlg{A}{\star}}{\rA}
    \proofstepwidestar[2]{\mathsf{UHGr}(U;A,\star)}{\rA}
    \proofstepwidestar[3]{x\in U}{\rA}
    \proofstepwidestar[4]{y\in U}{\rA}
    \proofstepwidestar[5]{z\in U}{\rA}
    \proofstepwidestar[1,2]{U\subseteq A}{%
      \FormulaRefAuto{SemigroupSubsemigroupSubset}[thm]{1,2}}
    \proofstepwidestar[1,2,3]{x\in A}{%
      \FormulaRefAuto{A\subseteq B,\,x\in A\vdash x\in B}[thm]{6,3}}
    \proofstepwidestar[1,2,4]{y\in A}{%
      \FormulaRefAuto{A\subseteq B,\,x\in A\vdash x\in B}[thm]{6,4}}
    \proofstepwidestar[1,2,5]{z\in A}{%
      \FormulaRefAuto{A\subseteq B,\,x\in A\vdash x\in B}[thm]{6,5}}
    \proofstepwidestar[1,2,3,4,5]{%
      (x\star y)\star z=x\star(y\star z)}{%
      \FormulaRefAuto{SemigroupAssociativityAxiom}[ax]{1,7,8,9}}
  \closeproofpartwideindR

  \proofpartwide{Schluss}
    \proofstepwidestar[1]{\SemigroupAlg{A}{\star}}{\rA}
    \proofstepwidestar[2]{\mathsf{UHGr}(U;A,\star)}{\rA}
    \proofstepwidestar[1,2]{\SemigroupAlg{U}{\star}}{%
      \FormulaRefAuto{SemigroupDef}[def]{%
        \FormulaRefAuto{SubsemigroupClosurePart}[thm]{1,2},
        \FormulaRefAuto{SubsemigroupAssociativityPart}[thm]{1,2}}}
  \closeproofpartwide
\end{tabproofsplitwide}

\FormulaThmDeltaK[Der Tr\"ager ist eine Unterhalbgruppe]{%
  \SemigroupAlg{A}{\star}\dsep A\neq\varnothing
  \vdash\mathsf{UHGr}(A;A,\star)%
}{SemigroupCarrierIsSubsemigroup}{%
  \DeltaRow{Mengen}{A\dsep x\dsep y}
  \DeltaRow{Funktionen}{\star}
}
\begin{tabproofwide}
  \proofstepwidestar[1]{\SemigroupAlg{A}{\star}}{\rA}
  \proofstepwidestar[2]{A\neq\varnothing}{\rA}
  \proofstepwidestar[]{A\subseteq A}{%
    \FormulaRefAuto{A\subseteq A}[thm]}
  \proofstepwidestar[4]{x\in A}{\rA}
  \proofstepwidestar[5]{y\in A}{\rA}
  \proofstepwidestar[1,4,5]{x\star y\in A}{%
    \FormulaRefAuto{SemigroupClosureAxiom}[ax]{1,4,5}}
  \proofstepwidestar[1]{%
    \forall x\in A\,\forall y\in A\,(x\star y\in A)}{%
    \rUI{\rRI{4,\rUI{\rRI{5,6}}}}}
  \proofstepwidestar[1,2]{%
    A\subseteq A\land A\neq\varnothing\land
    \forall x\in A\,\forall y\in A\,(x\star y\in A)}{%
    \rAI{3,\rAI{2,7}}}
  \proofstepwidestar[1,2]{\mathsf{UHGr}(A;A,\star)}{%
    \FormulaRefAuto{SemigroupSubsemigroupDef}[def]{1,8}}
\end{tabproofwide}

\subsection{Erzeugte Unterhalbgruppen}

\FormulaDefDeltaK[Von einer nichtleeren Menge erzeugte Unterhalbgruppe]{%
  \begin{aligned}[t]
    &\SemigroupAlg{A}{\star}\dsep X\subseteq A\dsep
      X\neq\varnothing\vdash{}\\[-2pt]
    &\langle X\rangle_{A,\star}\coloneqq
      \bigcap_{\mathsf{UHGr}(U;A,\star)\land X\subseteq U}U
  \end{aligned}
}{SemigroupGeneratedSubsemigroupDef}{%
  \DeltaRow{Mengen}{A\dsep U\dsep X}
  \DeltaRow{Funktionen}{\star}
  \DeltaRow{\textbf{Neues Symbol}}{}
  \DeltaPrem{Erzeugte Unterhalbgruppe}{\langle X\rangle_{A,\star}}
}

\FormulaThmDeltaK[Der Tr\"ager enth\"alt die Erzeuger]{%
  \begin{aligned}[t]
    &\SemigroupAlg{A}{\star}\dsep X\subseteq A\dsep
      X\neq\varnothing\vdash{}\\[-2pt]
    &X\subseteq\langle X\rangle_{A,\star}
  \end{aligned}
}{GeneratedSubsemigroupContainsGenerators}{%
  \DeltaRow{Mengen}{A\dsep U\dsep X\dsep x\dsep u}
  \DeltaRow{Funktionen}{\star}
}
\begin{tabproofwide}
  \proofstepwidestar[1]{\SemigroupAlg{A}{\star}}{\rA}
  \proofstepwidestar[2]{X\subseteq A}{\rA}
  \proofstepwidestar[3]{X\neq\varnothing}{\rA}
  \proofstepwidestar[3]{\exists u\,(u\in X)}{%
    \FormulaRefAuto{A\neq\varnothing\eqvdash\exists x\,(x\in A)}[thm]{3}}
  \proofstepwidestar[5]{u\in X}{\rA}
  \proofstepwidestar[2,5]{u\in A}{%
    \FormulaRefAuto{A\subseteq B,\,x\in A\vdash x\in B}[thm]{2,5}}
  \proofstepwidestar[2,5]{A\neq\varnothing}{%
    \FormulaRefAuto{a\in A\vdash A\neq\varnothing}[thm]{6}}
  \proofstepwidestar[2,3]{A\neq\varnothing}{\rEE{4,5,7}}
  \proofstepwidestar[1,2,3]{\mathsf{UHGr}(A;A,\star)}{%
    \FormulaRefAuto{SemigroupCarrierIsSubsemigroup}[thm]{1,8}}
  \proofstepwidestar[1,2,3]{%
    \mathsf{UHGr}(A;A,\star)\land X\subseteq A}{\rAI{9,2}}
  \proofstepwidestar[11]{x\in X}{\rA}
  \proofstepwidestar[12]{%
    \mathsf{UHGr}(U;A,\star)\land X\subseteq U}{\rA}
  \proofstepwidestar[12]{X\subseteq U}{\rAEb{12}}
  \proofstepwidestar[11,12]{x\in U}{%
    \FormulaRefAuto{A\subseteq B,\,x\in A\vdash x\in B}[thm]{13,11}}
  \proofstepwidestar[11]{%
    \forall U\,\Bigl(
      \mathsf{UHGr}(U;A,\star)\land X\subseteq U
      \rightarrow x\in U\Bigr)}{\rUI{\rRI{12,14}}}
  \proofstepwidestar[1,2,3,11]{%
    x\in\bigcap_{\mathsf{UHGr}(U;A,\star)\land X\subseteq U}U}{%
    \FormulaRefAuto{P(A)\vdash x\in\bigcap_{P(B)}B
      \leftrightarrow\forall C\,(P(C)\rightarrow x\in C)}[thm]{10,15}}
  \proofstepwidestar[1,2,3]{%
    \langle X\rangle_{A,\star}=
    \bigcap_{\mathsf{UHGr}(U;A,\star)\land X\subseteq U}U}{%
    \FormulaRefAuto{SemigroupGeneratedSubsemigroupDef}[def]{1,2,3}}
  \proofstepwidestar[1,2,3,11]{x\in\langle X\rangle_{A,\star}}{%
    \rIE{\FormulaRefAuto{a=b\vdash b=a}[thm]{17},16}}
  \proofstepwidestar[1,2,3]{X\subseteq\langle X\rangle_{A,\star}}{%
    \FormulaRefAuto{A\subseteq B\coloneq
      \forall x\,(x\in A\rightarrow x\in B)}[def]{\rUI{\rRI{11,18}}}}
\end{tabproofwide}

\FormulaThmDeltaK[Minimalit\"at der erzeugten Unterhalbgruppe]{%
  \begin{aligned}[t]
    &\SemigroupAlg{A}{\star}\dsep X\subseteq A\dsep X\neq\varnothing\dsep
      \mathsf{UHGr}(U;A,\star)\dsep X\subseteq U\vdash{}\\[-2pt]
    &\langle X\rangle_{A,\star}\subseteq U
  \end{aligned}
}{GeneratedSubsemigroupMinimal}{%
  \DeltaRow{Mengen}{A\dsep U\dsep V\dsep X}
  \DeltaRow{Funktionen}{\star}
}
\begin{tabproofwide}
  \proofstepwidestar[1]{\SemigroupAlg{A}{\star}}{\rA}
  \proofstepwidestar[2]{X\subseteq A}{\rA}
  \proofstepwidestar[3]{X\neq\varnothing}{\rA}
  \proofstepwidestar[4]{\mathsf{UHGr}(U;A,\star)}{\rA}
  \proofstepwidestar[5]{X\subseteq U}{\rA}
  \proofstepwidestar[4,5]{%
    \mathsf{UHGr}(U;A,\star)\land X\subseteq U}{\rAI{4,5}}
  \proofstepwidestar[4,5]{%
    \bigcap_{\mathsf{UHGr}(V;A,\star)\land X\subseteq V}V
    \subseteq U}{%
    \FormulaRefAuto{P(C)\vdash\bigcap_{P(A)}A\subseteq C}[thm]{6}}
  \proofstepwidestar[1,2,3]{%
    \langle X\rangle_{A,\star}=
    \bigcap_{\mathsf{UHGr}(V;A,\star)\land X\subseteq V}V}{%
    \FormulaRefAuto{SemigroupGeneratedSubsemigroupDef}[def]{1,2,3}}
  \proofstepwidestar[1,2,3,4,5]{\langle X\rangle_{A,\star}\subseteq U}{%
    \rIE{8,7}}
\end{tabproofwide}

\FormulaThmDeltaK[Produktabschluss der erzeugten Unterhalbgruppe]{%
  \begin{aligned}[t]
    &\SemigroupAlg{A}{\star}\dsep X\subseteq A\dsep X\neq\varnothing\dsep
      x,y\in\langle X\rangle_{A,\star}\vdash{}\\[-2pt]
    &x\star y\in\langle X\rangle_{A,\star}
  \end{aligned}
}{GeneratedSubsemigroupClosure}{%
  \DeltaRow{Mengen}{A\dsep U\dsep X\dsep x\dsep y\dsep u}
  \DeltaRow{Funktionen}{\star}
}
\begin{tabproofwide}
  \proofstepwidestar[1]{\SemigroupAlg{A}{\star}}{\rA}
  \proofstepwidestar[2]{X\subseteq A}{\rA}
  \proofstepwidestar[3]{X\neq\varnothing}{\rA}
  \proofstepwidestar[4]{x\in\langle X\rangle_{A,\star}}{\rA}
  \proofstepwidestar[5]{y\in\langle X\rangle_{A,\star}}{\rA}
  \proofstepwidestar[3]{\exists u\,(u\in X)}{%
    \FormulaRefAuto{A\neq\varnothing\eqvdash\exists u\,(u\in A)}[thm]{3}}
  \proofstepwidestar[7]{u\in X}{\rA}
  \proofstepwidestar[2,7]{u\in A}{%
    \FormulaRefAuto{A\subseteq B,\,x\in A\vdash x\in B}[thm]{2,7}}
  \proofstepwidestar[2,7]{A\neq\varnothing}{%
    \FormulaRefAuto{a\in A\vdash A\neq\varnothing}[thm]{8}}
  \proofstepwidestar[2,3]{A\neq\varnothing}{\rEE{6,7,9}}
  \proofstepwidestar[1,2,3]{\mathsf{UHGr}(A;A,\star)}{%
    \FormulaRefAuto{SemigroupCarrierIsSubsemigroup}[thm]{1,10}}
  \proofstepwidestar[1,2,3]{%
    \mathsf{UHGr}(A;A,\star)\land X\subseteq A}{\rAI{11,2}}
  \proofstepwidestar[1,2,3]{%
    \langle X\rangle_{A,\star}=
    \bigcap_{\mathsf{UHGr}(U;A,\star)\land X\subseteq U}U}{%
    \FormulaRefAuto{SemigroupGeneratedSubsemigroupDef}[def]{1,2,3}}
  \proofstepwidestar[1,2,3,4]{%
    x\in\bigcap_{\mathsf{UHGr}(U;A,\star)\land X\subseteq U}U}{\rIE{13,4}}
  \proofstepwidestar[1,2,3,5]{%
    y\in\bigcap_{\mathsf{UHGr}(U;A,\star)\land X\subseteq U}U}{\rIE{13,5}}
  \proofstepwidestar[1,2,3,4]{%
    \forall U\,\Bigl(
      \mathsf{UHGr}(U;A,\star)\land X\subseteq U
      \rightarrow x\in U\Bigr)}{%
    \FormulaRefAuto{P(A)\vdash x\in\bigcap_{P(B)}B
      \leftrightarrow\forall C\,(P(C)\rightarrow x\in C)}[thm]{12,14}}
  \proofstepwidestar[1,2,3,5]{%
    \forall U\,\Bigl(
      \mathsf{UHGr}(U;A,\star)\land X\subseteq U
      \rightarrow y\in U\Bigr)}{%
    \FormulaRefAuto{P(A)\vdash x\in\bigcap_{P(B)}B
      \leftrightarrow\forall C\,(P(C)\rightarrow x\in C)}[thm]{12,15}}
  \proofstepwidestar[18]{%
    \mathsf{UHGr}(U;A,\star)\land X\subseteq U}{\rA}
  \proofstepwidestar[18]{\mathsf{UHGr}(U;A,\star)}{\rAEa{18}}
  \proofstepwidestar[1,2,3,4,18]{x\in U}{\rRE{18,\rUE{16}}}
  \proofstepwidestar[1,2,3,5,18]{y\in U}{\rRE{18,\rUE{17}}}
  \proofstepwidestar[1,2,3,4,5,18]{x\star y\in U}{%
    \FormulaRefAuto{SemigroupSubsemigroupClosure}[thm]{1,19,20,21}}
  \proofstepwidestar[1,2,3,4,5]{%
    \forall U\,\Bigl(
      \mathsf{UHGr}(U;A,\star)\land X\subseteq U
      \rightarrow x\star y\in U\Bigr)}{\rUI{\rRI{18,22}}}
  \proofstepwidestar[1,2,3,4,5]{%
    x\star y\in
    \bigcap_{\mathsf{UHGr}(U;A,\star)\land X\subseteq U}U}{%
    \FormulaRefAuto{P(A)\vdash x\in\bigcap_{P(B)}B
      \leftrightarrow\forall C\,(P(C)\rightarrow x\in C)}[thm]{12,23}}
  \proofstepwidestar[1,2,3,4,5]{%
    x\star y\in\langle X\rangle_{A,\star}}{%
    \rIE{\FormulaRefAuto{a=b\vdash b=a}[thm]{13},24}}
\end{tabproofwide}

\FormulaThmDeltaK[Das Erzeugnis ist eine Unterhalbgruppe]{%
  \begin{aligned}[t]
    &\SemigroupAlg{A}{\star}\dsep X\subseteq A\dsep
      X\neq\varnothing\vdash{}\\[-2pt]
    &\mathsf{UHGr}(\langle X\rangle_{A,\star};A,\star)
  \end{aligned}
}{GeneratedSubsemigroupIsSubsemigroup}{%
  \DeltaRow{Mengen}{A\dsep X\dsep x\dsep y\dsep u}
  \DeltaRow{Funktionen}{\star}
}
\begin{tabproofsplitwide}
  \proofpartwideindR[Träger und Nichtleerheit des Erzeugnisses]{%
    \begin{aligned}[t]
      &\SemigroupAlg{A}{\star}\dsep X\subseteq A\dsep
        X\neq\varnothing\vdash{}\\[-2pt]
      &\langle X\rangle_{A,\star}\subseteq A\land
        \langle X\rangle_{A,\star}\neq\varnothing
    \end{aligned}
  }{%
    \begin{aligned}[t]
      &\SemigroupAlg{A}{\star}\dsep X\subseteq A\dsep
        X\neq\varnothing\vdash{}\\[-2pt]
      &\langle X\rangle_{A,\star}\subseteq A\land
        \langle X\rangle_{A,\star}\neq\varnothing
    \end{aligned}%
  }[GeneratedSubsemigroupCarrierAndNonempty]
    \proofstepwidestar[1]{\SemigroupAlg{A}{\star}}{\rA}
    \proofstepwidestar[2]{X\subseteq A}{\rA}
    \proofstepwidestar[3]{X\neq\varnothing}{\rA}
    \proofstepwidestar[3]{\exists u\,(u\in X)}{%
      \FormulaRefAuto{A\neq\varnothing\eqvdash
        \exists x\,(x\in A)}[thm]{3}}
    \proofstepwidestar[5]{u\in X}{\rA}
    \proofstepwidestar[2,5]{u\in A}{%
      \FormulaRefAuto{A\subseteq B,\,x\in A\vdash x\in B}[thm]{2,5}}
    \proofstepwidestar[2,5]{A\neq\varnothing}{%
      \FormulaRefAuto{a\in A\vdash A\neq\varnothing}[thm]{6}}
    \proofstepwidestar[2,3]{A\neq\varnothing}{\rEE{4,5,7}}
    \proofstepwidestar[1,2,3]{\mathsf{UHGr}(A;A,\star)}{%
      \FormulaRefAuto{SemigroupCarrierIsSubsemigroup}[thm]{1,8}}
    \proofstepwidestarbreak[1,2,3]{%
      \langle X\rangle_{A,\star}\subseteq A}{%
      \FormulaRefAuto{GeneratedSubsemigroupMinimal}[thm]{1,2,3,9,2}}
    \proofstepwidestar[1,2,3]{%
      X\subseteq\langle X\rangle_{A,\star}}{%
      \FormulaRefAuto{GeneratedSubsemigroupContainsGenerators}[thm]{1,2,3}}
    \proofstepwidestarbreak[2,3]{%
      \exists u\,(u\in\langle X\rangle_{A,\star})}{%
      \FormulaRefAuto{A\subseteq B\dsep A\neq\varnothing
        \vdash\exists a\in B\,a\in A}[thm]{11,3}}
    \proofstepwidestar[1,2,3]{%
      \langle X\rangle_{A,\star}\neq\varnothing}{%
      \FormulaRefAuto{A\neq\varnothing\eqvdash
        \exists x\,(x\in A)}[thm]{12}}
    \proofstepwidestarbreak[1,2,3]{%
      \langle X\rangle_{A,\star}\subseteq A\land
      \langle X\rangle_{A,\star}\neq\varnothing}{\rAI{10,13}}
  \closeproofpartwideindR

  \proofpartwideindR[Produktabschluss des Erzeugnisses]{%
    \begin{aligned}[t]
      &\SemigroupAlg{A}{\star}\dsep X\subseteq A\dsep
        X\neq\varnothing\vdash{}\\[-2pt]
      &\forall x\in\langle X\rangle_{A,\star}\,
        \forall y\in\langle X\rangle_{A,\star}\,
        (x\star y\in\langle X\rangle_{A,\star})
    \end{aligned}
  }{%
    \begin{aligned}[t]
      &\SemigroupAlg{A}{\star}\dsep X\subseteq A\dsep
        X\neq\varnothing\vdash{}\\[-2pt]
      &\forall x\in\langle X\rangle_{A,\star}\,
        \forall y\in\langle X\rangle_{A,\star}\,
        (x\star y\in\langle X\rangle_{A,\star})
    \end{aligned}%
  }[GeneratedSubsemigroupProductClosurePart]
    \proofstepwidestar[1]{\SemigroupAlg{A}{\star}}{\rA}
    \proofstepwidestar[2]{X\subseteq A}{\rA}
    \proofstepwidestar[3]{X\neq\varnothing}{\rA}
    \proofstepwidestar[4]{x\in\langle X\rangle_{A,\star}}{\rA}
    \proofstepwidestar[5]{y\in\langle X\rangle_{A,\star}}{\rA}
    \proofstepwidestarbreak[1,2,3,4,5]{%
      x\star y\in\langle X\rangle_{A,\star}}{%
      \FormulaRefAuto{GeneratedSubsemigroupClosure}[thm]{1,2,3,4,5}}
    \proofstepwidestarbreak[1,2,3]{%
      \forall x\in\langle X\rangle_{A,\star}\,
      \forall y\in\langle X\rangle_{A,\star}\,
      (x\star y\in\langle X\rangle_{A,\star})}{%
      \rUI{\rRI{4,\rUI{\rRI{5,6}}}}}
  \closeproofpartwideindR

  \proofpartwideindR[Schluss]{%
    \begin{aligned}[t]
      &\SemigroupAlg{A}{\star}\dsep X\subseteq A\dsep
        X\neq\varnothing\vdash{}\\[-2pt]
      &\mathsf{UHGr}(\langle X\rangle_{A,\star};A,\star)
    \end{aligned}
  }{%
    \begin{aligned}[t]
      &\SemigroupAlg{A}{\star}\dsep X\subseteq A\dsep
        X\neq\varnothing\vdash{}\\[-2pt]
      &\mathsf{UHGr}(\langle X\rangle_{A,\star};A,\star)
    \end{aligned}%
  }[GeneratedSubsemigroupConclusion]
    \proofstepwidestar[1]{\SemigroupAlg{A}{\star}}{\rA}
    \proofstepwidestar[2]{X\subseteq A}{\rA}
    \proofstepwidestar[3]{X\neq\varnothing}{\rA}
    \proofstepwidestarbreak[1,2,3]{%
      \begin{aligned}[t]
        &\langle X\rangle_{A,\star}\subseteq A\land{}\\[-2pt]
        &\langle X\rangle_{A,\star}\neq\varnothing
      \end{aligned}}{%
      \FormulaRefAuto{GeneratedSubsemigroupCarrierAndNonempty}[thm]{1,2,3}}
    \proofstepwidestar[1,2,3]{%
      \langle X\rangle_{A,\star}\subseteq A}{\rAEa{4}}
    \proofstepwidestar[1,2,3]{%
      \langle X\rangle_{A,\star}\neq\varnothing}{\rAEb{4}}
    \proofstepwidestarbreak[1,2,3]{%
      \begin{aligned}[t]
        &\forall x\in\langle X\rangle_{A,\star}\,\\[-2pt]
        &\quad\forall y\in\langle X\rangle_{A,\star}\,
          (x\star y\in\langle X\rangle_{A,\star})
      \end{aligned}}{%
      \FormulaRefAuto{GeneratedSubsemigroupProductClosurePart}[thm]{1,2,3}}
    \proofstepwidestarbreak[1,2,3]{%
      \begin{aligned}[t]
        &\langle X\rangle_{A,\star}\subseteq A\land
          \langle X\rangle_{A,\star}\neq\varnothing\land{}\\[-2pt]
        &\forall x\in\langle X\rangle_{A,\star}\,\\[-2pt]
        &\quad\forall y\in\langle X\rangle_{A,\star}\,
          (x\star y\in\langle X\rangle_{A,\star})
      \end{aligned}}{%
      \rAI{5,\rAI{6,7}}}
    \proofstepwidestarbreak[1,2,3]{%
      \mathsf{UHGr}(\langle X\rangle_{A,\star};A,\star)}{%
      \FormulaRefAuto{SemigroupSubsemigroupDef}[def]{1,8}}
  \closeproofpartwideindR
\end{tabproofsplitwide}

\subsection{Monogene Unterhalbgruppen}

\FormulaDefDeltaK[Von einem Element erzeugte Unterhalbgruppe]{%
  \SemigroupAlg{A}{\star}\dsep a\in A\vdash
  \langle a\rangle_{A,\star}\coloneqq
  \langle\{a\}\rangle_{A,\star}%
}{SemigroupMonogenicSubsemigroupDef}{%
  \DeltaRow{Mengen}{A\dsep a}
  \DeltaRow{Funktionen}{\star}
  \DeltaRow{\textbf{Neues Symbol}}{}
  \DeltaPrem{Monogene Unterhalbgruppe}{\langle a\rangle_{A,\star}}
}

\FormulaThmDeltaK[Das Monogen ist eine Unterhalbgruppe]{%
  \SemigroupAlg{A}{\star}\dsep a\in A\vdash
  \mathsf{UHGr}(\langle a\rangle_{A,\star};A,\star)%
}{MonogenicSubsemigroupIsSubsemigroup}{%
  \DeltaRow{Mengen}{A\dsep a}
  \DeltaRow{Funktionen}{\star}
}
\begin{tabproofwide}
  \proofstepwidestar[1]{\SemigroupAlg{A}{\star}}{\rA}
  \proofstepwidestar[2]{a\in A}{\rA}
  \proofstepwidestar[2]{\{a\}\subseteq A}{%
    \FormulaRefAuto{a\in A\vdash\{a\}\subseteq A}[thm]{2}}
  \proofstepwidestar[]{a\in\{a\}}{\FormulaRefAuto{a\in\{a\}}[thm]}
  \proofstepwidestar[]{\{a\}\neq\varnothing}{%
    \FormulaRefAuto{a\in A\vdash A\neq\varnothing}[thm]{4}}
  \proofstepwidestar[1,2]{%
    \mathsf{UHGr}(\langle\{a\}\rangle_{A,\star};A,\star)}{%
    \FormulaRefAuto{GeneratedSubsemigroupIsSubsemigroup}[thm]{1,3,5}}
  \proofstepwidestar[1,2]{%
    \langle a\rangle_{A,\star}=\langle\{a\}\rangle_{A,\star}}{%
    \FormulaRefAuto{SemigroupMonogenicSubsemigroupDef}[def]{1,2}}
  \proofstepwidestar[1,2]{%
    \mathsf{UHGr}(\langle a\rangle_{A,\star};A,\star)}{\rIE{7,6}}
\end{tabproofwide}

\FormulaDefDeltaK[Menge der positiven Potenzen]{%
  \SemigroupAlg{A}{\star}\dsep a\in A\vdash
  P_{A,\star}(a)\coloneqq
  \bigl\{x\in A\bigm|
    \exists n\in\mathbb N\,(n\neq0\land x=a^n)\bigr\}%
}{SemigroupPositivePowerSetDef}{%
  \DeltaRow{Mengen}{A\dsep a\dsep x\dsep n}
  \DeltaRow{Funktionen}{\star}
  \DeltaRow{\textbf{Neues Symbol}}{}
  \DeltaPrem{Positive Potenzmenge}{P_{A,\star}(a)}
}

\FormulaThmDeltaK[Elementkriterium der positiven Potenzmenge]{%
  \begin{aligned}[t]
    &\SemigroupAlg{A}{\star}\dsep a\in A\vdash{}\\[-2pt]
    &x\in P_{A,\star}(a)\leftrightarrow
      \Bigl(x\in A\land
        \exists n\in\mathbb N\,(n\neq0\land x=a^n)\Bigr)
  \end{aligned}
}{SemigroupPositivePowerSetMembership}{%
  \DeltaRow{Mengen}{A\dsep a\dsep x\dsep n}
  \DeltaRow{Funktionen}{\star}
}
\begin{tabproofwide}
  \proofstepwidestar[1]{\SemigroupAlg{A}{\star}}{\rA}
  \proofstepwidestar[2]{a\in A}{\rA}
  \proofstepwidestar[1,2]{%
    P_{A,\star}(a)=\{u\in A\mid
      \exists n\in\mathbb N\,(n\neq0\land u=a^n)\}}{%
    \FormulaRefAuto{SemigroupPositivePowerSetDef}[def]{1,2}}
  \proofstepwidestar[1,2]{%
    x\in P_{A,\star}(a)\leftrightarrow
    \Bigl(x\in A\land
      \exists n\in\mathbb N\,(n\neq0\land x=a^n)\Bigr)}{%
    \rIE{3,\FormulaRefAuto{x\in\{u\in A\mid P(u)\}\eqvdash
      x\in A\land P(x)}[thm]}}
\end{tabproofwide}

\FormulaThmDeltaK[Die positiven Potenzen bilden eine Unterhalbgruppe]{%
  \SemigroupAlg{A}{\star}\dsep a\in A\vdash
  \mathsf{UHGr}(P_{A,\star}(a);A,\star)%
}{SemigroupPositivePowersFormSubsemigroup}{%
  \DeltaRow{Mengen}{A\dsep a\dsep x\dsep y\dsep m\dsep n}
  \DeltaRow{Funktionen}{\star}
}
\begin{tabproofsplitwide}
  \proofpartwideindR[Teilmenge und Nichtleerheit]{%
    \begin{aligned}[t]
      &\SemigroupAlg{A}{\star}\dsep a\in A\vdash{}\\[-2pt]
      &P_{A,\star}(a)\subseteq A\land
        P_{A,\star}(a)\neq\varnothing
    \end{aligned}
  }{%
    \begin{aligned}[t]
      &\SemigroupAlg{A}{\star}\dsep a\in A\vdash{}\\[-2pt]
      &P_{A,\star}(a)\subseteq A\land
        P_{A,\star}(a)\neq\varnothing
    \end{aligned}%
  }[SemigroupPositivePowerSetCarrierAndNonempty]
    \proofstepwidestar[1]{\SemigroupAlg{A}{\star}}{\rA}
    \proofstepwidestar[2]{a\in A}{\rA}
    \proofstepwidestar[]{P_{A,\star}(a)\subseteq A}{%
      \FormulaRefAuto{\{x\in A\mid P(x)\}\subseteq A}[thm]}
    \proofstepwidestar[]{1\in\mathbb N}{%
      \FormulaRefAuto{PeanoOneInN}[thm]}
    \proofstepwidestar[]{1\neq0}{%
      \rIE{\FormulaRefAuto{PeanoOne}[def],
        \FormulaRefAuto{PeanoSuccNonzero}[ax]{%
          \FormulaRefAuto{PeanoZero}[ax]}}}
    \proofstepwidestar[1,2]{a^1=a}{%
      \FormulaRefAuto{SemigroupPowerOne}[thm]{1,2}}
    \proofstepwidestarbreak[1,2]{%
      a\in A\land
      \exists n\in\mathbb N\,(n\neq0\land a=a^n)}{%
      \rAI{2,\rEI{4,
        \rAI{5,\FormulaRefAuto{a=b\vdash b=a}[thm]{6}}}}}
    \proofstepwidestar[1,2]{a\in P_{A,\star}(a)}{%
      \FormulaRefAuto{SemigroupPositivePowerSetMembership}[thm]{1,2,7}}
    \proofstepwidestar[1,2]{P_{A,\star}(a)\neq\varnothing}{%
      \FormulaRefAuto{a\in A\vdash A\neq\varnothing}[thm]{8}}
    \proofstepwidestarbreak[1,2]{%
      P_{A,\star}(a)\subseteq A\land
      P_{A,\star}(a)\neq\varnothing}{\rAI{3,9}}
  \closeproofpartwideindR

  \proofpartwideindR[Produktabschluss der positiven Potenzen]{%
    \begin{aligned}[t]
      &\SemigroupAlg{A}{\star}\dsep a\in A\vdash{}\\[-2pt]
      &\forall x\in P_{A,\star}(a)\,
        \forall y\in P_{A,\star}(a)\,
        (x\star y\in P_{A,\star}(a))
    \end{aligned}
  }{%
    \begin{aligned}[t]
      &\SemigroupAlg{A}{\star}\dsep a\in A\vdash{}\\[-2pt]
      &\forall x\in P_{A,\star}(a)\,
        \forall y\in P_{A,\star}(a)\,
        (x\star y\in P_{A,\star}(a))
    \end{aligned}%
  }[SemigroupPositivePowerSetProductClosure]
    \proofstepwidestar[1]{\SemigroupAlg{A}{\star}}{\rA}
    \proofstepwidestar[2]{a\in A}{\rA}
    \proofstepwidestar[3]{x\in P_{A,\star}(a)}{\rA}
    \proofstepwidestar[4]{y\in P_{A,\star}(a)}{\rA}
    \proofstepwidestarbreak[1,2,3]{%
      \exists m\in\mathbb N\,(m\neq0\land x=a^m)}{%
      \rAEb{\FormulaRefAuto{SemigroupPositivePowerSetMembership}[thm]{%
        1,2,3}}}
    \proofstepwidestarbreak[1,2,4]{%
      \exists n\in\mathbb N\,(n\neq0\land y=a^n)}{%
      \rAEb{\FormulaRefAuto{SemigroupPositivePowerSetMembership}[thm]{%
        1,2,4}}}
    \proofstepwidestar[7]{%
      m\in\mathbb N\land(m\neq0\land x=a^m)}{\rA}
    \proofstepwidestar[8]{%
      n\in\mathbb N\land(n\neq0\land y=a^n)}{\rA}
    \proofstepwidestar[7]{m\in\mathbb N}{\rAEa{7}}
    \proofstepwidestar[7]{m\neq0}{\rAEa{\rAEb{7}}}
    \proofstepwidestar[7]{x=a^m}{\rAEb{\rAEb{7}}}
    \proofstepwidestar[8]{n\in\mathbb N}{\rAEa{8}}
    \proofstepwidestar[8]{n\neq0}{\rAEa{\rAEb{8}}}
    \proofstepwidestar[8]{y=a^n}{\rAEb{\rAEb{8}}}
    \proofstepwidestar[7,8]{m+n\in\mathbb N}{%
      \FormulaRefAuto{PeanoAddClosure}[thm]{9,12}}
    \proofstepwidestar[7]{\Pred(m)\in\mathbb N}{%
      \FormulaRefAuto{PeanoPredNatural}[thm]{9}}
    \proofstepwidestar[7]{m=\Succ(\Pred(m))}{%
      \FormulaRefAuto{PeanoPredNonzeroValue}[thm]{9,10}}
    \proofstepwidestarbreak[7,8]{%
      \Succ(\Pred(m))+n=\Succ(\Pred(m)+n)}{%
      \FormulaRefAuto{PeanoAddSuccLeft}[thm]{16,12}}
    \proofstepwidestar[7,8]{\Pred(m)+n\in\mathbb N}{%
      \FormulaRefAuto{PeanoAddClosure}[thm]{16,12}}
    \proofstepwidestar[7,8]{\Succ(\Pred(m)+n)\neq0}{%
      \FormulaRefAuto{PeanoSuccNonzero}[ax]{19}}
    \proofstepwidestar[7,8]{%
      m+n=\Succ(\Pred(m)+n)}{\rIE{17,18}}
    \proofstepwidestar[7,8]{m+n\neq0}{\rIE{21,20}}
    \proofstepwidestarbreak[1,2,7,8]{%
      a^{m+n}=a^m\star a^n}{%
      \FormulaRefAuto{SemigroupPowerAdditionLaw}[thm]{%
        1,2,9,12,10,13}}
    \proofstepwidestar[1,2,7,8]{%
      x\star y=a^{m+n}}{\rIE{11,14,23}}
    \proofstepwidestar[1,2,7,8]{x\star y\in A}{%
      \FormulaRefAuto{SemigroupPowerClosure}[thm]{1,2,15,22}}
    \proofstepwidestarbreak[1,2,7,8]{%
      x\star y\in A\land
      \exists k\in\mathbb N\,(k\neq0\land x\star y=a^k)}{%
      \rAI{25,\rEI{15,\rAI{22,24}}}}
    \proofstepwidestarbreak[1,2,7,8]{%
      x\star y\in P_{A,\star}(a)}{%
      \FormulaRefAuto{SemigroupPositivePowerSetMembership}[thm]{1,2,26}}
    \proofstepwidestar[1,2,3,4]{%
      x\star y\in P_{A,\star}(a)}{%
      \rEE{5,7,\rEE{6,8,27}}}
    \proofstepwidestarbreak[1,2]{%
      \forall x\in P_{A,\star}(a)\,
      \forall y\in P_{A,\star}(a)\,
      (x\star y\in P_{A,\star}(a))}{%
      \rUI{\rRI{3,\rUI{\rRI{4,28}}}}}
  \closeproofpartwideindR

  \proofpartwideindR[Schluss]{%
    \begin{aligned}[t]
      &\SemigroupAlg{A}{\star}\dsep a\in A\vdash{}\\[-2pt]
      &\mathsf{UHGr}(P_{A,\star}(a);A,\star)
    \end{aligned}
  }{%
    \begin{aligned}[t]
      &\SemigroupAlg{A}{\star}\dsep a\in A\vdash{}\\[-2pt]
      &\mathsf{UHGr}(P_{A,\star}(a);A,\star)
    \end{aligned}%
  }[SemigroupPositivePowerSetSubsemigroupConclusion]
    \proofstepwidestar[1]{\SemigroupAlg{A}{\star}}{\rA}
    \proofstepwidestar[2]{a\in A}{\rA}
    \proofstepwidestarbreak[1,2]{%
      \begin{aligned}[t]
        &P_{A,\star}(a)\subseteq A\land{}\\[-2pt]
        &P_{A,\star}(a)\neq\varnothing
      \end{aligned}}{%
      \FormulaRefAuto{SemigroupPositivePowerSetCarrierAndNonempty}[thm]{1,2}}
    \proofstepwidestar[1,2]{P_{A,\star}(a)\subseteq A}{\rAEa{3}}
    \proofstepwidestar[1,2]{P_{A,\star}(a)\neq\varnothing}{\rAEb{3}}
    \proofstepwidestarbreak[1,2]{%
      \begin{aligned}[t]
        &\forall x\in P_{A,\star}(a)\,\\[-2pt]
        &\quad\forall y\in P_{A,\star}(a)\,
          (x\star y\in P_{A,\star}(a))
      \end{aligned}}{%
      \FormulaRefAuto{SemigroupPositivePowerSetProductClosure}[thm]{1,2}}
    \proofstepwidestarbreak[1,2]{%
      \begin{aligned}[t]
        &P_{A,\star}(a)\subseteq A\land
          P_{A,\star}(a)\neq\varnothing\land{}\\[-2pt]
        &\forall x\in P_{A,\star}(a)\,\\[-2pt]
        &\quad\forall y\in P_{A,\star}(a)\,
          (x\star y\in P_{A,\star}(a))
      \end{aligned}}{%
      \rAI{4,\rAI{5,6}}}
    \proofstepwidestarbreak[1,2]{%
      \mathsf{UHGr}(P_{A,\star}(a);A,\star)}{%
      \FormulaRefAuto{SemigroupSubsemigroupDef}[def]{1,7}}
  \closeproofpartwideindR
\end{tabproofsplitwide}

\FormulaThmDeltaK[Beschreibung der monogenen Unterhalbgruppe durch Potenzen]{%
  \SemigroupAlg{A}{\star}\dsep a\in A\vdash
  \langle a\rangle_{A,\star}=P_{A,\star}(a)%
}{SemigroupMonogenicPowerDescription}{%
  \DeltaRow{Mengen}{A\dsep a\dsep x\dsep n\dsep k}
  \DeltaRow{Funktionen}{\star}
}
\begin{tabproofsplitwide}
  \proofpartwideindR[Positive Potenzen liegen im Monogen]{%
    \begin{aligned}[t]
      &\SemigroupAlg{A}{\star}\dsep a\in A\dsep k\in\mathbb N\vdash{}\\[-2pt]
      &a^{\Succ(k)}\in\langle a\rangle_{A,\star}
    \end{aligned}
  }{%
    \SemigroupAlg{A}{\star}\dsep a\in A\dsep k\in\mathbb N
    \vdash a^{\Succ(k)}\in\langle a\rangle_{A,\star}
  }[SemigroupPowerInMonogenicSubsemigroup]
    \proofstepwidestar[1]{\SemigroupAlg{A}{\star}}{\rA}
    \proofstepwidestar[2]{a\in A}{\rA}
    \proofstepwidestar[3]{k\in\mathbb N}{\rA}
    \proofstepwidestar[2]{\{a\}\subseteq A}{%
      \FormulaRefAuto{a\in A\vdash\{a\}\subseteq A}[thm]{2}}
    \proofstepwidestar[]{\{a\}\neq\varnothing}{%
      \FormulaRefAuto{a\in A\vdash A\neq\varnothing}[thm]{%
        \FormulaRefAuto{a\in\{a\}}[thm]}}
    \proofstepwidestar[1,2]{a\in\langle\{a\}\rangle_{A,\star}}{%
      \FormulaRefAuto{GeneratedSubsemigroupContainsGenerators}[thm]{%
        1,4,5,\FormulaRefAuto{a\in\{a\}}[thm]}}
    \proofstepwidestar[1,2]{a\in\langle a\rangle_{A,\star}}{%
      \rIE{\FormulaRefAuto{a=b\vdash b=a}[thm]{%
        \FormulaRefAuto{SemigroupMonogenicSubsemigroupDef}[def]{1,2}},6}}
    \proofstepwidestar[1,2]{a^{\Succ(0)}=a}{%
      \rIE{\FormulaRefAuto{PeanoOne}[def],
        \FormulaRefAuto{SemigroupPowerOne}[thm]{1,2}}}
    \proofstepwidestar[1,2]{a^{\Succ(0)}\in\langle a\rangle_{A,\star}}{%
      \rIE{\FormulaRefAuto{a=b\vdash b=a}[thm]{8},7}}
    \proofstepwidestar[1,2]{%
      \mathsf{UHGr}(\langle a\rangle_{A,\star};A,\star)}{%
      \FormulaRefAuto{MonogenicSubsemigroupIsSubsemigroup}[thm]{1,2}}
    \proofstepwidestar[1,2,3]{a^{\Succ(k)}\in A}{%
      \FormulaRefAuto{SemigroupPowerClosure}[thm]{%
        1,2,\FormulaRefAuto{PeanoSuccClosure}[ax]{3},
        \FormulaRefAuto{PeanoSuccNonzero}[ax]{3}}}
    \proofstepwidestar[12]{a^{\Succ(k)}\in\langle a\rangle_{A,\star}}{\rA}
    \proofstepwidestar[1,2,3,12]{%
      a^{\Succ(k)+1}=a^{\Succ(k)}\star a}{%
      \FormulaRefAuto{SemigroupPowerSuccessor}[thm]{%
        1,2,\FormulaRefAuto{PeanoSuccClosure}[ax]{3},
        \FormulaRefAuto{PeanoSuccNonzero}[ax]{3}}}
    \proofstepwidestar[1,2,3,12]{%
      a^{\Succ(k)}\star a\in\langle a\rangle_{A,\star}}{%
      \FormulaRefAuto{SemigroupSubsemigroupClosure}[thm]{1,10,12,7}}
    \proofstepwidestar[1,2,3,12]{%
      a^{\Succ(\Succ(k))}\in\langle a\rangle_{A,\star}}{%
      \rIE{\FormulaRefAuto{PeanoAddOneSucc}[thm]{%
          \FormulaRefAuto{PeanoSuccClosure}[ax]{3}},\rIE{13,14}}}
    \proofstepwidestar[1,2]{%
      \forall k\in\mathbb N\,
      \bigl(a^{\Succ(k)}\in\langle a\rangle_{A,\star}\bigr)}{%
      \FormulaRefAuto{PeanoInductionRule}[ax]{9,12,15}}
    \proofstepwidestar[1,2,3]{%
      a^{\Succ(k)}\in\langle a\rangle_{A,\star}}{\rRE{3,\rUE{16}}}
  \closeproofpartwideindR

  \proofpartwideindR[Monogen liegt in der Menge positiver Potenzen]{%
    \begin{aligned}[t]
      &\SemigroupAlg{A}{\star}\dsep a\in A\vdash{}\\[-2pt]
      &\langle a\rangle_{A,\star}\subseteq P_{A,\star}(a)
    \end{aligned}
  }{%
    \SemigroupAlg{A}{\star}\dsep a\in A\vdash
    \langle a\rangle_{A,\star}\subseteq P_{A,\star}(a)%
  }[SemigroupMonogenicSubsemigroupInPositivePowers]
    \proofstepwidestar[1]{\SemigroupAlg{A}{\star}}{\rA}
    \proofstepwidestar[2]{a\in A}{\rA}
    \proofstepwidestarbreak[1,2]{\mathsf{UHGr}(P_{A,\star}(a);A,\star)}{%
      \FormulaRefAuto{SemigroupPositivePowersFormSubsemigroup}[thm]{1,2}}
    \proofstepwidestar[]{1\in\mathbb N}{\FormulaRefAuto{PeanoOneInN}[thm]}
    \proofstepwidestar[]{1\neq0}{%
      \rIE{\FormulaRefAuto{PeanoOne}[def],
        \FormulaRefAuto{PeanoSuccNonzero}[ax]{\FormulaRefAuto{PeanoZero}[ax]}}}
    \proofstepwidestar[1,2]{a^1=a}{%
      \FormulaRefAuto{SemigroupPowerOne}[thm]{1,2}}
    \proofstepwidestarbreak[1,2]{%
      \exists n\in\mathbb N\,(n\neq0\land a=a^n)}{%
      \rEI{4,\rAI{5,\FormulaRefAuto{a=b\vdash b=a}[thm]{6}}}}
    \proofstepwidestar[1,2]{%
      a\in A\land\exists n\in\mathbb N\,(n\neq0\land a=a^n)}{%
      \rAI{2,7}}
    \proofstepwidestarbreak[1,2]{a\in P_{A,\star}(a)}{%
      \FormulaRefAuto{SemigroupPositivePowerSetMembership}[thm]{1,2,8}}
    \proofstepwidestarbreak[1,2]{\{a\}\subseteq P_{A,\star}(a)}{%
      \FormulaRefAuto{a\in A\vdash\{a\}\subseteq A}[thm]{9}}
    \proofstepwidestar[2]{\{a\}\subseteq A}{%
      \FormulaRefAuto{a\in A\vdash\{a\}\subseteq A}[thm]{2}}
    \proofstepwidestar[]{\{a\}\neq\varnothing}{%
      \FormulaRefAuto{a\in A\vdash A\neq\varnothing}[thm]{%
        \FormulaRefAuto{a\in\{a\}}[thm]}}
    \proofstepwidestarbreak[1,2]{%
      \langle\{a\}\rangle_{A,\star}\subseteq P_{A,\star}(a)}{%
      \FormulaRefAuto{GeneratedSubsemigroupMinimal}[thm]{1,11,12,3,10}}
    \proofstepwidestarbreak[1,2]{%
      \langle a\rangle_{A,\star}\subseteq P_{A,\star}(a)}{%
      \rIE{\FormulaRefAuto{SemigroupMonogenicSubsemigroupDef}[def]{1,2},13}}
  \closeproofpartwideindR

  \proofpartwideindR[Menge positiver Potenzen liegt im Monogen]{%
    \begin{aligned}[t]
      &\SemigroupAlg{A}{\star}\dsep a\in A\vdash{}\\[-2pt]
      &P_{A,\star}(a)\subseteq\langle a\rangle_{A,\star}
    \end{aligned}
  }{%
    \SemigroupAlg{A}{\star}\dsep a\in A\vdash
    P_{A,\star}(a)\subseteq\langle a\rangle_{A,\star}%
  }[SemigroupPositivePowersInMonogenicSubsemigroup]
    \proofstepwidestar[1]{\SemigroupAlg{A}{\star}}{\rA}
    \proofstepwidestar[2]{a\in A}{\rA}
    \proofstepwidestar[3]{x\in P_{A,\star}(a)}{\rA}
    \proofstepwidestarbreak[1,2,3]{%
      \exists n\in\mathbb N\,(n\neq0\land x=a^n)}{%
      \rAEb{\FormulaRefAuto{SemigroupPositivePowerSetMembership}[thm]{1,2,3}}}
    \proofstepwidestar[5]{n\in\mathbb N\land(n\neq0\land x=a^n)}{\rA}
    \proofstepwidestar[5]{n\in\mathbb N}{\rAEa{5}}
    \proofstepwidestar[5]{n\neq0}{\rAEa{\rAEb{5}}}
    \proofstepwidestar[5]{x=a^n}{\rAEb{\rAEb{5}}}
    \proofstepwidestar[5]{n=\Succ(\Pred(n))}{%
      \FormulaRefAuto{PeanoPredNonzeroValue}[thm]{6,7}}
    \proofstepwidestar[5]{\Pred(n)\in\mathbb N}{%
      \FormulaRefAuto{PeanoPredNatural}[thm]{6}}
    \proofstepwidestarbreak[1,2,5]{%
      a^{\Succ(\Pred(n))}\in\langle a\rangle_{A,\star}}{%
      \FormulaRefAuto{SemigroupPowerInMonogenicSubsemigroup}[thm]{1,2,10}}
    \proofstepwidestar[1,2,5]{a^n\in\langle a\rangle_{A,\star}}{\rIE{9,11}}
    \proofstepwidestar[1,2,5]{x\in\langle a\rangle_{A,\star}}{\rIE{8,12}}
    \proofstepwidestar[1,2,3]{x\in\langle a\rangle_{A,\star}}{\rEE{4,5,13}}
    \proofstepwidestarbreak[1,2]{%
      P_{A,\star}(a)\subseteq\langle a\rangle_{A,\star}}{%
      \FormulaRefAuto{A\subseteq B\coloneq
        \forall x\,(x\in A\rightarrow x\in B)}[def]{\rUI{\rRI{3,14}}}}
  \closeproofpartwideindR

  \proofpartwideindR[Gleichheit aus beiden Inklusionen]{%
    \begin{aligned}[t]
      &\SemigroupAlg{A}{\star}\dsep a\in A\vdash{}\\[-2pt]
      &\langle a\rangle_{A,\star}=P_{A,\star}(a)
    \end{aligned}
  }{%
    \SemigroupAlg{A}{\star}\dsep a\in A\vdash
    \langle a\rangle_{A,\star}=P_{A,\star}(a)%
  }[SemigroupMonogenicPowerDescriptionConclusion]
    \proofstepwidestar[1]{\SemigroupAlg{A}{\star}}{\rA}
    \proofstepwidestar[2]{a\in A}{\rA}
    \proofstepwidestarbreak[1,2]{%
      \langle a\rangle_{A,\star}\subseteq P_{A,\star}(a)}{%
      \FormulaRefAuto{SemigroupMonogenicSubsemigroupInPositivePowers}[thm]{1,2}}
    \proofstepwidestarbreak[1,2]{%
      P_{A,\star}(a)\subseteq\langle a\rangle_{A,\star}}{%
      \FormulaRefAuto{SemigroupPositivePowersInMonogenicSubsemigroup}[thm]{1,2}}
    \proofstepwidestarbreak[1,2]{%
      \langle a\rangle_{A,\star}=P_{A,\star}(a)}{%
      \FormulaRefAuto{A\subseteq B,\,B\subseteq A\vdash A=B}[thm]{3,4}}
  \closeproofpartwideindR
\end{tabproofsplitwide}

\section{Ideale}

Im Unterschied zu Unterhalbgruppen werden Ideale durch Multiplikation mit
beliebigen Elementen der umgebenden Halbgruppe absorbiert. Wir verlangen
ausdr\"ucklich Nichtleerheit; andernfalls w\"are die leere Menge stets das
kleinste Ideal und die sp\"atere Minimalidealtheorie w\"urde entarten.

\FormulaDefDeltaK[Linksideal]{%
  \begin{aligned}[t]
    &\SemigroupAlg{A}{\star}\vdash{}\\[-2pt]
    &\mathsf{LId}(I;A,\star)\coloneqq
      I\subseteq A\land I\neq\varnothing\land
      \forall a\in A\,\forall x\in I\,(a\star x\in I)
  \end{aligned}
}{SemigroupLeftIdealDef}{%
  \DeltaRow{Mengen}{A\dsep I\dsep a\dsep x}
  \DeltaRow{Funktionen}{\star}
  \DeltaRow{\textbf{Neues Symbol}}{}
  \DeltaPrem{Linksideal}{\mathsf{LId}(I;A,\star)}
}

\FormulaDefDeltaK[Rechtsideal]{%
  \begin{aligned}[t]
    &\SemigroupAlg{A}{\star}\vdash{}\\[-2pt]
    &\mathsf{RId}(I;A,\star)\coloneqq
      I\subseteq A\land I\neq\varnothing\land
      \forall x\in I\,\forall a\in A\,(x\star a\in I)
  \end{aligned}
}{SemigroupRightIdealDef}{%
  \DeltaRow{Mengen}{A\dsep I\dsep a\dsep x}
  \DeltaRow{Funktionen}{\star}
  \DeltaRow{\textbf{Neues Symbol}}{}
  \DeltaPrem{Rechtsideal}{\mathsf{RId}(I;A,\star)}
}

\FormulaDefDeltaK[Zweiseitiges Ideal]{%
  \SemigroupAlg{A}{\star}\vdash
  \mathsf{Id}(I;A,\star)\coloneqq
  \mathsf{LId}(I;A,\star)\land\mathsf{RId}(I;A,\star)%
}{SemigroupIdealDef}{%
  \DeltaRow{Mengen}{A\dsep I}
  \DeltaRow{Funktionen}{\star}
  \DeltaRow{\textbf{Neues Symbol}}{}
  \DeltaPrem{Zweiseitiges Ideal}{\mathsf{Id}(I;A,\star)}
}

\FormulaThmDeltaK[Tr\"agermenge eines Ideals]{%
  \SemigroupAlg{A}{\star}\dsep\mathsf{Id}(I;A,\star)
  \vdash I\subseteq A%
}{SemigroupIdealSubset}{%
  \DeltaRow{Mengen}{A\dsep I}
  \DeltaRow{Funktionen}{\star}
}
\begin{tabproofwide}
  \proofstepwidestar[1]{\SemigroupAlg{A}{\star}}{\rA}
  \proofstepwidestar[2]{\mathsf{Id}(I;A,\star)}{\rA}
  \proofstepwidestar[1,2]{%
    \mathsf{LId}(I;A,\star)\land\mathsf{RId}(I;A,\star)}{%
    \FormulaRefAuto{SemigroupIdealDef}[def]{1,2}}
  \proofstepwidestar[1,2]{\mathsf{LId}(I;A,\star)}{\rAEa{3}}
  \proofstepwidestar[1,2]{%
    I\subseteq A\land I\neq\varnothing\land
    \forall a\in A\,\forall x\in I\,(a\star x\in I)}{%
    \FormulaRefAuto{SemigroupLeftIdealDef}[def]{1,4}}
  \proofstepwidestar[1,2]{I\subseteq A}{\rAEa{5}}
\end{tabproofwide}

\FormulaThmDeltaK[Nichtleerheit eines Ideals]{%
  \SemigroupAlg{A}{\star}\dsep\mathsf{Id}(I;A,\star)
  \vdash I\neq\varnothing%
}{SemigroupIdealNonempty}{%
  \DeltaRow{Mengen}{A\dsep I}
  \DeltaRow{Funktionen}{\star}
}
\begin{tabproofwide}
  \proofstepwidestar[1]{\SemigroupAlg{A}{\star}}{\rA}
  \proofstepwidestar[2]{\mathsf{Id}(I;A,\star)}{\rA}
  \proofstepwidestar[1,2]{\mathsf{LId}(I;A,\star)}{%
    \rAEa{\FormulaRefAuto{SemigroupIdealDef}[def]{1,2}}}
  \proofstepwidestar[1,2]{%
    I\subseteq A\land I\neq\varnothing\land
    \forall a\in A\,\forall x\in I\,(a\star x\in I)}{%
    \FormulaRefAuto{SemigroupLeftIdealDef}[def]{1,3}}
  \proofstepwidestar[1,2]{I\neq\varnothing}{\rAEa{\rAEb{4}}}
\end{tabproofwide}

\FormulaThmDeltaK[Linksabsorption eines Ideals]{%
  \SemigroupAlg{A}{\star}\dsep\mathsf{Id}(I;A,\star)\dsep
  a\in A\dsep x\in I\vdash a\star x\in I%
}{SemigroupIdealLeftAbsorption}{%
  \DeltaRow{Mengen}{A\dsep I\dsep a\dsep x}
  \DeltaRow{Funktionen}{\star}
}
\begin{tabproofwide}
  \proofstepwidestar[1]{\SemigroupAlg{A}{\star}}{\rA}
  \proofstepwidestar[2]{\mathsf{Id}(I;A,\star)}{\rA}
  \proofstepwidestar[3]{a\in A}{\rA}
  \proofstepwidestar[4]{x\in I}{\rA}
  \proofstepwidestar[1,2]{\mathsf{LId}(I;A,\star)}{%
    \rAEa{\FormulaRefAuto{SemigroupIdealDef}[def]{1,2}}}
  \proofstepwidestar[1,2]{%
    \forall u\in A\,\forall v\in I\,(u\star v\in I)}{%
    \rAEb{\rAEb{\FormulaRefAuto{SemigroupLeftIdealDef}[def]{1,5}}}}
  \proofstepwidestar[1,2,3,4]{a\star x\in I}{\rRE{4,\rUE{\rRE{3,\rUE{6}}}}}
\end{tabproofwide}

\FormulaThmDeltaK[Rechtsabsorption eines Ideals]{%
  \SemigroupAlg{A}{\star}\dsep\mathsf{Id}(I;A,\star)\dsep
  x\in I\dsep a\in A\vdash x\star a\in I%
}{SemigroupIdealRightAbsorption}{%
  \DeltaRow{Mengen}{A\dsep I\dsep a\dsep x}
  \DeltaRow{Funktionen}{\star}
}
\begin{tabproofwide}
  \proofstepwidestar[1]{\SemigroupAlg{A}{\star}}{\rA}
  \proofstepwidestar[2]{\mathsf{Id}(I;A,\star)}{\rA}
  \proofstepwidestar[3]{x\in I}{\rA}
  \proofstepwidestar[4]{a\in A}{\rA}
  \proofstepwidestar[1,2]{\mathsf{RId}(I;A,\star)}{%
    \rAEb{\FormulaRefAuto{SemigroupIdealDef}[def]{1,2}}}
  \proofstepwidestar[1,2]{%
    \forall u\in I\,\forall v\in A\,(u\star v\in I)}{%
    \rAEb{\rAEb{\FormulaRefAuto{SemigroupRightIdealDef}[def]{1,5}}}}
  \proofstepwidestar[1,2,3,4]{x\star a\in I}{\rRE{4,\rUE{\rRE{3,\rUE{6}}}}}
\end{tabproofwide}

\FormulaThmDeltaK[Der nichtleere Tr\"ager ist ein Ideal]{%
  \SemigroupAlg{A}{\star}\dsep A\neq\varnothing
  \vdash\mathsf{Id}(A;A,\star)%
}{SemigroupCarrierIsIdeal}{%
  \DeltaRow{Mengen}{A\dsep a\dsep x}
  \DeltaRow{Funktionen}{\star}
}
\begin{tabproofwide}
  \proofstepwidestar[1]{\SemigroupAlg{A}{\star}}{\rA}
  \proofstepwidestar[2]{A\neq\varnothing}{\rA}
  \proofstepwidestar[]{A\subseteq A}{\FormulaRefAuto{A\subseteq A}[thm]}
  \proofstepwidestar[4]{a\in A}{\rA}
  \proofstepwidestar[5]{x\in A}{\rA}
  \proofstepwidestar[1,4,5]{a\star x\in A}{%
    \FormulaRefAuto{SemigroupClosureAxiom}[ax]{1,4,5}}
  \proofstepwidestar[1]{\forall a\in A\,\forall x\in A\,(a\star x\in A)}{%
    \rUI{\rRI{4,\rUI{\rRI{5,6}}}}}
  \proofstepwidestar[1,2]{\mathsf{LId}(A;A,\star)}{%
    \FormulaRefAuto{SemigroupLeftIdealDef}[def]{1,\rAI{3,\rAI{2,7}}}}
  \proofstepwidestar[1]{\forall x\in A\,\forall a\in A\,(x\star a\in A)}{%
    \rUI{\rRI{5,\rUI{\rRI{4,6}}}}}
  \proofstepwidestar[1,2]{\mathsf{RId}(A;A,\star)}{%
    \FormulaRefAuto{SemigroupRightIdealDef}[def]{1,\rAI{3,\rAI{2,9}}}}
  \proofstepwidestar[1,2]{\mathsf{Id}(A;A,\star)}{%
    \FormulaRefAuto{SemigroupIdealDef}[def]{1,\rAI{8,10}}}
\end{tabproofwide}

\FormulaThmDeltaK[Schnitt zweier Ideale]{%
  \begin{aligned}[t]
    &\SemigroupAlg{A}{\star}\dsep\mathsf{Id}(I;A,\star)\dsep
      \mathsf{Id}(J;A,\star)\vdash{}\\[-2pt]
    &\mathsf{Id}(I\cap J;A,\star)
  \end{aligned}
}{SemigroupIdealIntersection}{%
  \DeltaRow{Mengen}{A\dsep I\dsep J\dsep a\dsep x\dsep i\dsep j}
  \DeltaRow{Funktionen}{\star}
}
Der Schnitt ist nichtleer, weil das Produkt eines Elements von (I) mit
einem Element von (J) zugleich in beiden Idealen liegt.  Die linke und
die rechte Absorption werden danach komponentenweise vererbt.

\begin{tabproofsplitwide}
  \proofpartwideindR[Nichtleerheit]{%
    \SemigroupAlg{A}{\star}\dsep\mathsf{Id}(I;A,\star)\dsep
    \mathsf{Id}(J;A,\star)\vdash I\cap J\neq\varnothing
  }{%
    \SemigroupAlg{A}{\star}\dsep\mathsf{Id}(I;A,\star)\dsep
    \mathsf{Id}(J;A,\star)\vdash I\cap J\neq\varnothing
  }[SemigroupIdealIntersectionNonempty]
    \proofstepwidestar[1]{\SemigroupAlg{A}{\star}}{\rA}
    \proofstepwidestar[2]{\mathsf{Id}(I;A,\star)}{\rA}
    \proofstepwidestar[3]{\mathsf{Id}(J;A,\star)}{\rA}
    \proofstepwidestar[1,2]{I\subseteq A}{%
      \FormulaRefAuto{SemigroupIdealSubset}[thm]{1,2}}
    \proofstepwidestar[1,3]{J\subseteq A}{%
      \FormulaRefAuto{SemigroupIdealSubset}[thm]{1,3}}
    \proofstepwidestar[1,2]{I\neq\varnothing}{%
      \FormulaRefAuto{SemigroupIdealNonempty}[thm]{1,2}}
    \proofstepwidestar[1,3]{J\neq\varnothing}{%
      \FormulaRefAuto{SemigroupIdealNonempty}[thm]{1,3}}
    \proofstepwidestar[1,2]{\exists i\,(i\in I)}{%
      \FormulaRefAuto{A\neq\varnothing\eqvdash\exists x\,(x\in A)}[thm]{6}}
    \proofstepwidestar[1,3]{\exists j\,(j\in J)}{%
      \FormulaRefAuto{A\neq\varnothing\eqvdash\exists x\,(x\in A)}[thm]{7}}
    \proofstepwidestar[10]{i\in I}{\rA}
    \proofstepwidestar[11]{j\in J}{\rA}
    \proofstepwidestar[1,2,10]{i\in A}{%
      \FormulaRefAuto{A\subseteq B,\,x\in A\vdash x\in B}[thm]{4,10}}
    \proofstepwidestar[1,3,11]{j\in A}{%
      \FormulaRefAuto{A\subseteq B,\,x\in A\vdash x\in B}[thm]{5,11}}
    \proofstepwidestar[1,3,10,11]{i\star j\in J}{%
      \FormulaRefAuto{SemigroupIdealLeftAbsorption}[thm]{1,3,12,11}}
    \proofstepwidestar[1,2,10,11]{i\star j\in I}{%
      \FormulaRefAuto{SemigroupIdealRightAbsorption}[thm]{1,2,10,13}}
    \proofstepwidestar[1,2,3,10,11]{i\star j\in I\cap J}{%
      \FormulaRefAuto{x\in A\cap B\eqvdash x\in A\land x\in B}[thm]{%
        \rAI{15,14}}}
    \proofstepwidestar[1,2,3,10,11]{I\cap J\neq\varnothing}{%
      \FormulaRefAuto{a\in A\vdash A\neq\varnothing}[thm]{16}}
    \proofstepwidestar[1,2,3,10]{I\cap J\neq\varnothing}{%
      \rEE{9,11,17}}
    \proofstepwidestar[1,2,3]{I\cap J\neq\varnothing}{%
      \rEE{8,10,18}}
  \closeproofpartwideindR

  \proofpartwideindR[Linksidealeigenschaft]{%
    \SemigroupAlg{A}{\star}\dsep\mathsf{Id}(I;A,\star)\dsep
    \mathsf{Id}(J;A,\star)\vdash\mathsf{LId}(I\cap J;A,\star)
  }{%
    \SemigroupAlg{A}{\star}\dsep\mathsf{Id}(I;A,\star)\dsep
    \mathsf{Id}(J;A,\star)\vdash\mathsf{LId}(I\cap J;A,\star)
  }[SemigroupIdealIntersectionLeftIdeal]
    \proofstepwidestar[1]{\SemigroupAlg{A}{\star}}{\rA}
    \proofstepwidestar[2]{\mathsf{Id}(I;A,\star)}{\rA}
    \proofstepwidestar[3]{\mathsf{Id}(J;A,\star)}{\rA}
    \proofstepwidestar[1,2]{I\subseteq A}{%
      \FormulaRefAuto{SemigroupIdealSubset}[thm]{1,2}}
    \proofstepwidestar[]{I\cap J\subseteq I}{%
      \FormulaRefAuto{A\cap B\subseteq A}[thm]}
    \proofstepwidestar[1,2]{I\cap J\subseteq A}{%
      \FormulaRefAuto{A\subseteq B\dsep B\subseteq C\vdash
        A\subseteq C}[thm]{5,4}}
    \proofstepwidestar[1,2,3]{I\cap J\neq\varnothing}{%
      \FormulaRefAuto{SemigroupIdealIntersectionNonempty}[thm]{1,2,3}}
    \proofstepwidestar[8]{a\in A}{\rA}
    \proofstepwidestar[9]{x\in I\cap J}{\rA}
    \proofstepwidestar[9]{x\in I}{%
      \FormulaRefAuto{x\in A\cap B\vdash x\in A}[thm]{9}}
    \proofstepwidestar[9]{x\in J}{%
      \FormulaRefAuto{x\in A\cap B\vdash x\in B}[thm]{9}}
    \proofstepwidestar[1,2,8,9]{a\star x\in I}{%
      \FormulaRefAuto{SemigroupIdealLeftAbsorption}[thm]{1,2,8,10}}
    \proofstepwidestar[1,3,8,9]{a\star x\in J}{%
      \FormulaRefAuto{SemigroupIdealLeftAbsorption}[thm]{1,3,8,11}}
    \proofstepwidestar[1,2,3,8,9]{a\star x\in I\cap J}{%
      \FormulaRefAuto{x\in A\cap B\eqvdash x\in A\land x\in B}[thm]{%
        \rAI{12,13}}}
    \proofstepwidestar[1,2,3]{%
      \forall a\in A\,\forall x\in I\cap J\,(a\star x\in I\cap J)}{%
      \rUI{\rRI{8,\rUI{\rRI{9,14}}}}}
    \proofstepwidestar[1,2,3]{\mathsf{LId}(I\cap J;A,\star)}{%
      \FormulaRefAuto{SemigroupLeftIdealDef}[def]{1,\rAI{6,\rAI{7,15}}}}
  \closeproofpartwideindR

  \proofpartwideindR[Rechtsidealeigenschaft]{%
    \SemigroupAlg{A}{\star}\dsep\mathsf{Id}(I;A,\star)\dsep
    \mathsf{Id}(J;A,\star)\vdash\mathsf{RId}(I\cap J;A,\star)
  }{%
    \SemigroupAlg{A}{\star}\dsep\mathsf{Id}(I;A,\star)\dsep
    \mathsf{Id}(J;A,\star)\vdash\mathsf{RId}(I\cap J;A,\star)
  }[SemigroupIdealIntersectionRightIdeal]
    \proofstepwidestar[1]{\SemigroupAlg{A}{\star}}{\rA}
    \proofstepwidestar[2]{\mathsf{Id}(I;A,\star)}{\rA}
    \proofstepwidestar[3]{\mathsf{Id}(J;A,\star)}{\rA}
    \proofstepwidestar[1,2]{I\subseteq A}{%
      \FormulaRefAuto{SemigroupIdealSubset}[thm]{1,2}}
    \proofstepwidestar[]{I\cap J\subseteq I}{%
      \FormulaRefAuto{A\cap B\subseteq A}[thm]}
    \proofstepwidestar[1,2]{I\cap J\subseteq A}{%
      \FormulaRefAuto{A\subseteq B\dsep B\subseteq C\vdash
        A\subseteq C}[thm]{5,4}}
    \proofstepwidestar[1,2,3]{I\cap J\neq\varnothing}{%
      \FormulaRefAuto{SemigroupIdealIntersectionNonempty}[thm]{1,2,3}}
    \proofstepwidestar[8]{x\in I\cap J}{\rA}
    \proofstepwidestar[9]{a\in A}{\rA}
    \proofstepwidestar[8]{x\in I}{%
      \FormulaRefAuto{x\in A\cap B\vdash x\in A}[thm]{8}}
    \proofstepwidestar[8]{x\in J}{%
      \FormulaRefAuto{x\in A\cap B\vdash x\in B}[thm]{8}}
    \proofstepwidestar[1,2,8,9]{x\star a\in I}{%
      \FormulaRefAuto{SemigroupIdealRightAbsorption}[thm]{1,2,10,9}}
    \proofstepwidestar[1,3,8,9]{x\star a\in J}{%
      \FormulaRefAuto{SemigroupIdealRightAbsorption}[thm]{1,3,11,9}}
    \proofstepwidestar[1,2,3,8,9]{x\star a\in I\cap J}{%
      \FormulaRefAuto{x\in A\cap B\eqvdash x\in A\land x\in B}[thm]{%
        \rAI{12,13}}}
    \proofstepwidestar[1,2,3]{%
      \forall x\in I\cap J\,\forall a\in A\,(x\star a\in I\cap J)}{%
      \rUI{\rRI{8,\rUI{\rRI{9,14}}}}}
    \proofstepwidestar[1,2,3]{\mathsf{RId}(I\cap J;A,\star)}{%
      \FormulaRefAuto{SemigroupRightIdealDef}[def]{1,\rAI{6,\rAI{7,15}}}}
  \closeproofpartwideindR

  \proofpartwide{Schluss}
    \proofstepwidestar[1]{\SemigroupAlg{A}{\star}}{\rA}
    \proofstepwidestar[2]{\mathsf{Id}(I;A,\star)}{\rA}
    \proofstepwidestar[3]{\mathsf{Id}(J;A,\star)}{\rA}
    \proofstepwidestar[1,2,3]{\mathsf{LId}(I\cap J;A,\star)}{%
      \FormulaRefAuto{SemigroupIdealIntersectionLeftIdeal}[thm]{1,2,3}}
    \proofstepwidestar[1,2,3]{\mathsf{RId}(I\cap J;A,\star)}{%
      \FormulaRefAuto{SemigroupIdealIntersectionRightIdeal}[thm]{1,2,3}}
    \proofstepwidestar[1,2,3]{\mathsf{Id}(I\cap J;A,\star)}{%
      \FormulaRefAuto{SemigroupIdealDef}[def]{1,\rAI{4,5}}}
  \closeproofpartwide
\end{tabproofsplitwide}

\subsection{Hauptideale}

Statt die symbolische Einheitsadjunktion \(A^1\) informell zu benutzen,
definieren wir Hauptideale unmittelbar durch ihre Elemente. Damit sind auch
die F\"alle ohne linken oder rechten Faktor ausdr\"ucklich erfasst.

\FormulaDefDeltaK[Linkshauptideal]{%
  \SemigroupAlg{A}{\star}\dsep a\in A\vdash
  L_{A,\star}(a)\coloneqq
  \bigl\{b\in A\bigm|b=a\lor
    \exists x\in A\,(b=x\star a)\bigr\}%
}{SemigroupPrincipalLeftIdealDef}{%
  \DeltaRow{Mengen}{A\dsep a\dsep b\dsep x}
  \DeltaRow{Funktionen}{\star}
  \DeltaRow{\textbf{Neues Symbol}}{}
  \DeltaPrem{Linkshauptideal}{L_{A,\star}(a)}
}

\FormulaDefDeltaK[Rechtshauptideal]{%
  \SemigroupAlg{A}{\star}\dsep a\in A\vdash
  R_{A,\star}(a)\coloneqq
  \bigl\{b\in A\bigm|b=a\lor
    \exists y\in A\,(b=a\star y)\bigr\}%
}{SemigroupPrincipalRightIdealDef}{%
  \DeltaRow{Mengen}{A\dsep a\dsep b\dsep y}
  \DeltaRow{Funktionen}{\star}
  \DeltaRow{\textbf{Neues Symbol}}{}
  \DeltaPrem{Rechtshauptideal}{R_{A,\star}(a)}
}

\FormulaDefDeltaK[Zweiseitiges Hauptideal]{%
  \begin{aligned}[t]
    &\SemigroupAlg{A}{\star}\dsep a\in A\vdash{}\\[-2pt]
    &J_{A,\star}(a)\coloneqq
      \Bigl\{b\in A\Bigm|
        b=a\lor\exists x\in A\,(b=x\star a)\\[-2pt]
    &\hspace{93pt}{}\lor\exists y\in A\,(b=a\star y)
        \lor\exists x,y\in A\,(b=(x\star a)\star y)
      \Bigr\}
  \end{aligned}
}{SemigroupPrincipalTwoSidedIdealDef}{%
  \DeltaRow{Mengen}{A\dsep a\dsep b\dsep x\dsep y}
  \DeltaRow{Funktionen}{\star}
  \DeltaRow{\textbf{Neues Symbol}}{}
  \DeltaPrem{Zweiseitiges Hauptideal}{J_{A,\star}(a)}
}

\FormulaThmDeltaK[Elementkriterium des Linkshauptideals]{%
  \begin{aligned}[t]
    &\SemigroupAlg{A}{\star}\dsep a\in A\vdash{}\\[-2pt]
    &b\in L_{A,\star}(a)\leftrightarrow
      \bigl(b\in A\land(b=a\lor\exists x\in A\,(b=x\star a))\bigr)
  \end{aligned}
}{SemigroupPrincipalLeftIdealMembership}{%
  \DeltaRow{Mengen}{A\dsep a\dsep b\dsep x}
  \DeltaRow{Funktionen}{\star}
}
\begin{tabproofwide}
  \proofstepwidestar[1]{\SemigroupAlg{A}{\star}}{\rA}
  \proofstepwidestar[2]{a\in A}{\rA}
  \proofstepwidestar[1,2]{%
    L_{A,\star}(a)=\{u\in A\mid
      u=a\lor\exists x\in A\,(u=x\star a)\}}{%
    \FormulaRefAuto{SemigroupPrincipalLeftIdealDef}[def]{1,2}}
  \proofstepwidestar[1,2]{%
    b\in L_{A,\star}(a)\leftrightarrow
    \bigl(b\in A\land(b=a\lor\exists x\in A\,(b=x\star a))\bigr)}{%
    \rIE{3,\FormulaRefAuto{x\in\{u\in A\mid P(u)\}\eqvdash
      x\in A\land P(x)}[thm]}}
\end{tabproofwide}

\FormulaThmDeltaK[Elementkriterium des Rechtshauptideals]{%
  \begin{aligned}[t]
    &\SemigroupAlg{A}{\star}\dsep a\in A\vdash{}\\[-2pt]
    &b\in R_{A,\star}(a)\leftrightarrow
      \bigl(b\in A\land(b=a\lor\exists y\in A\,(b=a\star y))\bigr)
  \end{aligned}
}{SemigroupPrincipalRightIdealMembership}{%
  \DeltaRow{Mengen}{A\dsep a\dsep b\dsep y}
  \DeltaRow{Funktionen}{\star}
}
\begin{tabproofwide}
  \proofstepwidestar[1]{\SemigroupAlg{A}{\star}}{\rA}
  \proofstepwidestar[2]{a\in A}{\rA}
  \proofstepwidestar[1,2]{%
    R_{A,\star}(a)=\{u\in A\mid
      u=a\lor\exists y\in A\,(u=a\star y)\}}{%
    \FormulaRefAuto{SemigroupPrincipalRightIdealDef}[def]{1,2}}
  \proofstepwidestar[1,2]{%
    b\in R_{A,\star}(a)\leftrightarrow
    \bigl(b\in A\land(b=a\lor\exists y\in A\,(b=a\star y))\bigr)}{%
    \rIE{3,\FormulaRefAuto{x\in\{u\in A\mid P(u)\}\eqvdash
      x\in A\land P(x)}[thm]}}
\end{tabproofwide}

\FormulaThmDeltaK[Elementkriterium des zweiseitigen Hauptideals]{%
  \begin{aligned}[t]
    &\SemigroupAlg{A}{\star}\dsep a\in A\vdash{}\\[-2pt]
    &b\in J_{A,\star}(a)\leftrightarrow
      \Bigl(b\in A\land\bigl(
        b=a\lor\exists x\in A\,(b=x\star a)\\[-2pt]
    &\hspace{112pt}{}\lor\exists y\in A\,(b=a\star y)
        \lor\exists x,y\in A\,(b=(x\star a)\star y)
      \bigr)\Bigr)
  \end{aligned}
}{SemigroupPrincipalTwoSidedIdealMembership}{%
  \DeltaRow{Mengen}{A\dsep a\dsep b\dsep x\dsep y}
  \DeltaRow{Funktionen}{\star}
}
\begin{tabproofwide}
  \proofstepwidestar[1]{\SemigroupAlg{A}{\star}}{\rA}
  \proofstepwidestar[2]{a\in A}{\rA}
  \proofstepwidestarbreak[1,2]{%
    \begin{aligned}[t]
      J_{A,\star}(a)=\Bigl\{u\in A\Bigm|{}&
        u=a\lor\exists x\in A\,(u=x\star a)\\[-2pt]
        &{}\lor\exists y\in A\,(u=a\star y)\\[-2pt]
        &{}\lor\exists x,y\in A\,(u=(x\star a)\star y)\Bigr\}
    \end{aligned}}{%
    \FormulaRefAuto{SemigroupPrincipalTwoSidedIdealDef}[def]{1,2}}
  \proofstepwidestarbreak[1,2]{%
    \begin{aligned}[t]
      b\in J_{A,\star}(a)\leftrightarrow
      \Bigl(b\in A\land\bigl({}&
        b=a\lor\exists x\in A\,(b=x\star a)\\[-2pt]
        &{}\lor\exists y\in A\,(b=a\star y)\\[-2pt]
        &{}\lor\exists x,y\in A\,(b=(x\star a)\star y)\bigr)\Bigr)
    \end{aligned}}{%
    \rIE{3,\FormulaRefAuto{x\in\{u\in A\mid P(u)\}\eqvdash
      x\in A\land P(x)}[thm]}}
\end{tabproofwide}

\FormulaThmDeltaK[Der Erzeuger liegt in seinen Hauptidealen]{%
  \begin{aligned}[t]
    &\SemigroupAlg{A}{\star}\dsep a\in A\vdash{}\\[-2pt]
    &a\in L_{A,\star}(a)\land
      a\in R_{A,\star}(a)\land
      a\in J_{A,\star}(a)
  \end{aligned}
}{SemigroupPrincipalIdealsContainGenerator}{%
  \DeltaRow{Mengen}{A\dsep a}
  \DeltaRow{Funktionen}{\star}
}
\begin{tabproofwide}
  \proofstepwidestar[1]{\SemigroupAlg{A}{\star}}{\rA}
  \proofstepwidestar[2]{a\in A}{\rA}
  \proofstepwidestar[1,2]{a\in L_{A,\star}(a)}{%
    \FormulaRefAuto{SemigroupPrincipalLeftIdealMembership}[thm]{%
      1,2,\rAI{2,\rOIa{\rII}}}}
  \proofstepwidestar[1,2]{a\in R_{A,\star}(a)}{%
    \FormulaRefAuto{SemigroupPrincipalRightIdealMembership}[thm]{%
      1,2,\rAI{2,\rOIa{\rII}}}}
  \proofstepwidestar[1,2]{a\in J_{A,\star}(a)}{%
    \FormulaRefAuto{SemigroupPrincipalTwoSidedIdealMembership}[thm]{%
      1,2,\rAI{2,\rOIa{\rII}}}}
  \proofstepwidestar[1,2]{%
    a\in L_{A,\star}(a)\land a\in R_{A,\star}(a)\land
    a\in J_{A,\star}(a)}{\rAI{3,\rAI{4,5}}}
\end{tabproofwide}

\FormulaThmDeltaK[Das Linkshauptideal ist ein Linksideal]{%
  \SemigroupAlg{A}{\star}\dsep a\in A\vdash
  \mathsf{LId}(L_{A,\star}(a);A,\star)%
}{SemigroupPrincipalLeftIdealIsLeftIdeal}{%
  \DeltaRow{Mengen}{A\dsep a\dsep b\dsep s\dsep x}
  \DeltaRow{Funktionen}{\star}
}
\begin{tabproofwide}
  \proofstepwidestar[1]{\SemigroupAlg{A}{\star}}{\rA}
  \proofstepwidestar[2]{a\in A}{\rA}
  \proofstepwidestar[]{L_{A,\star}(a)\subseteq A}{%
    \FormulaRefAuto{\{x\in A\mid P(x)\}\subseteq A}[thm]}
  \proofstepwidestar[1,2]{a\in L_{A,\star}(a)}{%
    \rAEa{\FormulaRefAuto{SemigroupPrincipalIdealsContainGenerator}[thm]{1,2}}}
  \proofstepwidestar[1,2]{L_{A,\star}(a)\neq\varnothing}{%
    \FormulaRefAuto{a\in A\vdash A\neq\varnothing}[thm]{4}}
  \proofstepwidestar[6]{s\in A}{\rA}
  \proofstepwidestar[7]{b\in L_{A,\star}(a)}{\rA}
  \proofstepwidestar[1,2,7]{b\in A}{%
    \rAEa{\FormulaRefAuto{SemigroupPrincipalLeftIdealMembership}[thm]{1,2,7}}}
  \proofstepwidestar[1,2,6,7]{s\star b\in A}{%
    \FormulaRefAuto{SemigroupClosureAxiom}[ax]{1,6,8}}
  \proofstepwidestar[1,2,7]{b=a\lor\exists x\in A\,(b=x\star a)}{%
    \rAEb{\FormulaRefAuto{SemigroupPrincipalLeftIdealMembership}[thm]{1,2,7}}}
  \proofstepwidestar[11]{b=a}{\rA}
  \proofstepwidestar[6,11]{s\star b=s\star a}{\rIE{11}}
  \proofstepwidestar[6,11]{\exists u\in A\,(s\star b=u\star a)}{%
    \rEI{6,12}}
  \proofstepwidestar[1,2,6,7,11]{s\star b\in L_{A,\star}(a)}{%
    \FormulaRefAuto{SemigroupPrincipalLeftIdealMembership}[thm]{%
      1,2,\rAI{9,\rOIb{13}}}}
  \proofstepwidestar[15]{\exists x\in A\,(b=x\star a)}{\rA}
  \proofstepwidestar[16]{x\in A\land b=x\star a}{\rA}
  \proofstepwidestar[16]{x\in A}{\rAEa{16}}
  \proofstepwidestar[16]{b=x\star a}{\rAEb{16}}
  \proofstepwidestar[1,6,16]{s\star x\in A}{%
    \FormulaRefAuto{SemigroupClosureAxiom}[ax]{1,6,17}}
  \proofstepwidestar[1,2,6,16]{s\star b=(s\star x)\star a}{%
    \rIE{18,\FormulaRefAuto{SemigroupAssociativityAxiom}[ax]{1,6,17,2}}}
  \proofstepwidestar[1,2,6,16]{%
    \exists u\in A\,(s\star b=u\star a)}{\rEI{19,20}}
  \proofstepwidestar[1,2,6,7,16]{s\star b\in L_{A,\star}(a)}{%
    \FormulaRefAuto{SemigroupPrincipalLeftIdealMembership}[thm]{%
      1,2,\rAI{9,\rOIb{21}}}}
  \proofstepwidestar[1,2,6,7,15]{s\star b\in L_{A,\star}(a)}{%
    \rEE{15,16,22}}
  \proofstepwidestar[1,2,6,7]{s\star b\in L_{A,\star}(a)}{%
    \rOE{10,11,14,15,23}}
  \proofstepwidestar[1,2]{%
    \forall s\in A\,\forall b\in L_{A,\star}(a)\,
      (s\star b\in L_{A,\star}(a))}{%
    \rUI{\rRI{6,\rUI{\rRI{7,24}}}}}
  \proofstepwidestar[1,2]{\mathsf{LId}(L_{A,\star}(a);A,\star)}{%
    \FormulaRefAuto{SemigroupLeftIdealDef}[def]{1,\rAI{3,\rAI{5,25}}}}
\end{tabproofwide}

\FormulaThmDeltaK[Das Rechtshauptideal ist ein Rechtsideal]{%
  \SemigroupAlg{A}{\star}\dsep a\in A\vdash
  \mathsf{RId}(R_{A,\star}(a);A,\star)%
}{SemigroupPrincipalRightIdealIsRightIdeal}{%
  \DeltaRow{Mengen}{A\dsep a\dsep b\dsep s\dsep y}
  \DeltaRow{Funktionen}{\star}
}
\begin{tabproofwide}
  \proofstepwidestar[1]{\SemigroupAlg{A}{\star}}{\rA}
  \proofstepwidestar[2]{a\in A}{\rA}
  \proofstepwidestar[]{R_{A,\star}(a)\subseteq A}{%
    \FormulaRefAuto{\{x\in A\mid P(x)\}\subseteq A}[thm]}
  \proofstepwidestar[1,2]{a\in R_{A,\star}(a)}{%
    \rAEa{\rAEb{\FormulaRefAuto{SemigroupPrincipalIdealsContainGenerator}[thm]{1,2}}}}
  \proofstepwidestar[1,2]{R_{A,\star}(a)\neq\varnothing}{%
    \FormulaRefAuto{a\in A\vdash A\neq\varnothing}[thm]{4}}
  \proofstepwidestar[6]{b\in R_{A,\star}(a)}{\rA}
  \proofstepwidestar[7]{s\in A}{\rA}
  \proofstepwidestar[1,2,6]{b\in A}{%
    \rAEa{\FormulaRefAuto{SemigroupPrincipalRightIdealMembership}[thm]{1,2,6}}}
  \proofstepwidestar[1,2,6,7]{b\star s\in A}{%
    \FormulaRefAuto{SemigroupClosureAxiom}[ax]{1,8,7}}
  \proofstepwidestar[1,2,6]{b=a\lor\exists y\in A\,(b=a\star y)}{%
    \rAEb{\FormulaRefAuto{SemigroupPrincipalRightIdealMembership}[thm]{1,2,6}}}
  \proofstepwidestar[11]{b=a}{\rA}
  \proofstepwidestar[7,11]{b\star s=a\star s}{\rIE{11}}
  \proofstepwidestar[7,11]{\exists u\in A\,(b\star s=a\star u)}{\rEI{7,12}}
  \proofstepwidestar[1,2,6,7,11]{b\star s\in R_{A,\star}(a)}{%
    \FormulaRefAuto{SemigroupPrincipalRightIdealMembership}[thm]{%
      1,2,\rAI{9,\rOIb{13}}}}
  \proofstepwidestar[15]{\exists y\in A\,(b=a\star y)}{\rA}
  \proofstepwidestar[16]{y\in A\land b=a\star y}{\rA}
  \proofstepwidestar[16]{y\in A}{\rAEa{16}}
  \proofstepwidestar[16]{b=a\star y}{\rAEb{16}}
  \proofstepwidestar[1,7,16]{y\star s\in A}{%
    \FormulaRefAuto{SemigroupClosureAxiom}[ax]{1,17,7}}
  \proofstepwidestar[1,2,7,16]{b\star s=a\star(y\star s)}{%
    \rIE{18,\FormulaRefAuto{SemigroupAssociativityAxiom}[ax]{1,2,17,7}}}
  \proofstepwidestar[1,2,7,16]{\exists u\in A\,(b\star s=a\star u)}{\rEI{19,20}}
  \proofstepwidestar[1,2,6,7,16]{b\star s\in R_{A,\star}(a)}{%
    \FormulaRefAuto{SemigroupPrincipalRightIdealMembership}[thm]{%
      1,2,\rAI{9,\rOIb{21}}}}
  \proofstepwidestar[1,2,6,7,15]{b\star s\in R_{A,\star}(a)}{\rEE{15,16,22}}
  \proofstepwidestar[1,2,6,7]{b\star s\in R_{A,\star}(a)}{%
    \rOE{10,11,14,15,23}}
  \proofstepwidestar[1,2]{%
    \forall b\in R_{A,\star}(a)\,\forall s\in A\,
      (b\star s\in R_{A,\star}(a))}{%
    \rUI{\rRI{6,\rUI{\rRI{7,24}}}}}
  \proofstepwidestar[1,2]{\mathsf{RId}(R_{A,\star}(a);A,\star)}{%
    \FormulaRefAuto{SemigroupRightIdealDef}[def]{1,\rAI{3,\rAI{5,25}}}}
\end{tabproofwide}

\FormulaThmDeltaK[Minimalität des Linkshauptideals]{%
  \begin{aligned}[t]
    &\SemigroupAlg{A}{\star}\dsep a\in A\dsep
      \mathsf{LId}(I;A,\star)\dsep a\in I\vdash{}\\[-2pt]
    &L_{A,\star}(a)\subseteq I
  \end{aligned}
}{SemigroupPrincipalLeftIdealMinimal}{%
  \DeltaRow{Mengen}{A\dsep I\dsep a\dsep b\dsep x}
  \DeltaRow{Funktionen}{\star}
}
\begin{tabproofwide}
  \proofstepwidestar[1]{\SemigroupAlg{A}{\star}}{\rA}
  \proofstepwidestar[2]{a\in A}{\rA}
  \proofstepwidestar[3]{\mathsf{LId}(I;A,\star)}{\rA}
  \proofstepwidestar[4]{a\in I}{\rA}
  \proofstepwidestar[5]{b\in L_{A,\star}(a)}{\rA}
  \proofstepwidestar[1,2,5]{b=a\lor\exists x\in A\,(b=x\star a)}{%
    \rAEb{\FormulaRefAuto{SemigroupPrincipalLeftIdealMembership}[thm]{1,2,5}}}
  \proofstepwidestar[7]{b=a}{\rA}
  \proofstepwidestar[4,7]{b\in I}{\rIE{\FormulaRefAuto{a=b\vdash b=a}[thm]{7},4}}
  \proofstepwidestar[9]{\exists x\in A\,(b=x\star a)}{\rA}
  \proofstepwidestar[10]{x\in A\land b=x\star a}{\rA}
  \proofstepwidestar[10]{x\in A}{\rAEa{10}}
  \proofstepwidestar[10]{b=x\star a}{\rAEb{10}}
  \proofstepwidestar[1,3]{%
    \forall u\in A\,\forall v\in I\,(u\star v\in I)}{%
    \rAEb{\rAEb{\FormulaRefAuto{SemigroupLeftIdealDef}[def]{1,3}}}}
  \proofstepwidestar[1,3,4,10]{x\star a\in I}{\rRE{4,\rUE{\rRE{11,\rUE{13}}}}}
  \proofstepwidestar[1,3,4,10]{b\in I}{\rIE{12,14}}
  \proofstepwidestar[1,3,4,9]{b\in I}{\rEE{9,10,15}}
  \proofstepwidestar[1,2,3,4,5]{b\in I}{\rOE{6,7,8,9,16}}
  \proofstepwidestar[1,2,3,4]{L_{A,\star}(a)\subseteq I}{%
    \FormulaRefAuto{A\subseteq B\coloneq
      \forall x\,(x\in A\rightarrow x\in B)}[def]{\rUI{\rRI{5,17}}}}
\end{tabproofwide}

\FormulaThmDeltaK[Minimalität des Rechtshauptideals]{%
  \begin{aligned}[t]
    &\SemigroupAlg{A}{\star}\dsep a\in A\dsep
      \mathsf{RId}(I;A,\star)\dsep a\in I\vdash{}\\[-2pt]
    &R_{A,\star}(a)\subseteq I
  \end{aligned}
}{SemigroupPrincipalRightIdealMinimal}{%
  \DeltaRow{Mengen}{A\dsep I\dsep a\dsep b\dsep y}
  \DeltaRow{Funktionen}{\star}
}
\begin{tabproofwide}
  \proofstepwidestar[1]{\SemigroupAlg{A}{\star}}{\rA}
  \proofstepwidestar[2]{a\in A}{\rA}
  \proofstepwidestar[3]{\mathsf{RId}(I;A,\star)}{\rA}
  \proofstepwidestar[4]{a\in I}{\rA}
  \proofstepwidestar[5]{b\in R_{A,\star}(a)}{\rA}
  \proofstepwidestar[1,2,5]{b=a\lor\exists y\in A\,(b=a\star y)}{%
    \rAEb{\FormulaRefAuto{SemigroupPrincipalRightIdealMembership}[thm]{1,2,5}}}
  \proofstepwidestar[7]{b=a}{\rA}
  \proofstepwidestar[4,7]{b\in I}{\rIE{\FormulaRefAuto{a=b\vdash b=a}[thm]{7},4}}
  \proofstepwidestar[9]{\exists y\in A\,(b=a\star y)}{\rA}
  \proofstepwidestar[10]{y\in A\land b=a\star y}{\rA}
  \proofstepwidestar[10]{y\in A}{\rAEa{10}}
  \proofstepwidestar[10]{b=a\star y}{\rAEb{10}}
  \proofstepwidestar[1,3]{%
    \forall u\in I\,\forall v\in A\,(u\star v\in I)}{%
    \rAEb{\rAEb{\FormulaRefAuto{SemigroupRightIdealDef}[def]{1,3}}}}
  \proofstepwidestar[1,3,4,10]{a\star y\in I}{\rRE{11,\rUE{\rRE{4,\rUE{13}}}}}
  \proofstepwidestar[1,3,4,10]{b\in I}{\rIE{12,14}}
  \proofstepwidestar[1,3,4,9]{b\in I}{\rEE{9,10,15}}
  \proofstepwidestar[1,2,3,4,5]{b\in I}{\rOE{6,7,8,9,16}}
  \proofstepwidestar[1,2,3,4]{R_{A,\star}(a)\subseteq I}{%
    \FormulaRefAuto{A\subseteq B\coloneq
      \forall x\,(x\in A\rightarrow x\in B)}[def]{\rUI{\rRI{5,17}}}}
\end{tabproofwide}

\FormulaThmDeltaK[Das zweiseitige Hauptideal ist ein Ideal]{%
  \SemigroupAlg{A}{\star}\dsep a\in A\vdash
  \mathsf{Id}(J_{A,\star}(a);A,\star)%
}{SemigroupPrincipalTwoSidedIdealIsIdeal}{%
  \DeltaRow{Mengen}{A\dsep a\dsep b\dsep s\dsep x\dsep y}
  \DeltaRow{Funktionen}{\star}
}
Nach dem Elementkriterium besitzt jedes Element des Hauptideals eine von
vier Formen: Es ist der Erzeuger selbst, hat nur einen linken oder rechten
Faktor oder hat beide Faktoren.  Multiplikation von links beziehungsweise
rechts erh\"alt jede dieser Formen; die beiden Absorptionseigenschaften
werden deshalb getrennt bewiesen und erst im Schluss zusammengef\"uhrt.

\begin{tabproofsplitwide}
  \proofpartwideindR[Linksidealeigenschaft]{%
    \SemigroupAlg{A}{\star}\dsep a\in A\vdash
    \mathsf{LId}(J_{A,\star}(a);A,\star)
  }{%
    \SemigroupAlg{A}{\star}\dsep a\in A\vdash
    \mathsf{LId}(J_{A,\star}(a);A,\star)
  }[SemigroupPrincipalTwoSidedIdealLeftPart]
    \proofstepwidestar[1]{\SemigroupAlg{A}{\star}}{\rA}
    \proofstepwidestar[2]{a\in A}{\rA}
    \proofstepwidestar[]{J_{A,\star}(a)\subseteq A}{%
      \FormulaRefAuto{\{x\in A\mid P(x)\}\subseteq A}[thm]}
    \proofstepwidestar[1,2]{a\in J_{A,\star}(a)}{%
      \rAEb{\rAEb{\FormulaRefAuto{SemigroupPrincipalIdealsContainGenerator}[thm]{1,2}}}}
    \proofstepwidestar[1,2]{J_{A,\star}(a)\neq\varnothing}{%
      \FormulaRefAuto{a\in A\vdash A\neq\varnothing}[thm]{4}}
    \proofstepwidestar[6]{s\in A}{\rA}
    \proofstepwidestar[7]{b\in J_{A,\star}(a)}{\rA}
    \proofstepwidestar[1,2,7]{b\in A}{%
      \rAEa{\FormulaRefAuto{SemigroupPrincipalTwoSidedIdealMembership}[thm]{1,2,7}}}
    \proofstepwidestar[1,2,6,7]{s\star b\in A}{%
      \FormulaRefAuto{SemigroupClosureAxiom}[ax]{1,6,8}}
    \proofstepwidestarbreak[1,2,7]{%
      \begin{aligned}[t]
        b=a\lor{}&\exists x\in A\,(b=x\star a)\\[-2pt]
        &{}\lor\exists y\in A\,(b=a\star y)\\[-2pt]
        &{}\lor\exists x,y\in A\,(b=(x\star a)\star y)
      \end{aligned}}{%
      \rAEb{\FormulaRefAuto{SemigroupPrincipalTwoSidedIdealMembership}[thm]{1,2,7}}}
    \proofstepwidestar[11]{b=a}{\rA}
    \proofstepwidestar[6,11]{\exists u\in A\,(s\star b=u\star a)}{%
      \rEI{6,\rIE{11}}}
    \proofstepwidestar[1,2,6,7,11]{s\star b\in J_{A,\star}(a)}{%
      \FormulaRefAuto{SemigroupPrincipalTwoSidedIdealMembership}[thm]{%
        1,2,\rAI{9,\rOIb{12}}}}
    \proofstepwidestar[14]{\exists x\in A\,(b=x\star a)}{\rA}
    \proofstepwidestar[15]{x\in A\land b=x\star a}{\rA}
    \proofstepwidestar[15]{x\in A}{\rAEa{15}}
    \proofstepwidestar[15]{b=x\star a}{\rAEb{15}}
    \proofstepwidestar[1,6,15]{s\star x\in A}{%
      \FormulaRefAuto{SemigroupClosureAxiom}[ax]{1,6,16}}
    \proofstepwidestar[1,2,6,15]{s\star b=(s\star x)\star a}{%
      \rIE{17,\FormulaRefAuto{SemigroupAssociativityAxiom}[ax]{1,6,16,2}}}
    \proofstepwidestar[1,2,6,15]{\exists u\in A\,(s\star b=u\star a)}{\rEI{18,19}}
    \proofstepwidestar[1,2,6,7,15]{s\star b\in J_{A,\star}(a)}{%
      \FormulaRefAuto{SemigroupPrincipalTwoSidedIdealMembership}[thm]{%
        1,2,\rAI{9,\rOIb{20}}}}
    \proofstepwidestar[22]{\exists y\in A\,(b=a\star y)}{\rA}
    \proofstepwidestar[23]{y\in A\land b=a\star y}{\rA}
    \proofstepwidestar[23]{y\in A}{\rAEa{23}}
    \proofstepwidestar[23]{b=a\star y}{\rAEb{23}}
    \proofstepwidestar[1,2,6,23]{s\star b=(s\star a)\star y}{%
      \rIE{25,\FormulaRefAuto{SemigroupAssociativityAxiom}[ax]{1,6,2,24}}}
    \proofstepwidestar[1,2,6,23]{%
      \exists u,v\in A\,(s\star b=(u\star a)\star v)}{\rEI{6,\rEI{24,26}}}
    \proofstepwidestarbreak[1,2,6,7,23]{s\star b\in J_{A,\star}(a)}{%
      \FormulaRefAuto{SemigroupPrincipalTwoSidedIdealMembership}[thm]{%
        1,2,\rAI{9,\rOIb{\rOIb{\rOIb{27}}}}}}
    \proofstepwidestar[29]{\exists x,y\in A\,(b=(x\star a)\star y)}{\rA}
    \proofstepwidestar[30]{x\in A\land(y\in A\land b=(x\star a)\star y)}{\rA}
    \proofstepwidestar[30]{x\in A}{\rAEa{30}}
    \proofstepwidestar[30]{y\in A}{\rAEa{\rAEb{30}}}
    \proofstepwidestar[30]{b=(x\star a)\star y}{\rAEb{\rAEb{30}}}
    \proofstepwidestar[1,6,30]{s\star x\in A}{%
      \FormulaRefAuto{SemigroupClosureAxiom}[ax]{1,6,31}}
    \proofstepwidestarbreak[1,2,6,30]{%
      s\star b=((s\star x)\star a)\star y}{%
      \rIE{33,\FormulaRefAuto{a = b,\, b = c \vdash a = c}[thm]{%
        \FormulaRefAuto{SemigroupAssociativityAxiom}[ax]{%
          1,6,\FormulaRefAuto{SemigroupClosureAxiom}[ax]{1,31,2},32},
        \rIE{\FormulaRefAuto{SemigroupAssociativityAxiom}[ax]{1,6,31,2}}}}}
    \proofstepwidestar[1,2,6,30]{%
      \exists u,v\in A\,(s\star b=(u\star a)\star v)}{\rEI{34,\rEI{32,35}}}
    \proofstepwidestarbreak[1,2,6,7,30]{s\star b\in J_{A,\star}(a)}{%
      \FormulaRefAuto{SemigroupPrincipalTwoSidedIdealMembership}[thm]{%
        1,2,\rAI{9,\rOIb{\rOIb{\rOIb{36}}}}}}
    \proofstepwidestar[1,2,6,7,29]{s\star b\in J_{A,\star}(a)}{\rEE{29,30,37}}
    \proofstepwidestar[1,2,6,7,22]{s\star b\in J_{A,\star}(a)}{\rEE{22,23,28}}
    \proofstepwidestar[1,2,6,7,14]{s\star b\in J_{A,\star}(a)}{\rEE{14,15,21}}
    \proofstepwidestar[1,2,6,7]{s\star b\in J_{A,\star}(a)}{%
      \rOE{10,11,13,14,39,22,38,29,37}}
    \proofstepwidestar[1,2]{%
      \forall s\in A\,\forall b\in J_{A,\star}(a)\,
      (s\star b\in J_{A,\star}(a))}{%
      \rUI{\rRI{6,\rUI{\rRI{7,40}}}}}
    \proofstepwidestar[1,2]{\mathsf{LId}(J_{A,\star}(a);A,\star)}{%
      \FormulaRefAuto{SemigroupLeftIdealDef}[def]{1,\rAI{3,\rAI{5,41}}}}
  \closeproofpartwideindR

  \proofpartwideindR[Rechtsidealeigenschaft]{%
    \SemigroupAlg{A}{\star}\dsep a\in A\vdash
    \mathsf{RId}(J_{A,\star}(a);A,\star)
  }{%
    \SemigroupAlg{A}{\star}\dsep a\in A\vdash
    \mathsf{RId}(J_{A,\star}(a);A,\star)
  }[SemigroupPrincipalTwoSidedIdealRightPart]
    \proofstepwidestar[1]{\SemigroupAlg{A}{\star}}{\rA}
    \proofstepwidestar[2]{a\in A}{\rA}
    \proofstepwidestar[]{J_{A,\star}(a)\subseteq A}{%
      \FormulaRefAuto{\{x\in A\mid P(x)\}\subseteq A}[thm]}
    \proofstepwidestar[1,2]{a\in J_{A,\star}(a)}{%
      \rAEb{\rAEb{\FormulaRefAuto{SemigroupPrincipalIdealsContainGenerator}[thm]{1,2}}}}
    \proofstepwidestar[1,2]{J_{A,\star}(a)\neq\varnothing}{%
      \FormulaRefAuto{a\in A\vdash A\neq\varnothing}[thm]{4}}
    \proofstepwidestar[6]{b\in J_{A,\star}(a)}{\rA}
    \proofstepwidestar[7]{s\in A}{\rA}
    \proofstepwidestar[1,2,6]{b\in A}{%
      \rAEa{\FormulaRefAuto{SemigroupPrincipalTwoSidedIdealMembership}[thm]{1,2,6}}}
    \proofstepwidestar[1,2,6,7]{b\star s\in A}{%
      \FormulaRefAuto{SemigroupClosureAxiom}[ax]{1,8,7}}
    \proofstepwidestarbreak[1,2,6]{%
      \begin{aligned}[t]
        b=a\lor{}&\exists x\in A\,(b=x\star a)\\[-2pt]
        &{}\lor\exists y\in A\,(b=a\star y)\\[-2pt]
        &{}\lor\exists x,y\in A\,(b=(x\star a)\star y)
      \end{aligned}}{%
      \rAEb{\FormulaRefAuto{SemigroupPrincipalTwoSidedIdealMembership}[thm]{1,2,6}}}
    \proofstepwidestar[11]{b=a}{\rA}
    \proofstepwidestar[7,11]{\exists v\in A\,(b\star s=a\star v)}{\rEI{7,\rIE{11}}}
    \proofstepwidestarbreak[1,2,6,7,11]{b\star s\in J_{A,\star}(a)}{%
      \FormulaRefAuto{SemigroupPrincipalTwoSidedIdealMembership}[thm]{%
        1,2,\rAI{9,\rOIb{\rOIb{12}}}}}
    \proofstepwidestar[14]{\exists x\in A\,(b=x\star a)}{\rA}
    \proofstepwidestar[15]{x\in A\land b=x\star a}{\rA}
    \proofstepwidestar[15]{x\in A}{\rAEa{15}}
    \proofstepwidestar[15]{b=x\star a}{\rAEb{15}}
    \proofstepwidestar[1,2,7,15]{b\star s=(x\star a)\star s}{\rIE{17}}
    \proofstepwidestar[1,2,7,15]{%
      \exists u,v\in A\,(b\star s=(u\star a)\star v)}{\rEI{16,\rEI{7,18}}}
    \proofstepwidestarbreak[1,2,6,7,15]{b\star s\in J_{A,\star}(a)}{%
      \FormulaRefAuto{SemigroupPrincipalTwoSidedIdealMembership}[thm]{%
        1,2,\rAI{9,\rOIb{\rOIb{\rOIb{19}}}}}}
    \proofstepwidestar[21]{\exists y\in A\,(b=a\star y)}{\rA}
    \proofstepwidestar[22]{y\in A\land b=a\star y}{\rA}
    \proofstepwidestar[22]{y\in A}{\rAEa{22}}
    \proofstepwidestar[22]{b=a\star y}{\rAEb{22}}
    \proofstepwidestar[1,7,22]{y\star s\in A}{%
      \FormulaRefAuto{SemigroupClosureAxiom}[ax]{1,23,7}}
    \proofstepwidestar[1,2,7,22]{b\star s=a\star(y\star s)}{%
      \rIE{24,\FormulaRefAuto{SemigroupAssociativityAxiom}[ax]{1,2,23,7}}}
    \proofstepwidestar[1,2,7,22]{\exists v\in A\,(b\star s=a\star v)}{\rEI{25,26}}
    \proofstepwidestarbreak[1,2,6,7,22]{b\star s\in J_{A,\star}(a)}{%
      \FormulaRefAuto{SemigroupPrincipalTwoSidedIdealMembership}[thm]{%
        1,2,\rAI{9,\rOIb{\rOIb{27}}}}}
    \proofstepwidestar[29]{\exists x,y\in A\,(b=(x\star a)\star y)}{\rA}
    \proofstepwidestar[30]{x\in A\land(y\in A\land b=(x\star a)\star y)}{\rA}
    \proofstepwidestar[30]{x\in A}{\rAEa{30}}
    \proofstepwidestar[30]{y\in A}{\rAEa{\rAEb{30}}}
    \proofstepwidestar[30]{b=(x\star a)\star y}{\rAEb{\rAEb{30}}}
    \proofstepwidestar[1,7,30]{y\star s\in A}{%
      \FormulaRefAuto{SemigroupClosureAxiom}[ax]{1,32,7}}
    \proofstepwidestarbreak[1,2,7,30]{b\star s=(x\star a)\star(y\star s)}{%
      \rIE{33,\FormulaRefAuto{SemigroupAssociativityAxiom}[ax]{%
        1,\FormulaRefAuto{SemigroupClosureAxiom}[ax]{1,31,2},32,7}}}
    \proofstepwidestar[1,2,7,30]{%
      \exists u,v\in A\,(b\star s=(u\star a)\star v)}{\rEI{31,\rEI{34,35}}}
    \proofstepwidestarbreak[1,2,6,7,30]{b\star s\in J_{A,\star}(a)}{%
      \FormulaRefAuto{SemigroupPrincipalTwoSidedIdealMembership}[thm]{%
        1,2,\rAI{9,\rOIb{\rOIb{\rOIb{36}}}}}}
    \proofstepwidestar[1,2,6,7,29]{b\star s\in J_{A,\star}(a)}{\rEE{29,30,37}}
    \proofstepwidestar[1,2,6,7,21]{b\star s\in J_{A,\star}(a)}{\rEE{21,22,28}}
    \proofstepwidestar[1,2,6,7,14]{b\star s\in J_{A,\star}(a)}{\rEE{14,15,20}}
    \proofstepwidestar[1,2,6,7]{b\star s\in J_{A,\star}(a)}{%
      \rOE{10,11,13,14,39,21,38,29,37}}
    \proofstepwidestar[1,2]{%
      \forall b\in J_{A,\star}(a)\,\forall s\in A\,
      (b\star s\in J_{A,\star}(a))}{%
      \rUI{\rRI{6,\rUI{\rRI{7,40}}}}}
    \proofstepwidestar[1,2]{\mathsf{RId}(J_{A,\star}(a);A,\star)}{%
      \FormulaRefAuto{SemigroupRightIdealDef}[def]{1,\rAI{3,\rAI{5,41}}}}
  \closeproofpartwideindR

  \proofpartwide{Schluss}
    \proofstepwidestar[1]{\SemigroupAlg{A}{\star}}{\rA}
    \proofstepwidestar[2]{a\in A}{\rA}
    \proofstepwidestar[1,2]{\mathsf{LId}(J_{A,\star}(a);A,\star)}{%
      \FormulaRefAuto{SemigroupPrincipalTwoSidedIdealLeftPart}[thm]{1,2}}
    \proofstepwidestar[1,2]{\mathsf{RId}(J_{A,\star}(a);A,\star)}{%
      \FormulaRefAuto{SemigroupPrincipalTwoSidedIdealRightPart}[thm]{1,2}}
    \proofstepwidestar[1,2]{\mathsf{Id}(J_{A,\star}(a);A,\star)}{%
      \FormulaRefAuto{SemigroupIdealDef}[def]{1,\rAI{3,4}}}
  \closeproofpartwide
\end{tabproofsplitwide}

\FormulaThmDeltaK[Minimalit\"at des zweiseitigen Hauptideals]{%
  \begin{aligned}[t]
    &\SemigroupAlg{A}{\star}\dsep a\in A\dsep
      \mathsf{Id}(I;A,\star)\dsep a\in I\vdash{}\\[-2pt]
    &J_{A,\star}(a)\subseteq I
  \end{aligned}
}{SemigroupPrincipalTwoSidedIdealMinimal}{%
  \DeltaRow{Mengen}{A\dsep I\dsep a\dsep b\dsep x\dsep y}
  \DeltaRow{Funktionen}{\star}
}
Das Elementkriterium liefert vier und nur vier Formen.  In jedem Fall folgt
die Zugeh\"origkeit zu (I) unmittelbar aus (a\in I) und der linken
beziehungsweise rechten Idealeigenschaft.

\begin{tabproofsplitwide}
  \proofpartwideindR[Erzeugerfall]{%
    \SemigroupAlg{A}{\star}\dsep a\in A\dsep\mathsf{Id}(I;A,\star)\dsep
    a\in I\dsep b=a\vdash b\in I
  }{%
    \SemigroupAlg{A}{\star}\dsep a\in A\dsep\mathsf{Id}(I;A,\star)\dsep
    a\in I\dsep b=a\vdash b\in I
  }[SemigroupPrincipalTwoSidedIdealMinimalGeneratorCase]
    \proofstepwidestar[1]{\SemigroupAlg{A}{\star}}{\rA}
    \proofstepwidestar[2]{a\in A}{\rA}
    \proofstepwidestar[3]{\mathsf{Id}(I;A,\star)}{\rA}
    \proofstepwidestar[4]{a\in I}{\rA}
    \proofstepwidestar[5]{b=a}{\rA}
    \proofstepwidestar[4,5]{b\in I}{%
      \rIE{\FormulaRefAuto{a=b\vdash b=a}[thm]{5},4}}
  \closeproofpartwideindR

  \proofpartwideindR[Fall eines linken Faktors]{%
    \begin{aligned}[t]
      &\SemigroupAlg{A}{\star}\dsep a\in A\dsep
        \mathsf{Id}(I;A,\star)\dsep a\in I\dsep{}\\[-2pt]
      &\exists x\in A\,(b=x\star a)\vdash b\in I
    \end{aligned}
  }{%
    \SemigroupAlg{A}{\star}\dsep a\in A\dsep
    \mathsf{Id}(I;A,\star)\dsep a\in I\dsep
    \exists x\in A\,(b=x\star a)\vdash b\in I
  }[SemigroupPrincipalTwoSidedIdealMinimalLeftFactorCase]
    \proofstepwidestar[1]{\SemigroupAlg{A}{\star}}{\rA}
    \proofstepwidestar[2]{a\in A}{\rA}
    \proofstepwidestar[3]{\mathsf{Id}(I;A,\star)}{\rA}
    \proofstepwidestar[4]{a\in I}{\rA}
    \proofstepwidestar[5]{\exists x\in A\,(b=x\star a)}{\rA}
    \proofstepwidestar[6]{x\in A\land b=x\star a}{\rA}
    \proofstepwidestar[6]{x\in A}{\rAEa{6}}
    \proofstepwidestar[6]{b=x\star a}{\rAEb{6}}
    \proofstepwidestar[1,3,4,6]{x\star a\in I}{%
      \FormulaRefAuto{SemigroupIdealLeftAbsorption}[thm]{1,3,7,4}}
    \proofstepwidestar[1,3,4,6]{b\in I}{\rIE{8,9}}
    \proofstepwidestar[1,3,4,5]{b\in I}{\rEE{5,6,10}}
  \closeproofpartwideindR

  \proofpartwideindR[Fall eines rechten Faktors]{%
    \begin{aligned}[t]
      &\SemigroupAlg{A}{\star}\dsep a\in A\dsep
        \mathsf{Id}(I;A,\star)\dsep a\in I\dsep{}\\[-2pt]
      &\exists y\in A\,(b=a\star y)\vdash b\in I
    \end{aligned}
  }{%
    \SemigroupAlg{A}{\star}\dsep a\in A\dsep
    \mathsf{Id}(I;A,\star)\dsep a\in I\dsep
    \exists y\in A\,(b=a\star y)\vdash b\in I
  }[SemigroupPrincipalTwoSidedIdealMinimalRightFactorCase]
    \proofstepwidestar[1]{\SemigroupAlg{A}{\star}}{\rA}
    \proofstepwidestar[2]{a\in A}{\rA}
    \proofstepwidestar[3]{\mathsf{Id}(I;A,\star)}{\rA}
    \proofstepwidestar[4]{a\in I}{\rA}
    \proofstepwidestar[5]{\exists y\in A\,(b=a\star y)}{\rA}
    \proofstepwidestar[6]{y\in A\land b=a\star y}{\rA}
    \proofstepwidestar[6]{y\in A}{\rAEa{6}}
    \proofstepwidestar[6]{b=a\star y}{\rAEb{6}}
    \proofstepwidestar[1,3,4,6]{a\star y\in I}{%
      \FormulaRefAuto{SemigroupIdealRightAbsorption}[thm]{1,3,4,7}}
    \proofstepwidestar[1,3,4,6]{b\in I}{\rIE{8,9}}
    \proofstepwidestar[1,3,4,5]{b\in I}{\rEE{5,6,10}}
  \closeproofpartwideindR

  \proofpartwideindR[Fall zweier Faktoren]{%
    \begin{aligned}[t]
      &\SemigroupAlg{A}{\star}\dsep a\in A\dsep
        \mathsf{Id}(I;A,\star)\dsep a\in I\dsep{}\\[-2pt]
      &\exists x,y\in A\,(b=(x\star a)\star y)\vdash b\in I
    \end{aligned}
  }{%
    \begin{aligned}[t]
      &\SemigroupAlg{A}{\star}\dsep a\in A\dsep
        \mathsf{Id}(I;A,\star)\dsep a\in I\dsep{}\\[-2pt]
      &\exists x,y\in A\,(b=(x\star a)\star y)\vdash b\in I
    \end{aligned}
  }[SemigroupPrincipalTwoSidedIdealMinimalTwoFactorCase]
    \proofstepwidestar[1]{\SemigroupAlg{A}{\star}}{\rA}
    \proofstepwidestar[2]{a\in A}{\rA}
    \proofstepwidestar[3]{\mathsf{Id}(I;A,\star)}{\rA}
    \proofstepwidestar[4]{a\in I}{\rA}
    \proofstepwidestar[5]{\exists x,y\in A\,(b=(x\star a)\star y)}{\rA}
    \proofstepwidestar[6]{%
      x\in A\land\bigl(y\in A\land b=(x\star a)\star y\bigr)}{\rA}
    \proofstepwidestar[6]{x\in A}{\rAEa{6}}
    \proofstepwidestar[6]{y\in A}{\rAEa{\rAEb{6}}}
    \proofstepwidestar[6]{b=(x\star a)\star y}{\rAEb{\rAEb{6}}}
    \proofstepwidestar[1,3,4,6]{x\star a\in I}{%
      \FormulaRefAuto{SemigroupIdealLeftAbsorption}[thm]{1,3,7,4}}
    \proofstepwidestar[1,3,4,6]{(x\star a)\star y\in I}{%
      \FormulaRefAuto{SemigroupIdealRightAbsorption}[thm]{1,3,10,8}}
    \proofstepwidestar[1,3,4,6]{b\in I}{\rIE{9,11}}
    \proofstepwidestar[1,3,4,5]{b\in I}{\rEE{5,6,12}}
  \closeproofpartwideindR

  \proofpartwide{Fallunterscheidung und Schluss}
    \proofstepwidestar[1]{\SemigroupAlg{A}{\star}}{\rA}
    \proofstepwidestar[2]{a\in A}{\rA}
    \proofstepwidestar[3]{\mathsf{Id}(I;A,\star)}{\rA}
    \proofstepwidestar[4]{a\in I}{\rA}
    \proofstepwidestar[5]{b\in J_{A,\star}(a)}{\rA}
    \proofstepwidestarbreak[1,2,5]{%
      \begin{aligned}[t]
        b=a\lor{}&\exists x\in A\,(b=x\star a)\\[-2pt]
        &{}\lor\exists y\in A\,(b=a\star y)\\[-2pt]
        &{}\lor\exists x,y\in A\,(b=(x\star a)\star y)
      \end{aligned}}{%
      \rAEb{\FormulaRefAuto{SemigroupPrincipalTwoSidedIdealMembership}[thm]{1,2,5}}}
    \proofstepwidestar[7]{b=a}{\rA}
    \proofstepwidestar[1,2,3,4,7]{b\in I}{%
      \FormulaRefAuto{SemigroupPrincipalTwoSidedIdealMinimalGeneratorCase}[thm]{%
        1,2,3,4,7}}
    \proofstepwidestar[9]{\exists x\in A\,(b=x\star a)}{\rA}
    \proofstepwidestar[1,2,3,4,9]{b\in I}{%
      \FormulaRefAuto{SemigroupPrincipalTwoSidedIdealMinimalLeftFactorCase}[thm]{%
        1,2,3,4,9}}
    \proofstepwidestar[11]{\exists y\in A\,(b=a\star y)}{\rA}
    \proofstepwidestar[1,2,3,4,11]{b\in I}{%
      \FormulaRefAuto{SemigroupPrincipalTwoSidedIdealMinimalRightFactorCase}[thm]{%
        1,2,3,4,11}}
    \proofstepwidestar[13]{\exists x,y\in A\,(b=(x\star a)\star y)}{\rA}
    \proofstepwidestar[1,2,3,4,13]{b\in I}{%
      \FormulaRefAuto{SemigroupPrincipalTwoSidedIdealMinimalTwoFactorCase}[thm]{%
        1,2,3,4,13}}
    \proofstepwidestar[1,2,3,4,5]{b\in I}{%
      \rOE{6,7,8,9,10,11,12,13,14}}
    \proofstepwidestar[1,2,3,4]{J_{A,\star}(a)\subseteq I}{%
      \FormulaRefAuto{A\subseteq B\coloneq
        \forall x\,(x\in A\rightarrow x\in B)}[def]{\rUI{\rRI{5,15}}}}
  \closeproofpartwide
\end{tabproofsplitwide}

\section{Teilbarkeit in Halbgruppen}

Ohne neutrales Element wäre die Bedingung \(b=a\star c\) nicht notwendig
reflexiv. Deshalb nehmen die folgenden Definitionen den Fall \(a=b\)
ausdrücklich auf. Die zuvor definierten Hauptideale erfassen genau diese
fehlenden Faktoren; eine unbestimmte Schreibweise wie \(A^1\) wird daher
nicht benötigt.

\FormulaDefDeltaK[Linksteilbarkeit in einer Halbgruppe]{%
  \SemigroupLeftDivides{A,\star}{a}{b}%
}{SemigroupLeftDividesDef}{%
  \DeltaRow{Mengen}{A\dsep a\dsep b\dsep c}
  \DeltaRow{Funktionen}{\star}
  \DeltaPrem{Halbgruppe}{\SemigroupAlg{A}{\star}}
  \DeltaRow{Elemente}{a,b\in A}
  \DeltaRow{Linksteilbarkeit}{%
    b=a\lor\exists c\in A\,(b=a\star c)}
  \DeltaRow{\textbf{Neues Symbol}}{}
  \DeltaPrem{Linksteilbarkeit}{%
    \SemigroupLeftDivides{A,\star}{a}{b}}
}

\FormulaDefDeltaK[Rechtsteilbarkeit in einer Halbgruppe]{%
  \SemigroupRightDivides{A,\star}{a}{b}%
}{SemigroupRightDividesDef}{%
  \DeltaRow{Mengen}{A\dsep a\dsep b\dsep c}
  \DeltaRow{Funktionen}{\star}
  \DeltaPrem{Halbgruppe}{\SemigroupAlg{A}{\star}}
  \DeltaRow{Elemente}{a,b\in A}
  \DeltaRow{Rechtsteilbarkeit}{%
    b=a\lor\exists c\in A\,(b=c\star a)}
  \DeltaRow{\textbf{Neues Symbol}}{}
  \DeltaPrem{Rechtsteilbarkeit}{%
    \SemigroupRightDivides{A,\star}{a}{b}}
}

\FormulaDefDeltaK[Zweiseitige Teilbarkeit in einer Halbgruppe]{%
  \SemigroupDivides{A,\star}{a}{b}%
}{SemigroupTwoSidedDividesDef}{%
  \DeltaRow{Mengen}{A\dsep a\dsep b\dsep x\dsep y}
  \DeltaRow{Funktionen}{\star}
  \DeltaPrem{Halbgruppe}{\SemigroupAlg{A}{\star}}
  \DeltaRow{Elemente}{a,b\in A}
  \DeltaRow{Zweiseitige Teilbarkeit}{%
    \begin{aligned}[t]
      &b=a\lor\exists x\in A\,(b=x\star a)\\[-2pt]
      &{}\lor\exists y\in A\,(b=a\star y)\\[-2pt]
      &{}\lor\exists x,y\in A\,
        \bigl(b=(x\star a)\star y\bigr)
    \end{aligned}}
  \DeltaRow{\textbf{Neues Symbol}}{}
  \DeltaPrem{Zweiseitige Teilbarkeit}{%
    \SemigroupDivides{A,\star}{a}{b}}
}

\FormulaThmDeltaK[Linksteilbarkeit als Mitgliedschaft im Rechtshauptideal]{%
  \begin{aligned}[t]
    &\SemigroupAlg{A}{\star}\dsep a,b\in A\vdash{}\\[-2pt]
    &\SemigroupLeftDivides{A,\star}{a}{b}
      \leftrightarrow b\in R_{A,\star}(a)
  \end{aligned}
}{SemigroupLeftDividesPrincipalRightIdeal}{%
  \DeltaRow{Mengen}{A\dsep a\dsep b\dsep c}
  \DeltaRow{Funktionen}{\star}
}
\begin{tabproofwide}
  \proofstepwidestar[1]{\SemigroupAlg{A}{\star}}{\rA}
  \proofstepwidestar[2]{a\in A}{\rA}
  \proofstepwidestar[3]{b\in A}{\rA}
  \proofstepwidestar[4]{\SemigroupLeftDivides{A,\star}{a}{b}}{\rA}
  \proofstepwidestar[1,2,3,4]{b=a\lor\exists c\in A\,(b=a\star c)}{%
    \FormulaRefAuto{SemigroupLeftDividesDef}[def]{1,2,3,4}}
  \proofstepwidestar[1,2,3,4]{b\in R_{A,\star}(a)}{%
    \FormulaRefAuto{SemigroupPrincipalRightIdealMembership}[thm]{%
      1,2,\rAI{3,5}}}
  \proofstepwidestar[1,2,3]{%
    \SemigroupLeftDivides{A,\star}{a}{b}\rightarrow
    b\in R_{A,\star}(a)}{\rRI{4,6}}
  \proofstepwidestar[8]{b\in R_{A,\star}(a)}{\rA}
  \proofstepwidestar[1,2,8]{b=a\lor\exists c\in A\,(b=a\star c)}{%
    \rAEb{\FormulaRefAuto{SemigroupPrincipalRightIdealMembership}[thm]{1,2,8}}}
  \proofstepwidestar[1,2,3,8]{\SemigroupLeftDivides{A,\star}{a}{b}}{%
    \FormulaRefAuto{SemigroupLeftDividesDef}[def]{1,2,3,9}}
  \proofstepwidestar[1,2,3]{%
    b\in R_{A,\star}(a)\rightarrow
    \SemigroupLeftDivides{A,\star}{a}{b}}{\rRI{8,10}}
  \proofstepwidestar[1,2,3]{%
    \SemigroupLeftDivides{A,\star}{a}{b}
    \leftrightarrow b\in R_{A,\star}(a)}{\rLRI{7,11}}
\end{tabproofwide}

\FormulaThmDeltaK[Rechtsteilbarkeit als Mitgliedschaft im Linkshauptideal]{%
  \begin{aligned}[t]
    &\SemigroupAlg{A}{\star}\dsep a,b\in A\vdash{}\\[-2pt]
    &\SemigroupRightDivides{A,\star}{a}{b}
      \leftrightarrow b\in L_{A,\star}(a)
  \end{aligned}
}{SemigroupRightDividesPrincipalLeftIdeal}{%
  \DeltaRow{Mengen}{A\dsep a\dsep b\dsep c}
  \DeltaRow{Funktionen}{\star}
}
\begin{tabproofwide}
  \proofstepwidestar[1]{\SemigroupAlg{A}{\star}}{\rA}
  \proofstepwidestar[2]{a\in A}{\rA}
  \proofstepwidestar[3]{b\in A}{\rA}
  \proofstepwidestar[4]{\SemigroupRightDivides{A,\star}{a}{b}}{\rA}
  \proofstepwidestar[1,2,3,4]{b=a\lor\exists c\in A\,(b=c\star a)}{%
    \FormulaRefAuto{SemigroupRightDividesDef}[def]{1,2,3,4}}
  \proofstepwidestar[1,2,3,4]{b\in L_{A,\star}(a)}{%
    \FormulaRefAuto{SemigroupPrincipalLeftIdealMembership}[thm]{%
      1,2,\rAI{3,5}}}
  \proofstepwidestar[1,2,3]{%
    \SemigroupRightDivides{A,\star}{a}{b}\rightarrow
    b\in L_{A,\star}(a)}{\rRI{4,6}}
  \proofstepwidestar[8]{b\in L_{A,\star}(a)}{\rA}
  \proofstepwidestar[1,2,8]{b=a\lor\exists c\in A\,(b=c\star a)}{%
    \rAEb{\FormulaRefAuto{SemigroupPrincipalLeftIdealMembership}[thm]{1,2,8}}}
  \proofstepwidestar[1,2,3,8]{\SemigroupRightDivides{A,\star}{a}{b}}{%
    \FormulaRefAuto{SemigroupRightDividesDef}[def]{1,2,3,9}}
  \proofstepwidestar[1,2,3]{%
    b\in L_{A,\star}(a)\rightarrow
    \SemigroupRightDivides{A,\star}{a}{b}}{\rRI{8,10}}
  \proofstepwidestar[1,2,3]{%
    \SemigroupRightDivides{A,\star}{a}{b}
    \leftrightarrow b\in L_{A,\star}(a)}{\rLRI{7,11}}
\end{tabproofwide}

\FormulaThmDeltaK[Zweiseitige Teilbarkeit als Hauptidealmitgliedschaft]{%
  \begin{aligned}[t]
    &\SemigroupAlg{A}{\star}\dsep a,b\in A\vdash{}\\[-2pt]
    &\SemigroupDivides{A,\star}{a}{b}
      \leftrightarrow b\in J_{A,\star}(a)
  \end{aligned}
}{SemigroupTwoSidedDividesPrincipalIdeal}{%
  \DeltaRow{Mengen}{A\dsep a\dsep b\dsep x\dsep y}
  \DeltaRow{Funktionen}{\star}
}
\begin{tabproofwide}
  \proofstepwidestar[1]{\SemigroupAlg{A}{\star}}{\rA}
  \proofstepwidestar[2]{a\in A}{\rA}
  \proofstepwidestar[3]{b\in A}{\rA}
  \proofstepwidestar[4]{\SemigroupDivides{A,\star}{a}{b}}{\rA}
  \proofstepwidestarbreak[1,2,3,4]{%
    \begin{aligned}[t]
      b=a\lor{}&\exists x\in A\,(b=x\star a)\\[-2pt]
      &{}\lor\exists y\in A\,(b=a\star y)\\[-2pt]
      &{}\lor\exists x,y\in A\,(b=(x\star a)\star y)
    \end{aligned}}{%
    \FormulaRefAuto{SemigroupTwoSidedDividesDef}[def]{1,2,3,4}}
  \proofstepwidestar[1,2,3,4]{b\in J_{A,\star}(a)}{%
    \FormulaRefAuto{SemigroupPrincipalTwoSidedIdealMembership}[thm]{%
      1,2,\rAI{3,5}}}
  \proofstepwidestar[1,2,3]{%
    \SemigroupDivides{A,\star}{a}{b}\rightarrow
    b\in J_{A,\star}(a)}{\rRI{4,6}}
  \proofstepwidestar[8]{b\in J_{A,\star}(a)}{\rA}
  \proofstepwidestarbreak[1,2,8]{%
    \begin{aligned}[t]
      b=a\lor{}&\exists x\in A\,(b=x\star a)\\[-2pt]
      &{}\lor\exists y\in A\,(b=a\star y)\\[-2pt]
      &{}\lor\exists x,y\in A\,(b=(x\star a)\star y)
    \end{aligned}}{%
    \rAEb{\FormulaRefAuto{SemigroupPrincipalTwoSidedIdealMembership}[thm]{1,2,8}}}
  \proofstepwidestar[1,2,3,8]{\SemigroupDivides{A,\star}{a}{b}}{%
    \FormulaRefAuto{SemigroupTwoSidedDividesDef}[def]{1,2,3,9}}
  \proofstepwidestar[1,2,3]{%
    b\in J_{A,\star}(a)\rightarrow
    \SemigroupDivides{A,\star}{a}{b}}{\rRI{8,10}}
  \proofstepwidestar[1,2,3]{%
    \SemigroupDivides{A,\star}{a}{b}
    \leftrightarrow b\in J_{A,\star}(a)}{\rLRI{7,11}}
\end{tabproofwide}

\FormulaThmDeltaK[Reflexivität der drei Teilbarkeiten]{%
  \begin{aligned}[t]
    &\SemigroupAlg{A}{\star}\dsep a\in A\vdash{}\\[-2pt]
    &\SemigroupLeftDivides{A,\star}{a}{a}\land
      \SemigroupRightDivides{A,\star}{a}{a}\land
      \SemigroupDivides{A,\star}{a}{a}
  \end{aligned}
}{SemigroupDivisibilityReflexive}{%
  \DeltaRow{Mengen}{A\dsep a}
  \DeltaRow{Funktionen}{\star}
}
\begin{tabproofwide}
  \proofstepwidestar[1]{\SemigroupAlg{A}{\star}}{\rA}
  \proofstepwidestar[2]{a\in A}{\rA}
  \proofstepwidestar[1,2]{\SemigroupLeftDivides{A,\star}{a}{a}}{%
    \FormulaRefAuto{SemigroupLeftDividesDef}[def]{1,2,2,\rOIa{\rII}}}
  \proofstepwidestar[1,2]{\SemigroupRightDivides{A,\star}{a}{a}}{%
    \FormulaRefAuto{SemigroupRightDividesDef}[def]{1,2,2,\rOIa{\rII}}}
  \proofstepwidestar[1,2]{\SemigroupDivides{A,\star}{a}{a}}{%
    \FormulaRefAuto{SemigroupTwoSidedDividesDef}[def]{1,2,2,\rOIa{\rII}}}
  \proofstepwidestar[1,2]{%
    \SemigroupLeftDivides{A,\star}{a}{a}\land
    \SemigroupRightDivides{A,\star}{a}{a}\land
    \SemigroupDivides{A,\star}{a}{a}}{\rAI{3,\rAI{4,5}}}
\end{tabproofwide}

\FormulaThmDeltaK[Transitivität der Linksteilbarkeit]{%
  \begin{aligned}[t]
    &\SemigroupAlg{A}{\star}\dsep a,b,c\in A\dsep
      \SemigroupLeftDivides{A,\star}{a}{b}\dsep
      \SemigroupLeftDivides{A,\star}{b}{c}\vdash{}\\[-2pt]
    &\SemigroupLeftDivides{A,\star}{a}{c}
  \end{aligned}
}{SemigroupLeftDividesTransitive}{%
  \DeltaRow{Mengen}{A\dsep a\dsep b\dsep c}
  \DeltaRow{Funktionen}{\star}
}
\begin{tabproofwide}
  \proofstepwidestar[1]{\SemigroupAlg{A}{\star}}{\rA}
  \proofstepwidestar[2]{a\in A}{\rA}
  \proofstepwidestar[3]{b\in A}{\rA}
  \proofstepwidestar[4]{c\in A}{\rA}
  \proofstepwidestar[5]{\SemigroupLeftDivides{A,\star}{a}{b}}{\rA}
  \proofstepwidestar[6]{\SemigroupLeftDivides{A,\star}{b}{c}}{\rA}
  \proofstepwidestar[1,2,3,5]{b\in R_{A,\star}(a)}{%
    \FormulaRefAuto{SemigroupLeftDividesPrincipalRightIdeal}[thm]{1,2,3,5}}
  \proofstepwidestar[1,3,4,6]{c\in R_{A,\star}(b)}{%
    \FormulaRefAuto{SemigroupLeftDividesPrincipalRightIdeal}[thm]{1,3,4,6}}
  \proofstepwidestar[1,2]{\mathsf{RId}(R_{A,\star}(a);A,\star)}{%
    \FormulaRefAuto{SemigroupPrincipalRightIdealIsRightIdeal}[thm]{1,2}}
  \proofstepwidestar[1,2,3,5]{R_{A,\star}(b)\subseteq R_{A,\star}(a)}{%
    \FormulaRefAuto{SemigroupPrincipalRightIdealMinimal}[thm]{1,3,9,7}}
  \proofstepwidestar[1,2,3,4,5,6]{c\in R_{A,\star}(a)}{%
    \FormulaRefAuto{A\subseteq B,\,x\in A\vdash x\in B}[thm]{10,8}}
  \proofstepwidestar[1,2,3,4,5,6]{\SemigroupLeftDivides{A,\star}{a}{c}}{%
    \FormulaRefAuto{SemigroupLeftDividesPrincipalRightIdeal}[thm]{1,2,4,11}}
\end{tabproofwide}

\FormulaThmDeltaK[Transitivität der Rechtsteilbarkeit]{%
  \begin{aligned}[t]
    &\SemigroupAlg{A}{\star}\dsep a,b,c\in A\dsep
      \SemigroupRightDivides{A,\star}{a}{b}\dsep
      \SemigroupRightDivides{A,\star}{b}{c}\vdash{}\\[-2pt]
    &\SemigroupRightDivides{A,\star}{a}{c}
  \end{aligned}
}{SemigroupRightDividesTransitive}{%
  \DeltaRow{Mengen}{A\dsep a\dsep b\dsep c}
  \DeltaRow{Funktionen}{\star}
}
\begin{tabproofwide}
  \proofstepwidestar[1]{\SemigroupAlg{A}{\star}}{\rA}
  \proofstepwidestar[2]{a\in A}{\rA}
  \proofstepwidestar[3]{b\in A}{\rA}
  \proofstepwidestar[4]{c\in A}{\rA}
  \proofstepwidestar[5]{\SemigroupRightDivides{A,\star}{a}{b}}{\rA}
  \proofstepwidestar[6]{\SemigroupRightDivides{A,\star}{b}{c}}{\rA}
  \proofstepwidestar[1,2,3,5]{b\in L_{A,\star}(a)}{%
    \FormulaRefAuto{SemigroupRightDividesPrincipalLeftIdeal}[thm]{1,2,3,5}}
  \proofstepwidestar[1,3,4,6]{c\in L_{A,\star}(b)}{%
    \FormulaRefAuto{SemigroupRightDividesPrincipalLeftIdeal}[thm]{1,3,4,6}}
  \proofstepwidestar[1,2]{\mathsf{LId}(L_{A,\star}(a);A,\star)}{%
    \FormulaRefAuto{SemigroupPrincipalLeftIdealIsLeftIdeal}[thm]{1,2}}
  \proofstepwidestar[1,2,3,5]{L_{A,\star}(b)\subseteq L_{A,\star}(a)}{%
    \FormulaRefAuto{SemigroupPrincipalLeftIdealMinimal}[thm]{1,3,9,7}}
  \proofstepwidestar[1,2,3,4,5,6]{c\in L_{A,\star}(a)}{%
    \FormulaRefAuto{A\subseteq B,\,x\in A\vdash x\in B}[thm]{10,8}}
  \proofstepwidestar[1,2,3,4,5,6]{\SemigroupRightDivides{A,\star}{a}{c}}{%
    \FormulaRefAuto{SemigroupRightDividesPrincipalLeftIdeal}[thm]{1,2,4,11}}
\end{tabproofwide}

\FormulaThmDeltaK[Transitivität der zweiseitigen Teilbarkeit]{%
  \begin{aligned}[t]
    &\SemigroupAlg{A}{\star}\dsep a,b,c\in A\dsep
      \SemigroupDivides{A,\star}{a}{b}\dsep
      \SemigroupDivides{A,\star}{b}{c}\vdash{}\\[-2pt]
    &\SemigroupDivides{A,\star}{a}{c}
  \end{aligned}
}{SemigroupTwoSidedDividesTransitive}{%
  \DeltaRow{Mengen}{A\dsep a\dsep b\dsep c}
  \DeltaRow{Funktionen}{\star}
}
\begin{tabproofwide}
  \proofstepwidestar[1]{\SemigroupAlg{A}{\star}}{\rA}
  \proofstepwidestar[2]{a\in A}{\rA}
  \proofstepwidestar[3]{b\in A}{\rA}
  \proofstepwidestar[4]{c\in A}{\rA}
  \proofstepwidestar[5]{\SemigroupDivides{A,\star}{a}{b}}{\rA}
  \proofstepwidestar[6]{\SemigroupDivides{A,\star}{b}{c}}{\rA}
  \proofstepwidestar[1,2,3,5]{b\in J_{A,\star}(a)}{%
    \FormulaRefAuto{SemigroupTwoSidedDividesPrincipalIdeal}[thm]{1,2,3,5}}
  \proofstepwidestar[1,3,4,6]{c\in J_{A,\star}(b)}{%
    \FormulaRefAuto{SemigroupTwoSidedDividesPrincipalIdeal}[thm]{1,3,4,6}}
  \proofstepwidestar[1,2]{\mathsf{Id}(J_{A,\star}(a);A,\star)}{%
    \FormulaRefAuto{SemigroupPrincipalTwoSidedIdealIsIdeal}[thm]{1,2}}
  \proofstepwidestar[1,2,3,5]{J_{A,\star}(b)\subseteq J_{A,\star}(a)}{%
    \FormulaRefAuto{SemigroupPrincipalTwoSidedIdealMinimal}[thm]{1,3,9,7}}
  \proofstepwidestar[1,2,3,4,5,6]{c\in J_{A,\star}(a)}{%
    \FormulaRefAuto{A\subseteq B,\,x\in A\vdash x\in B}[thm]{10,8}}
  \proofstepwidestar[1,2,3,4,5,6]{\SemigroupDivides{A,\star}{a}{c}}{%
    \FormulaRefAuto{SemigroupTwoSidedDividesPrincipalIdeal}[thm]{1,2,4,11}}
\end{tabproofwide}

Linksteilbarkeit, Rechtsteilbarkeit und zweiseitige Teilbarkeit sind
reflexiv und transitiv.  In einer kommutativen Halbgruppe stimmen die drei
Relationen überein.  Für Monoide ist die ausdrückliche Alternative \(b=a\)
überflüssig, weil das neutrale Element als fehlender Faktor dient.

\section{Green-Relationen}

Die Green-Relationen vergleichen Elemente danach, welche Hauptideale sie
erzeugen. Die Indizes (A,\star) bleiben sichtbar, damit die Relation stets
eindeutig zu ihrer umgebenden Halbgruppe gehört.

\FormulaDefDeltaK[Greensche \(\mathcal L\)-Relation]{%
  \SemigroupAlg{A}{\star}\dsep a,b\in A\vdash
  a\mathrel{\mathcal L_{A,\star}}b\coloneqq
  L_{A,\star}(a)=L_{A,\star}(b)%
}{SemigroupGreenLDef}{%
  \DeltaRow{Mengen}{A\dsep a\dsep b}
  \DeltaRow{Funktionen}{\star}
  \DeltaRow{Relationsoperatoren}{\mathcal L_{A,\star}}
}

\FormulaDefDeltaK[Greensche \(\mathcal R\)-Relation]{%
  \SemigroupAlg{A}{\star}\dsep a,b\in A\vdash
  a\mathrel{\mathcal R_{A,\star}}b\coloneqq
  R_{A,\star}(a)=R_{A,\star}(b)%
}{SemigroupGreenRDef}{%
  \DeltaRow{Mengen}{A\dsep a\dsep b}
  \DeltaRow{Funktionen}{\star}
  \DeltaRow{Relationsoperatoren}{\mathcal R_{A,\star}}
}

\FormulaDefDeltaK[Greensche \(\mathcal J\)-Relation]{%
  \SemigroupAlg{A}{\star}\dsep a,b\in A\vdash
  a\mathrel{\mathcal J_{A,\star}}b\coloneqq
  J_{A,\star}(a)=J_{A,\star}(b)%
}{SemigroupGreenJDef}{%
  \DeltaRow{Mengen}{A\dsep a\dsep b}
  \DeltaRow{Funktionen}{\star}
  \DeltaRow{Relationsoperatoren}{\mathcal J_{A,\star}}
}

\FormulaDefDeltaK[Greensche \(\mathcal H\)-Relation]{%
  \begin{aligned}[t]
    &\SemigroupAlg{A}{\star}\dsep a,b\in A\vdash{}\\[-2pt]
    &a\mathrel{\mathcal H_{A,\star}}b\coloneqq
      \bigl(a\mathrel{\mathcal L_{A,\star}}b\land
            a\mathrel{\mathcal R_{A,\star}}b\bigr)
  \end{aligned}
}{SemigroupGreenHDef}{%
  \DeltaRow{Mengen}{A\dsep a\dsep b}
  \DeltaRow{Funktionen}{\star}
  \DeltaRow{Relationsoperatoren}{%
    \mathcal L_{A,\star}\dsep\mathcal R_{A,\star}\dsep
    \mathcal H_{A,\star}}
}

\FormulaThmDeltaK[\(\mathcal L\) als gegenseitige Rechtsteilbarkeit]{%
  \begin{aligned}[t]
    &\SemigroupAlg{A}{\star}\dsep a,b\in A\vdash{}\\[-2pt]
    &a\mathrel{\mathcal L_{A,\star}}b\leftrightarrow
      \Bigl(\SemigroupRightDivides{A,\star}{a}{b}\land
            \SemigroupRightDivides{A,\star}{b}{a}\Bigr)
  \end{aligned}
}{SemigroupGreenLMutualRightDivisibility}{%
  \DeltaRow{Mengen}{A\dsep a\dsep b}
  \DeltaRow{Funktionen}{\star}
  \DeltaRow{Relationsoperatoren}{\mathcal L_{A,\star}}
}
\begin{tabproofwide}
  \proofstepwidestar[1]{\SemigroupAlg{A}{\star}}{\rA}
  \proofstepwidestar[2]{a\in A}{\rA}
  \proofstepwidestar[3]{b\in A}{\rA}
  \proofstepwidestar[4]{a\mathrel{\mathcal L_{A,\star}}b}{\rA}
  \proofstepwidestar[1,2,3,4]{L_{A,\star}(a)=L_{A,\star}(b)}{%
    \FormulaRefAuto{SemigroupGreenLDef}[def]{1,2,3,4}}
  \proofstepwidestar[1,3]{b\in L_{A,\star}(b)}{%
    \rAEa{\FormulaRefAuto{SemigroupPrincipalIdealsContainGenerator}[thm]{1,3}}}
  \proofstepwidestar[1,2,3,4]{b\in L_{A,\star}(a)}{%
    \rIE{\FormulaRefAuto{a=b\vdash b=a}[thm]{5},6}}
  \proofstepwidestar[1,2,3,4]{\SemigroupRightDivides{A,\star}{a}{b}}{%
    \FormulaRefAuto{SemigroupRightDividesPrincipalLeftIdeal}[thm]{1,2,3,7}}
  \proofstepwidestar[1,2]{a\in L_{A,\star}(a)}{%
    \rAEa{\FormulaRefAuto{SemigroupPrincipalIdealsContainGenerator}[thm]{1,2}}}
  \proofstepwidestar[1,2,3,4]{a\in L_{A,\star}(b)}{\rIE{5,9}}
  \proofstepwidestar[1,2,3,4]{\SemigroupRightDivides{A,\star}{b}{a}}{%
    \FormulaRefAuto{SemigroupRightDividesPrincipalLeftIdeal}[thm]{1,3,2,10}}
  \proofstepwidestar[1,2,3,4]{%
    \SemigroupRightDivides{A,\star}{a}{b}\land
    \SemigroupRightDivides{A,\star}{b}{a}}{\rAI{8,11}}
  \proofstepwidestar[1,2,3]{%
    a\mathrel{\mathcal L_{A,\star}}b\rightarrow
    \bigl(\SemigroupRightDivides{A,\star}{a}{b}\land
          \SemigroupRightDivides{A,\star}{b}{a}\bigr)}{\rRI{4,12}}
  \proofstepwidestar[14]{%
    \SemigroupRightDivides{A,\star}{a}{b}\land
    \SemigroupRightDivides{A,\star}{b}{a}}{\rA}
  \proofstepwidestar[14]{\SemigroupRightDivides{A,\star}{a}{b}}{\rAEa{14}}
  \proofstepwidestar[14]{\SemigroupRightDivides{A,\star}{b}{a}}{\rAEb{14}}
  \proofstepwidestar[1,2,3,14]{b\in L_{A,\star}(a)}{%
    \FormulaRefAuto{SemigroupRightDividesPrincipalLeftIdeal}[thm]{1,2,3,15}}
  \proofstepwidestar[1,2,3,14]{a\in L_{A,\star}(b)}{%
    \FormulaRefAuto{SemigroupRightDividesPrincipalLeftIdeal}[thm]{1,3,2,16}}
  \proofstepwidestar[1,2]{\mathsf{LId}(L_{A,\star}(a);A,\star)}{%
    \FormulaRefAuto{SemigroupPrincipalLeftIdealIsLeftIdeal}[thm]{1,2}}
  \proofstepwidestar[1,3]{\mathsf{LId}(L_{A,\star}(b);A,\star)}{%
    \FormulaRefAuto{SemigroupPrincipalLeftIdealIsLeftIdeal}[thm]{1,3}}
  \proofstepwidestar[1,2,3,14]{L_{A,\star}(a)\subseteq L_{A,\star}(b)}{%
    \FormulaRefAuto{SemigroupPrincipalLeftIdealMinimal}[thm]{1,2,20,18}}
  \proofstepwidestar[1,2,3,14]{L_{A,\star}(b)\subseteq L_{A,\star}(a)}{%
    \FormulaRefAuto{SemigroupPrincipalLeftIdealMinimal}[thm]{1,3,19,17}}
  \proofstepwidestar[1,2,3,14]{L_{A,\star}(a)=L_{A,\star}(b)}{%
    \FormulaRefAuto{A\subseteq B,\,B\subseteq A\vdash A=B}[thm]{21,22}}
  \proofstepwidestar[1,2,3,14]{a\mathrel{\mathcal L_{A,\star}}b}{%
    \FormulaRefAuto{SemigroupGreenLDef}[def]{1,2,3,23}}
  \proofstepwidestar[1,2,3]{%
    \bigl(\SemigroupRightDivides{A,\star}{a}{b}\land
          \SemigroupRightDivides{A,\star}{b}{a}\bigr)
    \rightarrow a\mathrel{\mathcal L_{A,\star}}b}{\rRI{14,24}}
  \proofstepwidestar[1,2,3]{%
    a\mathrel{\mathcal L_{A,\star}}b\leftrightarrow
    \bigl(\SemigroupRightDivides{A,\star}{a}{b}\land
          \SemigroupRightDivides{A,\star}{b}{a}\bigr)}{\rLRI{13,25}}
\end{tabproofwide}

\FormulaThmDeltaK[\(\mathcal R\) als gegenseitige Linksteilbarkeit]{%
  \begin{aligned}[t]
    &\SemigroupAlg{A}{\star}\dsep a,b\in A\vdash{}\\[-2pt]
    &a\mathrel{\mathcal R_{A,\star}}b\leftrightarrow
      \Bigl(\SemigroupLeftDivides{A,\star}{a}{b}\land
            \SemigroupLeftDivides{A,\star}{b}{a}\Bigr)
  \end{aligned}
}{SemigroupGreenRMutualLeftDivisibility}{%
  \DeltaRow{Mengen}{A\dsep a\dsep b}
  \DeltaRow{Funktionen}{\star}
  \DeltaRow{Relationsoperatoren}{\mathcal R_{A,\star}}
}
\begin{tabproofwide}
  \proofstepwidestar[1]{\SemigroupAlg{A}{\star}}{\rA}
  \proofstepwidestar[2]{a\in A}{\rA}
  \proofstepwidestar[3]{b\in A}{\rA}
  \proofstepwidestar[4]{a\mathrel{\mathcal R_{A,\star}}b}{\rA}
  \proofstepwidestar[1,2,3,4]{R_{A,\star}(a)=R_{A,\star}(b)}{%
    \FormulaRefAuto{SemigroupGreenRDef}[def]{1,2,3,4}}
  \proofstepwidestar[1,3]{b\in R_{A,\star}(b)}{%
    \rAEa{\rAEb{\FormulaRefAuto{SemigroupPrincipalIdealsContainGenerator}[thm]{1,3}}}}
  \proofstepwidestar[1,2,3,4]{b\in R_{A,\star}(a)}{%
    \rIE{\FormulaRefAuto{a=b\vdash b=a}[thm]{5},6}}
  \proofstepwidestar[1,2,3,4]{\SemigroupLeftDivides{A,\star}{a}{b}}{%
    \FormulaRefAuto{SemigroupLeftDividesPrincipalRightIdeal}[thm]{1,2,3,7}}
  \proofstepwidestar[1,2]{a\in R_{A,\star}(a)}{%
    \rAEa{\rAEb{\FormulaRefAuto{SemigroupPrincipalIdealsContainGenerator}[thm]{1,2}}}}
  \proofstepwidestar[1,2,3,4]{a\in R_{A,\star}(b)}{\rIE{5,9}}
  \proofstepwidestar[1,2,3,4]{\SemigroupLeftDivides{A,\star}{b}{a}}{%
    \FormulaRefAuto{SemigroupLeftDividesPrincipalRightIdeal}[thm]{1,3,2,10}}
  \proofstepwidestar[1,2,3,4]{%
    \SemigroupLeftDivides{A,\star}{a}{b}\land
    \SemigroupLeftDivides{A,\star}{b}{a}}{\rAI{8,11}}
  \proofstepwidestar[1,2,3]{%
    a\mathrel{\mathcal R_{A,\star}}b\rightarrow
    \bigl(\SemigroupLeftDivides{A,\star}{a}{b}\land
          \SemigroupLeftDivides{A,\star}{b}{a}\bigr)}{\rRI{4,12}}
  \proofstepwidestar[14]{%
    \SemigroupLeftDivides{A,\star}{a}{b}\land
    \SemigroupLeftDivides{A,\star}{b}{a}}{\rA}
  \proofstepwidestar[14]{\SemigroupLeftDivides{A,\star}{a}{b}}{\rAEa{14}}
  \proofstepwidestar[14]{\SemigroupLeftDivides{A,\star}{b}{a}}{\rAEb{14}}
  \proofstepwidestar[1,2,3,14]{b\in R_{A,\star}(a)}{%
    \FormulaRefAuto{SemigroupLeftDividesPrincipalRightIdeal}[thm]{1,2,3,15}}
  \proofstepwidestar[1,2,3,14]{a\in R_{A,\star}(b)}{%
    \FormulaRefAuto{SemigroupLeftDividesPrincipalRightIdeal}[thm]{1,3,2,16}}
  \proofstepwidestar[1,2]{\mathsf{RId}(R_{A,\star}(a);A,\star)}{%
    \FormulaRefAuto{SemigroupPrincipalRightIdealIsRightIdeal}[thm]{1,2}}
  \proofstepwidestar[1,3]{\mathsf{RId}(R_{A,\star}(b);A,\star)}{%
    \FormulaRefAuto{SemigroupPrincipalRightIdealIsRightIdeal}[thm]{1,3}}
  \proofstepwidestar[1,2,3,14]{R_{A,\star}(a)\subseteq R_{A,\star}(b)}{%
    \FormulaRefAuto{SemigroupPrincipalRightIdealMinimal}[thm]{1,2,20,18}}
  \proofstepwidestar[1,2,3,14]{R_{A,\star}(b)\subseteq R_{A,\star}(a)}{%
    \FormulaRefAuto{SemigroupPrincipalRightIdealMinimal}[thm]{1,3,19,17}}
  \proofstepwidestar[1,2,3,14]{R_{A,\star}(a)=R_{A,\star}(b)}{%
    \FormulaRefAuto{A\subseteq B,\,B\subseteq A\vdash A=B}[thm]{21,22}}
  \proofstepwidestar[1,2,3,14]{a\mathrel{\mathcal R_{A,\star}}b}{%
    \FormulaRefAuto{SemigroupGreenRDef}[def]{1,2,3,23}}
  \proofstepwidestar[1,2,3]{%
    \bigl(\SemigroupLeftDivides{A,\star}{a}{b}\land
          \SemigroupLeftDivides{A,\star}{b}{a}\bigr)
    \rightarrow a\mathrel{\mathcal R_{A,\star}}b}{\rRI{14,24}}
  \proofstepwidestar[1,2,3]{%
    a\mathrel{\mathcal R_{A,\star}}b\leftrightarrow
    \bigl(\SemigroupLeftDivides{A,\star}{a}{b}\land
          \SemigroupLeftDivides{A,\star}{b}{a}\bigr)}{\rLRI{13,25}}
\end{tabproofwide}

\FormulaThmDeltaK[\(\mathcal J\) als gegenseitige zweiseitige Teilbarkeit]{%
  \begin{aligned}[t]
    &\SemigroupAlg{A}{\star}\dsep a,b\in A\vdash{}\\[-2pt]
    &a\mathrel{\mathcal J_{A,\star}}b\leftrightarrow
      \Bigl(\SemigroupDivides{A,\star}{a}{b}\land
            \SemigroupDivides{A,\star}{b}{a}\Bigr)
  \end{aligned}
}{SemigroupGreenJMutualTwoSidedDivisibility}{%
  \DeltaRow{Mengen}{A\dsep a\dsep b}
  \DeltaRow{Funktionen}{\star}
  \DeltaRow{Relationsoperatoren}{\mathcal J_{A,\star}}
}
\begin{tabproofwide}
  \proofstepwidestar[1]{\SemigroupAlg{A}{\star}}{\rA}
  \proofstepwidestar[2]{a\in A}{\rA}
  \proofstepwidestar[3]{b\in A}{\rA}
  \proofstepwidestar[4]{a\mathrel{\mathcal J_{A,\star}}b}{\rA}
  \proofstepwidestar[1,2,3,4]{J_{A,\star}(a)=J_{A,\star}(b)}{%
    \FormulaRefAuto{SemigroupGreenJDef}[def]{1,2,3,4}}
  \proofstepwidestar[1,3]{b\in J_{A,\star}(b)}{%
    \rAEb{\rAEb{\FormulaRefAuto{SemigroupPrincipalIdealsContainGenerator}[thm]{1,3}}}}
  \proofstepwidestar[1,2,3,4]{b\in J_{A,\star}(a)}{%
    \rIE{\FormulaRefAuto{a=b\vdash b=a}[thm]{5},6}}
  \proofstepwidestar[1,2,3,4]{\SemigroupDivides{A,\star}{a}{b}}{%
    \FormulaRefAuto{SemigroupTwoSidedDividesPrincipalIdeal}[thm]{1,2,3,7}}
  \proofstepwidestar[1,2]{a\in J_{A,\star}(a)}{%
    \rAEb{\rAEb{\FormulaRefAuto{SemigroupPrincipalIdealsContainGenerator}[thm]{1,2}}}}
  \proofstepwidestar[1,2,3,4]{a\in J_{A,\star}(b)}{\rIE{5,9}}
  \proofstepwidestar[1,2,3,4]{\SemigroupDivides{A,\star}{b}{a}}{%
    \FormulaRefAuto{SemigroupTwoSidedDividesPrincipalIdeal}[thm]{1,3,2,10}}
  \proofstepwidestar[1,2,3,4]{%
    \SemigroupDivides{A,\star}{a}{b}\land
    \SemigroupDivides{A,\star}{b}{a}}{\rAI{8,11}}
  \proofstepwidestar[1,2,3]{%
    a\mathrel{\mathcal J_{A,\star}}b\rightarrow
    \bigl(\SemigroupDivides{A,\star}{a}{b}\land
          \SemigroupDivides{A,\star}{b}{a}\bigr)}{\rRI{4,12}}
  \proofstepwidestar[14]{%
    \SemigroupDivides{A,\star}{a}{b}\land
    \SemigroupDivides{A,\star}{b}{a}}{\rA}
  \proofstepwidestar[14]{\SemigroupDivides{A,\star}{a}{b}}{\rAEa{14}}
  \proofstepwidestar[14]{\SemigroupDivides{A,\star}{b}{a}}{\rAEb{14}}
  \proofstepwidestar[1,2,3,14]{b\in J_{A,\star}(a)}{%
    \FormulaRefAuto{SemigroupTwoSidedDividesPrincipalIdeal}[thm]{1,2,3,15}}
  \proofstepwidestar[1,2,3,14]{a\in J_{A,\star}(b)}{%
    \FormulaRefAuto{SemigroupTwoSidedDividesPrincipalIdeal}[thm]{1,3,2,16}}
  \proofstepwidestar[1,2]{\mathsf{Id}(J_{A,\star}(a);A,\star)}{%
    \FormulaRefAuto{SemigroupPrincipalTwoSidedIdealIsIdeal}[thm]{1,2}}
  \proofstepwidestar[1,3]{\mathsf{Id}(J_{A,\star}(b);A,\star)}{%
    \FormulaRefAuto{SemigroupPrincipalTwoSidedIdealIsIdeal}[thm]{1,3}}
  \proofstepwidestar[1,2,3,14]{J_{A,\star}(a)\subseteq J_{A,\star}(b)}{%
    \FormulaRefAuto{SemigroupPrincipalTwoSidedIdealMinimal}[thm]{1,2,20,18}}
  \proofstepwidestar[1,2,3,14]{J_{A,\star}(b)\subseteq J_{A,\star}(a)}{%
    \FormulaRefAuto{SemigroupPrincipalTwoSidedIdealMinimal}[thm]{1,3,19,17}}
  \proofstepwidestar[1,2,3,14]{J_{A,\star}(a)=J_{A,\star}(b)}{%
    \FormulaRefAuto{A\subseteq B,\,B\subseteq A\vdash A=B}[thm]{21,22}}
  \proofstepwidestar[1,2,3,14]{a\mathrel{\mathcal J_{A,\star}}b}{%
    \FormulaRefAuto{SemigroupGreenJDef}[def]{1,2,3,23}}
  \proofstepwidestar[1,2,3]{%
    \bigl(\SemigroupDivides{A,\star}{a}{b}\land
          \SemigroupDivides{A,\star}{b}{a}\bigr)
    \rightarrow a\mathrel{\mathcal J_{A,\star}}b}{\rRI{14,24}}
  \proofstepwidestar[1,2,3]{%
    a\mathrel{\mathcal J_{A,\star}}b\leftrightarrow
    \bigl(\SemigroupDivides{A,\star}{a}{b}\land
          \SemigroupDivides{A,\star}{b}{a}\bigr)}{\rLRI{13,25}}
\end{tabproofwide}

\FormulaThmDeltaK[\(\mathcal H\) als gegenseitige einseitige Teilbarkeit]{%
  \begin{aligned}[t]
    &\SemigroupAlg{A}{\star}\dsep a,b\in A\vdash{}\\[-2pt]
    &a\mathrel{\mathcal H_{A,\star}}b\leftrightarrow
      \Bigl(
        \SemigroupRightDivides{A,\star}{a}{b}\land
        \SemigroupRightDivides{A,\star}{b}{a}\land
        \SemigroupLeftDivides{A,\star}{a}{b}\land
        \SemigroupLeftDivides{A,\star}{b}{a}\Bigr)
  \end{aligned}
}{SemigroupGreenHMutualOneSidedDivisibility}{%
  \DeltaRow{Mengen}{A\dsep a\dsep b}
  \DeltaRow{Funktionen}{\star}
  \DeltaRow{Relationsoperatoren}{\mathcal H_{A,\star}}
}
Die Definition von \(\mathcal H\) zerf\"allt in eine \(\mathcal L\)- und
eine \(\mathcal R\)-Bedingung.  Die beiden Richtungen bestehen daher nur
darin, diese Bedingungen in die zugeh\"origen Teilbarkeitsaussagen zu
\"ubersetzen beziehungsweise wieder zusammenzusetzen.

\begin{tabproofsplitwide}
  \proofpartwideindR[Hinrichtung]{%
    \begin{aligned}[t]
      &\SemigroupAlg{A}{\star}\dsep a,b\in A\dsep
        a\mathrel{\mathcal H_{A,\star}}b\vdash{}\\[-2pt]
      &\SemigroupRightDivides{A,\star}{a}{b}\land
        \SemigroupRightDivides{A,\star}{b}{a}\land
        \SemigroupLeftDivides{A,\star}{a}{b}\land
        \SemigroupLeftDivides{A,\star}{b}{a}
    \end{aligned}
  }{%
    \begin{aligned}[t]
      &\SemigroupAlg{A}{\star}\dsep a,b\in A\dsep
        a\mathrel{\mathcal H_{A,\star}}b\vdash{}\\[-2pt]
      &\SemigroupRightDivides{A,\star}{a}{b}\land
        \SemigroupRightDivides{A,\star}{b}{a}\land
        \SemigroupLeftDivides{A,\star}{a}{b}\land
        \SemigroupLeftDivides{A,\star}{b}{a}
    \end{aligned}
  }[SemigroupGreenHMutualOneSidedDivisibilityForward]
    \proofstepwidestar[1]{\SemigroupAlg{A}{\star}}{\rA}
    \proofstepwidestar[2]{a\in A}{\rA}
    \proofstepwidestar[3]{b\in A}{\rA}
    \proofstepwidestar[4]{a\mathrel{\mathcal H_{A,\star}}b}{\rA}
    \proofstepwidestar[1,2,3,4]{%
      a\mathrel{\mathcal L_{A,\star}}b\land
      a\mathrel{\mathcal R_{A,\star}}b}{%
      \FormulaRefAuto{SemigroupGreenHDef}[def]{1,2,3,4}}
    \proofstepwidestar[1,2,3,4]{%
      \SemigroupRightDivides{A,\star}{a}{b}\land
      \SemigroupRightDivides{A,\star}{b}{a}}{%
      \FormulaRefAuto{SemigroupGreenLMutualRightDivisibility}[thm]{%
        1,2,3,\rAEa{5}}}
    \proofstepwidestar[1,2,3,4]{%
      \SemigroupLeftDivides{A,\star}{a}{b}\land
      \SemigroupLeftDivides{A,\star}{b}{a}}{%
      \FormulaRefAuto{SemigroupGreenRMutualLeftDivisibility}[thm]{%
        1,2,3,\rAEb{5}}}
    \proofstepwidestarbreak[1,2,3,4]{%
      \begin{aligned}[t]
        &\SemigroupRightDivides{A,\star}{a}{b}\land
          \SemigroupRightDivides{A,\star}{b}{a}\land{}\\[-2pt]
        &\SemigroupLeftDivides{A,\star}{a}{b}\land
          \SemigroupLeftDivides{A,\star}{b}{a}
      \end{aligned}}{%
      \rAI{\rAEa{6},\rAI{\rAEb{6},\rAI{\rAEa{7},\rAEb{7}}}}}
  \closeproofpartwideindR

  \proofpartwideindR[R\"uckrichtung]{%
    \begin{aligned}[t]
      &\SemigroupAlg{A}{\star}\dsep a,b\in A\dsep{}\\[-2pt]
      &\SemigroupRightDivides{A,\star}{a}{b}\land
        \SemigroupRightDivides{A,\star}{b}{a}\land
        \SemigroupLeftDivides{A,\star}{a}{b}\land
        \SemigroupLeftDivides{A,\star}{b}{a}
        \vdash a\mathrel{\mathcal H_{A,\star}}b
    \end{aligned}
  }{%
    \begin{aligned}[t]
      &\SemigroupAlg{A}{\star}\dsep a,b\in A\dsep{}\\[-2pt]
      &\SemigroupRightDivides{A,\star}{a}{b}\land
        \SemigroupRightDivides{A,\star}{b}{a}\land
        \SemigroupLeftDivides{A,\star}{a}{b}\land
        \SemigroupLeftDivides{A,\star}{b}{a}
        \vdash a\mathrel{\mathcal H_{A,\star}}b
    \end{aligned}
  }[SemigroupGreenHMutualOneSidedDivisibilityBackward]
    \proofstepwidestar[1]{\SemigroupAlg{A}{\star}}{\rA}
    \proofstepwidestar[2]{a\in A}{\rA}
    \proofstepwidestar[3]{b\in A}{\rA}
    \proofstepwidestar[4]{%
      \begin{aligned}[t]
        &\SemigroupRightDivides{A,\star}{a}{b}\land
          \SemigroupRightDivides{A,\star}{b}{a}\land{}\\[-2pt]
        &\SemigroupLeftDivides{A,\star}{a}{b}\land
          \SemigroupLeftDivides{A,\star}{b}{a}
      \end{aligned}}{\rA}
    \proofstepwidestar[4]{%
      \SemigroupRightDivides{A,\star}{a}{b}\land
      \SemigroupRightDivides{A,\star}{b}{a}}{%
      \rAI{\rAEa{4},\rAEa{\rAEb{4}}}}
    \proofstepwidestar[4]{%
      \SemigroupLeftDivides{A,\star}{a}{b}\land
      \SemigroupLeftDivides{A,\star}{b}{a}}{%
      \rAI{\rAEa{\rAEb{\rAEb{4}}},\rAEb{\rAEb{\rAEb{4}}}}}
    \proofstepwidestar[1,2,3,4]{a\mathrel{\mathcal L_{A,\star}}b}{%
      \FormulaRefAuto{SemigroupGreenLMutualRightDivisibility}[thm]{1,2,3,5}}
    \proofstepwidestar[1,2,3,4]{a\mathrel{\mathcal R_{A,\star}}b}{%
      \FormulaRefAuto{SemigroupGreenRMutualLeftDivisibility}[thm]{1,2,3,6}}
    \proofstepwidestar[1,2,3,4]{a\mathrel{\mathcal H_{A,\star}}b}{%
      \FormulaRefAuto{SemigroupGreenHDef}[def]{1,2,3,\rAI{7,8}}}
  \closeproofpartwideindR

  \proofpartwide{Schluss}
    \proofstepwidestar[1]{\SemigroupAlg{A}{\star}}{\rA}
    \proofstepwidestar[2]{a\in A}{\rA}
    \proofstepwidestar[3]{b\in A}{\rA}
    \proofstepwidestar[4]{a\mathrel{\mathcal H_{A,\star}}b}{\rA}
    \proofstepwidestarbreak[1,2,3,4]{%
      \begin{aligned}[t]
        &\SemigroupRightDivides{A,\star}{a}{b}\land
          \SemigroupRightDivides{A,\star}{b}{a}\land{}\\[-2pt]
        &\SemigroupLeftDivides{A,\star}{a}{b}\land
          \SemigroupLeftDivides{A,\star}{b}{a}
      \end{aligned}}{%
      \FormulaRefAuto{SemigroupGreenHMutualOneSidedDivisibilityForward}[thm]{%
        1,2,3,4}}
    \proofstepwidestar[1,2,3]{%
      \begin{aligned}[t]
        a\mathrel{\mathcal H_{A,\star}}b\rightarrow\bigl(&
          \SemigroupRightDivides{A,\star}{a}{b}\land
          \SemigroupRightDivides{A,\star}{b}{a}\land{}\\[-2pt]
        &\SemigroupLeftDivides{A,\star}{a}{b}\land
          \SemigroupLeftDivides{A,\star}{b}{a}\bigr)
      \end{aligned}}{\rRI{4,5}}
    \proofstepwidestar[7]{%
      \begin{aligned}[t]
        &\SemigroupRightDivides{A,\star}{a}{b}\land
          \SemigroupRightDivides{A,\star}{b}{a}\land{}\\[-2pt]
        &\SemigroupLeftDivides{A,\star}{a}{b}\land
          \SemigroupLeftDivides{A,\star}{b}{a}
      \end{aligned}}{\rA}
    \proofstepwidestar[1,2,3,7]{a\mathrel{\mathcal H_{A,\star}}b}{%
      \FormulaRefAuto{SemigroupGreenHMutualOneSidedDivisibilityBackward}[thm]{%
        1,2,3,7}}
    \proofstepwidestar[1,2,3]{%
      \begin{aligned}[t]
        \bigl(&\SemigroupRightDivides{A,\star}{a}{b}\land
          \SemigroupRightDivides{A,\star}{b}{a}\land{}\\[-2pt]
        &\SemigroupLeftDivides{A,\star}{a}{b}\land
          \SemigroupLeftDivides{A,\star}{b}{a}\bigr)
          \rightarrow a\mathrel{\mathcal H_{A,\star}}b
      \end{aligned}}{\rRI{7,8}}
    \proofstepwidestarbreak[1,2,3]{%
      \begin{aligned}[t]
        a\mathrel{\mathcal H_{A,\star}}b\leftrightarrow\bigl(&
          \SemigroupRightDivides{A,\star}{a}{b}\land
          \SemigroupRightDivides{A,\star}{b}{a}\land{}\\[-2pt]
        &\SemigroupLeftDivides{A,\star}{a}{b}\land
          \SemigroupLeftDivides{A,\star}{b}{a}\bigr)
      \end{aligned}}{\rLRI{6,9}}
  \closeproofpartwide
\end{tabproofsplitwide}

\subsection{Äquivalenzeigenschaften}

\FormulaThmDeltaK[Reflexivität von \(\mathcal L\)]{%
  \SemigroupAlg{A}{\star}\dsep a\in A\vdash
  a\mathrel{\mathcal L_{A,\star}}a%
}{SemigroupGreenLReflexive}{%
  \DeltaRow{Mengen}{A\dsep a}
  \DeltaRow{Funktionen}{\star}
  \DeltaRow{Relationsoperatoren}{\mathcal L_{A,\star}}
}
\begin{tabproofwide}
  \proofstepwidestar[1]{\SemigroupAlg{A}{\star}}{\rA}
  \proofstepwidestar[2]{a\in A}{\rA}
  \proofstepwidestar[]{L_{A,\star}(a)=L_{A,\star}(a)}{\rII}
  \proofstepwidestar[1,2]{a\mathrel{\mathcal L_{A,\star}}a}{%
    \FormulaRefAuto{SemigroupGreenLDef}[def]{1,2,2,3}}
\end{tabproofwide}

\FormulaThmDeltaK[Symmetrie von \(\mathcal L\)]{%
  \begin{aligned}[t]
    &\SemigroupAlg{A}{\star}\dsep a,b\in A\dsep
      a\mathrel{\mathcal L_{A,\star}}b\vdash{}\\[-2pt]
    &b\mathrel{\mathcal L_{A,\star}}a
  \end{aligned}
}{SemigroupGreenLSymmetric}{%
  \DeltaRow{Mengen}{A\dsep a\dsep b}
  \DeltaRow{Funktionen}{\star}
  \DeltaRow{Relationsoperatoren}{\mathcal L_{A,\star}}
}
\begin{tabproofwide}
  \proofstepwidestar[1]{\SemigroupAlg{A}{\star}}{\rA}
  \proofstepwidestar[2]{a\in A}{\rA}
  \proofstepwidestar[3]{b\in A}{\rA}
  \proofstepwidestar[4]{a\mathrel{\mathcal L_{A,\star}}b}{\rA}
  \proofstepwidestar[1,2,3,4]{L_{A,\star}(a)=L_{A,\star}(b)}{%
    \FormulaRefAuto{SemigroupGreenLDef}[def]{1,2,3,4}}
  \proofstepwidestar[1,2,3,4]{L_{A,\star}(b)=L_{A,\star}(a)}{%
    \FormulaRefAuto{a=b\vdash b=a}[thm]{5}}
  \proofstepwidestar[1,2,3,4]{b\mathrel{\mathcal L_{A,\star}}a}{%
    \FormulaRefAuto{SemigroupGreenLDef}[def]{1,3,2,6}}
\end{tabproofwide}

\FormulaThmDeltaK[Transitivität von \(\mathcal L\)]{%
  \begin{aligned}[t]
    &\SemigroupAlg{A}{\star}\dsep a,b,c\in A\dsep
      a\mathrel{\mathcal L_{A,\star}}b\dsep
      b\mathrel{\mathcal L_{A,\star}}c\vdash{}\\[-2pt]
    &a\mathrel{\mathcal L_{A,\star}}c
  \end{aligned}
}{SemigroupGreenLTransitive}{%
  \DeltaRow{Mengen}{A\dsep a\dsep b\dsep c}
  \DeltaRow{Funktionen}{\star}
  \DeltaRow{Relationsoperatoren}{\mathcal L_{A,\star}}
}
\begin{tabproofwide}
  \proofstepwidestar[1]{\SemigroupAlg{A}{\star}}{\rA}
  \proofstepwidestar[2]{a\in A}{\rA}
  \proofstepwidestar[3]{b\in A}{\rA}
  \proofstepwidestar[4]{c\in A}{\rA}
  \proofstepwidestar[5]{a\mathrel{\mathcal L_{A,\star}}b}{\rA}
  \proofstepwidestar[6]{b\mathrel{\mathcal L_{A,\star}}c}{\rA}
  \proofstepwidestar[1,2,3,5]{L_{A,\star}(a)=L_{A,\star}(b)}{%
    \FormulaRefAuto{SemigroupGreenLDef}[def]{1,2,3,5}}
  \proofstepwidestar[1,3,4,6]{L_{A,\star}(b)=L_{A,\star}(c)}{%
    \FormulaRefAuto{SemigroupGreenLDef}[def]{1,3,4,6}}
  \proofstepwidestar[1,2,3,4,5,6]{L_{A,\star}(a)=L_{A,\star}(c)}{%
    \FormulaRefAuto{a=b\dsep b=c\vdash a=c}[thm]{7,8}}
  \proofstepwidestar[1,2,3,4,5,6]{a\mathrel{\mathcal L_{A,\star}}c}{%
    \FormulaRefAuto{SemigroupGreenLDef}[def]{1,2,4,9}}
\end{tabproofwide}

\FormulaThmDeltaK[\(\mathcal L\) ist eine Äquivalenzrelation]{%
  \SemigroupAlg{A}{\star}\vdash\EqRel{A,\mathcal L_{A,\star}}%
}{SemigroupGreenLIsEquivalence}{%
  \DeltaRow{Mengen}{A\dsep a\dsep b\dsep c}
  \DeltaRow{Funktionen}{\star}
  \DeltaRow{Relationsoperatoren}{\mathcal L_{A,\star}}
}
\begin{tabproofwide}
  \proofstepwidestar[1]{\SemigroupAlg{A}{\star}}{\rA}
  \proofstepwidestar[1]{a\in A\rightarrow a\mathrel{\mathcal L_{A,\star}}a}{%
    \FormulaRefAuto{SemigroupGreenLReflexive}[thm]{1}}
  \proofstepwidestar[1]{%
    a\mathrel{\mathcal L_{A,\star}}b\rightarrow
    b\mathrel{\mathcal L_{A,\star}}a}{%
    \FormulaRefAuto{SemigroupGreenLSymmetric}[thm]{1}}
  \proofstepwidestar[1]{%
    \bigl(a\mathrel{\mathcal L_{A,\star}}b\land
          b\mathrel{\mathcal L_{A,\star}}c\bigr)
    \rightarrow a\mathrel{\mathcal L_{A,\star}}c}{%
    \FormulaRefAuto{SemigroupGreenLTransitive}[thm]{1}}
  \proofstepwidestar[1]{\EqRel{A,\mathcal L_{A,\star}}}{%
    \FormulaRefAuto{EqRelDef}[def]{2,3,4}}
\end{tabproofwide}

\FormulaThmDeltaK[Reflexivität von \(\mathcal R\)]{%
  \SemigroupAlg{A}{\star}\dsep a\in A\vdash
  a\mathrel{\mathcal R_{A,\star}}a%
}{SemigroupGreenRReflexive}{%
  \DeltaRow{Mengen}{A\dsep a}
  \DeltaRow{Funktionen}{\star}
  \DeltaRow{Relationsoperatoren}{\mathcal R_{A,\star}}
}
\begin{tabproofwide}
  \proofstepwidestar[1]{\SemigroupAlg{A}{\star}}{\rA}
  \proofstepwidestar[2]{a\in A}{\rA}
  \proofstepwidestar[]{R_{A,\star}(a)=R_{A,\star}(a)}{\rII}
  \proofstepwidestar[1,2]{a\mathrel{\mathcal R_{A,\star}}a}{%
    \FormulaRefAuto{SemigroupGreenRDef}[def]{1,2,2,3}}
\end{tabproofwide}

\FormulaThmDeltaK[Symmetrie von \(\mathcal R\)]{%
  \SemigroupAlg{A}{\star}\dsep a,b\in A\dsep
  a\mathrel{\mathcal R_{A,\star}}b\vdash
  b\mathrel{\mathcal R_{A,\star}}a%
}{SemigroupGreenRSymmetric}{%
  \DeltaRow{Mengen}{A\dsep a\dsep b}
  \DeltaRow{Funktionen}{\star}
  \DeltaRow{Relationsoperatoren}{\mathcal R_{A,\star}}
}
\begin{tabproofwide}
  \proofstepwidestar[1]{\SemigroupAlg{A}{\star}}{\rA}
  \proofstepwidestar[2]{a\in A}{\rA}
  \proofstepwidestar[3]{b\in A}{\rA}
  \proofstepwidestar[4]{a\mathrel{\mathcal R_{A,\star}}b}{\rA}
  \proofstepwidestar[1,2,3,4]{R_{A,\star}(a)=R_{A,\star}(b)}{%
    \FormulaRefAuto{SemigroupGreenRDef}[def]{1,2,3,4}}
  \proofstepwidestar[1,2,3,4]{R_{A,\star}(b)=R_{A,\star}(a)}{%
    \FormulaRefAuto{a=b\vdash b=a}[thm]{5}}
  \proofstepwidestar[1,2,3,4]{b\mathrel{\mathcal R_{A,\star}}a}{%
    \FormulaRefAuto{SemigroupGreenRDef}[def]{1,3,2,6}}
\end{tabproofwide}

\FormulaThmDeltaK[Transitivität von \(\mathcal R\)]{%
  \begin{aligned}[t]
    &\SemigroupAlg{A}{\star}\dsep a,b,c\in A\dsep
      a\mathrel{\mathcal R_{A,\star}}b\dsep
      b\mathrel{\mathcal R_{A,\star}}c\vdash{}\\[-2pt]
    &a\mathrel{\mathcal R_{A,\star}}c
  \end{aligned}
}{SemigroupGreenRTransitive}{%
  \DeltaRow{Mengen}{A\dsep a\dsep b\dsep c}
  \DeltaRow{Funktionen}{\star}
  \DeltaRow{Relationsoperatoren}{\mathcal R_{A,\star}}
}
\begin{tabproofwide}
  \proofstepwidestar[1]{\SemigroupAlg{A}{\star}}{\rA}
  \proofstepwidestar[2]{a\in A}{\rA}
  \proofstepwidestar[3]{b\in A}{\rA}
  \proofstepwidestar[4]{c\in A}{\rA}
  \proofstepwidestar[5]{a\mathrel{\mathcal R_{A,\star}}b}{\rA}
  \proofstepwidestar[6]{b\mathrel{\mathcal R_{A,\star}}c}{\rA}
  \proofstepwidestar[1,2,3,5]{R_{A,\star}(a)=R_{A,\star}(b)}{%
    \FormulaRefAuto{SemigroupGreenRDef}[def]{1,2,3,5}}
  \proofstepwidestar[1,3,4,6]{R_{A,\star}(b)=R_{A,\star}(c)}{%
    \FormulaRefAuto{SemigroupGreenRDef}[def]{1,3,4,6}}
  \proofstepwidestar[1,2,3,4,5,6]{R_{A,\star}(a)=R_{A,\star}(c)}{%
    \FormulaRefAuto{a=b\dsep b=c\vdash a=c}[thm]{7,8}}
  \proofstepwidestar[1,2,3,4,5,6]{a\mathrel{\mathcal R_{A,\star}}c}{%
    \FormulaRefAuto{SemigroupGreenRDef}[def]{1,2,4,9}}
\end{tabproofwide}

\FormulaThmDeltaK[\(\mathcal R\) ist eine Äquivalenzrelation]{%
  \SemigroupAlg{A}{\star}\vdash\EqRel{A,\mathcal R_{A,\star}}%
}{SemigroupGreenRIsEquivalence}{%
  \DeltaRow{Mengen}{A\dsep a\dsep b\dsep c}
  \DeltaRow{Funktionen}{\star}
  \DeltaRow{Relationsoperatoren}{\mathcal R_{A,\star}}
}
\begin{tabproofwide}
  \proofstepwidestar[1]{\SemigroupAlg{A}{\star}}{\rA}
  \proofstepwidestar[1]{a\in A\rightarrow a\mathrel{\mathcal R_{A,\star}}a}{%
    \FormulaRefAuto{SemigroupGreenRReflexive}[thm]{1}}
  \proofstepwidestar[1]{%
    a\mathrel{\mathcal R_{A,\star}}b\rightarrow
    b\mathrel{\mathcal R_{A,\star}}a}{%
    \FormulaRefAuto{SemigroupGreenRSymmetric}[thm]{1}}
  \proofstepwidestar[1]{%
    \bigl(a\mathrel{\mathcal R_{A,\star}}b\land
          b\mathrel{\mathcal R_{A,\star}}c\bigr)
    \rightarrow a\mathrel{\mathcal R_{A,\star}}c}{%
    \FormulaRefAuto{SemigroupGreenRTransitive}[thm]{1}}
  \proofstepwidestar[1]{\EqRel{A,\mathcal R_{A,\star}}}{%
    \FormulaRefAuto{EqRelDef}[def]{2,3,4}}
\end{tabproofwide}

\FormulaThmDeltaK[Reflexivität von \(\mathcal J\)]{%
  \SemigroupAlg{A}{\star}\dsep a\in A\vdash
  a\mathrel{\mathcal J_{A,\star}}a%
}{SemigroupGreenJReflexive}{%
  \DeltaRow{Mengen}{A\dsep a}
  \DeltaRow{Funktionen}{\star}
  \DeltaRow{Relationsoperatoren}{\mathcal J_{A,\star}}
}
\begin{tabproofwide}
  \proofstepwidestar[1]{\SemigroupAlg{A}{\star}}{\rA}
  \proofstepwidestar[2]{a\in A}{\rA}
  \proofstepwidestar[]{J_{A,\star}(a)=J_{A,\star}(a)}{\rII}
  \proofstepwidestar[1,2]{a\mathrel{\mathcal J_{A,\star}}a}{%
    \FormulaRefAuto{SemigroupGreenJDef}[def]{1,2,2,3}}
\end{tabproofwide}

\FormulaThmDeltaK[Symmetrie von \(\mathcal J\)]{%
  \SemigroupAlg{A}{\star}\dsep a,b\in A\dsep
  a\mathrel{\mathcal J_{A,\star}}b\vdash
  b\mathrel{\mathcal J_{A,\star}}a%
}{SemigroupGreenJSymmetric}{%
  \DeltaRow{Mengen}{A\dsep a\dsep b}
  \DeltaRow{Funktionen}{\star}
  \DeltaRow{Relationsoperatoren}{\mathcal J_{A,\star}}
}
\begin{tabproofwide}
  \proofstepwidestar[1]{\SemigroupAlg{A}{\star}}{\rA}
  \proofstepwidestar[2]{a\in A}{\rA}
  \proofstepwidestar[3]{b\in A}{\rA}
  \proofstepwidestar[4]{a\mathrel{\mathcal J_{A,\star}}b}{\rA}
  \proofstepwidestar[1,2,3,4]{J_{A,\star}(a)=J_{A,\star}(b)}{%
    \FormulaRefAuto{SemigroupGreenJDef}[def]{1,2,3,4}}
  \proofstepwidestar[1,2,3,4]{J_{A,\star}(b)=J_{A,\star}(a)}{%
    \FormulaRefAuto{a=b\vdash b=a}[thm]{5}}
  \proofstepwidestar[1,2,3,4]{b\mathrel{\mathcal J_{A,\star}}a}{%
    \FormulaRefAuto{SemigroupGreenJDef}[def]{1,3,2,6}}
\end{tabproofwide}

\FormulaThmDeltaK[Transitivität von \(\mathcal J\)]{%
  \begin{aligned}[t]
    &\SemigroupAlg{A}{\star}\dsep a,b,c\in A\dsep
      a\mathrel{\mathcal J_{A,\star}}b\dsep
      b\mathrel{\mathcal J_{A,\star}}c\vdash{}\\[-2pt]
    &a\mathrel{\mathcal J_{A,\star}}c
  \end{aligned}
}{SemigroupGreenJTransitive}{%
  \DeltaRow{Mengen}{A\dsep a\dsep b\dsep c}
  \DeltaRow{Funktionen}{\star}
  \DeltaRow{Relationsoperatoren}{\mathcal J_{A,\star}}
}
\begin{tabproofwide}
  \proofstepwidestar[1]{\SemigroupAlg{A}{\star}}{\rA}
  \proofstepwidestar[2]{a\in A}{\rA}
  \proofstepwidestar[3]{b\in A}{\rA}
  \proofstepwidestar[4]{c\in A}{\rA}
  \proofstepwidestar[5]{a\mathrel{\mathcal J_{A,\star}}b}{\rA}
  \proofstepwidestar[6]{b\mathrel{\mathcal J_{A,\star}}c}{\rA}
  \proofstepwidestar[1,2,3,5]{J_{A,\star}(a)=J_{A,\star}(b)}{%
    \FormulaRefAuto{SemigroupGreenJDef}[def]{1,2,3,5}}
  \proofstepwidestar[1,3,4,6]{J_{A,\star}(b)=J_{A,\star}(c)}{%
    \FormulaRefAuto{SemigroupGreenJDef}[def]{1,3,4,6}}
  \proofstepwidestar[1,2,3,4,5,6]{J_{A,\star}(a)=J_{A,\star}(c)}{%
    \FormulaRefAuto{a=b\dsep b=c\vdash a=c}[thm]{7,8}}
  \proofstepwidestar[1,2,3,4,5,6]{a\mathrel{\mathcal J_{A,\star}}c}{%
    \FormulaRefAuto{SemigroupGreenJDef}[def]{1,2,4,9}}
\end{tabproofwide}

\FormulaThmDeltaK[\(\mathcal J\) ist eine Äquivalenzrelation]{%
  \SemigroupAlg{A}{\star}\vdash\EqRel{A,\mathcal J_{A,\star}}%
}{SemigroupGreenJIsEquivalence}{%
  \DeltaRow{Mengen}{A\dsep a\dsep b\dsep c}
  \DeltaRow{Funktionen}{\star}
  \DeltaRow{Relationsoperatoren}{\mathcal J_{A,\star}}
}
\begin{tabproofwide}
  \proofstepwidestar[1]{\SemigroupAlg{A}{\star}}{\rA}
  \proofstepwidestar[1]{a\in A\rightarrow a\mathrel{\mathcal J_{A,\star}}a}{%
    \FormulaRefAuto{SemigroupGreenJReflexive}[thm]{1}}
  \proofstepwidestar[1]{%
    a\mathrel{\mathcal J_{A,\star}}b\rightarrow
    b\mathrel{\mathcal J_{A,\star}}a}{%
    \FormulaRefAuto{SemigroupGreenJSymmetric}[thm]{1}}
  \proofstepwidestar[1]{%
    \bigl(a\mathrel{\mathcal J_{A,\star}}b\land
          b\mathrel{\mathcal J_{A,\star}}c\bigr)
    \rightarrow a\mathrel{\mathcal J_{A,\star}}c}{%
    \FormulaRefAuto{SemigroupGreenJTransitive}[thm]{1}}
  \proofstepwidestar[1]{\EqRel{A,\mathcal J_{A,\star}}}{%
    \FormulaRefAuto{EqRelDef}[def]{2,3,4}}
\end{tabproofwide}

\FormulaThmDeltaK[Reflexivität von \(\mathcal H\)]{%
  \SemigroupAlg{A}{\star}\dsep a\in A\vdash
  a\mathrel{\mathcal H_{A,\star}}a%
}{SemigroupGreenHReflexive}{%
  \DeltaRow{Mengen}{A\dsep a}
  \DeltaRow{Funktionen}{\star}
  \DeltaRow{Relationsoperatoren}{\mathcal H_{A,\star}}
}
\begin{tabproofwide}
  \proofstepwidestar[1]{\SemigroupAlg{A}{\star}}{\rA}
  \proofstepwidestar[2]{a\in A}{\rA}
  \proofstepwidestar[1,2]{a\mathrel{\mathcal L_{A,\star}}a}{%
    \FormulaRefAuto{SemigroupGreenLReflexive}[thm]{1,2}}
  \proofstepwidestar[1,2]{a\mathrel{\mathcal R_{A,\star}}a}{%
    \FormulaRefAuto{SemigroupGreenRReflexive}[thm]{1,2}}
  \proofstepwidestar[1,2]{a\mathrel{\mathcal H_{A,\star}}a}{%
    \FormulaRefAuto{SemigroupGreenHDef}[def]{1,2,2,\rAI{3,4}}}
\end{tabproofwide}

\FormulaThmDeltaK[Symmetrie von \(\mathcal H\)]{%
  \SemigroupAlg{A}{\star}\dsep a,b\in A\dsep
  a\mathrel{\mathcal H_{A,\star}}b\vdash
  b\mathrel{\mathcal H_{A,\star}}a%
}{SemigroupGreenHSymmetric}{%
  \DeltaRow{Mengen}{A\dsep a\dsep b}
  \DeltaRow{Funktionen}{\star}
  \DeltaRow{Relationsoperatoren}{\mathcal H_{A,\star}}
}
\begin{tabproofwide}
  \proofstepwidestar[1]{\SemigroupAlg{A}{\star}}{\rA}
  \proofstepwidestar[2]{a\in A}{\rA}
  \proofstepwidestar[3]{b\in A}{\rA}
  \proofstepwidestar[4]{a\mathrel{\mathcal H_{A,\star}}b}{\rA}
  \proofstepwidestar[1,2,3,4]{%
    a\mathrel{\mathcal L_{A,\star}}b\land
    a\mathrel{\mathcal R_{A,\star}}b}{%
    \FormulaRefAuto{SemigroupGreenHDef}[def]{1,2,3,4}}
  \proofstepwidestar[1,2,3,4]{b\mathrel{\mathcal L_{A,\star}}a}{%
    \FormulaRefAuto{SemigroupGreenLSymmetric}[thm]{1,2,3,\rAEa{5}}}
  \proofstepwidestar[1,2,3,4]{b\mathrel{\mathcal R_{A,\star}}a}{%
    \FormulaRefAuto{SemigroupGreenRSymmetric}[thm]{1,2,3,\rAEb{5}}}
  \proofstepwidestar[1,2,3,4]{b\mathrel{\mathcal H_{A,\star}}a}{%
    \FormulaRefAuto{SemigroupGreenHDef}[def]{1,3,2,\rAI{6,7}}}
\end{tabproofwide}

\FormulaThmDeltaK[Transitivität von \(\mathcal H\)]{%
  \begin{aligned}[t]
    &\SemigroupAlg{A}{\star}\dsep a,b,c\in A\dsep
      a\mathrel{\mathcal H_{A,\star}}b\dsep
      b\mathrel{\mathcal H_{A,\star}}c\vdash{}\\[-2pt]
    &a\mathrel{\mathcal H_{A,\star}}c
  \end{aligned}
}{SemigroupGreenHTransitive}{%
  \DeltaRow{Mengen}{A\dsep a\dsep b\dsep c}
  \DeltaRow{Funktionen}{\star}
  \DeltaRow{Relationsoperatoren}{\mathcal H_{A,\star}}
}
\begin{tabproofwide}
  \proofstepwidestar[1]{\SemigroupAlg{A}{\star}}{\rA}
  \proofstepwidestar[2]{a\in A}{\rA}
  \proofstepwidestar[3]{b\in A}{\rA}
  \proofstepwidestar[4]{c\in A}{\rA}
  \proofstepwidestar[5]{a\mathrel{\mathcal H_{A,\star}}b}{\rA}
  \proofstepwidestar[6]{b\mathrel{\mathcal H_{A,\star}}c}{\rA}
  \proofstepwidestar[1,2,3,5]{%
    a\mathrel{\mathcal L_{A,\star}}b\land
    a\mathrel{\mathcal R_{A,\star}}b}{%
    \FormulaRefAuto{SemigroupGreenHDef}[def]{1,2,3,5}}
  \proofstepwidestar[1,3,4,6]{%
    b\mathrel{\mathcal L_{A,\star}}c\land
    b\mathrel{\mathcal R_{A,\star}}c}{%
    \FormulaRefAuto{SemigroupGreenHDef}[def]{1,3,4,6}}
  \proofstepwidestar[1,2,3,4,5,6]{a\mathrel{\mathcal L_{A,\star}}c}{%
    \FormulaRefAuto{SemigroupGreenLTransitive}[thm]{1,2,3,4,\rAEa{7},\rAEa{8}}}
  \proofstepwidestar[1,2,3,4,5,6]{a\mathrel{\mathcal R_{A,\star}}c}{%
    \FormulaRefAuto{SemigroupGreenRTransitive}[thm]{1,2,3,4,\rAEb{7},\rAEb{8}}}
  \proofstepwidestar[1,2,3,4,5,6]{a\mathrel{\mathcal H_{A,\star}}c}{%
    \FormulaRefAuto{SemigroupGreenHDef}[def]{1,2,4,\rAI{9,10}}}
\end{tabproofwide}

\FormulaThmDeltaK[\(\mathcal H\) ist eine Äquivalenzrelation]{%
  \SemigroupAlg{A}{\star}\vdash\EqRel{A,\mathcal H_{A,\star}}%
}{SemigroupGreenHIsEquivalence}{%
  \DeltaRow{Mengen}{A\dsep a\dsep b\dsep c}
  \DeltaRow{Funktionen}{\star}
  \DeltaRow{Relationsoperatoren}{\mathcal H_{A,\star}}
}
\begin{tabproofwide}
  \proofstepwidestar[1]{\SemigroupAlg{A}{\star}}{\rA}
  \proofstepwidestar[1]{a\in A\rightarrow a\mathrel{\mathcal H_{A,\star}}a}{%
    \FormulaRefAuto{SemigroupGreenHReflexive}[thm]{1}}
  \proofstepwidestar[1]{%
    a\mathrel{\mathcal H_{A,\star}}b\rightarrow
    b\mathrel{\mathcal H_{A,\star}}a}{%
    \FormulaRefAuto{SemigroupGreenHSymmetric}[thm]{1}}
  \proofstepwidestar[1]{%
    \bigl(a\mathrel{\mathcal H_{A,\star}}b\land
          b\mathrel{\mathcal H_{A,\star}}c\bigr)
    \rightarrow a\mathrel{\mathcal H_{A,\star}}c}{%
    \FormulaRefAuto{SemigroupGreenHTransitive}[thm]{1}}
  \proofstepwidestar[1]{\EqRel{A,\mathcal H_{A,\star}}}{%
    \FormulaRefAuto{EqRelDef}[def]{2,3,4}}
\end{tabproofwide}

\section{Morphismen von Halbgruppen}

\FormulaDefDeltaK[Begriff des Halbgruppenhomomorphismus]{%
  \SemigroupHom{f}{A,\star}{B,\diamond}%
}{SemigroupHomDef}{%
  \DeltaRow{Mengen}{A\dsep B\dsep x\dsep y}
  \DeltaRow{Funktionen}{f\dsep\star\dsep\diamond}
  \DeltaRow{\textbf{Axiome}}{}
  \DeltaPrem{Quellhalbgruppe}{\SemigroupAlg{A}{\star}}
  \DeltaPrem{Zielhalbgruppe}{\SemigroupAlg{B}{\diamond}}
  \DeltaPrem{Abbildung}{f\colon A\to B}
  \DeltaRow{Operationserhaltung}{%
    x,y\in A\vdash
    f(x\star y)=f(x)\diamond f(y)}
  \DeltaRow{\textbf{Neues Symbol}}{}
  \DeltaPrem{Halbgruppenhomomorphismus}{%
    \SemigroupHom{f}{A,\star}{B,\diamond}}
}

\FormulaAxiomDeltaK[Quellhalbgruppe]{%
  \SemigroupHom{f}{A,\star}{B,\diamond}
  \vdash\SemigroupAlg{A}{\star}%
}{SemigroupHomSourceAxiom}{%
  \DeltaRow{Mengen}{A\dsep B}
  \DeltaRow{Funktionen}{f\dsep\star\dsep\diamond}
}

\FormulaAxiomDeltaK[Zielhalbgruppe]{%
  \SemigroupHom{f}{A,\star}{B,\diamond}
  \vdash\SemigroupAlg{B}{\diamond}%
}{SemigroupHomTargetAxiom}{%
  \DeltaRow{Mengen}{A\dsep B}
  \DeltaRow{Funktionen}{f\dsep\star\dsep\diamond}
}

\FormulaAxiomDeltaK[Operationserhaltung eines Halbgruppenhomomorphismus]{%
  \SemigroupHom{f}{A,\star}{B,\diamond}\dsep x,y\in A
  \vdash f(x\star y)=f(x)\diamond f(y)%
}{SemigroupHomOperationAxiom}{%
  \DeltaRow{Mengen}{A\dsep B\dsep x\dsep y}
  \DeltaRow{Funktionen}{f\dsep\star\dsep\diamond}
}

\FormulaAxiomDeltaK[Trägerabbildung eines Halbgruppenhomomorphismus]{%
  \SemigroupHom{f}{A,\star}{B,\diamond}
  \vdash f\colon A\to B%
}{SemigroupHomFunctionAxiom}{%
  \DeltaRow{Mengen}{A\dsep B}
  \DeltaRow{Funktionen}{f\dsep\star\dsep\diamond}
}

\FormulaDefDeltaK[Begriff des Halbgruppenisomorphismus]{%
  \SemigroupIso{f}{A,\star}{B,\diamond}%
}{SemigroupIsoDef}{%
  \DeltaRow{Mengen}{A\dsep B}
  \DeltaRow{Funktionen}{f\dsep\star\dsep\diamond}
  \DeltaRow{\textbf{Axiome}}{}
  \DeltaPrem{Halbgruppenhomomorphismus}{%
    \SemigroupHom{f}{A,\star}{B,\diamond}}
  \DeltaPrem{Bijektivität}{f\colon A\bij B}
  \DeltaRow{\textbf{Neues Symbol}}{}
  \DeltaPrem{Halbgruppenisomorphismus}{%
    \SemigroupIso{f}{A,\star}{B,\diamond}}
}

\FormulaAxiomDeltaK[Homomorphismuszugriff]{%
  \SemigroupIso{f}{A,\star}{B,\diamond}
  \vdash\SemigroupHom{f}{A,\star}{B,\diamond}%
}{SemigroupIsoHomAxiom}{%
  \DeltaRow{Mengen}{A\dsep B}
  \DeltaRow{Funktionen}{f\dsep\star\dsep\diamond}
}

\FormulaAxiomDeltaK[Bijektivitätszugriff]{%
  \SemigroupIso{f}{A,\star}{B,\diamond}
  \vdash f\colon A\bij B%
}{SemigroupIsoBijectiveAxiom}{%
  \DeltaRow{Mengen}{A\dsep B}
  \DeltaRow{Funktionen}{f\dsep\star\dsep\diamond}
}

\subsection{Grundlegende Eigenschaften}

\FormulaThmDeltaK[Komposition von Halbgruppenhomomorphismen]{%
  \begin{aligned}[t]
    &\SemigroupHom{f}{A,\star}{B,\diamond}\dsep
      \SemigroupHom{g}{B,\diamond}{C,\bullet}\\
    &\vdash\SemigroupHom{g\circ f}{A,\star}{C,\bullet}
  \end{aligned}
}{SemigroupHomComposition}{%
  \DeltaRow{Mengen}{A\dsep B\dsep C\dsep x\dsep y}
  \DeltaRow{Funktionen}{f\dsep g\dsep\star\dsep\diamond\dsep\bullet}
}
\begin{tabproofsplitwide}
  \proofpartwideindR[Komposition der Trägerabbildungen]{%
    \begin{aligned}[t]
      &\SemigroupHom{f}{A,\star}{B,\diamond}\dsep
        \SemigroupHom{g}{B,\diamond}{C,\bullet}\vdash{}\\[-2pt]
      &g\circ f\colon A\to C
    \end{aligned}
  }{%
    \SemigroupHom{f}{A,\star}{B,\diamond}\dsep
    \SemigroupHom{g}{B,\diamond}{C,\bullet}
    \vdash g\circ f\colon A\to C%
  }[SemigroupHomCompositionFunction]
    \proofstepwidestar[1]{\SemigroupHom{f}{A,\star}{B,\diamond}}{\rA}
    \proofstepwidestar[2]{\SemigroupHom{g}{B,\diamond}{C,\bullet}}{\rA}
    \proofstepwidestar[1]{f\colon A\to B}{%
      \FormulaRefAuto{SemigroupHomFunctionAxiom}[ax]{1}}
    \proofstepwidestar[2]{g\colon B\to C}{%
      \FormulaRefAuto{SemigroupHomFunctionAxiom}[ax]{2}}
    \proofstepwidestar[1,2]{g\circ f\colon A\to C}{%
      \FormulaRefAuto{G\circ F\colon A\to C}[thm]{3,4}}
  \closeproofpartwideindR

  \proofpartwideindR[Einsetzen der ersten Homomorphie]{%
    \begin{aligned}[t]
      &\SemigroupHom{f}{A,\star}{B,\diamond}\dsep
        \SemigroupHom{g}{B,\diamond}{C,\bullet}\dsep x,y\in A\vdash{}\\[-2pt]
      &(g\circ f)(x\star y)=g\bigl(f(x)\diamond f(y)\bigr)
    \end{aligned}
  }{%
    \SemigroupHom{f}{A,\star}{B,\diamond}\dsep
    \SemigroupHom{g}{B,\diamond}{C,\bullet}\dsep x,y\in A
    \vdash (g\circ f)(x\star y)=g\bigl(f(x)\diamond f(y)\bigr)%
  }[SemigroupHomCompositionFirstOperation]
    \proofstepwidestar[1]{\SemigroupHom{f}{A,\star}{B,\diamond}}{\rA}
    \proofstepwidestar[2]{\SemigroupHom{g}{B,\diamond}{C,\bullet}}{\rA}
    \proofstepwidestar[3]{x\in A}{\rA}
    \proofstepwidestar[4]{y\in A}{\rA}
    \proofstepwidestar[1]{f\colon A\to B}{%
      \FormulaRefAuto{SemigroupHomFunctionAxiom}[ax]{1}}
    \proofstepwidestar[2]{g\colon B\to C}{%
      \FormulaRefAuto{SemigroupHomFunctionAxiom}[ax]{2}}
    \proofstepwidestar[1]{\SemigroupAlg{A}{\star}}{%
      \FormulaRefAuto{SemigroupHomSourceAxiom}[ax]{1}}
    \proofstepwidestar[1,3,4]{x\star y\in A}{%
      \FormulaRefAuto{SemigroupClosureAxiom}[ax]{7,3,4}}
    \proofstepwidestar[1,3,4]{f(x\star y)=f(x)\diamond f(y)}{%
      \FormulaRefAuto{SemigroupHomOperationAxiom}[ax]{1,3,4}}
    \proofstepwide[1,2,3,4]{(g\circ f)(x\star y)}{=}{g(f(x\star y))}{%
      \FormulaRefAuto{%
        x\in A \vdash (F\circ G)(x)=F(G(x))%
      }[thm]{5,6,8}}
    \proofstepwide[1,2,3,4]{}{=}{g\bigl(f(x)\diamond f(y)\bigr)}{%
      \rIE{9}}
    \proofstepwidestar[1,2,3,4]{%
      (g\circ f)(x\star y)=g\bigl(f(x)\diamond f(y)\bigr)}{%
      \rChain{10,11}}
  \closeproofpartwideindR

  \proofpartwideindR[Einsetzen der zweiten Homomorphie]{%
    \begin{aligned}[t]
      &\SemigroupHom{f}{A,\star}{B,\diamond}\dsep
        \SemigroupHom{g}{B,\diamond}{C,\bullet}\dsep x,y\in A\vdash{}\\[-2pt]
      &g\bigl(f(x)\diamond f(y)\bigr)
        =(g\circ f)(x)\bullet(g\circ f)(y)
    \end{aligned}
  }{%
    \SemigroupHom{f}{A,\star}{B,\diamond}\dsep
    \SemigroupHom{g}{B,\diamond}{C,\bullet}\dsep x,y\in A
    \vdash g\bigl(f(x)\diamond f(y)\bigr)
      =(g\circ f)(x)\bullet(g\circ f)(y)%
  }[SemigroupHomCompositionSecondOperation]
    \proofstepwidestar[1]{\SemigroupHom{f}{A,\star}{B,\diamond}}{\rA}
    \proofstepwidestar[2]{\SemigroupHom{g}{B,\diamond}{C,\bullet}}{\rA}
    \proofstepwidestar[3]{x\in A}{\rA}
    \proofstepwidestar[4]{y\in A}{\rA}
    \proofstepwidestar[1]{f\colon A\to B}{%
      \FormulaRefAuto{SemigroupHomFunctionAxiom}[ax]{1}}
    \proofstepwidestar[2]{g\colon B\to C}{%
      \FormulaRefAuto{SemigroupHomFunctionAxiom}[ax]{2}}
    \proofstepwidestar[1,3]{f(x)\in B}{%
      \FormulaRefAuto{%
        F\colon A\to B\dsep x\in A\vdash F(x)\in B%
      }[thm]{5,3}}
    \proofstepwidestar[1,4]{f(y)\in B}{%
      \FormulaRefAuto{%
        F\colon A\to B\dsep x\in A\vdash F(x)\in B%
      }[thm]{5,4}}
    \proofstepwidestar[1,2,3]{g(f(x))=(g\circ f)(x)}{%
      \FormulaRefAuto{%
        x\in A \vdash F(G(x))=(F\circ G)(x)%
      }[thm]{5,6,3}}
    \proofstepwidestar[1,2,4]{g(f(y))=(g\circ f)(y)}{%
      \FormulaRefAuto{%
        x\in A \vdash F(G(x))=(F\circ G)(x)%
      }[thm]{5,6,4}}
    \proofstepwide[1,2,3,4]{g\bigl(f(x)\diamond f(y)\bigr)}{=}{%
      g(f(x))\bullet g(f(y))}{%
      \FormulaRefAuto{SemigroupHomOperationAxiom}[ax]{2,7,8}}
    \proofstepwide[1,2,3,4]{}{=}{(g\circ f)(x)\bullet g(f(y))}{\rIE{9}}
    \proofstepwide[1,2,3,4]{}{=}{%
      (g\circ f)(x)\bullet(g\circ f)(y)}{\rIE{10}}
    \proofstepwidestar[1,2,3,4]{%
      g\bigl(f(x)\diamond f(y)\bigr)
      =(g\circ f)(x)\bullet(g\circ f)(y)}{%
      \rChain{11,13}}
  \closeproofpartwideindR

  \proofpartwideindR[Zusammenführung]{%
    \begin{aligned}[t]
      &\SemigroupHom{f}{A,\star}{B,\diamond}\dsep
        \SemigroupHom{g}{B,\diamond}{C,\bullet}\vdash{}\\[-2pt]
      &\SemigroupHom{g\circ f}{A,\star}{C,\bullet}
    \end{aligned}
  }{%
    \SemigroupHom{f}{A,\star}{B,\diamond}\dsep
    \SemigroupHom{g}{B,\diamond}{C,\bullet}
    \vdash \SemigroupHom{g\circ f}{A,\star}{C,\bullet}%
  }[SemigroupHomCompositionConclusion]
    \proofstepwidestar[1]{\SemigroupHom{f}{A,\star}{B,\diamond}}{\rA}
    \proofstepwidestar[2]{\SemigroupHom{g}{B,\diamond}{C,\bullet}}{\rA}
    \proofstepwidestar[1]{\SemigroupAlg{A}{\star}}{%
      \FormulaRefAuto{SemigroupHomSourceAxiom}[ax]{1}}
    \proofstepwidestar[2]{\SemigroupAlg{C}{\bullet}}{%
      \FormulaRefAuto{SemigroupHomTargetAxiom}[ax]{2}}
    \proofstepwidestar[1,2]{g\circ f\colon A\to C}{%
      \FormulaRefAuto{SemigroupHomCompositionFunction}[thm]{1,2}}
    \proofstepwidestar[6]{x\in A}{\rA}
    \proofstepwidestar[7]{y\in A}{\rA}
    \proofstepwidestar[1,2,6,7]{%
      (g\circ f)(x\star y)=g\bigl(f(x)\diamond f(y)\bigr)}{%
      \FormulaRefAuto{SemigroupHomCompositionFirstOperation}[thm]{1,2,6,7}}
    \proofstepwidestar[1,2,6,7]{%
      g\bigl(f(x)\diamond f(y)\bigr)
      =(g\circ f)(x)\bullet(g\circ f)(y)}{%
      \FormulaRefAuto{SemigroupHomCompositionSecondOperation}[thm]{1,2,6,7}}
    \proofstepwidestar[1,2,6,7]{%
      (g\circ f)(x\star y)=(g\circ f)(x)\bullet(g\circ f)(y)}{%
      \rChain{8,9}}
    \proofstepwidestar[1,2]{%
      \SemigroupHom{g\circ f}{A,\star}{C,\bullet}}{%
      \FormulaRefAuto{SemigroupHomDef}[def]{3,4,5,10}}
  \closeproofpartwideindR
\end{tabproofsplitwide}

\FormulaThmDeltaK[Die Identität ist ein Halbgruppenisomorphismus]{%
  \SemigroupAlg{A}{\star}
  \vdash\SemigroupIso{\Id_A}{A,\star}{A,\star}%
}{SemigroupIdentityIsomorphism}{%
  \DeltaRow{Mengen}{A\dsep x\dsep y}
  \DeltaRow{Funktionen}{\star}
}
\begin{tabproofwide}
  \proofstepwidestar[1]{\SemigroupAlg{A}{\star}}{\rA}
  \proofstepwidestar[]{\Id_A\colon A\to A}{%
    \FormulaRefAuto{IdA}[thm]}
  \proofstepwidestar[2]{x\in A}{\rA}
  \proofstepwidestar[3]{y\in A}{\rA}
  \proofstepwidestar[1,2,3]{x\star y\in A}{%
    \FormulaRefAuto{SemigroupClosureAxiom}[ax]{1,2,3}}
  \proofstepwidestar[2]{\Id_A(x)=x}{%
    \FormulaRefAuto{x\in A\vdash \Id_A(x)=x}{2}}
  \proofstepwidestar[2]{x=\Id_A(x)}{%
    \FormulaRefAuto{a=b\vdash b=a}{6}}
  \proofstepwidestar[3]{\Id_A(y)=y}{%
    \FormulaRefAuto{x\in A\vdash \Id_A(x)=x}{3}}
  \proofstepwidestar[3]{y=\Id_A(y)}{%
    \FormulaRefAuto{a=b\vdash b=a}{8}}
  \proofstepwide[1,2,3]{\Id_A(x\star y)}{=}{x\star y}{%
    \FormulaRefAuto{x\in A\vdash \Id_A(x)=x}{5}}
  \proofstepwide[1,2,3]{}{=}{\Id_A(x)\star y}{\rIE{7}}
  \proofstepwide[1,2,3]{}{=}{\Id_A(x)\star\Id_A(y)}{\rIE{9}}
  \proofstepwidestar[1,2,3]{%
    \Id_A(x\star y)=\Id_A(x)\star\Id_A(y)}{\rChain{10,12}}
  \proofstepwidestar[1]{%
    \SemigroupHom{\Id_A}{A,\star}{A,\star}}{%
    \FormulaRefAuto{SemigroupHomDef}[def]{1,1,2,13}}
  \proofstepwidestar[]{\Id_A\colon A\bij A}{%
    \FormulaRefAuto{\Id_A\colon A\bij A}[thm]}
  \proofstepwidestar[1]{%
    \SemigroupIso{\Id_A}{A,\star}{A,\star}}{%
    \FormulaRefAuto{SemigroupIsoDef}[def]{14,15}}
\end{tabproofwide}

\FormulaThmDeltaK[Komposition von Halbgruppenisomorphismen]{%
  \begin{aligned}[t]
    &\SemigroupIso{f}{A,\star}{B,\diamond}\dsep
      \SemigroupIso{g}{B,\diamond}{C,\bullet}\\
    &\vdash\SemigroupIso{g\circ f}{A,\star}{C,\bullet}
  \end{aligned}
}{SemigroupIsoComposition}{%
  \DeltaRow{Mengen}{A\dsep B\dsep C}
  \DeltaRow{Funktionen}{f\dsep g\dsep\star\dsep\diamond\dsep\bullet}
}
\begin{tabproof}
  \proofstep{1}{\SemigroupIso{f}{A,\star}{B,\diamond}}{\rA}
  \proofstep{2}{\SemigroupIso{g}{B,\diamond}{C,\bullet}}{\rA}
  \proofstep{1}{\SemigroupHom{f}{A,\star}{B,\diamond}}{%
    \FormulaRefAuto{SemigroupIsoHomAxiom}[ax]{1}}
  \proofstep{2}{\SemigroupHom{g}{B,\diamond}{C,\bullet}}{%
    \FormulaRefAuto{SemigroupIsoHomAxiom}[ax]{2}}
  \proofstep{1,2}{\SemigroupHom{g\circ f}{A,\star}{C,\bullet}}{%
    \FormulaRefAuto{SemigroupHomComposition}[thm]{3,4}}
  \proofstep{1}{f\colon A\bij B}{%
    \FormulaRefAuto{SemigroupIsoBijectiveAxiom}[ax]{1}}
  \proofstep{2}{g\colon B\bij C}{%
    \FormulaRefAuto{SemigroupIsoBijectiveAxiom}[ax]{2}}
  \proofstep{1,2}{g\circ f\colon A\bij C}{%
    \FormulaRefAuto{F\colon A\bij B \dsep G\colon B\bij C\vdash
      G\circ F\colon A\bij C}[thm]{6,7}}
  \proofstep{1,2}{\SemigroupIso{g\circ f}{A,\star}{C,\bullet}}{%
    \FormulaRefAuto{SemigroupIsoDef}[def]{5,8}}
\end{tabproof}

\section{Morphismen von Potenzhalbgruppen}

\FormulaThmDeltaK[Ein Homomorphismus erhält das Mengenprodukt]{%
  \begin{aligned}[t]
    &\SemigroupHom{f}{A,\star}{B,\diamond}\dsep
      X,Y\in\PowNE{A}\vdash{}\\[-2pt]
    &f[X\star_{\mathcal P}Y]
      =f[X]\diamond_{\mathcal P}f[Y]
  \end{aligned}%
}{SemigroupHomPreservesPowerProduct}{%
  \DeltaRow{Mengen}{A\dsep B\dsep X\dsep Y\dsep
    u\dsep v\dsep w\dsep x\dsep y\dsep z}
  \DeltaRow{Funktionen}{f\dsep\star\dsep\diamond\dsep
    \star_{\mathcal P}\dsep\diamond_{\mathcal P}}
}
\begin{tabproofsplitwide}
  \proofpartwideindR[Bilder fester Faktoren]{%
    \begin{aligned}[t]
      &\SemigroupHom{f}{A,\star}{B,\diamond}\dsep
        X,Y\in\PowNE{A}\dsep x\in X\dsep y\in Y\vdash{}\\[-2pt]
      &f(x)\diamond f(y)\in f[X]\diamond_{\mathcal P}f[Y]
    \end{aligned}
  }{%
    \begin{aligned}[t]
      &\SemigroupHom{f}{A,\star}{B,\diamond}\dsep
        X,Y\in\PowNE{A}\dsep x\in X\dsep y\in Y\vdash{}\\[-2pt]
      &f(x)\diamond f(y)\in f[X]\diamond_{\mathcal P}f[Y]
    \end{aligned}
  }[PowerHomImagesOfFactors]
    \proofstepwidestar[1]{\SemigroupHom{f}{A,\star}{B,\diamond}}{\rA}
    \proofstepwidestar[2]{X\in\PowNE{A}}{\rA}
    \proofstepwidestar[3]{Y\in\PowNE{A}}{\rA}
    \proofstepwidestar[4]{x\in X}{\rA}
    \proofstepwidestar[5]{y\in Y}{\rA}
    \proofstepwidestar[1]{\SemigroupAlg{B}{\diamond}}{%
      \FormulaRefAuto{SemigroupHomTargetAxiom}[ax]{1}}
    \proofstepwidestar[1]{f\colon A\to B}{%
      \FormulaRefAuto{SemigroupHomFunctionAxiom}[ax]{1}}
    \proofstepwidestar[1,2]{f[X]\in\PowNE{B}}{%
      \FormulaRefAuto{NonemptyDirectImage}[thm]{7,2}}
    \proofstepwidestar[1,3]{f[Y]\in\PowNE{B}}{%
      \FormulaRefAuto{NonemptyDirectImage}[thm]{7,3}}
    \proofstepwidestar[2]{X\subseteq A}{%
      \rAEa{\FormulaRefAuto{NonemptyPowersetMembership}[thm]{2}}}
    \proofstepwidestar[3]{Y\subseteq A}{%
      \rAEa{\FormulaRefAuto{NonemptyPowersetMembership}[thm]{3}}}
    \proofstepwidestar[1,2,4]{f(x)\in f[X]}{%
      \FormulaRefAuto{C\subseteq A\dsep F\colon A\to B\dsep
        x\in C\vdash F(x)\in F[C]}[thm]{10,7,4}}
    \proofstepwidestar[1,3,5]{f(y)\in f[Y]}{%
      \FormulaRefAuto{C\subseteq A\dsep F\colon A\to B\dsep
        x\in C\vdash F(x)\in F[C]}[thm]{11,7,5}}
    \proofstepwidestar[1,2,3,4,5]{%
      f(x)\diamond f(y)\in f[X]\diamond_{\mathcal P}f[Y]}{%
      \FormulaRefAuto{PowerSemigroupProductContains}[thm]{6,8,9,12,13}}
  \closeproofpartwideindR

  \proofpartwideindR[Vorwärtsrichtung mit festen Zeugen]{%
    \begin{aligned}[t]
      &\SemigroupHom{f}{A,\star}{B,\diamond}\dsep X,Y\in\PowNE{A}
        \dsep x\in X\dsep y\in Y\dsep{}\\[-2pt]
      &w=x\star y\dsep z=f(w)
        \vdash z\in f[X]\diamond_{\mathcal P}f[Y]
    \end{aligned}
  }{%
    \begin{aligned}[t]
      &\SemigroupHom{f}{A,\star}{B,\diamond}\dsep X,Y\in\PowNE{A}
        \dsep x\in X\dsep y\in Y\dsep{}\\[-2pt]
      &w=x\star y\dsep z=f(w)
        \vdash z\in f[X]\diamond_{\mathcal P}f[Y]
    \end{aligned}
  }[PowerHomForwardFixedWitnesses]
    \proofstepwidestar[1]{\SemigroupHom{f}{A,\star}{B,\diamond}}{\rA}
    \proofstepwidestar[2]{X\in\PowNE{A}}{\rA}
    \proofstepwidestar[3]{Y\in\PowNE{A}}{\rA}
    \proofstepwidestar[4]{x\in X}{\rA}
    \proofstepwidestar[5]{y\in Y}{\rA}
    \proofstepwidestar[6]{w=x\star y}{\rA}
    \proofstepwidestar[7]{z=f(w)}{\rA}
    \proofstepwidestar[2,4]{x\in A}{%
      \FormulaRefAuto{PowerSemigroupCarrierElement}[thm]{2,4}}
    \proofstepwidestar[3,5]{y\in A}{%
      \FormulaRefAuto{PowerSemigroupCarrierElement}[thm]{3,5}}
    \proofstepwidestar[6,7]{z=f(x\star y)}{\rIE{6,7}}
    \proofstepwidestar[1,2,3,4,5]{%
      f(x\star y)=f(x)\diamond f(y)}{%
      \FormulaRefAuto{SemigroupHomOperationAxiom}[ax]{1,8,9}}
    \proofstepwidestar[1,2,3,4,5,6,7]{z=f(x)\diamond f(y)}{%
      \rChain{10,11}}
    \proofstepwidestar[1,2,3,4,5]{%
      f(x)\diamond f(y)\in f[X]\diamond_{\mathcal P}f[Y]}{%
      \FormulaRefAuto{PowerHomImagesOfFactors}[thm]{1,2,3,4,5}}
    \proofstepwidestar[1,2,3,4,5,6,7]{%
      z\in f[X]\diamond_{\mathcal P}f[Y]}{\rIE{12,13}}
  \closeproofpartwideindR

  \proofpartwideindR[Vorwärtsrichtung mit einem Produktzeugen]{%
    \begin{aligned}[t]
      &\SemigroupHom{f}{A,\star}{B,\diamond}\dsep X,Y\in\PowNE{A}
        \dsep w\in X\star_{\mathcal P}Y\dsep z=f(w)\vdash{}\\[-2pt]
      &z\in f[X]\diamond_{\mathcal P}f[Y]
    \end{aligned}
  }{%
    \begin{aligned}[t]
      &\SemigroupHom{f}{A,\star}{B,\diamond}\dsep X,Y\in\PowNE{A}
        \dsep w\in X\star_{\mathcal P}Y\dsep z=f(w)\vdash{}\\[-2pt]
      &z\in f[X]\diamond_{\mathcal P}f[Y]
    \end{aligned}
  }[PowerHomForwardProductWitness]
    \proofstepwidestar[1]{\SemigroupHom{f}{A,\star}{B,\diamond}}{\rA}
    \proofstepwidestar[2]{X\in\PowNE{A}}{\rA}
    \proofstepwidestar[3]{Y\in\PowNE{A}}{\rA}
    \proofstepwidestar[4]{w\in X\star_{\mathcal P}Y}{\rA}
    \proofstepwidestar[5]{z=f(w)}{\rA}
    \proofstepwidestar[1]{\SemigroupAlg{A}{\star}}{%
      \FormulaRefAuto{SemigroupHomSourceAxiom}[ax]{1}}
    \proofstepwidestar[1,2,3,4]{%
      \exists x\in X\,\exists y\in Y\,(w=x\star y)}{%
      \FormulaRefAuto{PowerSemigroupProductWitnesses}[thm]{6,2,3,4}}
    \proofstepwidestar[8]{x\in X\land\exists y\in Y\,(w=x\star y)}{\rA}
    \proofstepwidestar[8]{x\in X}{\rAEa{8}}
    \proofstepwidestar[8]{\exists y\in Y\,(w=x\star y)}{\rAEb{8}}
    \proofstepwidestar[11]{y\in Y\land w=x\star y}{\rA}
    \proofstepwidestar[11]{y\in Y}{\rAEa{11}}
    \proofstepwidestar[11]{w=x\star y}{\rAEb{11}}
    \proofstepwidestar[1,2,3,5,8,11]{%
      z\in f[X]\diamond_{\mathcal P}f[Y]}{%
      \FormulaRefAuto{PowerHomForwardFixedWitnesses}[thm]{%
        1,2,3,9,12,13,5}}
    \proofstepwidestar[1,2,3,5,8]{%
      z\in f[X]\diamond_{\mathcal P}f[Y]}{\rEE{10,11,14}}
    \proofstepwidestar[1,2,3,4,5]{%
      z\in f[X]\diamond_{\mathcal P}f[Y]}{\rEE{7,8,15}}
  \closeproofpartwideindR

  \proofpartwideindR[Vorwärtsrichtung für Bildelemente]{%
    \begin{aligned}[t]
      &\SemigroupHom{f}{A,\star}{B,\diamond}\dsep X,Y\in\PowNE{A}
        \dsep z\in f[X\star_{\mathcal P}Y]\vdash{}\\[-2pt]
      &z\in f[X]\diamond_{\mathcal P}f[Y]
    \end{aligned}
  }{%
    \begin{aligned}[t]
      &\SemigroupHom{f}{A,\star}{B,\diamond}\dsep X,Y\in\PowNE{A}
        \dsep z\in f[X\star_{\mathcal P}Y]\vdash{}\\[-2pt]
      &z\in f[X]\diamond_{\mathcal P}f[Y]
    \end{aligned}
  }[PowerHomForwardElement]
    \proofstepwidestar[1]{\SemigroupHom{f}{A,\star}{B,\diamond}}{\rA}
    \proofstepwidestar[2]{X\in\PowNE{A}}{\rA}
    \proofstepwidestar[3]{Y\in\PowNE{A}}{\rA}
    \proofstepwidestar[4]{z\in f[X\star_{\mathcal P}Y]}{\rA}
    \proofstepwidestar[1]{\SemigroupAlg{A}{\star}}{%
      \FormulaRefAuto{SemigroupHomSourceAxiom}[ax]{1}}
    \proofstepwidestar[1]{f\colon A\to B}{%
      \FormulaRefAuto{SemigroupHomFunctionAxiom}[ax]{1}}
    \proofstepwidestar[1,2,3]{X\star_{\mathcal P}Y\subseteq A}{%
      \FormulaRefAuto{PowerSemigroupProductSubset}[thm]{5,2,3}}
    \proofstepwidestar[1,2,3,4]{%
      \exists w\in X\star_{\mathcal P}Y\,(z=f(w))}{%
      \FormulaRefAuto{C\subseteq A\dsep F\colon A\to B\dsep
        y\in F[C]\vdash\exists x\in C\,y=F(x)}[thm]{7,6,4}}
    \proofstepwidestar[9]{w\in X\star_{\mathcal P}Y\land z=f(w)}{\rA}
    \proofstepwidestar[9]{w\in X\star_{\mathcal P}Y}{\rAEa{9}}
    \proofstepwidestar[9]{z=f(w)}{\rAEb{9}}
    \proofstepwidestar[1,2,3,9]{z\in f[X]\diamond_{\mathcal P}f[Y]}{%
      \FormulaRefAuto{PowerHomForwardProductWitness}[thm]{1,2,3,10,11}}
    \proofstepwidestar[1,2,3,4]{z\in f[X]\diamond_{\mathcal P}f[Y]}{%
      \rEE{8,9,12}}
  \closeproofpartwideindR

  \proofpartwideindR[Erste Teilmengenrichtung]{%
    \begin{aligned}[t]
      &\SemigroupHom{f}{A,\star}{B,\diamond}\dsep X,Y\in\PowNE{A}
        \vdash{}\\[-2pt]
      &f[X\star_{\mathcal P}Y]
        \subseteq f[X]\diamond_{\mathcal P}f[Y]
    \end{aligned}
  }{%
    \begin{aligned}[t]
      &\SemigroupHom{f}{A,\star}{B,\diamond}\dsep X,Y\in\PowNE{A}
        \vdash{}\\[-2pt]
      &f[X\star_{\mathcal P}Y]
        \subseteq f[X]\diamond_{\mathcal P}f[Y]
    \end{aligned}
  }[PowerHomForwardSubset]
    \proofstepwidestar[1]{\SemigroupHom{f}{A,\star}{B,\diamond}}{\rA}
    \proofstepwidestar[2]{X\in\PowNE{A}}{\rA}
    \proofstepwidestar[3]{Y\in\PowNE{A}}{\rA}
    \proofstepwidestar[4]{z\in f[X\star_{\mathcal P}Y]}{\rA}
    \proofstepwidestar[1,2,3,4]{z\in f[X]\diamond_{\mathcal P}f[Y]}{%
      \FormulaRefAuto{PowerHomForwardElement}[thm]{1,2,3,4}}
    \proofstepwidestar[1,2,3]{%
      \forall z\bigl(z\in f[X\star_{\mathcal P}Y]\rightarrow
      z\in f[X]\diamond_{\mathcal P}f[Y]\bigr)}{\rUI{\rRI{4,5}}}
    \proofstepwidestar[1,2,3]{%
      f[X\star_{\mathcal P}Y]\subseteq
      f[X]\diamond_{\mathcal P}f[Y]}{%
      \FormulaRefAuto{A\subseteq B\coloneq
        \forall x\,(x\in A\rightarrow x\in B)}[def]{6}}
  \closeproofpartwideindR

  \proofpartwideindR[Rückwärtsrichtung der Gleichheit]{%
    \begin{aligned}[t]
      &\SemigroupHom{f}{A,\star}{B,\diamond}\dsep x,y\in A\dsep
        u=f(x)\dsep v=f(y)\dsep z=u\diamond v\vdash{}\\[-2pt]
      &z=f(x\star y)
    \end{aligned}
  }{%
    \begin{aligned}[t]
      &\SemigroupHom{f}{A,\star}{B,\diamond}\dsep x,y\in A\dsep
        u=f(x)\dsep v=f(y)\dsep z=u\diamond v\vdash{}\\[-2pt]
      &z=f(x\star y)
    \end{aligned}
  }[PowerHomBackwardEquality]
    \proofstepwidestar[1]{\SemigroupHom{f}{A,\star}{B,\diamond}}{\rA}
    \proofstepwidestar[2]{x\in A}{\rA}
    \proofstepwidestar[3]{y\in A}{\rA}
    \proofstepwidestar[4]{u=f(x)}{\rA}
    \proofstepwidestar[5]{v=f(y)}{\rA}
    \proofstepwidestar[6]{z=u\diamond v}{\rA}
    \proofstepwidestar[1,2,3]{f(x\star y)=f(x)\diamond f(y)}{%
      \FormulaRefAuto{SemigroupHomOperationAxiom}[ax]{1,2,3}}
    \proofstepwidestar[1,2,3]{f(x)\diamond f(y)=f(x\star y)}{%
      \FormulaRefAuto{a=b\vdash b=a}[thm]{7}}
    \proofstepwidestar[4,6]{z=f(x)\diamond v}{\rIE{4,6}}
    \proofstepwidestar[4,5,6]{z=f(x)\diamond f(y)}{\rIE{5,9}}
    \proofstepwidestar[1,2,3,4,5,6]{z=f(x\star y)}{\rChain{10,8}}
  \closeproofpartwideindR

  \proofpartwideindR[Rückwärtsrichtung mit festen Zeugen]{%
    \begin{aligned}[t]
      &\SemigroupHom{f}{A,\star}{B,\diamond}\dsep X,Y\in\PowNE{A}
        \dsep x\in X\dsep y\in Y\dsep{}\\[-2pt]
      &u=f(x)\dsep v=f(y)\dsep z=u\diamond v
        \vdash z\in f[X\star_{\mathcal P}Y]
    \end{aligned}
  }{%
    \begin{aligned}[t]
      &\SemigroupHom{f}{A,\star}{B,\diamond}\dsep X,Y\in\PowNE{A}
        \dsep x\in X\dsep y\in Y\dsep{}\\[-2pt]
      &u=f(x)\dsep v=f(y)\dsep z=u\diamond v
        \vdash z\in f[X\star_{\mathcal P}Y]
    \end{aligned}
  }[PowerHomBackwardFixedWitnesses]
    \proofstepwidestar[1]{\SemigroupHom{f}{A,\star}{B,\diamond}}{\rA}
    \proofstepwidestar[2]{X\in\PowNE{A}}{\rA}
    \proofstepwidestar[3]{Y\in\PowNE{A}}{\rA}
    \proofstepwidestar[4]{x\in X}{\rA}
    \proofstepwidestar[5]{y\in Y}{\rA}
    \proofstepwidestar[6]{u=f(x)}{\rA}
    \proofstepwidestar[7]{v=f(y)}{\rA}
    \proofstepwidestar[8]{z=u\diamond v}{\rA}
    \proofstepwidestar[1]{\SemigroupAlg{A}{\star}}{%
      \FormulaRefAuto{SemigroupHomSourceAxiom}[ax]{1}}
    \proofstepwidestar[1]{f\colon A\to B}{%
      \FormulaRefAuto{SemigroupHomFunctionAxiom}[ax]{1}}
    \proofstepwidestar[2,4]{x\in A}{%
      \FormulaRefAuto{PowerSemigroupCarrierElement}[thm]{2,4}}
    \proofstepwidestar[3,5]{y\in A}{%
      \FormulaRefAuto{PowerSemigroupCarrierElement}[thm]{3,5}}
    \proofstepwidestar[1,2,3,4,5]{x\star y\in X\star_{\mathcal P}Y}{%
      \FormulaRefAuto{PowerSemigroupProductContains}[thm]{9,2,3,4,5}}
    \proofstepwidestar[1,2,3]{X\star_{\mathcal P}Y\subseteq A}{%
      \FormulaRefAuto{PowerSemigroupProductSubset}[thm]{9,2,3}}
    \proofstepwidestar[1,2,3,4,5]{f(x\star y)\in f[X\star_{\mathcal P}Y]}{%
      \FormulaRefAuto{C\subseteq A\dsep F\colon A\to B\dsep
        x\in C\vdash F(x)\in F[C]}[thm]{14,10,13}}
    \proofstepwidestar[1,2,3,4,5,6,7,8]{z=f(x\star y)}{%
      \FormulaRefAuto{PowerHomBackwardEquality}[thm]{1,11,12,6,7,8}}
    \proofstepwidestar[1,2,3,4,5,6,7,8]{z\in f[X\star_{\mathcal P}Y]}{%
      \rIE{16,15}}
  \closeproofpartwideindR

  \proofpartwideindR[Rückwärtsrichtung mit rechtem Bildzeugen]{%
    \begin{aligned}[t]
      &\SemigroupHom{f}{A,\star}{B,\diamond}\dsep X,Y\in\PowNE{A}
        \dsep x\in X\dsep u=f(x)\dsep v\in f[Y]\dsep{}\\[-2pt]
      &z=u\diamond v\vdash z\in f[X\star_{\mathcal P}Y]
    \end{aligned}
  }{%
    \begin{aligned}[t]
      &\SemigroupHom{f}{A,\star}{B,\diamond}\dsep X,Y\in\PowNE{A}
        \dsep x\in X\dsep u=f(x)\dsep v\in f[Y]\dsep{}\\[-2pt]
      &z=u\diamond v\vdash z\in f[X\star_{\mathcal P}Y]
    \end{aligned}
  }[PowerHomBackwardRightWitness]
    \proofstepwidestar[1]{\SemigroupHom{f}{A,\star}{B,\diamond}}{\rA}
    \proofstepwidestar[2]{X\in\PowNE{A}}{\rA}
    \proofstepwidestar[3]{Y\in\PowNE{A}}{\rA}
    \proofstepwidestar[4]{x\in X}{\rA}
    \proofstepwidestar[5]{u=f(x)}{\rA}
    \proofstepwidestar[6]{v\in f[Y]}{\rA}
    \proofstepwidestar[7]{z=u\diamond v}{\rA}
    \proofstepwidestar[1]{f\colon A\to B}{%
      \FormulaRefAuto{SemigroupHomFunctionAxiom}[ax]{1}}
    \proofstepwidestar[3]{Y\subseteq A}{%
      \rAEa{\FormulaRefAuto{NonemptyPowersetMembership}[thm]{3}}}
    \proofstepwidestar[1,3,6]{\exists y\in Y\,(v=f(y))}{%
      \FormulaRefAuto{C\subseteq A\dsep F\colon A\to B\dsep
        y\in F[C]\vdash\exists x\in C\,y=F(x)}[thm]{9,8,6}}
    \proofstepwidestar[11]{y\in Y\land v=f(y)}{\rA}
    \proofstepwidestar[11]{y\in Y}{\rAEa{11}}
    \proofstepwidestar[11]{v=f(y)}{\rAEb{11}}
    \proofstepwidestar[1,2,3,4,5,7,11]{z\in f[X\star_{\mathcal P}Y]}{%
      \FormulaRefAuto{PowerHomBackwardFixedWitnesses}[thm]{%
        1,2,3,4,12,5,13,7}}
    \proofstepwidestar[1,2,3,4,5,6,7]{z\in f[X\star_{\mathcal P}Y]}{%
      \rEE{10,11,14}}
  \closeproofpartwideindR

  \proofpartwideindR[Rückwärtsrichtung mit zwei Bildzeugen]{%
    \begin{aligned}[t]
      &\SemigroupHom{f}{A,\star}{B,\diamond}\dsep X,Y\in\PowNE{A}
        \dsep u\in f[X]\dsep v\in f[Y]\dsep z=u\diamond v\vdash{}\\[-2pt]
      &z\in f[X\star_{\mathcal P}Y]
    \end{aligned}
  }{%
    \begin{aligned}[t]
      &\SemigroupHom{f}{A,\star}{B,\diamond}\dsep X,Y\in\PowNE{A}
        \dsep u\in f[X]\dsep v\in f[Y]\dsep z=u\diamond v\vdash{}\\[-2pt]
      &z\in f[X\star_{\mathcal P}Y]
    \end{aligned}
  }[PowerHomBackwardImageWitnesses]
    \proofstepwidestar[1]{\SemigroupHom{f}{A,\star}{B,\diamond}}{\rA}
    \proofstepwidestar[2]{X\in\PowNE{A}}{\rA}
    \proofstepwidestar[3]{Y\in\PowNE{A}}{\rA}
    \proofstepwidestar[4]{u\in f[X]}{\rA}
    \proofstepwidestar[5]{v\in f[Y]}{\rA}
    \proofstepwidestar[6]{z=u\diamond v}{\rA}
    \proofstepwidestar[1]{f\colon A\to B}{%
      \FormulaRefAuto{SemigroupHomFunctionAxiom}[ax]{1}}
    \proofstepwidestar[2]{X\subseteq A}{%
      \rAEa{\FormulaRefAuto{NonemptyPowersetMembership}[thm]{2}}}
    \proofstepwidestar[1,2,4]{\exists x\in X\,(u=f(x))}{%
      \FormulaRefAuto{C\subseteq A\dsep F\colon A\to B\dsep
        y\in F[C]\vdash\exists x\in C\,y=F(x)}[thm]{8,7,4}}
    \proofstepwidestar[10]{x\in X\land u=f(x)}{\rA}
    \proofstepwidestar[10]{x\in X}{\rAEa{10}}
    \proofstepwidestar[10]{u=f(x)}{\rAEb{10}}
    \proofstepwidestar[1,2,3,5,6,10]{z\in f[X\star_{\mathcal P}Y]}{%
      \FormulaRefAuto{PowerHomBackwardRightWitness}[thm]{%
        1,2,3,11,12,5,6}}
    \proofstepwidestar[1,2,3,4,5,6]{z\in f[X\star_{\mathcal P}Y]}{%
      \rEE{9,10,13}}
  \closeproofpartwideindR

  \proofpartwideindR[Auflösung des zweiten Produktzeugen]{%
    \begin{aligned}[t]
      &\SemigroupHom{f}{A,\star}{B,\diamond}\dsep X,Y\in\PowNE{A}
        \dsep u\in f[X]\dsep{}\\[-2pt]
      &\exists v\in f[Y]\,(z=u\diamond v)
        \vdash z\in f[X\star_{\mathcal P}Y]
    \end{aligned}
  }{%
    \begin{aligned}[t]
      &\SemigroupHom{f}{A,\star}{B,\diamond}\dsep X,Y\in\PowNE{A}
        \dsep u\in f[X]\dsep{}\\[-2pt]
      &\exists v\in f[Y]\,(z=u\diamond v)
        \vdash z\in f[X\star_{\mathcal P}Y]
    \end{aligned}
  }[PowerHomBackwardSecondProductWitness]
    \proofstepwidestar[1]{\SemigroupHom{f}{A,\star}{B,\diamond}}{\rA}
    \proofstepwidestar[2]{X\in\PowNE{A}}{\rA}
    \proofstepwidestar[3]{Y\in\PowNE{A}}{\rA}
    \proofstepwidestar[4]{u\in f[X]}{\rA}
    \proofstepwidestar[5]{\exists v\in f[Y]\,(z=u\diamond v)}{\rA}
    \proofstepwidestar[6]{v\in f[Y]\land z=u\diamond v}{\rA}
    \proofstepwidestar[6]{v\in f[Y]}{\rAEa{6}}
    \proofstepwidestar[6]{z=u\diamond v}{\rAEb{6}}
    \proofstepwidestar[1,2,3,4,6]{z\in f[X\star_{\mathcal P}Y]}{%
      \FormulaRefAuto{PowerHomBackwardImageWitnesses}[thm]{1,2,3,4,7,8}}
    \proofstepwidestar[1,2,3,4,5]{z\in f[X\star_{\mathcal P}Y]}{%
      \rEE{5,6,9}}
  \closeproofpartwideindR

  \proofpartwideindR[Rückwärtsrichtung für Produktelemente]{%
    \begin{aligned}[t]
      &\SemigroupHom{f}{A,\star}{B,\diamond}\dsep X,Y\in\PowNE{A}
        \dsep z\in f[X]\diamond_{\mathcal P}f[Y]\vdash{}\\[-2pt]
      &z\in f[X\star_{\mathcal P}Y]
    \end{aligned}
  }{%
    \begin{aligned}[t]
      &\SemigroupHom{f}{A,\star}{B,\diamond}\dsep X,Y\in\PowNE{A}
        \dsep z\in f[X]\diamond_{\mathcal P}f[Y]\vdash{}\\[-2pt]
      &z\in f[X\star_{\mathcal P}Y]
    \end{aligned}
  }[PowerHomBackwardElement]
    \proofstepwidestar[1]{\SemigroupHom{f}{A,\star}{B,\diamond}}{\rA}
    \proofstepwidestar[2]{X\in\PowNE{A}}{\rA}
    \proofstepwidestar[3]{Y\in\PowNE{A}}{\rA}
    \proofstepwidestar[4]{z\in f[X]\diamond_{\mathcal P}f[Y]}{\rA}
    \proofstepwidestar[1]{\SemigroupAlg{B}{\diamond}}{%
      \FormulaRefAuto{SemigroupHomTargetAxiom}[ax]{1}}
    \proofstepwidestar[1]{f\colon A\to B}{%
      \FormulaRefAuto{SemigroupHomFunctionAxiom}[ax]{1}}
    \proofstepwidestar[1,2]{f[X]\in\PowNE{B}}{%
      \FormulaRefAuto{NonemptyDirectImage}[thm]{6,2}}
    \proofstepwidestar[1,3]{f[Y]\in\PowNE{B}}{%
      \FormulaRefAuto{NonemptyDirectImage}[thm]{6,3}}
    \proofstepwidestar[1,2,3,4]{%
      \exists u\in f[X]\,\exists v\in f[Y]\,(z=u\diamond v)}{%
      \FormulaRefAuto{PowerSemigroupProductWitnesses}[thm]{5,7,8,4}}
    \proofstepwidestar[10]{%
      u\in f[X]\land\exists v\in f[Y]\,(z=u\diamond v)}{\rA}
    \proofstepwidestar[10]{u\in f[X]}{\rAEa{10}}
    \proofstepwidestar[10]{\exists v\in f[Y]\,(z=u\diamond v)}{\rAEb{10}}
    \proofstepwidestar[1,2,3,10]{z\in f[X\star_{\mathcal P}Y]}{%
      \FormulaRefAuto{PowerHomBackwardSecondProductWitness}[thm]{%
        1,2,3,11,12}}
    \proofstepwidestar[1,2,3,4]{z\in f[X\star_{\mathcal P}Y]}{%
      \rEE{9,10,13}}
  \closeproofpartwideindR

  \proofpartwideindR[Zweite Teilmengenrichtung]{%
    \begin{aligned}[t]
      &\SemigroupHom{f}{A,\star}{B,\diamond}\dsep X,Y\in\PowNE{A}
        \vdash{}\\[-2pt]
      &f[X]\diamond_{\mathcal P}f[Y]
        \subseteq f[X\star_{\mathcal P}Y]
    \end{aligned}
  }{%
    \begin{aligned}[t]
      &\SemigroupHom{f}{A,\star}{B,\diamond}\dsep X,Y\in\PowNE{A}
        \vdash{}\\[-2pt]
      &f[X]\diamond_{\mathcal P}f[Y]
        \subseteq f[X\star_{\mathcal P}Y]
    \end{aligned}
  }[PowerHomBackwardSubset]
    \proofstepwidestar[1]{\SemigroupHom{f}{A,\star}{B,\diamond}}{\rA}
    \proofstepwidestar[2]{X\in\PowNE{A}}{\rA}
    \proofstepwidestar[3]{Y\in\PowNE{A}}{\rA}
    \proofstepwidestar[4]{z\in f[X]\diamond_{\mathcal P}f[Y]}{\rA}
    \proofstepwidestar[1,2,3,4]{z\in f[X\star_{\mathcal P}Y]}{%
      \FormulaRefAuto{PowerHomBackwardElement}[thm]{1,2,3,4}}
    \proofstepwidestar[1,2,3]{%
      \forall z\bigl(z\in f[X]\diamond_{\mathcal P}f[Y]\rightarrow
      z\in f[X\star_{\mathcal P}Y]\bigr)}{\rUI{\rRI{4,5}}}
    \proofstepwidestar[1,2,3]{%
      f[X]\diamond_{\mathcal P}f[Y]\subseteq
      f[X\star_{\mathcal P}Y]}{%
      \FormulaRefAuto{A\subseteq B\coloneq
        \forall x\,(x\in A\rightarrow x\in B)}[def]{6}}
  \closeproofpartwideindR

  \proofpartwide{Schluss}
    \proofstepwidestar[1]{\SemigroupHom{f}{A,\star}{B,\diamond}}{\rA}
    \proofstepwidestar[2]{X\in\PowNE{A}}{\rA}
    \proofstepwidestar[3]{Y\in\PowNE{A}}{\rA}
    \proofstepwidestar[1,2,3]{%
      f[X\star_{\mathcal P}Y]\subseteq
      f[X]\diamond_{\mathcal P}f[Y]}{%
      \FormulaRefAuto{PowerHomForwardSubset}[thm]{1,2,3}}
    \proofstepwidestar[1,2,3]{%
      f[X]\diamond_{\mathcal P}f[Y]\subseteq
      f[X\star_{\mathcal P}Y]}{%
      \FormulaRefAuto{PowerHomBackwardSubset}[thm]{1,2,3}}
    \proofstepwidestar[1,2,3]{%
      f[X\star_{\mathcal P}Y]=f[X]\diamond_{\mathcal P}f[Y]}{%
      \FormulaRefAuto{A\subseteq B, B\subseteq A\vdash A=B}[thm]{4,5}}
  \closeproofpartwide
\end{tabproofsplitwide}

\FormulaThmDeltaK[Operationserhaltung der induzierten Potenzabbildung]{%
  \begin{aligned}[t]
    &\SemigroupHom{f}{A,\star}{B,\diamond}\dsep
      X,Y\in\PowNE{A}\vdash{}\\[-2pt]
    &\PowMapNE{f}(X\star_{\mathcal P}Y)
      =\PowMapNE{f}(X)\diamond_{\mathcal P}\PowMapNE{f}(Y)
  \end{aligned}%
}{PowerSemigroupInducedMapOperation}{%
  \DeltaRow{Mengen}{A\dsep B\dsep X\dsep Y}
  \DeltaRow{Funktionen}{f\dsep\star\dsep\diamond\dsep
    \star_{\mathcal P}\dsep\diamond_{\mathcal P}\dsep
    \PowMap{f}\dsep\PowMapNE{f}}
}
\begin{tabproofwide}
  \proofstepwidestar[1]{\SemigroupHom{f}{A,\star}{B,\diamond}}{\rA}
  \proofstepwidestar[2]{X\in\PowNE{A}}{\rA}
  \proofstepwidestar[3]{Y\in\PowNE{A}}{\rA}
  \proofstepwidestar[1]{\SemigroupAlg{A}{\star}}{%
    \FormulaRefAuto{SemigroupHomSourceAxiom}[ax]{1}}
  \proofstepwidestar[1]{f\colon A\to B}{%
    \FormulaRefAuto{SemigroupHomFunctionAxiom}[ax]{1}}
  \proofstepwidestar[1,2,3]{X\star_{\mathcal P}Y\in\PowNE{A}}{%
    \FormulaRefAuto{PowerSemigroupProductClosure}[thm]{4,2,3}}
  \proofstepwidestar[1,2]{%
    \PowMapNE{f}(X)=f[X]}{%
    \FormulaRefAuto{NonemptyPowerMapValue}[thm]{5,2}}
  \proofstepwidestar[1,2]{%
    f[X]=\PowMapNE{f}(X)}{%
    \FormulaRefAuto{a=b\vdash b=a}[thm]{7}}
  \proofstepwidestar[1,3]{%
    \PowMapNE{f}(Y)=f[Y]}{%
    \FormulaRefAuto{NonemptyPowerMapValue}[thm]{5,3}}
  \proofstepwidestar[1,3]{%
    f[Y]=\PowMapNE{f}(Y)}{%
    \FormulaRefAuto{a=b\vdash b=a}[thm]{9}}
  \proofstepwide[1,2,3]{%
    \PowMapNE{f}(X\star_{\mathcal P}Y)}{=}{f[X\star_{\mathcal P}Y]}{%
    \FormulaRefAuto{NonemptyPowerMapValue}[thm]{5,6}}
  \proofstepwide[1,2,3]{}{=}{f[X]\diamond_{\mathcal P}f[Y]}{%
    \FormulaRefAuto{SemigroupHomPreservesPowerProduct}[thm]{1,2,3}}
  \proofstepwide[1,2,3]{}{=}{%
    \PowMapNE{f}(X)\diamond_{\mathcal P}f[Y]}{\rLRS{8}}
  \proofstepwide[1,2,3]{}{=}{%
    \PowMapNE{f}(X)\diamond_{\mathcal P}\PowMapNE{f}(Y)}{\rLRS{10}}
  \proofstepwide[1,2,3]{%
    \PowMapNE{f}(X\star_{\mathcal P}Y)}{=}{%
    \PowMapNE{f}(X)\diamond_{\mathcal P}\PowMapNE{f}(Y)}{\rChain{11,14}}
\end{tabproofwide}

\FormulaThmDeltaK[Isomorphismen induzieren Potenzhalbgruppenisomorphismen]{%
  \begin{aligned}[t]
    &\SemigroupIso{f}{A,\star}{B,\diamond}\vdash{}\\[-2pt]
    &\SemigroupIso{\PowMapNE{f}}
      {\PowNE{A},\star_{\mathcal P}}
      {\PowNE{B},\diamond_{\mathcal P}}
  \end{aligned}%
}{SemigroupIsoInducesPowerSemigroupIso}{%
  \DeltaRow{Mengen}{A\dsep B\dsep X\dsep Y}
  \DeltaRow{Funktionen}{f\dsep\star\dsep\diamond\dsep
    \star_{\mathcal P}\dsep\diamond_{\mathcal P}\dsep
    \PowMap{f}\dsep\PowMapNE{f}}
}
\begin{tabproofwide}
  \proofstepwidestar[1]{\SemigroupIso{f}{A,\star}{B,\diamond}}{\rA}
  \proofstepwidestar[1]{\SemigroupHom{f}{A,\star}{B,\diamond}}{%
    \FormulaRefAuto{SemigroupIsoHomAxiom}[ax]{1}}
  \proofstepwidestar[1]{f\colon A\bij B}{%
    \FormulaRefAuto{SemigroupIsoBijectiveAxiom}[ax]{1}}
  \proofstepwidestar[1]{\SemigroupAlg{A}{\star}}{%
    \FormulaRefAuto{SemigroupHomSourceAxiom}[ax]{2}}
  \proofstepwidestar[1]{\SemigroupAlg{B}{\diamond}}{%
    \FormulaRefAuto{SemigroupHomTargetAxiom}[ax]{2}}
  \proofstepwidestar[1]{\SemigroupAlg{\PowNE{A}}{\star_{\mathcal P}}}{%
    \FormulaRefAuto{NonemptyPowerSemigroup}[thm]{4}}
  \proofstepwidestar[1]{\SemigroupAlg{\PowNE{B}}{\diamond_{\mathcal P}}}{%
    \FormulaRefAuto{NonemptyPowerSemigroup}[thm]{5}}
  \proofstepwidestar[1]{%
    \PowMapNE{f}\colon\PowNE{A}\bij\PowNE{B}}{%
    \FormulaRefAuto{NonemptyPowerMapBijective}[thm]{3}}
  \proofstepwidestar[9]{X\in\PowNE{A}}{\rA}
  \proofstepwidestar[10]{Y\in\PowNE{A}}{\rA}
  \proofstepwidestar[1,9,10]{%
    \PowMapNE{f}(X\star_{\mathcal P}Y)
      =\PowMapNE{f}(X)\diamond_{\mathcal P}\PowMapNE{f}(Y)}{%
    \FormulaRefAuto{PowerSemigroupInducedMapOperation}[thm]{2,9,10}}
  \proofstepwidestar[1]{%
    \PowMapNE{f}\colon\PowNE{A}\to\PowNE{B}}{%
    \FormulaRefAuto{Bijektive Funktion}[def]{8}}
  \proofstepwidestar[1]{%
    \SemigroupHom{\PowMapNE{f}}
      {\PowNE{A},\star_{\mathcal P}}
      {\PowNE{B},\diamond_{\mathcal P}}}{%
    \FormulaRefAuto{SemigroupHomDef}[def]{6,7,12,11}}
  \proofstepwidestar[1]{%
    \SemigroupIso{\PowMapNE{f}}
      {\PowNE{A},\star_{\mathcal P}}
      {\PowNE{B},\diamond_{\mathcal P}}}{%
    \FormulaRefAuto{SemigroupIsoDef}[def]{13,8}}
\end{tabproofwide}

\section{Das Isomorphieproblem für Potenzhalbgruppen}

Zwei Halbgruppen heißen \emph{global isomorph}, wenn ihre
Potenzhalbgruppen isomorph sind.  Die folgende Richtung gilt ohne jede
Endlichkeitsvoraussetzung.

\FormulaThmDeltaK[Die isomorphe Richtung des Potenzhalbgruppenproblems]{%
  \begin{aligned}[t]
    &\exists f\,\SemigroupIso{f}{A,\star}{B,\diamond}\vdash{}\\[-2pt]
    &\exists F\,\SemigroupIso{F}
      {\PowNE{A},\star_{\mathcal P}}
      {\PowNE{B},\diamond_{\mathcal P}}
  \end{aligned}%
}{SemigroupIsoImpliesPowerSemigroupIso}{%
  \DeltaRow{Mengen}{A\dsep B}
  \DeltaRow{Funktionen}{F\dsep f\dsep\star\dsep\diamond\dsep
    \star_{\mathcal P}\dsep\diamond_{\mathcal P}\dsep
    \PowMap{f}\dsep\PowMapNE{f}}
}
\begin{tabproofwide}
  \proofstepwidestar[1]{\exists f\,\SemigroupIso{f}{A,\star}{B,\diamond}}{\rA}
  \proofstepwidestar[2]{\SemigroupIso{f}{A,\star}{B,\diamond}}{\rA}
  \proofstepwidestar[2]{%
    \SemigroupIso{\PowMapNE{f}}
      {\PowNE{A},\star_{\mathcal P}}
      {\PowNE{B},\diamond_{\mathcal P}}}{%
    \FormulaRefAuto{SemigroupIsoInducesPowerSemigroupIso}[thm]{2}}
  \proofstepwidestar[2]{%
    \exists F\,\SemigroupIso{F}
      {\PowNE{A},\star_{\mathcal P}}
      {\PowNE{B},\diamond_{\mathcal P}}}{\rEI{3}}
  \proofstepwidestar[1]{%
    \exists F\,\SemigroupIso{F}
      {\PowNE{A},\star_{\mathcal P}}
      {\PowNE{B},\diamond_{\mathcal P}}}{\rEE{1,2,4}}
\end{tabproofwide}

\subsection{Inflationen und translationelle Ununterscheidbarkeit}

Eine Erweiterung kann neue Elemente enthalten, ohne neue Multiplikationswerte
zu erzeugen.  Der folgende Begriff fasst diese Situation zusammen.  Die
Abbildung \(r\) zieht jedes neue Element auf ein altes Element zurück; die
Multiplikation der Erweiterung hängt nur von diesen Rückzugswerten ab.

\FormulaDefDeltaK[Inflation einer Halbgruppe]{%
  \mathsf{Infl}_{r}\bigl((A,\star),(B,\diamond)\bigr)%
}{SemigroupInflationDef}{%
  \DeltaRow{Mengen}{A\dsep B\dsep x\dsep y\dsep a}
  \DeltaRow{Funktionen}{r\dsep\star\dsep\diamond}
  \DeltaRow{\textbf{Axiome}}{}
  \DeltaPrem{Quellhalbgruppe}{\SemigroupAlg{A}{\star}}
  \DeltaPrem{Zielhalbgruppe}{\SemigroupAlg{B}{\diamond}}
  \DeltaPrem{Erweiterung}{A\subseteq B}
  \DeltaPrem{Rückzug}{r\colon B\to A}
  \DeltaRow{Retraktion}{a\in A\vdash r(a)=a}
  \DeltaRow{Produkt über dem Rückzug}{%
    x,y\in B\vdash x\diamond y=r(x)\star r(y)}
  \DeltaRow{\textbf{Neues Symbol}}{}
  \DeltaPrem{Inflation}{%
    \mathsf{Infl}_{r}\bigl((A,\star),(B,\diamond)\bigr)}
}

Für die nachfolgende Anwendung ist nicht die Gleichheit zweier Elemente
entscheidend, sondern die Gleichheit ihrer Links- und Rechtstranslationen.

\FormulaDefDeltaK[Translationelle Ununterscheidbarkeit]{%
  \begin{aligned}[t]
    &\SemigroupAlg{B}{\diamond}\dsep x,y\in B\vdash{}\\[-2pt]
    &x\mathrel{\equiv^{\mathrm{tr}}_{B,\diamond}}y
      \coloneqq
      \forall z\in B\,\bigl(
        x\diamond z=y\diamond z\land
        z\diamond x=z\diamond y\bigr)
  \end{aligned}%
}{SemigroupTranslationIndistinguishableDef}{%
  \DeltaRow{Mengen}{B\dsep x\dsep y\dsep z}
  \DeltaRow{Funktionen}{\diamond}
  \DeltaRow{Relationsoperatoren}{%
    \equiv^{\mathrm{tr}}_{B,\diamond}}
}

\FormulaThmDeltaK[Zugriff auf die gemeinsamen Translationen]{%
  \begin{aligned}[t]
    &\SemigroupAlg{B}{\diamond}\dsep
      x,y\in B\dsep
      x\mathrel{\equiv^{\mathrm{tr}}_{B,\diamond}}y\dsep z\in B
      \vdash{}\\[-2pt]
    &x\diamond z=y\diamond z\land z\diamond x=z\diamond y
  \end{aligned}%
}{SemigroupTranslationIndistinguishableActions}{%
  \DeltaRow{Mengen}{B\dsep x\dsep y\dsep z}
  \DeltaRow{Funktionen}{\diamond}
}
\begin{tabproofwide}
  \proofstepwidestar[1]{\SemigroupAlg{B}{\diamond}}{\rA}
  \proofstepwidestar[2]{x\in B}{\rA}
  \proofstepwidestar[3]{y\in B}{\rA}
  \proofstepwidestar[4]{%
    x\mathrel{\equiv^{\mathrm{tr}}_{B,\diamond}}y}{\rA}
  \proofstepwidestar[5]{z\in B}{\rA}
  \proofstepwidestar[1,2,3,4]{%
    \forall u\in B\,\bigl(
      x\diamond u=y\diamond u\land
      u\diamond x=u\diamond y\bigr)}{%
    \FormulaRefAuto{SemigroupTranslationIndistinguishableDef}[def]{1,2,3,4}}
  \proofstepwidestar[1,2,3,4,5]{%
    x\diamond z=y\diamond z\land z\diamond x=z\diamond y}{%
    \rRE{5,\rUE{6}}}
\end{tabproofwide}

Zunächst halten wir die für die Anwendung benötigte Reflexivität des
Translationsprofils fest.

\FormulaThmDeltaK[Reflexivität der translationellen Ununterscheidbarkeit]{%
  \SemigroupAlg{B}{\diamond}\dsep x\in B\vdash
  x\mathrel{\equiv^{\mathrm{tr}}_{B,\diamond}}x%
}{SemigroupTranslationIndistinguishableReflexive}{%
  \DeltaRow{Mengen}{B\dsep x\dsep z}
  \DeltaRow{Funktionen}{\diamond}
  \DeltaRow{Relationsoperatoren}{\equiv^{\mathrm{tr}}_{B,\diamond}}
}
\begin{tabproofwide}
  \proofstepwidestar[1]{\SemigroupAlg{B}{\diamond}}{\rA}
  \proofstepwidestar[2]{x\in B}{\rA}
  \proofstepwidestar[3]{z\in B}{\rA}
  \proofstepwidestar[3]{x\diamond z=x\diamond z}{\rII}
  \proofstepwidestar[3]{z\diamond x=z\diamond x}{\rII}
  \proofstepwidestar[3]{%
    x\diamond z=x\diamond z\land z\diamond x=z\diamond x}{\rAI{4,5}}
  \proofstepwidestar[]{%
    \forall z\in B\,\bigl(
      x\diamond z=x\diamond z\land z\diamond x=z\diamond x\bigr)}{%
    \rUI{\rRI{3,6}}}
  \proofstepwidestar[1,2]{%
    x\mathrel{\equiv^{\mathrm{tr}}_{B,\diamond}}x}{%
    \FormulaRefAuto{SemigroupTranslationIndistinguishableDef}[def]{1,2,2,7}}
\end{tabproofwide}

Der folgende Satz ist das algebraische Bindeglied zur späteren
Reservefamilie.  Er verlangt nicht, dass die Bijektion die einzelnen
Elemente festhält.  Es genügt, sie nur innerhalb ihrer Translationsprofile
zu verschieben und sämtliche wirklichen Produktwerte festzulassen.

\FormulaThmDeltaK[Isomorphiekriterium durch Fixierung der Produktwerte]{%
  \begin{aligned}[t]
    &\SemigroupAlg{A}{\star}\dsep\SemigroupAlg{B}{\diamond}\dsep
      A\subseteq B\dsep F\colon A\bij B\dsep{}\\[-2pt]
    &\forall x,y\in A\,(x\star y=x\diamond y)\dsep{}\\[-2pt]
    &\forall x\in A\,
      \bigl(x\mathrel{\equiv^{\mathrm{tr}}_{B,\diamond}}F(x)\bigr)
      \dsep{}\\[-2pt]
    &\forall x,y\in A\,F(x\star y)=x\star y
      \vdash \SemigroupIso{F}{A,\star}{B,\diamond}
  \end{aligned}%
}{SemigroupFixedProductsIsomorphismCriterion}{%
  \DeltaRow{Mengen}{A\dsep B\dsep x\dsep y}
  \DeltaRow{Funktionen}{F\dsep\star\dsep\diamond}
}
\begin{tabproofwide}
  \proofstepwidestar[1]{\SemigroupAlg{A}{\star}}{\rA}
  \proofstepwidestar[2]{\SemigroupAlg{B}{\diamond}}{\rA}
  \proofstepwidestar[3]{A\subseteq B}{\rA}
  \proofstepwidestar[4]{F\colon A\bij B}{\rA}
  \proofstepwidestar[5]{\forall u,v\in A\,(u\star v=u\diamond v)}{\rA}
  \proofstepwidestar[6]{\forall u\in A\,
    \bigl(u\mathrel{\equiv^{\mathrm{tr}}_{B,\diamond}}F(u)\bigr)}{\rA}
  \proofstepwidestar[7]{\forall u,v\in A\,F(u\star v)=u\star v}{\rA}
  \proofstepwidestar[4]{F\colon A\to B}{%
    \FormulaRefAuto{Bijektive Funktion}[def]{4}}
  \proofstepwidestar[9]{x\in A}{\rA}
  \proofstepwidestar[10]{y\in A}{\rA}
  \proofstepwidestar[3,9]{x\in B}{%
    \FormulaRefAuto{A\subseteq B\dsep x\in A\vdash x\in B}[thm]{3,9}}
  \proofstepwidestar[3,10]{y\in B}{%
    \FormulaRefAuto{A\subseteq B\dsep x\in A\vdash x\in B}[thm]{3,10}}
  \proofstepwidestar[4,9]{F(x)\in B}{%
    \FormulaRefAuto{F\colon A\to B\dsep x\in A\vdash F(x)\in B}[thm]{8,9}}
  \proofstepwidestar[4,10]{F(y)\in B}{%
    \FormulaRefAuto{F\colon A\to B\dsep x\in A\vdash F(x)\in B}[thm]{8,10}}
  \proofstepwidestar[6,9]{%
    x\mathrel{\equiv^{\mathrm{tr}}_{B,\diamond}}F(x)}{\rRE{9,\rUE{6}}}
  \proofstepwidestar[6,10]{%
    y\mathrel{\equiv^{\mathrm{tr}}_{B,\diamond}}F(y)}{\rRE{10,\rUE{6}}}
  \proofstepwidestar[2,3,4,6,9,10]{%
    x\diamond F(y)=F(x)\diamond F(y)}{%
    \rAEa{\FormulaRefAuto{SemigroupTranslationIndistinguishableActions}[thm]{%
      2,11,13,15,14}}}
  \proofstepwidestar[2,3,4,6,9,10]{%
    F(x)\diamond F(y)=x\diamond F(y)}{%
    \FormulaRefAuto{a=b\vdash b=a}[thm]{17}}
  \proofstepwidestar[2,3,4,6,9,10]{x\diamond y=x\diamond F(y)}{%
    \rAEb{\FormulaRefAuto{SemigroupTranslationIndistinguishableActions}[thm]{%
      2,12,14,16,11}}}
  \proofstepwidestar[2,3,4,6,9,10]{x\diamond F(y)=x\diamond y}{%
    \FormulaRefAuto{a=b\vdash b=a}[thm]{19}}
  \proofstepwidestar[5,9,10]{x\star y=x\diamond y}{%
    \rRE{10,\rUE{\rRE{9,\rUE{5}}}}}
  \proofstepwidestar[5,9,10]{x\diamond y=x\star y}{%
    \FormulaRefAuto{a=b\vdash b=a}[thm]{21}}
  \proofstepwidestar[7,9,10]{F(x\star y)=x\star y}{%
    \rRE{10,\rUE{\rRE{9,\rUE{7}}}}}
  \proofstepwidestar[7,9,10]{x\star y=F(x\star y)}{%
    \FormulaRefAuto{a=b\vdash b=a}[thm]{23}}
  \proofstepwidestar[2,3,4,5,6,7,9,10]{%
    F(x\star y)=F(x)\diamond F(y)}{%
    \rChain{23,21,19,17}}
  \proofstepwidestar[1,2,3,4,5,6,7]{%
    \SemigroupHom{F}{A,\star}{B,\diamond}}{%
    \FormulaRefAuto{SemigroupHomDef}[def]{1,2,8,25}}
  \proofstepwidestar[1,2,3,4,5,6,7]{%
    \SemigroupIso{F}{A,\star}{B,\diamond}}{%
    \FormulaRefAuto{SemigroupIsoDef}[def]{26,4}}
\end{tabproofwide}

\FormulaThmDeltaK[Gleiche Translationsprofile in der Potenzhalbgruppe]{%
  \begin{aligned}[t]
    &\SemigroupAlg{B}{\diamond}\dsep X,Y\in\PowNE{B}\dsep{}\\[-2pt]
    &\forall x\in X\,\exists y\in Y\,
      \bigl(x\mathrel{\equiv^{\mathrm{tr}}_{B,\diamond}}y\bigr)
      \dsep{}\\[-2pt]
    &\forall y\in Y\,\exists x\in X\,
      \bigl(y\mathrel{\equiv^{\mathrm{tr}}_{B,\diamond}}x\bigr)
      \vdash
      X\mathrel{\equiv^{\mathrm{tr}}_{\PowNE{B},
        \diamond_{\mathcal P}}}Y
  \end{aligned}%
}{PowerSemigroupTranslationSupportCriterion}{%
  \DeltaRow{Mengen}{B\dsep X\dsep Y\dsep Z\dsep
    x\dsep y\dsep z\dsep w}
  \DeltaRow{Funktionen}{\diamond\dsep\diamond_{\mathcal P}}
}
\begin{tabproofsplitwide}
  \proofpartwideindR[Rechte Mengenprodukte]{%
    \begin{aligned}[t]
      &\SemigroupAlg{B}{\diamond}\dsep X,Y,Z\in\PowNE{B}\dsep{}\\[-2pt]
      &\forall x\in X\,\exists y\in Y\,
        \bigl(x\mathrel{\equiv^{\mathrm{tr}}_{B,\diamond}}y\bigr)
        \dsep{}\\[-2pt]
      &\forall y\in Y\,\exists x\in X\,
        \bigl(y\mathrel{\equiv^{\mathrm{tr}}_{B,\diamond}}x\bigr)
        \vdash X\diamond_{\mathcal P}Z=Y\diamond_{\mathcal P}Z
    \end{aligned}%
  }{%
    \begin{aligned}[t]
      &\SemigroupAlg{B}{\diamond}\dsep X,Y,Z\in\PowNE{B}\dsep{}\\[-2pt]
      &\forall x\in X\,\exists y\in Y\,
        \bigl(x\mathrel{\equiv^{\mathrm{tr}}_{B,\diamond}}y\bigr)
        \dsep{}\\[-2pt]
      &\forall y\in Y\,\exists x\in X\,
        \bigl(y\mathrel{\equiv^{\mathrm{tr}}_{B,\diamond}}x\bigr)
        \vdash X\diamond_{\mathcal P}Z=Y\diamond_{\mathcal P}Z
    \end{aligned}%
  }[PowerSemigroupTranslationSupportRightProduct]
    \proofstepwidestar[1]{\SemigroupAlg{B}{\diamond}}{\rA}
    \proofstepwidestar[2]{X\in\PowNE{B}}{\rA}
    \proofstepwidestar[3]{Y\in\PowNE{B}}{\rA}
    \proofstepwidestar[4]{Z\in\PowNE{B}}{\rA}
    \proofstepwidestar[5]{\forall u\in X\,\exists v\in Y\,
      \bigl(u\mathrel{\equiv^{\mathrm{tr}}_{B,\diamond}}v\bigr)}{\rA}
    \proofstepwidestar[6]{\forall v\in Y\,\exists u\in X\,
      \bigl(v\mathrel{\equiv^{\mathrm{tr}}_{B,\diamond}}u\bigr)}{\rA}

    \proofstepwidestar[7]{w\in X\diamond_{\mathcal P}Z}{\rA}
    \proofstepwidestar[1,2,4,7]{%
      \exists x\in X\,\exists z\in Z\,(w=x\diamond z)}{%
      \FormulaRefAuto{PowerSemigroupProductWitnesses}[thm]{1,2,4,7}}
    \proofstepwidestar[8]{x\in X}{\rA}
    \proofstepwidestar[8]{z\in Z}{\rA}
    \proofstepwidestar[8]{w=x\diamond z}{\rA}
    \proofstepwidestar[5,8]{\exists y\in Y\,
      \bigl(x\mathrel{\equiv^{\mathrm{tr}}_{B,\diamond}}y\bigr)}{%
      \rRE{9,\rUE{5}}}
    \proofstepwidestar[12]{y\in Y}{\rA}
    \proofstepwidestar[12]{%
      x\mathrel{\equiv^{\mathrm{tr}}_{B,\diamond}}y}{\rA}
    \proofstepwidestar[2,8]{x\in B}{%
      \FormulaRefAuto{PowerSemigroupCarrierElement}[thm]{2,9}}
    \proofstepwidestar[3,12]{y\in B}{%
      \FormulaRefAuto{PowerSemigroupCarrierElement}[thm]{3,13}}
    \proofstepwidestar[4,8]{z\in B}{%
      \FormulaRefAuto{PowerSemigroupCarrierElement}[thm]{4,10}}
    \proofstepwidestar[1,2,3,4,8,12]{x\diamond z=y\diamond z}{%
      \rAEa{\FormulaRefAuto{SemigroupTranslationIndistinguishableActions}[thm]{%
        1,15,16,14,17}}}
    \proofstepwidestar[1,2,3,4,8,12]{w=y\diamond z}{\rChain{11,18}}
    \proofstepwidestar[1,3,4,8,12]{%
      y\diamond z\in Y\diamond_{\mathcal P}Z}{%
      \FormulaRefAuto{PowerSemigroupProductContains}[thm]{1,3,4,13,10}}
    \proofstepwidestar[1,2,3,4,8,12]{%
      w\in Y\diamond_{\mathcal P}Z}{\rIE{19,20}}
    \proofstepwidestar[1,2,3,4,5,8]{%
      w\in Y\diamond_{\mathcal P}Z}{\rEEStar{12\text{--}21}}
    \proofstepwidestar[1,2,3,4,5,7]{%
      w\in Y\diamond_{\mathcal P}Z}{\rEEStar{8\text{--}22}}
    \proofstepwidestar[1,2,3,4,5]{%
      \forall w\bigl(w\in X\diamond_{\mathcal P}Z\rightarrow
        w\in Y\diamond_{\mathcal P}Z\bigr)}{\rUI{\rRI{7,23}}}
    \proofstepwidestar[1,2,3,4,5]{%
      X\diamond_{\mathcal P}Z\subseteq Y\diamond_{\mathcal P}Z}{%
      \FormulaRefAuto{A\subseteq B\coloneq
        \forall x\,(x\in A\rightarrow x\in B)}[def]{24}}

    \proofstepwidestar[26]{w\in Y\diamond_{\mathcal P}Z}{\rA}
    \proofstepwidestar[1,3,4,26]{%
      \exists y\in Y\,\exists z\in Z\,(w=y\diamond z)}{%
      \FormulaRefAuto{PowerSemigroupProductWitnesses}[thm]{1,3,4,26}}
    \proofstepwidestar[27]{y\in Y}{\rA}
    \proofstepwidestar[27]{z\in Z}{\rA}
    \proofstepwidestar[27]{w=y\diamond z}{\rA}
    \proofstepwidestar[6,27]{\exists x\in X\,
      \bigl(y\mathrel{\equiv^{\mathrm{tr}}_{B,\diamond}}x\bigr)}{%
      \rRE{28,\rUE{6}}}
    \proofstepwidestar[31]{x\in X}{\rA}
    \proofstepwidestar[31]{%
      y\mathrel{\equiv^{\mathrm{tr}}_{B,\diamond}}x}{\rA}
    \proofstepwidestar[3,27]{y\in B}{%
      \FormulaRefAuto{PowerSemigroupCarrierElement}[thm]{3,28}}
    \proofstepwidestar[2,31]{x\in B}{%
      \FormulaRefAuto{PowerSemigroupCarrierElement}[thm]{2,32}}
    \proofstepwidestar[4,27]{z\in B}{%
      \FormulaRefAuto{PowerSemigroupCarrierElement}[thm]{4,29}}
    \proofstepwidestar[1,2,3,4,27,31]{y\diamond z=x\diamond z}{%
      \rAEa{\FormulaRefAuto{SemigroupTranslationIndistinguishableActions}[thm]{%
        1,34,35,33,36}}}
    \proofstepwidestar[1,2,3,4,27,31]{w=x\diamond z}{\rChain{30,37}}
    \proofstepwidestar[1,2,4,27,31]{%
      x\diamond z\in X\diamond_{\mathcal P}Z}{%
      \FormulaRefAuto{PowerSemigroupProductContains}[thm]{1,2,4,32,29}}
    \proofstepwidestar[1,2,3,4,27,31]{%
      w\in X\diamond_{\mathcal P}Z}{\rIE{38,39}}
    \proofstepwidestar[1,2,3,4,6,27]{%
      w\in X\diamond_{\mathcal P}Z}{\rEEStar{31\text{--}40}}
    \proofstepwidestar[1,2,3,4,6,26]{%
      w\in X\diamond_{\mathcal P}Z}{\rEEStar{27\text{--}41}}
    \proofstepwidestar[1,2,3,4,6]{%
      \forall w\bigl(w\in Y\diamond_{\mathcal P}Z\rightarrow
        w\in X\diamond_{\mathcal P}Z\bigr)}{\rUI{\rRI{26,42}}}
    \proofstepwidestar[1,2,3,4,6]{%
      Y\diamond_{\mathcal P}Z\subseteq X\diamond_{\mathcal P}Z}{%
      \FormulaRefAuto{A\subseteq B\coloneq
        \forall x\,(x\in A\rightarrow x\in B)}[def]{43}}
    \proofstepwidestar[1,2,3,4,5,6]{%
      X\diamond_{\mathcal P}Z=Y\diamond_{\mathcal P}Z}{%
      \FormulaRefAuto{A\subseteq B\dsep B\subseteq A\vdash A=B}[thm]{25,44}}
  \closeproofpartwideindR

  \proofpartwideindR[Linke Mengenprodukte]{%
    \begin{aligned}[t]
      &\SemigroupAlg{B}{\diamond}\dsep X,Y,Z\in\PowNE{B}\dsep{}\\[-2pt]
      &\forall x\in X\,\exists y\in Y\,
        \bigl(x\mathrel{\equiv^{\mathrm{tr}}_{B,\diamond}}y\bigr)
        \dsep{}\\[-2pt]
      &\forall y\in Y\,\exists x\in X\,
        \bigl(y\mathrel{\equiv^{\mathrm{tr}}_{B,\diamond}}x\bigr)
        \vdash Z\diamond_{\mathcal P}X=Z\diamond_{\mathcal P}Y
    \end{aligned}%
  }{%
    \begin{aligned}[t]
      &\SemigroupAlg{B}{\diamond}\dsep X,Y,Z\in\PowNE{B}\dsep{}\\[-2pt]
      &\forall x\in X\,\exists y\in Y\,
        \bigl(x\mathrel{\equiv^{\mathrm{tr}}_{B,\diamond}}y\bigr)
        \dsep{}\\[-2pt]
      &\forall y\in Y\,\exists x\in X\,
        \bigl(y\mathrel{\equiv^{\mathrm{tr}}_{B,\diamond}}x\bigr)
        \vdash Z\diamond_{\mathcal P}X=Z\diamond_{\mathcal P}Y
    \end{aligned}%
  }[PowerSemigroupTranslationSupportLeftProduct]
    \proofstepwidestar[1]{\SemigroupAlg{B}{\diamond}}{\rA}
    \proofstepwidestar[2]{X\in\PowNE{B}}{\rA}
    \proofstepwidestar[3]{Y\in\PowNE{B}}{\rA}
    \proofstepwidestar[4]{Z\in\PowNE{B}}{\rA}
    \proofstepwidestar[5]{\forall u\in X\,\exists v\in Y\,
      \bigl(u\mathrel{\equiv^{\mathrm{tr}}_{B,\diamond}}v\bigr)}{\rA}
    \proofstepwidestar[6]{\forall v\in Y\,\exists u\in X\,
      \bigl(v\mathrel{\equiv^{\mathrm{tr}}_{B,\diamond}}u\bigr)}{\rA}

    \proofstepwidestar[7]{w\in Z\diamond_{\mathcal P}X}{\rA}
    \proofstepwidestar[1,2,4,7]{%
      \exists z\in Z\,\exists x\in X\,(w=z\diamond x)}{%
      \FormulaRefAuto{PowerSemigroupProductWitnesses}[thm]{1,4,2,7}}
    \proofstepwidestar[8]{z\in Z}{\rA}
    \proofstepwidestar[8]{x\in X}{\rA}
    \proofstepwidestar[8]{w=z\diamond x}{\rA}
    \proofstepwidestar[5,8]{\exists y\in Y\,
      \bigl(x\mathrel{\equiv^{\mathrm{tr}}_{B,\diamond}}y\bigr)}{%
      \rRE{10,\rUE{5}}}
    \proofstepwidestar[12]{y\in Y}{\rA}
    \proofstepwidestar[12]{%
      x\mathrel{\equiv^{\mathrm{tr}}_{B,\diamond}}y}{\rA}
    \proofstepwidestar[2,8]{x\in B}{%
      \FormulaRefAuto{PowerSemigroupCarrierElement}[thm]{2,10}}
    \proofstepwidestar[3,12]{y\in B}{%
      \FormulaRefAuto{PowerSemigroupCarrierElement}[thm]{3,13}}
    \proofstepwidestar[4,8]{z\in B}{%
      \FormulaRefAuto{PowerSemigroupCarrierElement}[thm]{4,9}}
    \proofstepwidestar[1,2,3,4,8,12]{z\diamond x=z\diamond y}{%
      \rAEb{\FormulaRefAuto{SemigroupTranslationIndistinguishableActions}[thm]{%
        1,15,16,14,17}}}
    \proofstepwidestar[1,2,3,4,8,12]{w=z\diamond y}{\rChain{11,18}}
    \proofstepwidestar[1,3,4,8,12]{%
      z\diamond y\in Z\diamond_{\mathcal P}Y}{%
      \FormulaRefAuto{PowerSemigroupProductContains}[thm]{1,4,3,9,13}}
    \proofstepwidestar[1,2,3,4,8,12]{%
      w\in Z\diamond_{\mathcal P}Y}{\rIE{19,20}}
    \proofstepwidestar[1,2,3,4,5,8]{%
      w\in Z\diamond_{\mathcal P}Y}{\rEEStar{12\text{--}21}}
    \proofstepwidestar[1,2,3,4,5,7]{%
      w\in Z\diamond_{\mathcal P}Y}{\rEEStar{8\text{--}22}}
    \proofstepwidestar[1,2,3,4,5]{%
      \forall w\bigl(w\in Z\diamond_{\mathcal P}X\rightarrow
        w\in Z\diamond_{\mathcal P}Y\bigr)}{\rUI{\rRI{7,23}}}
    \proofstepwidestar[1,2,3,4,5]{%
      Z\diamond_{\mathcal P}X\subseteq Z\diamond_{\mathcal P}Y}{%
      \FormulaRefAuto{A\subseteq B\coloneq
        \forall x\,(x\in A\rightarrow x\in B)}[def]{24}}

    \proofstepwidestar[26]{w\in Z\diamond_{\mathcal P}Y}{\rA}
    \proofstepwidestar[1,3,4,26]{%
      \exists z\in Z\,\exists y\in Y\,(w=z\diamond y)}{%
      \FormulaRefAuto{PowerSemigroupProductWitnesses}[thm]{1,4,3,26}}
    \proofstepwidestar[27]{z\in Z}{\rA}
    \proofstepwidestar[27]{y\in Y}{\rA}
    \proofstepwidestar[27]{w=z\diamond y}{\rA}
    \proofstepwidestar[6,27]{\exists x\in X\,
      \bigl(y\mathrel{\equiv^{\mathrm{tr}}_{B,\diamond}}x\bigr)}{%
      \rRE{29,\rUE{6}}}
    \proofstepwidestar[31]{x\in X}{\rA}
    \proofstepwidestar[31]{%
      y\mathrel{\equiv^{\mathrm{tr}}_{B,\diamond}}x}{\rA}
    \proofstepwidestar[3,27]{y\in B}{%
      \FormulaRefAuto{PowerSemigroupCarrierElement}[thm]{3,29}}
    \proofstepwidestar[2,31]{x\in B}{%
      \FormulaRefAuto{PowerSemigroupCarrierElement}[thm]{2,32}}
    \proofstepwidestar[4,27]{z\in B}{%
      \FormulaRefAuto{PowerSemigroupCarrierElement}[thm]{4,28}}
    \proofstepwidestar[1,2,3,4,27,31]{z\diamond y=z\diamond x}{%
      \rAEb{\FormulaRefAuto{SemigroupTranslationIndistinguishableActions}[thm]{%
        1,34,35,33,36}}}
    \proofstepwidestar[1,2,3,4,27,31]{w=z\diamond x}{\rChain{30,37}}
    \proofstepwidestar[1,2,4,27,31]{%
      z\diamond x\in Z\diamond_{\mathcal P}X}{%
      \FormulaRefAuto{PowerSemigroupProductContains}[thm]{1,4,2,28,32}}
    \proofstepwidestar[1,2,3,4,27,31]{%
      w\in Z\diamond_{\mathcal P}X}{\rIE{38,39}}
    \proofstepwidestar[1,2,3,4,6,27]{%
      w\in Z\diamond_{\mathcal P}X}{\rEEStar{31\text{--}40}}
    \proofstepwidestar[1,2,3,4,6,26]{%
      w\in Z\diamond_{\mathcal P}X}{\rEEStar{27\text{--}41}}
    \proofstepwidestar[1,2,3,4,6]{%
      \forall w\bigl(w\in Z\diamond_{\mathcal P}Y\rightarrow
        w\in Z\diamond_{\mathcal P}X\bigr)}{\rUI{\rRI{26,42}}}
    \proofstepwidestar[1,2,3,4,6]{%
      Z\diamond_{\mathcal P}Y\subseteq Z\diamond_{\mathcal P}X}{%
      \FormulaRefAuto{A\subseteq B\coloneq
        \forall x\,(x\in A\rightarrow x\in B)}[def]{43}}
    \proofstepwidestar[1,2,3,4,5,6]{%
      Z\diamond_{\mathcal P}X=Z\diamond_{\mathcal P}Y}{%
      \FormulaRefAuto{A\subseteq B\dsep B\subseteq A\vdash A=B}[thm]{25,44}}
  \closeproofpartwideindR

  \proofpartwide{Schluss}
    \proofstepwidestar[1]{\SemigroupAlg{B}{\diamond}}{\rA}
    \proofstepwidestar[2]{X\in\PowNE{B}}{\rA}
    \proofstepwidestar[3]{Y\in\PowNE{B}}{\rA}
    \proofstepwidestar[4]{\forall x\in X\,\exists y\in Y\,
      \bigl(x\mathrel{\equiv^{\mathrm{tr}}_{B,\diamond}}y\bigr)}{\rA}
    \proofstepwidestar[5]{\forall y\in Y\,\exists x\in X\,
      \bigl(y\mathrel{\equiv^{\mathrm{tr}}_{B,\diamond}}x\bigr)}{\rA}
    \proofstepwidestar[1]{%
      \SemigroupAlg{\PowNE{B}}{\diamond_{\mathcal P}}}{%
      \FormulaRefAuto{NonemptyPowerSemigroup}[thm]{1}}
    \proofstepwidestar[7]{Z\in\PowNE{B}}{\rA}
    \proofstepwidestarbreak[1,2,3,4,5,7]{%
      X\diamond_{\mathcal P}Z=Y\diamond_{\mathcal P}Z}{%
      \FormulaRefAuto{PowerSemigroupTranslationSupportRightProduct}[thm]{%
        1,2,3,7,4,5}}
    \proofstepwidestarbreak[1,2,3,4,5,7]{%
      Z\diamond_{\mathcal P}X=Z\diamond_{\mathcal P}Y}{%
      \FormulaRefAuto{PowerSemigroupTranslationSupportLeftProduct}[thm]{%
        1,2,3,7,4,5}}
    \proofstepwidestar[1,2,3,4,5,7]{%
      \begin{aligned}[t]
        &X\diamond_{\mathcal P}Z=Y\diamond_{\mathcal P}Z\land{}\\[-2pt]
        &Z\diamond_{\mathcal P}X=Z\diamond_{\mathcal P}Y
      \end{aligned}}{\rAI{8,9}}
    \proofstepwidestarbreak[1,2,3,4,5]{%
      \forall Z\in\PowNE{B}\,\left(
        \begin{aligned}[t]
          &X\diamond_{\mathcal P}Z=Y\diamond_{\mathcal P}Z\land{}\\[-2pt]
          &Z\diamond_{\mathcal P}X=Z\diamond_{\mathcal P}Y
        \end{aligned}\right)}{%
      \rUI{\rRI{7,10}}}
    \proofstepwidestarbreak[1,2,3,4,5]{%
      X\mathrel{\equiv^{\mathrm{tr}}_{\PowNE{B},
        \diamond_{\mathcal P}}}Y}{%
      \FormulaRefAuto{SemigroupTranslationIndistinguishableDef}[def]{%
        6,2,3,11}}
  \closeproofpartwide
\end{tabproofsplitwide}

% Bei besonders langen Formeln steht die Begründung in einer eigenen
% Tabellenzeile.  Abhängigkeiten und Schrittzählung bleiben unverändert.
\NewDocumentCommand{\MogProofStepWideStarBreak}{O{} m m}{%
  \stepcounter{proofstepnr}%
  #1 & \theproofstepnr
     & \multicolumn{3}{@{}>{$}l<{$}@{}}{#2} & \\[-.2ex]%
     & & \multicolumn{4}{@{}r@{}}{%
       \makebox[0pt][r]{\ensuremath{#3}}}\\%
}

\newcommand{\MogStar}{\mathbin{\ast}}
\newcommand{\MogDiamond}{\mathbin{\diamond}}

\subsection{Ein unendliches Gegenbeispiel zur Umkehrung}

Die folgende Konstruktion geht auf
\href{https://doi.org/10.1007/BF02389140}%
  {Mogiljanskaja, \emph{Non-isomorphic Semigroups with Isomorphic
  Semigroups of Subsets}, 1973, S.~330--333}
zurück. Die Konstruktion lässt sich rückblickend von ihrem Ziel her
folgendermaßen verstehen.

\begin{remark}[Kernidee und Beweisskizze]
Gesucht sind zwei nicht isomorphe unendliche Halbgruppen, deren
Potenzhalbgruppen dennoch isomorph sind. Dazu wird die Multiplikation zunächst
so gewählt, dass sie möglichst wenig Information trägt. Die Elemente \(a_i\)
bilden eine erste Schicht \(L\); ihre paarweise verschiedenen Produkte
\(a_i\MogStar a_j=b_{ij}\) bilden eine zweite Schicht \(D\). Alle übrigen
Produkte sind \(0\). Insbesondere gilt für alle \(x,y,z\in S\)
\[
  (x\MogStar y)\MogStar z
  =\mathbf0
  =x\MogStar(y\MogStar z).
\]

Nun wird \(S\) um ein neues Element \(k\) erweitert, wobei jedes Produkt mit
\(k\) gleich \(0\) gesetzt wird. Das Element \(k\) ist selbst kein Produkt.
Wäre \(f\) ein Isomorphismus von \(S\) nach \(T\), so gäbe es wegen der
Surjektivität ein \(s\in S\) mit \(f(s)=k\). Jedes Element von
\(S\setminus L\) ist ein Produkt; daher müsste \(s=a_i\) für ein
\(i\in\mathbb N\) gelten. Dann
wäre jedoch
\[
  f(b_{i0})=f(a_i\MogStar a_0)
    =k\MogDiamond f(a_0)=\mathbf0,
  \qquad
  f(b_{i1})=f(a_i\MogStar a_1)
    =k\MogDiamond f(a_1)=\mathbf0,
\]
obwohl \(b_{i0}\neq b_{i1}\). Dies widerspricht der Injektivität von \(f\).

Auf der Ebene nichtleerer Teilmengen geht diese Unterscheidung verloren. Für
\(X,Y\in\PowNE{S}\) setze
\[
  R(X,Y)\coloneqq
  \bigl\{z\in D\bigm|
    \exists i,j\in\mathbb N\,
      (z=b_{ij}\land a_i\in X\cap L\land a_j\in Y\cap L)\bigr\}
\]
gilt
\[
  X\MogStar_{\mathcal P}Y
  =R(X,Y)\cup
  \begin{cases}
    \varnothing,&X\subseteq L\land Y\subseteq L,\\
    \{\mathbf0\},&\text{sonst}.
  \end{cases}
\]
Für \(X,Y\in\PowNE{T}\) gilt dieselbe Beschreibung mit
\(\MogDiamond_{\mathcal P}\). Ein Mengenprodukt hängt somit nur von den
beiden \(L\)-Anteilen und davon ab, ob beide Faktoren vollständig in \(L\)
liegen; einzelne Elemente außerhalb von \(L\) werden nicht unterschieden.

Die neue Möglichkeit, \(k\) in eine Teilmenge aufzunehmen, wird deshalb in
einer unendlichen Reservefamilie \(\mathcal C\) von Teilmengen der
Produktschicht codiert. Jedes \(K\in\mathcal C\) enthält \(b_{00}\) und
\(b_{11}\), liegt aber in der in Band~17 definierten Menge \(D'\). Wäre ein
solches \(K\) ein Rechteck \(R(X,Y)\), so folgte aus
\(b_{00},b_{11}\in R(X,Y)\) stets \(b_{01}\in R(X,Y)\). Dies ist unmöglich,
weil \(b_{01}\notin D'\) gilt. Daher enthält \(\mathcal C\) kein
Produktrechteck.
Die Unendlichkeit ermöglicht dann eine Bijektion
\(\psi\colon\mathcal C\bij\mathcal C\cup\mathcal E\), welche die neuen
\(k\)-haltigen \(D\)-Anteile aufnimmt. Die daraus gebildete Abbildung \(\Phi\)
erhält die \(L\)-Anteile sowie die Eigenschaft, vollständig in \(L\) zu liegen,
und lässt alle tatsächlich auftretenden Mengenprodukte fest. Daher gilt
\[
  \Phi(X\MogStar_{\mathcal P}Y)
  =X\MogStar_{\mathcal P}Y
  =\Phi(X)\MogDiamond_{\mathcal P}\Phi(Y).
\]
Genau dieser Absorptionsschritt ist ohne die Unendlichkeit nicht möglich.
\end{remark}

Für die benötigten beiden Schichten und die beiden Sonderwerte genügen
vier verschiedene natürliche Tags. Die Tags \(0\) und \(1\) trennen die
einfach und die zweifach indizierte Schicht; die Tags \(2\) und \(3\) halten
\(\mathbf0\) und \(k\) außerhalb beider Schichten und voneinander verschieden.

Die familien- und mengentheoretischen Daten werden aus Band~17 übernommen:
\begin{align*}
  a&\colon\mathbb N\to L,
  &a_i&=(0,i),
  &a[\mathbb N]&=L,\\[-2pt]
  b&\colon\mathbb N\times\mathbb N\to D,
  &b_{ij}&=(1,(i,j)),
  &b[\mathbb N\times\mathbb N]&=D.
\end{align*}
Dabei sind die Indizes eindeutig und \(L\cap D=\varnothing\); dies sind
\FormulaRefAuto{MogiljanskajaAIndexUnique}[thm],
\FormulaRefAuto{MogiljanskajaBIndexUnique}[thm] und
\FormulaRefAuto{MogiljanskajaLayersDisjoint}[thm]. Die beiden
Bildmengengleichheiten sind in
\FormulaRefAuto{MogiljanskajaLayerLImage}[thm] und
\FormulaRefAuto{MogiljanskajaLayerDImage}[thm] registriert. Im Folgenden werden nur
noch die beiden algebraischen Sonderwerte und die Operationen ergänzt.

\FormulaDefDeltaK[Das absorbierende Symbol]{%
  \mathbf0\coloneqq(2,0)%
}{MogiljanskajaZeroDef}{%
  \DeltaRow{Mengen}{\mathbf0}
}

\FormulaDefDeltaK[Das Zusatzsymbol]{%
  k\coloneqq(3,0)%
}{MogiljanskajaKDef}{%
  \DeltaRow{Mengen}{k}
}

\FormulaThmDeltaK[Das absorbierende Symbol liegt nicht in der (a)-Schicht]{%
  \mathbf0\notin L%
}{MogiljanskajaZeroNotInL}{%
  \DeltaRow{Mengen}{L\dsep\mathbf0}
}
\begin{tabproof}
  \proofstep{}{\mathbf0=(2,0)}{%
    \FormulaRefAuto{MogiljanskajaZeroDef}[def]}
  \proofstep{}{L=\{0\}\times\mathbb N}{%
    \FormulaRefAuto{MogiljanskajaLayerLDef}[def]}
  \proofstep{}{0\neq2}{\FormulaRefAuto{PeanoZeroNeTwo}[thm]}
  \proofstep{}{2\neq0}{%
    \FormulaRefAuto{a\neq b\vdash b\neq a}[thm]{3}}
  \proofstep{}{(2,0)\notin\{0\}\times\mathbb N}{%
    \FormulaRefAuto{TaggedPointOutsideProductLayer}[thm]{4}}
  \proofstep{}{(2,0)\notin L}{\rIE{2,5}}
  \proofstep{}{\mathbf0\notin L}{\rIE{1,6}}
\end{tabproof}

\FormulaThmDeltaK[Das absorbierende Symbol liegt nicht in der Produktschicht]{%
  \mathbf0\notin L\cup D\vdash\mathbf0\notin D%
}{MogiljanskajaZeroNotInD}{%
  \DeltaRow{Mengen}{L\dsep D\dsep\mathbf0}
}
\begin{tabproof}
  \proofstep{1}{\mathbf0\notin L\cup D}{\rA}
  \proofstep{1}{\mathbf0\notin L\land\mathbf0\notin D}{%
    \FormulaRefAuto{OutsideUnionComponents}[thm]{1}}
  \proofstep{1}{\mathbf0\notin D}{\rAEb{2}}
\end{tabproof}

\FormulaThmDeltaK[Das absorbierende Symbol liegt in keiner Grundschicht]{%
  \mathbf0\notin L\cup D%
}{MogiljanskajaZeroOutsideLayers}{%
  \DeltaRow{Mengen}{L\dsep D\dsep\mathbf0}
}
\begin{tabproof}
  \proofstep{}{\mathbf0\notin L}{%
    \FormulaRefAuto{MogiljanskajaZeroNotInL}[thm]}
  \proofstep{}{\mathbf0=(2,0)}{%
    \FormulaRefAuto{MogiljanskajaZeroDef}[def]}
  \proofstep{}{D=\{1\}\times(\mathbb N\times\mathbb N)}{%
    \FormulaRefAuto{MogiljanskajaLayerDDef}[def]}
  \proofstep{}{1\neq2}{\FormulaRefAuto{PeanoOneNeTwo}[thm]}
  \proofstep{}{2\neq1}{%
    \FormulaRefAuto{a\neq b\vdash b\neq a}[thm]{4}}
  \proofstep{}{(2,0)\notin\{1\}\times(\mathbb N\times\mathbb N)}{%
    \FormulaRefAuto{TaggedPointOutsideProductLayer}[thm]{5}}
  \proofstep{}{(2,0)\notin D}{\rIE{3,6}}
  \proofstep{}{\mathbf0\notin D}{\rIE{2,7}}
  \proofstep{}{\mathbf0\notin L\cup D}{%
    \FormulaRefAuto{NotInUnion}[thm]{1,8}}
\end{tabproof}

\FormulaThmDeltaK[Zusatzsymbol außerhalb der (a)-Schicht]{%
  k\notin L%
}{MogiljanskajaKNotInL}{%
  \DeltaRow{Mengen}{L\dsep k}
}
\begin{tabproof}
  \proofstep{}{k=(3,0)}{\FormulaRefAuto{MogiljanskajaKDef}[def]}
  \proofstep{}{L=\{0\}\times\mathbb N}{%
    \FormulaRefAuto{MogiljanskajaLayerLDef}[def]}
  \proofstep{}{0\neq3}{\FormulaRefAuto{PeanoZeroNeThree}[thm]}
  \proofstep{}{3\neq0}{%
    \FormulaRefAuto{a\neq b\vdash b\neq a}[thm]{3}}
  \proofstep{}{(3,0)\notin\{0\}\times\mathbb N}{%
    \FormulaRefAuto{TaggedPointOutsideProductLayer}[thm]{4}}
  \proofstep{}{(3,0)\notin L}{\rIE{2,5}}
  \proofstep{}{k\notin L}{\rIE{1,6}}
\end{tabproof}

\FormulaThmDeltaK[Das Zusatzsymbol liegt nicht in der Produktschicht]{%
  k\notin L\cup D\vdash k\notin D%
}{MogiljanskajaKNotInD}{%
  \DeltaRow{Mengen}{L\dsep D\dsep k}
}
\begin{tabproof}
  \proofstep{1}{k\notin L\cup D}{\rA}
  \proofstep{1}{k\notin L\land k\notin D}{%
    \FormulaRefAuto{OutsideUnionComponents}[thm]{1}}
  \proofstep{1}{k\notin D}{\rAEb{2}}
\end{tabproof}

\FormulaThmDeltaK[Das Zusatzsymbol liegt in keiner Grundschicht]{%
  k\notin L\cup D%
}{MogiljanskajaKOutsideLayers}{%
  \DeltaRow{Mengen}{L\dsep D\dsep k}
}
\begin{tabproof}
  \proofstep{}{k\notin L}{\FormulaRefAuto{MogiljanskajaKNotInL}[thm]}
  \proofstep{}{k=(3,0)}{\FormulaRefAuto{MogiljanskajaKDef}[def]}
  \proofstep{}{D=\{1\}\times(\mathbb N\times\mathbb N)}{%
    \FormulaRefAuto{MogiljanskajaLayerDDef}[def]}
  \proofstep{}{1\neq3}{\FormulaRefAuto{PeanoOneNeThree}[thm]}
  \proofstep{}{3\neq1}{%
    \FormulaRefAuto{a\neq b\vdash b\neq a}[thm]{4}}
  \proofstep{}{(3,0)\notin\{1\}\times(\mathbb N\times\mathbb N)}{%
    \FormulaRefAuto{TaggedPointOutsideProductLayer}[thm]{5}}
  \proofstep{}{(3,0)\notin D}{\rIE{3,6}}
  \proofstep{}{k\notin D}{\rIE{2,7}}
  \proofstep{}{k\notin L\cup D}{%
    \FormulaRefAuto{NotInUnion}[thm]{1,8}}
\end{tabproof}

\FormulaThmDeltaK[Die beiden Sonderwerte sind verschieden]{%
  \mathbf0\neq k%
}{MogiljanskajaZeroNeK}{%
  \DeltaRow{Mengen}{\mathbf0\dsep k}
}
\begin{tabproof}
  \proofstep{}{\mathbf0=(2,0)}{%
    \FormulaRefAuto{MogiljanskajaZeroDef}[def]}
  \proofstep{}{k=(3,0)}{\FormulaRefAuto{MogiljanskajaKDef}[def]}
  \proofstep{}{2\neq3}{\FormulaRefAuto{PeanoTwoNeThree}[thm]}
  \proofstep{}{(2,0)\neq(3,0)}{%
    \FormulaRefAuto{TaggedPointsDifferent}[thm]{3}}
  \proofstep{}{(2,0)\neq k}{\rIE{2,4}}
  \proofstep{}{\mathbf0\neq k}{\rIE{1,5}}
\end{tabproof}

\FormulaThmDeltaK[Das Zusatzsymbol liegt nicht in der Nullmenge]{%
  k\notin\{\mathbf0\}%
}{MogiljanskajaKNotInZeroSingleton}{%
  \DeltaRow{Mengen}{\mathbf0\dsep k}
}
\begin{tabproof}
  \proofstep{}{\mathbf0\neq k}{%
    \FormulaRefAuto{MogiljanskajaZeroNeK}[thm]}
  \proofstep{}{k\neq\mathbf0}{%
    \FormulaRefAuto{a\neq b\vdash b\neq a}[thm]{1}}
  \proofstep{}{k\notin\{\mathbf0\}}{%
    \FormulaRefAuto{x\notin\{a\}\eqvdash x\neq a}[thm]{2}}
\end{tabproof}

\FormulaDefDeltaK[Der erste Träger]{%
  S\coloneqq L\cup D\cup\{\mathbf0\}%
}{MogiljanskajaCarrierSDef}{%
  \DeltaRow{Mengen}{S\dsep L\dsep D\dsep\mathbf0}
}

\FormulaDefDeltaK[Der zweite Träger]{%
  T\coloneqq S\cup\{k\}%
}{MogiljanskajaCarrierTDef}{%
  \DeltaRow{Mengen}{T\dsep S\dsep k}
}

\FormulaThmDeltaK[Elementkriterium des ersten Trägers]{%
  s\in S\leftrightarrow(s\in L\lor s\in D\lor s=\mathbf0)%
}{MogiljanskajaCarrierSMembership}{%
  \DeltaRow{Mengen}{S\dsep L\dsep D\dsep\mathbf0\dsep s}
}
\begin{tabproofwide}
  \proofstepwidestar[]{S=L\cup D\cup\{\mathbf0\}}{%
    \FormulaRefAuto{MogiljanskajaCarrierSDef}[def]}
  \proofstepwidestar[]{s\in S\leftrightarrow
    s\in L\cup D\cup\{\mathbf0\}}{\rLRS{1}}
  \proofstepwidestar[]{s\in L\cup D\cup\{\mathbf0\}
    \leftrightarrow (s\in L\lor s\in D)\lor s\in\{\mathbf0\}}{%
    \FormulaRefAuto{x\in A\cup B\eqvdash x\in A\lor x\in B}[thm]}
  \proofstepwidestar[]{(s\in L\lor s\in D)\lor s\in\{\mathbf0\}
    \leftrightarrow s\in L\lor s\in D\lor s=\mathbf0}{%
    \FormulaRefAuto{x\in\{a\}\eqvdash x=a}[thm]}
  \proofstepwidestar[]{s\in S\leftrightarrow
    (s\in L\lor s\in D\lor s=\mathbf0)}{%
    \FormulaRefAuto{P\leftrightarrow Q\dsep Q\leftrightarrow R
      \vdash P\leftrightarrow R}[thm]{2,3,4}}
\end{tabproofwide}

\FormulaThmDeltaK[Elementkriterium des zweiten Trägers]{%
  t\in T\leftrightarrow(t\in S\lor t=k)%
}{MogiljanskajaCarrierTMembership}{%
  \DeltaRow{Mengen}{S\dsep T\dsep k\dsep t}
}
\begin{tabproofwide}
  \proofstepwidestar[]{T=S\cup\{k\}}{%
    \FormulaRefAuto{MogiljanskajaCarrierTDef}[def]}
  \proofstepwidestar[]{t\in T\leftrightarrow t\in S\cup\{k\}}{\rLRS{1}}
  \proofstepwidestar[]{t\in S\cup\{k\}\leftrightarrow
    t\in S\lor t\in\{k\}}{%
    \FormulaRefAuto{x\in A\cup B\eqvdash x\in A\lor x\in B}[thm]}
  \proofstepwidestar[]{t\in S\lor t\in\{k\}\leftrightarrow
    t\in S\lor t=k}{\FormulaRefAuto{x\in\{a\}\eqvdash x=a}[thm]}
  \proofstepwidestar[]{t\in T\leftrightarrow(t\in S\lor t=k)}{%
    \FormulaRefAuto{P\leftrightarrow Q\dsep Q\leftrightarrow R
      \vdash P\leftrightarrow R}[thm]{2,3,4}}
\end{tabproofwide}

\FormulaThmDeltaK[Der erste Träger liegt im zweiten]{%
  S\subseteq T%
}{MogiljanskajaSSubsetT}{%
  \DeltaRow{Mengen}{S\dsep T\dsep k}
}
\begin{tabproof}
  \proofstep{}{T=S\cup\{k\}}{\FormulaRefAuto{MogiljanskajaCarrierTDef}[def]}
  \proofstep{}{S\subseteq S\cup\{k\}}{\FormulaRefAuto{A\subseteq A\cup B}[thm]}
  \proofstep{}{S\subseteq T}{\rIE{1,2}}
\end{tabproof}

\FormulaThmDeltaK[Das Zusatzsymbol gehört zum zweiten Träger]{%
  k\in T%
}{MogiljanskajaKInT}{%
  \DeltaRow{Mengen}{T\dsep k}
}
\begin{tabproof}
  \proofstep{}{k=k}{\rII}
  \proofstep{}{k\in S\lor k=k}{\rOIb{1}}
  \proofstep{}{k\in T}{\FormulaRefAuto{MogiljanskajaCarrierTMembership}[thm]{2}}
\end{tabproof}

\FormulaThmDeltaK[Das Zusatzsymbol liegt außerhalb der drei Trägerteile]{%
  k\notin(L\cup D)\cup\{\mathbf0\}%
}{MogiljanskajaKOutsideCarrierPieces}{%
  \DeltaRow{Mengen}{L\dsep D\dsep\mathbf0\dsep k}
}
\begin{tabproof}
  \proofstep{}{k\notin L\cup D}{%
    \FormulaRefAuto{MogiljanskajaKOutsideLayers}[thm]}
  \proofstep{}{k\notin\{\mathbf0\}}{%
    \FormulaRefAuto{MogiljanskajaKNotInZeroSingleton}[thm]}
  \proofstep{}{k\notin(L\cup D)\cup\{\mathbf0\}}{%
    \FormulaRefAuto{NotInUnion}[thm]{1,2}}
\end{tabproof}

\FormulaThmDeltaK[Das Zusatzsymbol gehört nicht zum ersten Träger]{%
  k\notin(L\cup D)\cup\{\mathbf0\}\vdash k\notin S%
}{MogiljanskajaKNotInS}{%
  \DeltaRow{Mengen}{S\dsep L\dsep D\dsep\mathbf0\dsep k}
}
\begin{tabproofwide}
  \proofstepwidestar[1]{k\notin(L\cup D)\cup\{\mathbf0\}}{\rA}
  \proofstepwidestar[]{S=(L\cup D)\cup\{\mathbf0\}}{%
    \FormulaRefAuto{MogiljanskajaCarrierSDef}[def]}
  \proofstepwidestar[]{(L\cup D)\cup\{\mathbf0\}=S}{%
    \FormulaRefAuto{a=b\vdash b=a}[thm]{2}}
  \proofstepwidestar[1]{k\notin S}{\rIE{3,1}}
\end{tabproofwide}

Die folgenden beiden Vorschriften sind eindeutig: Liegen beide Faktoren in
\(L\), so bestimmen die Eindeutigkeitssätze ihre Indizes; andernfalls greift
der zweite Fall.  Für \(T=S\cup\{k\}\) sind ebenso die beiden Fälle der
zweiten Vorschrift disjunkt und vollständig.

\FormulaDefDeltaK[Die Multiplikation auf (S)]{%
  \begin{aligned}[t]
    \MogStar&\colon S\times S\to S,\\[-2pt]
    i,j\in\mathbb N&\vdash a_i\MogStar a_j\coloneqq b_{ij},\\[-2pt]
    x,y\in S\dsep(x\notin L\lor y\notin L)
      &\vdash x\MogStar y\coloneqq\mathbf0
  \end{aligned}
}{MogiljanskajaStarDef}{%
  \DeltaRow{Mengen}{S\dsep L\dsep\mathbf0\dsep a_i\dsep a_j\dsep b_{ij}
    \dsep x\dsep y\dsep i\dsep j}
  \DeltaRow{Funktionen}{\MogStar}
}

\FormulaDefDeltaK[Die Multiplikation auf (T)]{%
  \begin{aligned}[t]
    \MogDiamond&\colon T\times T\to T,\\[-2pt]
    x,y\in S&\vdash x\MogDiamond y\coloneqq x\MogStar y,\\[-2pt]
    x,y\in T\dsep k\in\{x,y\}
      &\vdash x\MogDiamond y\coloneqq\mathbf0
  \end{aligned}
}{MogiljanskajaDiamondDef}{%
  \DeltaRow{Mengen}{S\dsep T\dsep k\dsep\mathbf0\dsep x\dsep y}
  \DeltaRow{Funktionen}{\MogStar\dsep\MogDiamond}
}

\FormulaThmDeltaK[Wertebereich der ersten Multiplikation]{%
  x,y\in S\vdash x\MogStar y\in D\cup\{\mathbf0\}%
}{MogiljanskajaStarRange}{%
  \DeltaRow{Mengen}{S\dsep L\dsep D\dsep\mathbf0\dsep x\dsep y}
  \DeltaRow{Funktionen}{\MogStar}
}
\begin{tabproof}
  \proofstep{1}{x\in S}{\rA}
  \proofstep{2}{y\in S}{\rA}
  \proofstep{1,2}{%
    (x\in L\land y\in L)\lor(x\notin L\lor y\notin L)}{%
    \FormulaRefAuto{(P\land Q)\lor(\neg P\lor\neg Q)}[thm]}
  \proofstep{4}{x\in L\land y\in L}{\rA}
  \proofstep{4}{\exists i,j\in\mathbb N\,(x=a_i\land y=a_j)}{%
    \FormulaRefAuto{MogiljanskajaLayerLMembership}[thm]{4}}
  \proofstep{6}{i,j\in\mathbb N\land x=a_i\land y=a_j}{\rA}
  \proofstep{6}{x\MogStar y=b_{ij}}{%
    \FormulaRefAuto{MogiljanskajaStarDef}[def]{6}}
  \proofstep{6}{b_{ij}\in D}{%
    \FormulaRefAuto{MogiljanskajaLayerDMembership}[thm]{\rEI{6}}}
  \proofstep{6}{b_{ij}\in D\cup\{\mathbf0\}}{%
    \FormulaRefAuto{z\in A\vdash z\in A\cup B}[thm]{8}}
  \proofstep{6}{x\MogStar y\in D\cup\{\mathbf0\}}{\rIE{7,9}}
  \proofstep{4}{x\MogStar y\in D\cup\{\mathbf0\}}{\rEE{5,6,10}}
  \proofstep{12}{x\notin L\lor y\notin L}{\rA}
  \proofstep{1,2,12}{x\MogStar y=\mathbf0}{%
    \FormulaRefAuto{MogiljanskajaStarDef}[def]{1,2,12}}
  \proofstep{}{\mathbf0\in\{\mathbf0\}}{%
    \FormulaRefAuto{a\in\{a\}}[thm]}
  \proofstep{}{\{\mathbf0\}\subseteq D\cup\{\mathbf0\}}{%
    \FormulaRefAuto{B\subseteq A\cup B}[thm]}
  \proofstep{}{\mathbf0\in D\cup\{\mathbf0\}}{%
    \FormulaRefAuto{A\subseteq B\dsep x\in A\vdash x\in B}[thm]{15,14}}
  \proofstep{1,2,12}{x\MogStar y\in D\cup\{\mathbf0\}}{%
    \rIE{13,16}}
  \proofstep{1,2}{x\MogStar y\in D\cup\{\mathbf0\}}{%
    \rOE{3,4,11,12,17}}
\end{tabproof}

\FormulaThmDeltaK[Produktschicht und Null liegen außerhalb der (a)-Schicht]{%
  u\in D\cup\{\mathbf0\}\vdash u\notin L%
}{MogiljanskajaProductValuesOutsideL}{%
  \DeltaRow{Mengen}{L\dsep D\dsep\mathbf0\dsep u}
}
\begin{tabproof}
  \proofstep{1}{u\in D\cup\{\mathbf0\}}{\rA}
  \proofstep{1}{u\in D\lor u\in\{\mathbf0\}}{%
    \FormulaRefAuto{x\in A\cup B\eqvdash x\in A\lor x\in B}[thm]{1}}
  \proofstep{3}{u\in D}{\rA}
  \proofstep{}{L\cap D=\varnothing}{%
    \FormulaRefAuto{MogiljanskajaLayersDisjoint}[thm]}
  \proofstep{3}{u\notin L}{%
    \FormulaRefAuto{A\cap B=\varnothing\dsep x\in B\vdash x\notin A}[thm]{%
      4,3}}
  \proofstep{6}{u\in\{\mathbf0\}}{\rA}
  \proofstep{6}{u=\mathbf0}{\FormulaRefAuto{x\in\{a\}\eqvdash x=a}[thm]{6}}
  \proofstep{}{\mathbf0\notin L}{\FormulaRefAuto{MogiljanskajaZeroNotInL}[thm]}
  \proofstep{6}{u\notin L}{\rIE{7,8}}
  \proofstep{1}{u\notin L}{\rOE{2,3,5,6,9}}
\end{tabproof}

\FormulaThmDeltaK[Produktschicht und Null liegen im ersten Träger]{%
  u\in D\cup\{\mathbf0\}\vdash u\in S%
}{MogiljanskajaProductValuesInS}{%
  \DeltaRow{Mengen}{S\dsep D\dsep\mathbf0\dsep u}
}
\begin{tabproof}
  \proofstep{1}{u\in D\cup\{\mathbf0\}}{\rA}
  \proofstep{1}{u\in D\lor u\in\{\mathbf0\}}{%
    \FormulaRefAuto{x\in A\cup B\eqvdash x\in A\lor x\in B}[thm]{1}}
  \proofstep{3}{u\in D}{\rA}
  \proofstep{3}{u\in L\lor u\in D\lor u=\mathbf0}{\rOIb{\rOIa{3}}}
  \proofstep{3}{u\in S}{\FormulaRefAuto{MogiljanskajaCarrierSMembership}[thm]{4}}
  \proofstep{6}{u\in\{\mathbf0\}}{\rA}
  \proofstep{6}{u=\mathbf0}{\FormulaRefAuto{x\in\{a\}\eqvdash x=a}[thm]{6}}
  \proofstep{6}{u\in L\lor u\in D\lor u=\mathbf0}{\rOIb{\rOIb{7}}}
  \proofstep{6}{u\in S}{\FormulaRefAuto{MogiljanskajaCarrierSMembership}[thm]{8}}
  \proofstep{1}{u\in S}{\rOE{2,3,5,6,9}}
\end{tabproof}

\FormulaThmDeltaK[Linksabsorption der Produkte der ersten Multiplikation]{%
  u\in D\cup\{\mathbf0\}\dsep z\in S\vdash
  u\MogStar z=\mathbf0%
}{MogiljanskajaStarLeftAbsorption}{%
  \DeltaRow{Mengen}{S\dsep L\dsep D\dsep\mathbf0\dsep u\dsep z}
  \DeltaRow{Funktionen}{\MogStar}
}
\begin{tabproof}
  \proofstep{1}{u\in D\cup\{\mathbf0\}}{\rA}
  \proofstep{2}{z\in S}{\rA}
  \proofstep{1}{u\notin L}{%
    \FormulaRefAuto{MogiljanskajaProductValuesOutsideL}[thm]{1}}
  \proofstep{1}{u\in S}{%
    \FormulaRefAuto{MogiljanskajaProductValuesInS}[thm]{1}}
  \proofstep{1}{u\notin L\lor z\notin L}{\rOIa{3}}
  \proofstep{1,2}{u\MogStar z=\mathbf0}{%
    \FormulaRefAuto{MogiljanskajaStarDef}[def]{4,2,5}}
\end{tabproof}

\FormulaThmDeltaK[Rechtsabsorption der Produkte der ersten Multiplikation]{%
  u\in D\cup\{\mathbf0\}\dsep z\in S\vdash
  z\MogStar u=\mathbf0%
}{MogiljanskajaStarRightAbsorption}{%
  \DeltaRow{Mengen}{S\dsep L\dsep D\dsep\mathbf0\dsep u\dsep z}
  \DeltaRow{Funktionen}{\MogStar}
}
\begin{tabproof}
  \proofstep{1}{u\in D\cup\{\mathbf0\}}{\rA}
  \proofstep{2}{z\in S}{\rA}
  \proofstep{1}{u\notin L}{%
    \FormulaRefAuto{MogiljanskajaProductValuesOutsideL}[thm]{1}}
  \proofstep{1}{u\in S}{%
    \FormulaRefAuto{MogiljanskajaProductValuesInS}[thm]{1}}
  \proofstep{1}{z\notin L\lor u\notin L}{\rOIb{3}}
  \proofstep{1,2}{z\MogStar u=\mathbf0}{%
    \FormulaRefAuto{MogiljanskajaStarDef}[def]{2,4,5}}
\end{tabproof}

\FormulaThmDeltaK[Wertebereich der zweiten Multiplikation]{%
  x,y\in T\vdash x\MogDiamond y\in D\cup\{\mathbf0\}%
}{MogiljanskajaDiamondRange}{%
  \DeltaRow{Mengen}{S\dsep T\dsep D\dsep k\dsep\mathbf0\dsep x\dsep y}
  \DeltaRow{Funktionen}{\MogStar\dsep\MogDiamond}
}
\begin{tabproof}
  \proofstep{1}{x\in T}{\rA}
  \proofstep{2}{y\in T}{\rA}
  \proofstep{1,2}{(x,y\in S)\lor k\in\{x,y\}}{%
    \FormulaRefAuto{MogiljanskajaCarrierTMembership}[thm]{1,2}}
  \proofstep{4}{x,y\in S}{\rA}
  \proofstep{4}{x\MogStar y\in D\cup\{\mathbf0\}}{%
    \FormulaRefAuto{MogiljanskajaStarRange}[thm]{4}}
  \proofstep{4}{x\MogDiamond y=x\MogStar y}{%
    \FormulaRefAuto{MogiljanskajaDiamondDef}[def]{4}}
  \proofstep{4}{x\MogDiamond y\in D\cup\{\mathbf0\}}{\rIE{6,5}}
  \proofstep{8}{k\in\{x,y\}}{\rA}
  \proofstep{1,2,8}{x\MogDiamond y=\mathbf0}{%
    \FormulaRefAuto{MogiljanskajaDiamondDef}[def]{1,2,8}}
  \proofstep{}{\mathbf0\in\{\mathbf0\}}{%
    \FormulaRefAuto{x\in\{a\}\eqvdash x=a}[thm]{\rII}}
  \proofstep{}{\{\mathbf0\}\subseteq D\cup\{\mathbf0\}}{%
    \FormulaRefAuto{B\subseteq A\cup B}[thm]}
  \proofstep{}{\mathbf0\in D\cup\{\mathbf0\}}{%
    \FormulaRefAuto{A\subseteq B\dsep x\in A\vdash x\in B}[thm]{11,10}}
  \proofstep{1,2,8}{x\MogDiamond y\in D\cup\{\mathbf0\}}{\rIE{9,12}}
  \proofstep{1,2}{x\MogDiamond y\in D\cup\{\mathbf0\}}{%
    \rOE{3,4,7,8,13}}
\end{tabproof}

\FormulaThmDeltaK[Linksabsorption der Produkte der zweiten Multiplikation]{%
  u\in D\cup\{\mathbf0\}\dsep z\in T\vdash
  u\MogDiamond z=\mathbf0%
}{MogiljanskajaDiamondLeftAbsorption}{%
  \DeltaRow{Mengen}{S\dsep T\dsep D\dsep k\dsep\mathbf0\dsep u\dsep z}
  \DeltaRow{Funktionen}{\MogStar\dsep\MogDiamond}
}
\begin{tabproof}
  \proofstep{1}{u\in D\cup\{\mathbf0\}}{\rA}
  \proofstep{2}{z\in T}{\rA}
  \proofstep{1}{u\in S}{%
    \FormulaRefAuto{MogiljanskajaProductValuesInS}[thm]{1}}
  \proofstep{2}{z\in S\lor z=k}{%
    \FormulaRefAuto{MogiljanskajaCarrierTMembership}[thm]{2}}
  \proofstep{5}{z\in S}{\rA}
  \proofstep{1,5}{u\MogStar z=\mathbf0}{%
    \FormulaRefAuto{MogiljanskajaStarLeftAbsorption}[thm]{1,5}}
  \proofstep{1,5}{u\MogDiamond z=u\MogStar z}{%
    \FormulaRefAuto{MogiljanskajaDiamondDef}[def]{3,5}}
  \proofstep{1,5}{u\MogDiamond z=\mathbf0}{%
    \FormulaRefAuto{a=b\dsep b=c\vdash a=c}[thm]{7,6}}
  \proofstep{9}{z=k}{\rA}
  \proofstep{}{S\subseteq T}{%
    \FormulaRefAuto{MogiljanskajaSSubsetT}[thm]}
  \proofstep{1}{u\in T}{%
    \FormulaRefAuto{A\subseteq B\dsep x\in A\vdash x\in B}[thm]{10,3}}
  \proofstep{9}{k=z}{%
    \FormulaRefAuto{a=b\vdash b=a}[thm]{9}}
  \proofstep{9}{k=u\lor k=z}{\rOIb{12}}
  \proofstep{9}{k\in\{u,z\}}{%
    \FormulaRefAuto{x\in\{a,b\}\eqvdash (x=a\lor x=b)}[thm]{13}}
  \proofstep{1,2,9}{u\MogDiamond z=\mathbf0}{%
    \FormulaRefAuto{MogiljanskajaDiamondDef}[def]{11,2,14}}
  \proofstep{1,2}{u\MogDiamond z=\mathbf0}{%
    \rOE{4,5,8,9,15}}
\end{tabproof}

\FormulaThmDeltaK[Rechtsabsorption der Produkte der zweiten Multiplikation]{%
  u\in D\cup\{\mathbf0\}\dsep z\in T\vdash
  z\MogDiamond u=\mathbf0%
}{MogiljanskajaDiamondRightAbsorption}{%
  \DeltaRow{Mengen}{S\dsep T\dsep D\dsep k\dsep\mathbf0\dsep u\dsep z}
  \DeltaRow{Funktionen}{\MogStar\dsep\MogDiamond}
}
\begin{tabproof}
  \proofstep{1}{u\in D\cup\{\mathbf0\}}{\rA}
  \proofstep{2}{z\in T}{\rA}
  \proofstep{1}{u\in S}{%
    \FormulaRefAuto{MogiljanskajaProductValuesInS}[thm]{1}}
  \proofstep{2}{z\in S\lor z=k}{%
    \FormulaRefAuto{MogiljanskajaCarrierTMembership}[thm]{2}}
  \proofstep{5}{z\in S}{\rA}
  \proofstep{1,5}{z\MogStar u=\mathbf0}{%
    \FormulaRefAuto{MogiljanskajaStarRightAbsorption}[thm]{1,5}}
  \proofstep{1,5}{z\MogDiamond u=z\MogStar u}{%
    \FormulaRefAuto{MogiljanskajaDiamondDef}[def]{5,3}}
  \proofstep{1,5}{z\MogDiamond u=\mathbf0}{%
    \FormulaRefAuto{a=b\dsep b=c\vdash a=c}[thm]{7,6}}
  \proofstep{9}{z=k}{\rA}
  \proofstep{}{S\subseteq T}{%
    \FormulaRefAuto{MogiljanskajaSSubsetT}[thm]}
  \proofstep{1}{u\in T}{%
    \FormulaRefAuto{A\subseteq B\dsep x\in A\vdash x\in B}[thm]{10,3}}
  \proofstep{9}{k=z}{%
    \FormulaRefAuto{a=b\vdash b=a}[thm]{9}}
  \proofstep{9}{k=z\lor k=u}{\rOIa{12}}
  \proofstep{9}{k\in\{z,u\}}{%
    \FormulaRefAuto{x\in\{a,b\}\eqvdash (x=a\lor x=b)}[thm]{13}}
  \proofstep{1,2,9}{z\MogDiamond u=\mathbf0}{%
    \FormulaRefAuto{MogiljanskajaDiamondDef}[def]{2,11,14}}
  \proofstep{1,2}{z\MogDiamond u=\mathbf0}{%
    \rOE{4,5,8,9,15}}
\end{tabproof}

\FormulaThmDeltaK[Der erste Träger bildet eine Halbgruppe]{%
  \SemigroupAlg{S}{\MogStar}%
}{MogiljanskajaFirstSemigroup}{%
  \DeltaRow{Mengen}{S\dsep D\dsep\mathbf0\dsep x\dsep y\dsep z}
  \DeltaRow{Funktionen}{\MogStar}
}
\begin{tabproofwide}
  \proofstepwidestar[1]{x\in S}{\rA}
  \proofstepwidestar[2]{y\in S}{\rA}
  \proofstepwidestar[3]{z\in S}{\rA}
  \proofstepwidestar[]{\MogStar\colon S\times S\to S}{%
    \FormulaRefAuto{MogiljanskajaStarDef}[def]}
  \proofstepwidestar[1,2]{(x,y)\in S\times S}{%
    \FormulaRefAuto{a\in A\dsep b\in B\vdash(a,b)\in A\times B}[thm]{1,2}}
  \proofstepwidestar[1,2]{x\MogStar y\in S}{%
    \FormulaRefAuto{F\colon A\to B\dsep x\in A\vdash F(x)\in B}[thm]{%
      4,5}}
  \proofstepwidestar[2,3]{y\MogStar z\in D\cup\{\mathbf0\}}{%
    \FormulaRefAuto{MogiljanskajaStarRange}[thm]{2,3}}
  \proofstepwidestar[1,2]{x\MogStar y\in D\cup\{\mathbf0\}}{%
    \FormulaRefAuto{MogiljanskajaStarRange}[thm]{1,2}}
  \proofstepwide[1,2,3]{(x\MogStar y)\MogStar z}{=}{\mathbf0}{%
    \FormulaRefAuto{MogiljanskajaStarLeftAbsorption}[thm]{8,3}}
  \proofstepwidestar[1,2,3]{x\MogStar(y\MogStar z)=\mathbf0}{%
    \FormulaRefAuto{MogiljanskajaStarRightAbsorption}[thm]{7,1}}
  \proofstepwide[1,2,3]{\mathbf0}{=}{x\MogStar(y\MogStar z)}{%
    \FormulaRefAuto{a=b\vdash b=a}[thm]{10}}
  \proofstepwide[1,2,3]{(x\MogStar y)\MogStar z}{=}{%
    x\MogStar(y\MogStar z)}{\rChain{9,11}}
  \proofstepwidestar[]{\SemigroupAlg{S}{\MogStar}}{%
    \FormulaRefAuto{SemigroupDef}[def]{6,12}}
\end{tabproofwide}

\FormulaThmDeltaK[Der zweite Träger bildet eine Halbgruppe]{%
  \SemigroupAlg{T}{\MogDiamond}%
}{MogiljanskajaSecondSemigroup}{%
  \DeltaRow{Mengen}{T\dsep D\dsep\mathbf0\dsep x\dsep y\dsep z}
  \DeltaRow{Funktionen}{\MogDiamond}
}
\begin{tabproofwide}
  \proofstepwidestar[1]{x\in T}{\rA}
  \proofstepwidestar[2]{y\in T}{\rA}
  \proofstepwidestar[3]{z\in T}{\rA}
  \proofstepwidestar[]{\MogDiamond\colon T\times T\to T}{%
    \FormulaRefAuto{MogiljanskajaDiamondDef}[def]}
  \proofstepwidestar[1,2]{(x,y)\in T\times T}{%
    \FormulaRefAuto{a\in A\dsep b\in B\vdash(a,b)\in A\times B}[thm]{1,2}}
  \proofstepwidestar[1,2]{x\MogDiamond y\in T}{%
    \FormulaRefAuto{F\colon A\to B\dsep x\in A\vdash F(x)\in B}[thm]{%
      4,5}}
  \proofstepwidestar[2,3]{y\MogDiamond z\in D\cup\{\mathbf0\}}{%
    \FormulaRefAuto{MogiljanskajaDiamondRange}[thm]{2,3}}
  \proofstepwidestar[1,2]{x\MogDiamond y\in D\cup\{\mathbf0\}}{%
    \FormulaRefAuto{MogiljanskajaDiamondRange}[thm]{1,2}}
  \proofstepwide[1,2,3]{(x\MogDiamond y)\MogDiamond z}{=}{\mathbf0}{%
    \FormulaRefAuto{MogiljanskajaDiamondLeftAbsorption}[thm]{8,3}}
  \proofstepwidestar[1,2,3]{x\MogDiamond(y\MogDiamond z)=\mathbf0}{%
    \FormulaRefAuto{MogiljanskajaDiamondRightAbsorption}[thm]{7,1}}
  \proofstepwide[1,2,3]{\mathbf0}{=}{x\MogDiamond(y\MogDiamond z)}{%
    \FormulaRefAuto{a=b\vdash b=a}[thm]{10}}
  \proofstepwide[1,2,3]{(x\MogDiamond y)\MogDiamond z}{=}{%
    x\MogDiamond(y\MogDiamond z)}{\rChain{9,11}}
  \proofstepwidestar[]{\SemigroupAlg{T}{\MogDiamond}}{%
    \FormulaRefAuto{SemigroupDef}[def]{6,12}}
\end{tabproofwide}

\FormulaThmDeltaK[Der erste Träger ist nicht endlich]{%
  \neg\Finite{S}%
}{MogiljanskajaCarrierSInfinite}{%
  \DeltaRow{Mengen}{L\dsep S}
  \DeltaRow{Funktionen}{a\dsep F}
}
\begin{proof}
Nach \FormulaRefAuto{MogiljanskajaAElementsDef}[def] und
\FormulaRefAuto{MogiljanskajaAIndexUnique}[thm] ist
\(a\colon\mathbb N\inj L\). Aus
\FormulaRefAuto{MogiljanskajaCarrierSDef}[def] folgt \(L\subseteq S\);
die Zielmengenerweiterung liefert daher
\(a\colon\mathbb N\inj S\). Dies widerspricht im Falle
\(\Finite{S}\) dem Satz
\FormulaRefAuto{FiniteNoNatInjection}[thm].
\end{proof}

\FormulaThmDeltaK[Beide Träger sind nicht endlich]{%
  \neg\Finite{S}\land\neg\Finite{T}%
}{MogiljanskajaCarrierTInfinite}{%
  \DeltaRow{Mengen}{L\dsep S\dsep T}
  \DeltaRow{Funktionen}{a\dsep F}
}
\begin{proof}
Der erste Konjunkt ist
\FormulaRefAuto{MogiljanskajaCarrierSInfinite}[thm]. Wie zuvor ist
\(a\colon\mathbb N\inj L\). Aus
\FormulaRefAuto{MogiljanskajaCarrierSDef}[def] und
\FormulaRefAuto{MogiljanskajaCarrierTDef}[def] folgt
\(L\subseteq S\subseteq T\). Mit
\FormulaRefAuto{InjectiveCodomainExtension}[thm] erhalten wir
\(a\colon\mathbb N\inj T\), was nach
\FormulaRefAuto{FiniteNoNatInjection}[thm] die Endlichkeit von \(T\)
ausschließt.
\end{proof}

\FormulaThmDeltaK[Elemente außerhalb von (L) sind Produkte]{%
  s\in S\dsep s\notin L\vdash
  \exists x,y\in S\,(s=x\MogStar y)%
}{MogiljanskajaNonLIsProduct}{%
  \DeltaRow{Mengen}{L\dsep S\dsep D\dsep\mathbf0\dsep s\dsep x\dsep y
    \dsep i\dsep j}
  \DeltaRow{Funktionen}{\MogStar}
}
\begin{tabproof}
  \proofstep{1}{s\in S}{\rA}
  \proofstep{2}{s\notin L}{\rA}
  \proofstep{1,2}{s\in D\lor s=\mathbf0}{%
    \FormulaRefAuto{MogiljanskajaCarrierSMembership}[thm]{1,2}}
  \proofstep{4}{s\in D}{\rA}
  \proofstep{4}{\exists i,j\in\mathbb N\,(s=b_{ij})}{%
    \FormulaRefAuto{MogiljanskajaLayerDMembership}[thm]{4}}
  \proofstep{6}{i,j\in\mathbb N\land s=b_{ij}}{\rA}
  \proofstep{6}{i,j\in\mathbb N}{\rAEa{6}}
  \proofstep{6}{a_i,a_j\in L}{%
    \FormulaRefAuto{MogiljanskajaLayerLMembership}[thm]{7}}
  \proofstep{6}{a_i,a_j\in S}{%
    \FormulaRefAuto{MogiljanskajaCarrierSMembership}[thm]{8}}
  \proofstep{6}{s=a_i\MogStar a_j}{%
    \FormulaRefAuto{MogiljanskajaStarDef}[def]{6}}
  \proofstep{6}{\exists x,y\in S\,(s=x\MogStar y)}{\rEI{9,10}}
  \proofstep{4}{\exists x,y\in S\,(s=x\MogStar y)}{\rEE{5,6,11}}
  \proofstep{13}{s=\mathbf0}{\rA}
  \proofstep{13}{\mathbf0\in S}{%
    \FormulaRefAuto{MogiljanskajaCarrierSMembership}[thm]{\rOIb{\rOIb{\rII}}}}
  \proofstep{}{\mathbf0\in\{\mathbf0\}}{%
    \FormulaRefAuto{a\in\{a\}}[thm]}
  \proofstep{}{\{\mathbf0\}\subseteq D\cup\{\mathbf0\}}{%
    \FormulaRefAuto{B\subseteq A\cup B}[thm]}
  \proofstep{}{\mathbf0\in D\cup\{\mathbf0\}}{%
    \FormulaRefAuto{A\subseteq B\dsep x\in A\vdash x\in B}[thm]{16,15}}
  \proofstep{13}{\mathbf0\MogStar\mathbf0=\mathbf0}{%
    \FormulaRefAuto{MogiljanskajaStarLeftAbsorption}[thm]{17,14}}
  \proofstep{13}{\mathbf0=\mathbf0\MogStar\mathbf0}{%
    \FormulaRefAuto{a=b\vdash b=a}[thm]{18}}
  \proofstep{13}{s=\mathbf0\MogStar\mathbf0}{\rChain{13,19}}
  \proofstep{13}{\exists x,y\in S\,(s=x\MogStar y)}{\rEI{14,20}}
  \proofstep{1,2}{\exists x,y\in S\,(s=x\MogStar y)}{%
    \rOE{3,4,12,13,21}}
\end{tabproof}

\FormulaThmDeltaK[Kein Produkt im zweiten Träger ist (k)]{%
  x,y\in T\vdash x\MogDiamond y\neq k%
}{MogiljanskajaKNotProduct}{%
  \DeltaRow{Mengen}{T\dsep D\dsep\mathbf0\dsep k\dsep x\dsep y}
  \DeltaRow{Funktionen}{\MogDiamond}
}
\begin{tabproof}
  \proofstep{1}{x\in T}{\rA}
  \proofstep{2}{y\in T}{\rA}
  \proofstep{1,2}{x\MogDiamond y\in D\cup\{\mathbf0\}}{%
    \FormulaRefAuto{MogiljanskajaDiamondRange}[thm]{1,2}}
  \proofstep{}{k\notin L\cup D}{%
    \FormulaRefAuto{MogiljanskajaKOutsideLayers}[thm]}
  \proofstep{5}{k\in D}{\rA}
  \proofstep{5}{k\in L\cup D}{%
    \FormulaRefAuto{x\in A\cup B\eqvdash x\in A\lor x\in B}[thm]{\rOIb{5}}}
  \proofstep{5}{\bot}{\rBI{6,4}}
  \proofstep{}{k\notin D}{\rCI{5,7}}
  \proofstep{}{\mathbf0\neq k}{\FormulaRefAuto{MogiljanskajaZeroNeK}[thm]}
  \proofstep{}{k\neq\mathbf0}{\FormulaRefAuto{a\neq b\vdash b\neq a}[thm]{9}}
  \proofstep{}{k\notin\{\mathbf0\}}{%
    \FormulaRefAuto{x\notin\{a\}\eqvdash x\neq a}[thm]{10}}
  \proofstep{12}{k\in D\cup\{\mathbf0\}}{\rA}
  \proofstep{12}{k\in D\lor k\in\{\mathbf0\}}{%
    \FormulaRefAuto{x\in A\cup B\eqvdash x\in A\lor x\in B}[thm]{12}}
  \proofstep{14}{k\in D}{\rA}
  \proofstep{14}{\bot}{\rBI{14,8}}
  \proofstep{16}{k\in\{\mathbf0\}}{\rA}
  \proofstep{16}{\bot}{\rBI{16,11}}
  \proofstep{12}{\bot}{\rOE{13,14,15,16,17}}
  \proofstep{}{k\notin D\cup\{\mathbf0\}}{\rCI{12,18}}
  \proofstep{20}{x\MogDiamond y=k}{\rA}
  \proofstep{1,2,20}{k\in D\cup\{\mathbf0\}}{\rIE{20,3}}
  \proofstep{1,2,20}{\bot}{\rBI{19,21}}
  \proofstep{1,2}{x\MogDiamond y\neq k}{\rCI{20,22}}
\end{tabproof}

\FormulaThmDeltaK[Das neue Element absorbiert links]{%
  t\in T\vdash k\MogDiamond t=\mathbf0%
}{MogiljanskajaKAbsorbsLeft}{%
  \DeltaRow{Mengen}{T\dsep k\dsep\mathbf0\dsep t}
  \DeltaRow{Funktionen}{\MogDiamond}
}
\begin{tabproof}
  \proofstep{1}{t\in T}{\rA}
  \proofstep{}{k\in T}{\FormulaRefAuto{MogiljanskajaKInT}[thm]}
  \proofstep{}{k=k}{\rII}
  \proofstep{}{k=k\lor k=t}{\rOIa{3}}
  \proofstep{}{k\in\{k,t\}}{%
    \FormulaRefAuto{x\in\{a,b\}\eqvdash (x=a\lor x=b)}[thm]{4}}
  \proofstep{1}{k\MogDiamond t=\mathbf0}{%
    \FormulaRefAuto{MogiljanskajaDiamondDef}[def]{2,1,5}}
\end{tabproof}

\FormulaThmDeltaK[Das neue Element absorbiert rechts]{%
  t\in T\vdash t\MogDiamond k=\mathbf0%
}{MogiljanskajaKAbsorbsRight}{%
  \DeltaRow{Mengen}{T\dsep k\dsep\mathbf0\dsep t}
  \DeltaRow{Funktionen}{\MogDiamond}
}
\begin{tabproof}
  \proofstep{1}{t\in T}{\rA}
  \proofstep{}{k\in T}{\FormulaRefAuto{MogiljanskajaKInT}[thm]}
  \proofstep{}{k=k}{\rII}
  \proofstep{}{k=t\lor k=k}{\rOIb{3}}
  \proofstep{}{k\in\{t,k\}}{%
    \FormulaRefAuto{x\in\{a,b\}\eqvdash (x=a\lor x=b)}[thm]{4}}
  \proofstep{1}{t\MogDiamond k=\mathbf0}{%
    \FormulaRefAuto{MogiljanskajaDiamondDef}[def]{1,2,5}}
\end{tabproof}

\FormulaThmDeltaK[Ein angenommener Grundhalbgruppenisomorphismus
führt zum Widerspruch]{%
  \SemigroupIso{f}{S,\MogStar}{T,\MogDiamond}\vdash\bot%
}{MogiljanskajaIsoContradiction}{%
  \DeltaRow{Mengen}{L\dsep S\dsep T\dsep\mathbf0\dsep k\dsep s
    \dsep x\dsep y\dsep a_i\dsep a_j\dsep b_{ij}\dsep b_{i0}
    \dsep b_{i1}\dsep i\dsep j}
  \DeltaRow{Funktionen}{f\dsep\MogStar\dsep\MogDiamond}
}
\begin{tabproofsplitwide}
  \proofpartwideindR[Ein Urbild von (k) liegt in der (a)-Schicht]{%
    \begin{aligned}[t]
      &\SemigroupIso{f}{S,\MogStar}{T,\MogDiamond}\dsep
        s\in S\dsep f(s)=k\vdash{}\\[-2pt]
      &s\in L
    \end{aligned}
  }{%
    \SemigroupIso{f}{S,\MogStar}{T,\MogDiamond}\dsep
    s\in S\dsep f(s)=k\vdash s\in L%
  }[MogiljanskajaIsoKPreimageInL]
    \proofstepwidestar[1]{%
      \SemigroupIso{f}{S,\MogStar}{T,\MogDiamond}}{\rA}
    \proofstepwidestar[2]{s\in S}{\rA}
    \proofstepwidestar[3]{f(s)=k}{\rA}
    \proofstepwidestar[1]{%
      \SemigroupHom{f}{S,\MogStar}{T,\MogDiamond}}{%
      \FormulaRefAuto{SemigroupIsoHomAxiom}[ax]{1}}
    \proofstepwidestar[5]{\exists x,y\in S\,(s=x\MogStar y)}{\rA}
    \proofstepwidestar[6]{x,y\in S\land s=x\MogStar y}{\rA}
    \proofstepwidestar[6]{x,y\in S}{\rAEa{6}}
    \proofstepwidestar[6]{s=x\MogStar y}{\rAEb{6}}
    \proofstepwidestar[1,6]{%
      f(x\MogStar y)=f(x)\MogDiamond f(y)}{%
      \FormulaRefAuto{SemigroupHomOperationAxiom}[ax]{4,7}}
    \proofstepwidestar[6]{f(s)=f(x\MogStar y)}{\rLRS{8}}
    \proofstepwidestar[1,6]{f(s)=f(x)\MogDiamond f(y)}{%
      \rChain{10,9}}
    \proofstepwidestar[1,3,6]{f(x)\MogDiamond f(y)=k}{%
      \FormulaRefAuto{a=b\dsep a=c\vdash b=c}[thm]{11,3}}
    \proofstepwidestar[1]{f\colon S\to T}{%
      \FormulaRefAuto{SemigroupHomFunctionAxiom}[ax]{4}}
    \proofstepwidestar[1,6]{f(x)\in T}{%
      \FormulaRefAuto{F\colon A\to B\dsep x\in A\vdash F(x)\in B}[thm]{%
        13,7}}
    \proofstepwidestar[1,6]{f(y)\in T}{%
      \FormulaRefAuto{F\colon A\to B\dsep x\in A\vdash F(x)\in B}[thm]{%
        13,7}}
    \proofstepwidestar[1,6]{f(x)\MogDiamond f(y)\neq k}{%
      \FormulaRefAuto{MogiljanskajaKNotProduct}[thm]{14,15}}
    \proofstepwidestar[1,3,6]{\bot}{\rBI{16,12}}
    \proofstepwidestar[1,3,5]{\bot}{\rEE{5,6,17}}
    \proofstepwidestar[1,3]{%
      \neg\exists x,y\in S\,(s=x\MogStar y)}{\rCI{5,18}}
    \proofstepwidestar[20]{s\notin L}{\rA}
    \proofstepwidestar[2,20]{\exists x,y\in S\,(s=x\MogStar y)}{%
      \FormulaRefAuto{MogiljanskajaNonLIsProduct}[thm]{2,20}}
    \proofstepwidestar[1,2,3,20]{\bot}{\rBI{19,21}}
    \proofstepwidestar[1,2,3]{s\in L}{\rCI{20,22}}
  \closeproofpartwideindR

  \proofpartwideindR[Produktbilder unter einem Urbild von (k)]{%
    \begin{aligned}[t]
      &\SemigroupIso{f}{S,\MogStar}{T,\MogDiamond}\dsep
        i,j\in\mathbb N\dsep f(a_i)=k\vdash{}\\[-2pt]
      &f(b_{ij})=\mathbf0
    \end{aligned}
  }{%
    \SemigroupIso{f}{S,\MogStar}{T,\MogDiamond}\dsep
    i,j\in\mathbb N\dsep f(a_i)=k\vdash f(b_{ij})=\mathbf0%
  }[MogiljanskajaIsoRowProductImageZero]
    \proofstepwidestar[1]{%
      \SemigroupIso{f}{S,\MogStar}{T,\MogDiamond}}{\rA}
    \proofstepwidestar[2]{i\in\mathbb N}{\rA}
    \proofstepwidestar[3]{j\in\mathbb N}{\rA}
    \proofstepwidestar[4]{f(a_i)=k}{\rA}
    \proofstepwidestar[1]{%
      \SemigroupHom{f}{S,\MogStar}{T,\MogDiamond}}{%
      \FormulaRefAuto{SemigroupIsoHomAxiom}[ax]{1}}
    \proofstepwidestar[1]{f\colon S\to T}{%
      \FormulaRefAuto{SemigroupHomFunctionAxiom}[ax]{5}}
    \proofstepwidestar[2]{a_i\in L}{%
      \FormulaRefAuto{MogiljanskajaLayerLMembership}[thm]{2}}
    \proofstepwidestar[2]{a_i\in S}{%
      \FormulaRefAuto{MogiljanskajaCarrierSMembership}[thm]{7}}
    \proofstepwidestar[3]{a_j\in L}{%
      \FormulaRefAuto{MogiljanskajaLayerLMembership}[thm]{3}}
    \proofstepwidestar[3]{a_j\in S}{%
      \FormulaRefAuto{MogiljanskajaCarrierSMembership}[thm]{9}}
    \proofstepwidestar[1,3]{f(a_j)\in T}{%
      \FormulaRefAuto{F\colon A\to B\dsep x\in A\vdash F(x)\in B}[thm]{%
        6,10}}
    \proofstepwidestar[2,3]{a_i\MogStar a_j=b_{ij}}{%
      \FormulaRefAuto{MogiljanskajaStarDef}[def]{2,3}}
    \proofstepwidestar[2,3]{b_{ij}=a_i\MogStar a_j}{%
      \FormulaRefAuto{a=b\vdash b=a}[thm]{12}}
    \proofstepwide[2,3]{f(b_{ij})}{=}{f(a_i\MogStar a_j)}{\rLRS{13}}
    \proofstepwide[1,2,3]{}{=}{f(a_i)\MogDiamond f(a_j)}{%
      \FormulaRefAuto{SemigroupHomOperationAxiom}[ax]{5,8,10}}
    \proofstepwide[1,2,3,4]{}{=}{k\MogDiamond f(a_j)}{\rLRS{4}}
    \proofstepwide[1,3]{}{=}{\mathbf0}{%
      \FormulaRefAuto{MogiljanskajaKAbsorbsLeft}[thm]{11}}
    \proofstepwidestar[1,2,3,4]{f(b_{ij})=\mathbf0}{\rChain{14,17}}
  \closeproofpartwideindR

  \proofpartwide{Schluss}
    \proofstepwidestar[1]{%
      \SemigroupIso{f}{S,\MogStar}{T,\MogDiamond}}{\rA}
    \proofstepwidestar[1]{f\colon S\bij T}{%
      \FormulaRefAuto{SemigroupIsoBijectiveAxiom}[ax]{1}}
    \proofstepwidestar[1]{f\colon S\sur T}{%
      \FormulaRefAuto{F\colon A\bij B\vdash F\colon A\sur B}[thm]{2}}
    \proofstepwidestar[]{k\in T}{%
      \FormulaRefAuto{MogiljanskajaKInT}[thm]}
    \proofstepwidestar[1]{\exists s\in S\,(f(s)=k)}{%
      \FormulaRefAuto{Surjektivität}[ax]{3,4}}
    \proofstepwidestar[6]{s\in S\land f(s)=k}{\rA}
    \proofstepwidestar[6]{s\in S}{\rAEa{6}}
    \proofstepwidestar[6]{f(s)=k}{\rAEb{6}}
    \proofstepwidestar[1,6]{s\in L}{%
      \FormulaRefAuto{MogiljanskajaIsoKPreimageInL}[thm]{1,7,8}}
    \proofstepwidestar[1,6]{\exists i\in\mathbb N\,(s=a_i)}{%
      \FormulaRefAuto{MogiljanskajaLayerLMembership}[thm]{9}}
    \proofstepwidestar[11]{i\in\mathbb N\land s=a_i}{\rA}
    \proofstepwidestar[11]{i\in\mathbb N}{\rAEa{11}}
    \proofstepwidestar[11]{s=a_i}{\rAEb{11}}
    \proofstepwidestar[6,11]{f(a_i)=k}{\rIE{13,8}}
    \proofstepwidestar[]{0\in\mathbb N}{\FormulaRefAuto{PeanoZero}[ax]}
    \proofstepwidestar[]{1\in\mathbb N}{%
      \FormulaRefAuto{PeanoOneInN}[thm]}
    \proofstepwidestar[1,6,11]{f(b_{i0})=\mathbf0}{%
      \FormulaRefAuto{MogiljanskajaIsoRowProductImageZero}[thm]{%
        1,12,15,14}}
    \proofstepwidestar[1,6,11]{f(b_{i1})=\mathbf0}{%
      \FormulaRefAuto{MogiljanskajaIsoRowProductImageZero}[thm]{%
        1,12,16,14}}
    \proofstepwidestar[1,6,11]{f(b_{i0})=f(b_{i1})}{%
      \FormulaRefAuto{a=b\dsep c=b\vdash a=c}[thm]{17,18}}
    \proofstepwidestar[1]{f\colon S\inj T}{%
      \FormulaRefAuto{F\colon A\bij B\vdash F\colon A\inj B}[thm]{2}}
    \proofstepwidestar[11]{b_{i0}\in D}{%
      \FormulaRefAuto{MogiljanskajaLayerDMembership}[thm]{12,15}}
    \proofstepwidestar[11]{b_{i0}\in S}{%
      \FormulaRefAuto{MogiljanskajaCarrierSMembership}[thm]{21}}
    \proofstepwidestar[11]{b_{i1}\in D}{%
      \FormulaRefAuto{MogiljanskajaLayerDMembership}[thm]{12,16}}
    \proofstepwidestar[11]{b_{i1}\in S}{%
      \FormulaRefAuto{MogiljanskajaCarrierSMembership}[thm]{23}}
    \proofstepwidestar[1,6,11]{b_{i0}=b_{i1}}{%
      \FormulaRefAuto{Injektivität}[ax]{20,22,24,19}}
    \proofstepwidestar[1,6,11]{i=i\land0=1}{%
      \FormulaRefAuto{MogiljanskajaBIndexUnique}[thm]{12,15,12,16,25}}
    \proofstepwidestar[1,6,11]{0=1}{\rAEb{26}}
    \proofstepwidestar[]{0\neq1}{%
      \FormulaRefAuto{PeanoZeroNeOne}[thm]}
    \proofstepwidestar[1,6,11]{\bot}{\rBI{28,27}}
    \proofstepwidestar[1,6]{\bot}{\rEE{10,11,29}}
    \proofstepwidestar[1]{\bot}{\rEE{5,6,30}}
  \closeproofpartwide
\end{tabproofsplitwide}

\FormulaThmDeltaK[Kein Grundhalbgruppenisomorphismus existiert]{%
  \neg\exists f\,\SemigroupIso{f}{S,\MogStar}{T,\MogDiamond}%
}{MogiljanskajaNotIsomorphic}{%
  \DeltaRow{Mengen}{S\dsep T}
  \DeltaRow{Funktionen}{f\dsep\MogStar\dsep\MogDiamond}
}
\begin{tabproof}
  \proofstep{1}{\exists f\,\SemigroupIso{f}{S,\MogStar}{T,\MogDiamond}}{\rA}
  \proofstep{2}{\SemigroupIso{f}{S,\MogStar}{T,\MogDiamond}}{\rA}
  \proofstep{2}{\bot}{%
    \FormulaRefAuto{MogiljanskajaIsoContradiction}[thm]{2}}
  \proofstep{1}{\bot}{\rEE{1,2,3}}
  \proofstep{}{\neg\exists f\,\SemigroupIso{f}{S,\MogStar}{T,\MogDiamond}}{%
    \rCI{1,4}}
\end{tabproof}

Die Folgen- und Mengenkonstruktion der Reserve ist bereits in Band~17
vollständig ausgeführt. Metasprachlich verwenden wir
\[
  c_0=b_{00},\quad c_1=b_{11},\quad c_2=b_{10},\qquad
  U=\{b_{\Succ(\Succ(n)),j}\mid n,j\in\mathbb N\},
\]
sowie \(V=U\cup\{c_2\}\), \(P=\{c_0,c_1\}\) und
\(D'=P\cup V\); dies ist
\FormulaRefAuto{MogiljanskajaDPrimeDef}[def].
Die benötigten Markerfakten
\[
  c_0,c_1,c_2\in D\setminus U,\qquad
  D'\subseteq D,\qquad b_{01}\notin D',\qquad
  c_2\notin P,\qquad P\cap U=\varnothing
\]
stehen in \FormulaRefAuto{MogiljanskajaReserveMarkerFacts}[thm] und
\FormulaRefAuto{MogiljanskajaReserveCoreFacts}[thm].

Ebenfalls aus Band~17 übernehmen wir die Bijektionen
\[
  \theta\colon U\bij D,
  \qquad
  \sigma\colon U\bij V,
\]
vgl. \FormulaRefAuto{MogiljanskajaThetaBijective}[thm] und
\FormulaRefAuto{MogiljanskajaSigmaBijective}[thm].
Die der zweiten Bijektion zugrunde liegende Folge
\(H_n=b_{\Succ(n),0}\) ist injektiv in \(V\), aber keine Aufzählung von
\(V\): Nach \FormulaRefAuto{MogiljanskajaShiftSequenceImage}[thm] gilt
\[
  H[\mathbb N]
  =\{b_{\Succ(n),0}\mid n\in\mathbb N\}
  \subsetneq V.
\]
Für die folgenden halbgruppentheoretischen Argumente werden nur die so
bereitgestellten Mengen, Marker und Bijektionen benötigt.

\FormulaDefDeltaK[Reservefamilien mit festen Kernen]{%
  \begin{aligned}[t]
    \mathcal C&\coloneqq\CoreFamily{P}{V},&
    \mathcal C_0&\coloneqq\CoreFamily{P}{U},\\[-2pt]
    \mathcal C_1&\coloneqq\CoreFamily{P\cup\{c_2\}}{U},&
    \mathcal E&\coloneqq\CoreFamily{\{k\}}{D}
  \end{aligned}%
}{MogiljanskajaReserveFamiliesDef}{%
  \DeltaRow{Mengen}{D\dsep U\dsep V\dsep P\dsep c_2\dsep k\dsep
    \mathcal C\dsep\mathcal C_0\dsep\mathcal C_1\dsep\mathcal E}
}

\FormulaThmDeltaK[Disjunkte Kerne und Zerlegung der Reservefamilie]{%
  \begin{aligned}[t]
    &P\cap U=\varnothing\dsep P\cap V=\varnothing\dsep
      (P\cup\{c_2\})\cap U=\varnothing\dsep\{k\}\cap D=\varnothing,\\[-2pt]
    &\mathcal C=\mathcal C_0\cup\mathcal C_1\dsep
      \mathcal C_0\cap\mathcal C_1=\varnothing\dsep
      \mathcal C\subseteq\powerset(D)\dsep
      \mathcal E\subseteq\powerset(D\cup\{k\})\dsep
      \mathcal C\cap\mathcal E=\varnothing
  \end{aligned}%
}{MogiljanskajaReserveFamilyFacts}{%
  \DeltaRow{Mengen}{D\dsep D'\dsep U\dsep V\dsep P\dsep c_2\dsep k
    \dsep\mathcal C\dsep\mathcal C_0\dsep\mathcal C_1\dsep\mathcal E
    \dsep K}
}
\begin{tabproofwide}
  \proofstepwidestar[]{c_0,c_1,c_2\notin U}{%
    \rAEa{\rAEb{\FormulaRefAuto{MogiljanskajaReserveMarkerFacts}[thm]}}}
  \MogProofStepWideStarBreak[]{c_2\notin P\cup U}{%
    \FormulaRefAuto{NotInUnion}[thm]{%
      \rAEa{\FormulaRefAuto{MogiljanskajaReserveCoreFacts}[thm]},
      \rAEb{\rAEb{1}}}}
  \proofstepwidestar[]{P\cap U=\varnothing}{%
    \rAEb{\FormulaRefAuto{MogiljanskajaReserveCoreFacts}[thm]}}
  \proofstepwidestar[]{P\cap\{c_2\}=\varnothing}{%
    \FormulaRefAuto{AvoidedSingletonDisjoint}[thm]{%
      \rAEa{\FormulaRefAuto{MogiljanskajaReserveCoreFacts}[thm]}}}
  \MogProofStepWideStarBreak[]{P\cap V=\varnothing}{%
    \FormulaRefAuto{A\cap(B\cup C)=(A\cap B)\cup(A\cap C)}[thm]{%
      \FormulaRefAuto{MogiljanskajaDPrimeDef}[def],3,4}}
  \MogProofStepWideStarBreak[]{(P\cup\{c_2\})\cap U=\varnothing}{%
    \FormulaRefAuto{(A\cup B)\cap C=(A\cap C)\cup(B\cap C)}[thm]{%
      3,\FormulaRefAuto{SingletonDisjointFromAvoidedSet}[thm]{\rAEb{\rAEb{1}}}}}
  \proofstepwidestar[]{k\notin D}{%
    \FormulaRefAuto{MogiljanskajaKNotInD}[thm]{%
      \FormulaRefAuto{MogiljanskajaKOutsideLayers}[thm]}}
  \proofstepwidestar[]{\{k\}\cap D=\varnothing}{%
    \FormulaRefAuto{SingletonDisjointFromAvoidedSet}[thm]{7}}
  \MogProofStepWideStarBreak[]{\mathcal C=\mathcal C_0\cup\mathcal C_1}{%
    \rIE{%
      \FormulaRefAuto{MogiljanskajaReserveFamiliesDef}[def],
      \FormulaRefAuto{MogiljanskajaDPrimeDef}[def],
      \rAEa{\FormulaRefAuto{CoreFamilyAdjoinedSplit}[thm]{2}}}}
  \MogProofStepWideStarBreak[]{\mathcal C_0\cap\mathcal C_1=\varnothing}{%
    \rIE{%
      \FormulaRefAuto{MogiljanskajaReserveFamiliesDef}[def],
      \rAEb{\FormulaRefAuto{CoreFamilyAdjoinedSplit}[thm]{2}}}}
  \MogProofStepWideStarBreak[]{D'\subseteq D}{%
    \rAEa{\rAEb{\rAEb{\rAEb{\FormulaRefAuto{MogiljanskajaReserveMarkerFacts}[thm]}}}}}
  \MogProofStepWideStarBreak[]{\mathcal C\subseteq\powerset(D)}{%
    \FormulaRefAuto{A\subseteq B\dsep B\subseteq C\vdash A\subseteq C}[thm]{%
      \rIE{\FormulaRefAuto{MogiljanskajaReserveFamiliesDef}[def],
        \FormulaRefAuto{CoreFamilySubsetPowerset}[thm]},
      \FormulaRefAuto{A \subseteq B \vdash
        \powerset(A) \subseteq \powerset(B)}[thm]{%
        \rIE{\FormulaRefAuto{MogiljanskajaDPrimeDef}[def],11}}}}
  \MogProofStepWideStarBreak[]{\mathcal E\subseteq\powerset(D\cup\{k\})}{%
    \rIE{\FormulaRefAuto{MogiljanskajaReserveFamiliesDef}[def],
      \FormulaRefAuto{A\cup B=B\cup A}[thm],
      \FormulaRefAuto{CoreFamilySubsetPowerset}[thm]}}
  \proofstepwidestar[14]{K\in\mathcal C\cap\mathcal E}{\rA}
  \proofstepwidestar[14]{K\in\mathcal C\land K\in\mathcal E}{%
    \FormulaRefAuto{x\in A\cap B\eqvdash x\in A\land x\in B}[thm]{14}}
  \proofstepwidestar[14]{K\subseteq D}{%
    \FormulaRefAuto{x\in\powerset(A)\eqvdash x\subseteq A}[thm]{%
      \FormulaRefAuto{A\subseteq B\dsep x\in A\vdash x\in B}[thm]{%
        12,\rAEa{15}}}}
  \MogProofStepWideStarBreak[14]{\{k\}\subseteq K}{%
    \rAEa{\FormulaRefAuto{CoreFamilyMembership}[thm]{%
      \rAEb{15},\FormulaRefAuto{MogiljanskajaReserveFamiliesDef}[def]}}}
  \proofstepwidestar[14]{k\in K}{%
    \FormulaRefAuto{A\subseteq B\dsep x\in A\vdash x\in B}[thm]{%
      17,\FormulaRefAuto{a\in\{a\}}[thm]}}
  \proofstepwidestar[14]{k\in D}{%
    \FormulaRefAuto{A\subseteq B\dsep x\in A\vdash x\in B}[thm]{16,18}}
  \proofstepwidestar[14]{\bot}{\rBI{7,19}}
  \MogProofStepWideStarBreak[]{\mathcal C\cap\mathcal E=\varnothing}{%
    \FormulaRefAuto{\forall x\,x\notin A\vdash A=\varnothing}[thm]{%
      \rUI{\rCI{14,20}}}}
  \MogProofStepWideStarBreak[]{%
    \begin{aligned}[t]
      &P\cap U=\varnothing\dsep P\cap V=\varnothing\dsep{}\\[-2pt]
      &(P\cup\{c_2\})\cap U=\varnothing\dsep
        \{k\}\cap D=\varnothing,\\[-2pt]
      &\mathcal C=\mathcal C_0\cup\mathcal C_1\dsep
        \mathcal C_0\cap\mathcal C_1=\varnothing\dsep{}\\[-2pt]
      &\mathcal C\subseteq\powerset(D)\dsep
        \mathcal E\subseteq\powerset(D\cup\{k\})\dsep{}\\[-2pt]
      &\mathcal C\cap\mathcal E=\varnothing
    \end{aligned}}{%
    \rAI{3,\rAI{5,\rAI{6,\rAI{8,
      \rAI{9,\rAI{10,\rAI{12,\rAI{13,21}}}}}}}}}
\end{tabproofwide}

\FormulaDefDeltaK[Die beiden Kerntransporte und ihre Verklebung]{%
  \begin{aligned}[t]
    \kappa_0&\coloneqq\CoreTransport{P}{P}{\sigma},\\[-2pt]
    \kappa_1&\coloneqq\CoreTransport{P\cup\{c_2\}}{\{k\}}{\theta},\\[-2pt]
    \psi&\coloneqq\CaseFun{\mathcal C_0}{\mathcal C_1}{\kappa_0}{\kappa_1}
  \end{aligned}%
}{MogiljanskajaPsiDef}{%
  \DeltaRow{Mengen}{P\dsep U\dsep V\dsep D\dsep c_2\dsep k\dsep
    \mathcal C\dsep\mathcal C_0\dsep\mathcal C_1\dsep\mathcal E}
  \DeltaRow{Funktionen}{\sigma\dsep\theta\dsep\kappa_0\dsep\kappa_1\dsep\psi}
}

\FormulaThmDeltaK[Die Reserve absorbiert die neue Mengenfamilie]{%
  \psi\colon\mathcal C\bij\mathcal C\cup\mathcal E%
}{MogiljanskajaPsiBijective}{%
  \DeltaRow{Mengen}{P\dsep U\dsep V\dsep D\dsep c_2\dsep k\dsep
    \mathcal C\dsep\mathcal C_0\dsep\mathcal C_1\dsep\mathcal E}
  \DeltaRow{Funktionen}{\sigma\dsep\theta\dsep\kappa_0\dsep\kappa_1\dsep\psi}
}
\begin{tabproofwide}
  \proofstepwidestar[]{P\cap U=\varnothing}{%
    \rAEa{\FormulaRefAuto{MogiljanskajaReserveFamilyFacts}[thm]}}
  \proofstepwidestar[]{c_0,c_1,c_2\notin U}{%
    \rAEa{\rAEb{\FormulaRefAuto{MogiljanskajaReserveMarkerFacts}[thm]}}}
  \proofstepwidestar[]{c_2\notin P\cup U}{%
    \FormulaRefAuto{NotInUnion}[thm]{%
      \rAEa{\FormulaRefAuto{MogiljanskajaReserveCoreFacts}[thm]},
      \rAEb{\rAEb{2}}}}
  \proofstepwidestar[]{\{k\}\cap D=\varnothing}{%
    \rAEa{\rAEb{\rAEb{\rAEb{%
      \FormulaRefAuto{MogiljanskajaReserveFamilyFacts}[thm]}}}}}
  \proofstepwidestar[]{\sigma\colon U\bij U\cup\{c_2\}}{%
    \rIE{\FormulaRefAuto{MogiljanskajaDPrimeDef}[def],
      \FormulaRefAuto{MogiljanskajaSigmaBijective}[thm]}}
  \proofstepwidestar[]{\theta\colon U\bij D}{%
    \FormulaRefAuto{MogiljanskajaThetaBijective}[thm]}
  \MogProofStepWideStarBreak[]{%
    \begin{aligned}[t]
      &\CoreFamily{P}{U\cup\{c_2\}}\cap{}\\[-2pt]
      &\quad\CoreFamily{\{k\}}{D}=\varnothing
    \end{aligned}}{%
    \rIE{%
      \FormulaRefAuto{MogiljanskajaReserveFamiliesDef}[def],
      \FormulaRefAuto{MogiljanskajaDPrimeDef}[def],
      \rAEb{\rAEb{\rAEb{\rAEb{\rAEb{\rAEb{\rAEb{\rAEb{%
        \FormulaRefAuto{MogiljanskajaReserveFamilyFacts}[thm]}}}}}}}}}}
  \MogProofStepWideStarBreak[]{\psi\colon\mathcal C\bij\mathcal C\cup\mathcal E}{%
    \rIE{%
      \FormulaRefAuto{MogiljanskajaPsiDef}[def],
      \FormulaRefAuto{MogiljanskajaReserveFamiliesDef}[def],
      \FormulaRefAuto{MogiljanskajaDPrimeDef}[def],
      \FormulaRefAuto{CoreFamilyAbsorptionBijective}[thm]{1,3,4,5,6,7}}}
\end{tabproofwide}

\FormulaDefDeltaK[Der unveränderte Teil und die lokale Potenzmengenbijektion]{%
  \begin{aligned}[t]
    \mathcal O&\coloneqq\powerset(D)\setminus\mathcal C,\\[-2pt]
    \eta&\coloneqq\CaseFun{\mathcal O}{\mathcal C}{\Id_{\mathcal O}}{\psi}
  \end{aligned}%
}{MogiljanskajaEtaDef}{%
  \DeltaRow{Mengen}{D\dsep k\dsep\mathcal O\dsep\mathcal C\dsep\mathcal E}
  \DeltaRow{Funktionen}{\psi\dsep\eta}
}

\FormulaThmDeltaK[Geschützte Eigenschaften der lokalen Potenzmengenabbildung]{%
  \begin{aligned}[t]
    &\bigl(\eta\colon\powerset(D)\bij\powerset(D\cup\{k\})\bigr)
      \land{}\\[-2pt]
    &\quad\Bigl(\eta(\varnothing)=\varnothing\land{}\\[-2pt]
    &\qquad\forall K\in\powerset(D)\,\\[-2pt]
    &\qquad\quad
      \bigl(K\notin\mathcal C\rightarrow\eta(K)=K\bigr)\Bigr)
  \end{aligned}%
}{MogiljanskajaEtaProperties}{%
  \DeltaRow{Mengen}{D\dsep P\dsep k\dsep K\dsep
    \mathcal O\dsep\mathcal C\dsep\mathcal E}
  \DeltaRow{Funktionen}{\psi\dsep\eta}
}
\begin{tabproofwide}
  \proofstepwidestar[]{k\notin D}{%
    \FormulaRefAuto{MogiljanskajaKNotInD}[thm]{%
      \FormulaRefAuto{MogiljanskajaKOutsideLayers}[thm]}}
  \MogProofStepWideStarBreak[]{\mathcal C\subseteq\powerset(D)}{%
    \rAEa{\rAEb{\rAEb{\rAEb{\rAEb{\rAEb{\rAEb{%
      \FormulaRefAuto{MogiljanskajaReserveFamilyFacts}[thm]}}}}}}}}
  \proofstepwidestar[]{P\neq\varnothing}{%
    \FormulaRefAuto{a\in\{a,b\}}[thm]{%
      \FormulaRefAuto{MogiljanskajaDPrimeDef}[def]}}
  \proofstepwidestar[]{\varnothing\notin\mathcal C}{%
    \FormulaRefAuto{CoreFamilyMembership}[thm]{%
      \FormulaRefAuto{MogiljanskajaReserveFamiliesDef}[def],3}}
  \proofstepwidestar[]{%
    \mathcal E=\mathfrak{I}[\{k\},D\cup\{k\}]}{%
    \rIE{%
      \FormulaRefAuto{MogiljanskajaReserveFamiliesDef}[def],
      \FormulaRefAuto{CoreFamilyBooleanInterval}[thm],
      \FormulaRefAuto{A\cup B=B\cup A}[thm]}}
  \proofstepwidestar[]{%
    \psi\colon\mathcal C\bij
      \mathcal C\cup\mathfrak{I}[\{k\},D\cup\{k\}]}{%
    \rIE{\FormulaRefAuto{MogiljanskajaPsiBijective}[thm],5}}
  \proofstepwidestar[]{\mathcal O=\powerset(D)\setminus\mathcal C}{%
    \FormulaRefAuto{MogiljanskajaEtaDef}[def]}
  \proofstepwidestar[]{%
    \eta=\CaseFun{\mathcal O}{\mathcal C}{\Id_{\mathcal O}}{\psi}}{%
    \FormulaRefAuto{MogiljanskajaEtaDef}[def]}
  \MogProofStepWideStarBreak[]{%
    \begin{aligned}[t]
      &\bigl(\eta\colon\powerset(D)\bij\powerset(D\cup\{k\})\bigr)
        \land{}\\[-2pt]
      &\quad\Bigl(\eta(\varnothing)=\varnothing\land{}\\[-2pt]
      &\qquad\forall K\in\powerset(D)\,\\[-2pt]
      &\qquad\quad
        \bigl(K\notin\mathcal C\rightarrow\eta(K)=K\bigr)\Bigr)
    \end{aligned}}{%
    \FormulaRefAuto{ProtectedPowerSetExtension}[thm]{1,2,4,6,7,8}}
\end{tabproofwide}

\FormulaThmDeltaK[Die lokale Potenzmengenabbildung ist bijektiv und nulltreu]{%
  \bigl(\eta\colon\powerset(D)\bij\powerset(D\cup\{k\})\bigr)\land
  \eta(\varnothing)=\varnothing%
}{MogiljanskajaEtaBijective}{%
  \DeltaRow{Mengen}{D\dsep k\dsep\mathcal O\dsep\mathcal C\dsep\mathcal E}
  \DeltaRow{Funktionen}{\psi\dsep\eta}
}
\begin{tabproofwide}
  \proofstepwidestar[]{%
    \begin{aligned}[t]
      &\bigl(\eta\colon\powerset(D)\bij\powerset(D\cup\{k\})\bigr)
        \land\Bigl(\eta(\varnothing)=\varnothing\land{}\\[-2pt]
      &\qquad\forall K\in\powerset(D)\,
        \bigl(K\notin\mathcal C\rightarrow\eta(K)=K\bigr)\Bigr)
    \end{aligned}}{%
    \FormulaRefAuto{MogiljanskajaEtaProperties}[thm]}
  \proofstepwidestar[]{%
    \eta\colon\powerset(D)\bij\powerset(D\cup\{k\})}{\rAEa{1}}
  \proofstepwidestar[]{\eta(\varnothing)=\varnothing}{\rAEa{\rAEb{1}}}
  \proofstepwidestar[]{%
    \bigl(\eta\colon\powerset(D)\bij\powerset(D\cup\{k\})\bigr)\land
    \eta(\varnothing)=\varnothing}{\rAI{2,3}}
\end{tabproofwide}

\FormulaThmDeltaK[Außerhalb der Reserve wirkt die lokale Abbildung identisch]{%
  K\in\powerset(D)\dsep K\notin\mathcal C\vdash\eta(K)=K%
}{MogiljanskajaEtaFixesOutsideReserve}{%
  \DeltaRow{Mengen}{D\dsep\mathcal O\dsep\mathcal C\dsep\mathcal E\dsep K}
  \DeltaRow{Funktionen}{\psi\dsep\eta}
}
\begin{tabproof}
  \proofstep{1}{K\in\powerset(D)}{\rA}
  \proofstep{2}{K\notin\mathcal C}{\rA}
  \proofstep{}{%
    \forall J\in\powerset(D)\,
      \bigl(J\notin\mathcal C\rightarrow\eta(J)=J\bigr)}{%
    \rAEb{\rAEb{\FormulaRefAuto{MogiljanskajaEtaProperties}[thm]}}}
  \proofstep{1}{K\notin\mathcal C\rightarrow\eta(K)=K}{%
    \rRE{\rUE{3},1}}
  \proofstep{1,2}{\eta(K)=K}{\rRE{4,2}}
\end{tabproof}
\FormulaDefDeltaK[Geschützte Schicht und globale Abbildung]{%
  \begin{aligned}[t]
    A_*&\coloneqq L\cup\{\mathbf0\},\\[-2pt]
    \Phi&\coloneqq
      \LayerPowerMap{A_*}{D}{D\cup\{k\}}{\eta}
      \restriction_{\PowNE{S}}
  \end{aligned}%
}{MogiljanskajaPhiDef}{%
  \DeltaRow{Mengen}{A_*\dsep L\dsep D\dsep S\dsep T\dsep\mathbf0\dsep k}
  \DeltaRow{Funktionen}{\eta\dsep\Phi}
}

\FormulaThmDeltaK[Trägerzerlegung für die Schichtfortsetzung]{%
  \begin{aligned}[t]
    &A_*\cap D=\varnothing\dsep A_*\cap(D\cup\{k\})=\varnothing,\\[-2pt]
    &S=A_*\cup D\dsep T=A_*\cup(D\cup\{k\})
  \end{aligned}%
}{MogiljanskajaLayerCarrierFacts}{%
  \DeltaRow{Mengen}{A_*\dsep L\dsep D\dsep S\dsep T\dsep\mathbf0\dsep k}
}
\begin{tabproofwide}
  \proofstepwidestar[]{L\cap D=\varnothing}{%
    \FormulaRefAuto{MogiljanskajaLayersDisjoint}[thm]}
  \proofstepwidestar[]{\mathbf0\notin D}{%
    \FormulaRefAuto{MogiljanskajaZeroNotInD}[thm]{%
      \FormulaRefAuto{MogiljanskajaZeroOutsideLayers}[thm]}}
  \proofstepwidestar[]{\{\mathbf0\}\cap D=\varnothing}{%
    \FormulaRefAuto{SingletonDisjointFromAvoidedSet}[thm]{2}}
  \proofstepwidestar[]{A_*\cap D=\varnothing}{%
    \FormulaRefAuto{(A\cup B)\cap C=(A\cap C)\cup(B\cap C)}[thm]{%
      \FormulaRefAuto{MogiljanskajaPhiDef}[def],1,3}}
  \proofstepwidestar[]{k\notin A_*}{%
    \FormulaRefAuto{MogiljanskajaKOutsideCarrierPieces}[thm]{%
      \FormulaRefAuto{MogiljanskajaPhiDef}[def]}}
  \proofstepwidestar[]{A_*\cap\{k\}=\varnothing}{%
    \FormulaRefAuto{AvoidedSingletonDisjoint}[thm]{5}}
  \proofstepwidestar[]{A_*\cap(D\cup\{k\})=\varnothing}{%
    \FormulaRefAuto{A\cap(B\cup C)=(A\cap B)\cup(A\cap C)}[thm]{4,6}}
  \MogProofStepWideStarBreak[]{S=A_*\cup D}{%
    \FormulaRefAuto{MogiljanskajaCarrierSDef}[def]{%
      \FormulaRefAuto{MogiljanskajaPhiDef}[def],
      \FormulaRefAuto{A\cup(B\cup C)=(A\cup B)\cup C}[thm]}}
  \MogProofStepWideStarBreak[]{T=A_*\cup(D\cup\{k\})}{%
    \FormulaRefAuto{MogiljanskajaCarrierTDef}[def]{8,
      \FormulaRefAuto{A\cup(B\cup C)=(A\cup B)\cup C}[thm]}}
  \MogProofStepWideStarBreak[]{%
    \begin{aligned}[t]
      &A_*\cap D=\varnothing\dsep
        A_*\cap(D\cup\{k\})=\varnothing,\\[-2pt]
      &S=A_*\cup D\dsep T=A_*\cup(D\cup\{k\})
    \end{aligned}}{%
    \rAI{4,\rAI{7,\rAI{8,9}}}}
\end{tabproofwide}

\FormulaThmDeltaK[Die globale Abbildung ist bijektiv]{%
  \Phi\colon\PowNE{S}\bij\PowNE{T}%
}{MogiljanskajaPhiBijective}{%
  \DeltaRow{Mengen}{A_*\dsep D\dsep S\dsep T\dsep k}
  \DeltaRow{Funktionen}{\eta\dsep\Phi}
}
\begin{tabproofwide}
  \proofstepwidestar[]{A_*\cap D=\varnothing}{%
    \rAEa{\FormulaRefAuto{MogiljanskajaLayerCarrierFacts}[thm]}}
  \proofstepwidestar[]{A_*\cap(D\cup\{k\})=\varnothing}{%
    \rAEa{\rAEb{\FormulaRefAuto{MogiljanskajaLayerCarrierFacts}[thm]}}}
  \proofstepwidestar[]{S=A_*\cup D}{%
    \rAEa{\rAEb{\rAEb{\FormulaRefAuto{MogiljanskajaLayerCarrierFacts}[thm]}}}}
  \proofstepwidestar[]{T=A_*\cup(D\cup\{k\})}{%
    \rAEb{\rAEb{\rAEb{\FormulaRefAuto{MogiljanskajaLayerCarrierFacts}[thm]}}}}
  \proofstepwidestar[]{\eta\colon\powerset(D)\bij\powerset(D\cup\{k\})}{%
    \rAEa{\FormulaRefAuto{MogiljanskajaEtaBijective}[thm]}}
  \proofstepwidestar[]{\eta(\varnothing)=\varnothing}{%
    \rAEb{\FormulaRefAuto{MogiljanskajaEtaBijective}[thm]}}
  \MogProofStepWideStarBreak[]{\Phi\colon\PowNE{S}\bij\PowNE{T}}{%
    \rIE{\FormulaRefAuto{MogiljanskajaPhiDef}[def],3,4,
      \FormulaRefAuto{LayerPowerMapNonemptyBijective}[thm]{1,2,5,6}}}
\end{tabproofwide}

\FormulaThmDeltaK[Grundgleichung und Invarianten der globalen Abbildung]{%
  \begin{aligned}[t]
    &X\in\PowNE{S}\vdash
      \Phi(X)=(X\cap A_*)\cup\eta(X\cap D),\\[-2pt]
    &X\in\PowNE{S}\vdash\Phi(X)\cap L=X\cap L,\\[-2pt]
    &X\in\PowNE{S}\vdash
      \bigl(X\subseteq L\leftrightarrow\Phi(X)\subseteq L\bigr)
  \end{aligned}%
}{MogiljanskajaPhiInvariants}{%
  \DeltaRow{Mengen}{A_*\dsep L\dsep D\dsep S\dsep T\dsep k\dsep X}
  \DeltaRow{Funktionen}{\eta\dsep\Phi}
}
\begin{tabproofwide}
  \proofstepwidestar[1]{X\in\PowNE{S}}{\rA}
  \proofstepwidestar[1]{X\in\powerset(S)}{%
    \FormulaRefAuto{A\setminus B\subseteq A}[thm]{1}}
  \proofstepwidestar[]{S=A_*\cup D}{%
    \rAEa{\rAEb{\rAEb{\FormulaRefAuto{MogiljanskajaLayerCarrierFacts}[thm]}}}}
  \proofstepwidestar[]{\PowNE{S}\subseteq\powerset(A_*\cup D)}{%
    \rIE{3,\FormulaRefAuto{A\setminus B\subseteq A}[thm]}}
  \proofstepwidestar[1]{X\in\powerset(A_*\cup D)}{\rIE{3,2}}
  \proofstepwidestar[]{\eta\colon\powerset(D)\to\powerset(D\cup\{k\})}{%
    \FormulaRefAuto{Bijektive Funktion}[def]{%
      \rAEa{\FormulaRefAuto{MogiljanskajaEtaBijective}[thm]}}}
  \MogProofStepWideStarBreak[1]{\Phi(X)=
    \LayerPowerMap{A_*}{D}{D\cup\{k\}}{\eta}(X)}{%
    \FormulaRefAuto{C\subseteq A\dsep x\in C\vdash F\restriction_C(x)=F(x)}[thm]{%
      \FormulaRefAuto{MogiljanskajaPhiDef}[def],
      4,1}}
  \proofstepwidestar[1]{\LayerPowerMap{A_*}{D}{D\cup\{k\}}{\eta}(X)
    =(X\cap A_*)\cup\eta(X\cap D)}{%
    \FormulaRefAuto{LayerPowerMapValue}[thm]{6,5}}
  \proofstepwidestar[1]{\Phi(X)=(X\cap A_*)\cup\eta(X\cap D)}{%
    \rChain{7,8}}
  \proofstepwidestar[]{A_*\cap(D\cup\{k\})=\varnothing}{%
    \rAEa{\rAEb{\FormulaRefAuto{MogiljanskajaLayerCarrierFacts}[thm]}}}
  \proofstepwidestar[1]{%
    \LayerPowerMap{A_*}{D}{D\cup\{k\}}{\eta}(X)\cap A_*=X\cap A_*}{%
    \FormulaRefAuto{LayerPowerMapStablePart}[thm]{10,6,5}}
  \proofstepwidestar[]{L\subseteq A_*}{%
    \FormulaRefAuto{A\subseteq A\cup B}[thm]{%
      \FormulaRefAuto{MogiljanskajaPhiDef}[def]}}
  \proofstepwidestar[]{A_*\cap L=L}{%
    \FormulaRefAuto{B\subseteq A\eqvdash A\cap B=B}[thm]{12}}
  \MogProofStepWideStarBreak[1]{%
    \Phi(X)\cap L=
      \LayerPowerMap{A_*}{D}{D\cup\{k\}}{\eta}(X)\cap L}{%
    \rLRS{7}}
  \MogProofStepWideStarBreak[1]{%
    \begin{aligned}[t]
      &\LayerPowerMap{A_*}{D}{D\cup\{k\}}{\eta}(X)\cap L\\[-2pt]
      &\quad=\bigl(\LayerPowerMap{A_*}{D}{D\cup\{k\}}{\eta}(X)
        \cap A_*\bigr)\cap L
    \end{aligned}}{%
    \rChain{%
      \rLRS{\FormulaRefAuto{a=b\vdash b=a}[thm]{%
        13}},
      \FormulaRefAuto{a=b\vdash b=a}[thm]{%
        \FormulaRefAuto{(A\cap B)\cap C=A\cap(B\cap C)}[thm]}}}
  \MogProofStepWideStarBreak[1]{%
    \bigl(\LayerPowerMap{A_*}{D}{D\cup\{k\}}{\eta}(X)
      \cap A_*\bigr)\cap L=(X\cap A_*)\cap L}{%
    \rLRS{11}}
  \proofstepwidestar[1]{(X\cap A_*)\cap L=X\cap L}{%
    \rChain{%
      \FormulaRefAuto{(A\cap B)\cap C=A\cap(B\cap C)}[thm],
      \rLRS{13}}}
  \proofstepwidestar[1]{\Phi(X)\cap L=X\cap L}{%
    \rChain{14,15,16,17}}
  \proofstepwidestar[]{A_*\cap D=\varnothing}{%
    \rAEa{\FormulaRefAuto{MogiljanskajaLayerCarrierFacts}[thm]}}
  \proofstepwidestar[]{\eta\colon\powerset(D)\bij\powerset(D\cup\{k\})}{%
    \rAEa{\FormulaRefAuto{MogiljanskajaEtaBijective}[thm]}}
  \proofstepwidestar[]{\eta(\varnothing)=\varnothing}{%
    \rAEb{\FormulaRefAuto{MogiljanskajaEtaBijective}[thm]}}
  \MogProofStepWideStarBreak[1]{X\subseteq L\leftrightarrow
    \LayerPowerMap{A_*}{D}{D\cup\{k\}}{\eta}(X)\subseteq L}{%
    \FormulaRefAuto{LayerPowerMapPreservesSubsets}[thm]{%
      19,10,20,21,12,5}}
  \proofstepwidestar[1]{X\subseteq L\leftrightarrow\Phi(X)\subseteq L}{%
    \rIE{7,22}}
  \MogProofStepWideStarBreak[1]{%
    \begin{aligned}[t]
      &\Phi(X)=(X\cap A_*)\cup\eta(X\cap D),\\[-2pt]
      &\Phi(X)\cap L=X\cap L,\\[-2pt]
      &X\subseteq L\leftrightarrow\Phi(X)\subseteq L
    \end{aligned}}{\rAI{9,\rAI{18,23}}}
\end{tabproofwide}


\FormulaDefDeltaK[Rechteckanteil eines Mengenprodukts]{%
  R(X,Y)\coloneqq\bigl\{z\in D\bigm|
    \exists i,j\in\mathbb N\,
    (z=b_{ij}\land a_i\in X\cap L\land a_j\in Y\cap L)\bigr\}%
}{MogiljanskajaProductPartsDef}{%
  \DeltaRow{Mengen}{D\dsep L\dsep X\dsep Y\dsep z
    \dsep a_i\dsep a_j\dsep b_{ij}\dsep i\dsep j}
}

\FormulaThmDeltaK[Produkte mit einem Faktor au\ss erhalb von (L)]{%
  x,y\in T\dsep(x\notin L\lor y\notin L)
  \vdash x\MogDiamond y=\mathbf0%
}{MogiljanskajaDiamondElementOutsideL}{%
  \DeltaRow{Mengen}{L\dsep S\dsep T\dsep\mathbf0\dsep k\dsep x\dsep y}
  \DeltaRow{Funktionen}{\MogStar\dsep\MogDiamond}
}
\begin{tabproofwide}
  \proofstepwidestar[1]{x\in T}{\rA}
  \proofstepwidestar[2]{y\in T}{\rA}
  \proofstepwidestar[3]{x\notin L\lor y\notin L}{\rA}
  \proofstepwidestar[1]{x\in S\lor x=k}{%
    \FormulaRefAuto{MogiljanskajaCarrierTMembership}[thm]{1}}
  \proofstepwidestar[2]{y\in S\lor y=k}{%
    \FormulaRefAuto{MogiljanskajaCarrierTMembership}[thm]{2}}
  \proofstepwidestar[6]{x\in S}{\rA}
  \proofstepwidestar[7]{y\in S}{\rA}
  \proofstepwidestar[3,6,7]{x\MogStar y=\mathbf0}{%
    \FormulaRefAuto{MogiljanskajaStarDef}[def]{6,7,3}}
  \proofstepwidestar[6,7]{x\MogDiamond y=x\MogStar y}{%
    \FormulaRefAuto{MogiljanskajaDiamondDef}[def]{6,7}}
  \proofstepwidestar[3,6,7]{x\MogDiamond y=\mathbf0}{\rChain{9,8}}
  \proofstepwidestar[11]{y=k}{\rA}
  \proofstepwidestar[11]{k=y}{\FormulaRefAuto{a=b\vdash b=a}[thm]{11}}
  \proofstepwidestar[11]{k\in\{x,y\}}{%
    \FormulaRefAuto{x\in\{a,b\}\eqvdash(x=a\lor x=b)}[thm]{\rOIb{12}}}
  \proofstepwidestar[1,2,11]{x\MogDiamond y=\mathbf0}{%
    \FormulaRefAuto{MogiljanskajaDiamondDef}[def]{1,2,13}}
  \proofstepwidestar[1,2,3,6]{x\MogDiamond y=\mathbf0}{%
    \rOE{5,7,10,11,14}}
  \proofstepwidestar[16]{x=k}{\rA}
  \proofstepwidestar[16]{k=x}{\FormulaRefAuto{a=b\vdash b=a}[thm]{16}}
  \proofstepwidestar[16]{k\in\{x,y\}}{%
    \FormulaRefAuto{x\in\{a,b\}\eqvdash(x=a\lor x=b)}[thm]{\rOIa{17}}}
  \proofstepwidestar[1,2,16]{x\MogDiamond y=\mathbf0}{%
    \FormulaRefAuto{MogiljanskajaDiamondDef}[def]{1,2,18}}
  \proofstepwidestar[1,2,3]{x\MogDiamond y=\mathbf0}{\rOE{4,6,15,16,19}}
\end{tabproofwide}

\FormulaThmDeltaK[Au\ss erhalb von (L) sind alle Elemente
translationell ununterscheidbar]{%
  x,y\in T\dsep x\notin L\dsep y\notin L\vdash
  x\mathrel{\equiv^{\mathrm{tr}}_{T,\MogDiamond}}y%
}{MogiljanskajaOutsideLTranslationIndistinguishable}{%
  \DeltaRow{Mengen}{L\dsep T\dsep x\dsep y\dsep z}
  \DeltaRow{Funktionen}{\MogDiamond}
  \DeltaRow{Relationsoperatoren}{\equiv^{\mathrm{tr}}_{T,\MogDiamond}}
}
\begin{tabproofwide}
  \proofstepwidestar[1]{x\in T}{\rA}
  \proofstepwidestar[2]{y\in T}{\rA}
  \proofstepwidestar[3]{x\notin L}{\rA}
  \proofstepwidestar[4]{y\notin L}{\rA}
  \proofstepwidestar[]{\SemigroupAlg{T}{\MogDiamond}}{%
    \FormulaRefAuto{MogiljanskajaSecondSemigroup}[thm]}
  \proofstepwidestar[6]{z\in T}{\rA}
  \proofstepwidestar[1,3,6]{x\MogDiamond z=\mathbf0}{%
    \FormulaRefAuto{MogiljanskajaDiamondElementOutsideL}[thm]{1,6,\rOIa{3}}}
  \proofstepwidestar[2,4,6]{y\MogDiamond z=\mathbf0}{%
    \FormulaRefAuto{MogiljanskajaDiamondElementOutsideL}[thm]{2,6,\rOIa{4}}}
  \proofstepwidestar[1,2,3,4,6]{x\MogDiamond z=y\MogDiamond z}{%
    \rChain{7,\FormulaRefAuto{a=b\vdash b=a}[thm]{8}}}
  \proofstepwidestar[1,3,6]{z\MogDiamond x=\mathbf0}{%
    \FormulaRefAuto{MogiljanskajaDiamondElementOutsideL}[thm]{6,1,\rOIb{3}}}
  \proofstepwidestar[2,4,6]{z\MogDiamond y=\mathbf0}{%
    \FormulaRefAuto{MogiljanskajaDiamondElementOutsideL}[thm]{6,2,\rOIb{4}}}
  \proofstepwidestar[1,2,3,4,6]{z\MogDiamond x=z\MogDiamond y}{%
    \rChain{10,\FormulaRefAuto{a=b\vdash b=a}[thm]{11}}}
  \proofstepwidestar[1,2,3,4,6]{%
    x\MogDiamond z=y\MogDiamond z\land z\MogDiamond x=z\MogDiamond y}{%
    \rAI{9,12}}
  \proofstepwidestar[1,2,3,4]{%
    \forall z\in T\,\bigl(x\MogDiamond z=y\MogDiamond z\land
      z\MogDiamond x=z\MogDiamond y\bigr)}{\rUI{\rRI{6,13}}}
  \proofstepwidestar[1,2,3,4]{%
    x\mathrel{\equiv^{\mathrm{tr}}_{T,\MogDiamond}}y}{%
    \FormulaRefAuto{SemigroupTranslationIndistinguishableDef}[def]{5,1,2,14}}
\end{tabproofwide}

\FormulaThmDeltaK[Translationsprofile vor und nach der globalen Abbildung]{%
  \begin{aligned}[t]
    X\in\PowNE{S}\vdash{}&
      \Bigl(\forall x\in X\,\exists y\in\Phi(X)\,
      \bigl(x\mathrel{\equiv^{\mathrm{tr}}_{T,\MogDiamond}}y\bigr)\Bigr)\\[-2pt]
    &\land\Bigl(\forall y\in\Phi(X)\,\exists x\in X\,
      \bigl(y\mathrel{\equiv^{\mathrm{tr}}_{T,\MogDiamond}}x\bigr)\Bigr)
  \end{aligned}%
}{MogiljanskajaPhiTranslationSupport}{%
  \DeltaRow{Mengen}{L\dsep S\dsep T\dsep X\dsep x\dsep y}
  \DeltaRow{Funktionen}{\Phi\dsep\MogDiamond}
  \DeltaRow{Relationsoperatoren}{\equiv^{\mathrm{tr}}_{T,\MogDiamond}}
}
\begin{tabproofsplitwide}
  \proofpartwideindR[Von (X) nach (\Phi(X))]{%
    X\in\PowNE{S}\vdash\forall x\in X\,\exists y\in\Phi(X)\,
      \bigl(x\mathrel{\equiv^{\mathrm{tr}}_{T,\MogDiamond}}y\bigr)
  }{%
    X\in\PowNE{S}\vdash\forall x\in X\,\exists y\in\Phi(X)\,
      \bigl(x\mathrel{\equiv^{\mathrm{tr}}_{T,\MogDiamond}}y\bigr)
  }[MogiljanskajaPhiTranslationSupportForward]
    \proofstepwidestar[1]{X\in\PowNE{S}}{\rA}
    \proofstepwidestar[2]{x\in X}{\rA}
    \proofstepwidestar[]{\SemigroupAlg{T}{\MogDiamond}}{%
      \FormulaRefAuto{MogiljanskajaSecondSemigroup}[thm]}
    \proofstepwidestar[]{S\subseteq T}{\FormulaRefAuto{MogiljanskajaSSubsetT}[thm]}
    \proofstepwidestar[1,2]{x\in S}{\FormulaRefAuto{PowerSemigroupCarrierElement}[thm]{1,2}}
    \proofstepwidestar[1,2]{x\in T}{%
      \FormulaRefAuto{A\subseteq B\dsep x\in A\vdash x\in B}[thm]{4,5}}
    \proofstepwidestar[]{\Phi\colon\PowNE{S}\to\PowNE{T}}{%
      \FormulaRefAuto{Bijektive Funktion}[def]{\FormulaRefAuto{MogiljanskajaPhiBijective}[thm]}}
    \proofstepwidestar[1]{\Phi(X)\in\PowNE{T}}{%
      \FormulaRefAuto{F\colon A\to B\dsep x\in A\vdash F(x)\in B}[thm]{7,1}}
    \proofstepwidestar[1]{\Phi(X)\cap L=X\cap L}{%
      \rAEa{\rAEb{\FormulaRefAuto{MogiljanskajaPhiInvariants}[thm]{1}}}}
    \proofstepwidestar[1]{X\subseteq L\leftrightarrow\Phi(X)\subseteq L}{%
      \rAEb{\rAEb{\FormulaRefAuto{MogiljanskajaPhiInvariants}[thm]{1}}}}
    \proofstepwidestar[]{x\in L\lor x\notin L}{\FormulaRefAuto{P\lor\neg P}[thm]}
    \proofstepwidestar[12]{x\in L}{\rA}
    \proofstepwidestar[2,12]{x\in X\cap L}{%
      \FormulaRefAuto{x\in A\cap B\eqvdash x\in A\land x\in B}[thm]{\rAI{2,12}}}
    \proofstepwidestar[1,2,12]{x\in\Phi(X)\cap L}{%
      \rIE{\FormulaRefAuto{a=b\vdash b=a}[thm]{9},13}}
    \proofstepwidestar[1,2,12]{x\in\Phi(X)}{%
      \rAEa{\FormulaRefAuto{x\in A\cap B\eqvdash x\in A\land x\in B}[thm]{14}}}
    \proofstepwidestar[1,2]{x\mathrel{\equiv^{\mathrm{tr}}_{T,\MogDiamond}}x}{%
      \FormulaRefAuto{SemigroupTranslationIndistinguishableReflexive}[thm]{3,6}}
    \proofstepwidestar[1,2,12]{\exists y\in\Phi(X)\,
      \bigl(x\mathrel{\equiv^{\mathrm{tr}}_{T,\MogDiamond}}y\bigr)}{%
      \rEI{\rAI{15,16}}}
    \proofstepwidestar[18]{x\notin L}{\rA}
    \proofstepwidestar[2,18]{X\not\subseteq L}{%
      \FormulaRefAuto{OutsideWitnessBlocksSubset}[thm]{2,18}}
    \proofstepwidestar[1,2,18]{\Phi(X)\not\subseteq L}{%
      \FormulaRefAuto{P\leftrightarrow Q,\neg P\vdash\neg Q}[thm]{10,19}}
    \proofstepwidestar[1,2,18]{\exists y\in\Phi(X)\,(y\notin L)}{%
      \FormulaRefAuto{NonSubsetWitness}[thm]{20}}
    \proofstepwidestar[22]{y\in\Phi(X)\land y\notin L}{\rA}
    \proofstepwidestar[22]{y\in\Phi(X)}{\rAEa{22}}
    \proofstepwidestar[22]{y\notin L}{\rAEb{22}}
    \proofstepwidestar[1,22]{y\in T}{\FormulaRefAuto{PowerSemigroupCarrierElement}[thm]{8,23}}
    \proofstepwidestar[1,2,18,22]{x\mathrel{\equiv^{\mathrm{tr}}_{T,\MogDiamond}}y}{%
      \FormulaRefAuto{MogiljanskajaOutsideLTranslationIndistinguishable}[thm]{6,25,18,24}}
    \proofstepwidestar[1,2,18,22]{\exists y\in\Phi(X)\,
      \bigl(x\mathrel{\equiv^{\mathrm{tr}}_{T,\MogDiamond}}y\bigr)}{\rEI{\rAI{23,26}}}
    \proofstepwidestar[1,2,18]{\exists y\in\Phi(X)\,
      \bigl(x\mathrel{\equiv^{\mathrm{tr}}_{T,\MogDiamond}}y\bigr)}{\rEE{21,22,27}}
    \proofstepwidestar[1,2]{\exists y\in\Phi(X)\,
      \bigl(x\mathrel{\equiv^{\mathrm{tr}}_{T,\MogDiamond}}y\bigr)}{\rOE{11,12,17,18,28}}
    \proofstepwidestar[1]{\forall x\in X\,\exists y\in\Phi(X)\,
      \bigl(x\mathrel{\equiv^{\mathrm{tr}}_{T,\MogDiamond}}y\bigr)}{\rUI{\rRI{2,29}}}
  \closeproofpartwideindR

  \proofpartwideindR[Von (\Phi(X)) nach (X)]{%
    X\in\PowNE{S}\vdash\forall y\in\Phi(X)\,\exists x\in X\,
      \bigl(y\mathrel{\equiv^{\mathrm{tr}}_{T,\MogDiamond}}x\bigr)
  }{%
    X\in\PowNE{S}\vdash\forall y\in\Phi(X)\,\exists x\in X\,
      \bigl(y\mathrel{\equiv^{\mathrm{tr}}_{T,\MogDiamond}}x\bigr)
  }[MogiljanskajaPhiTranslationSupportBackward]
    \proofstepwidestar[1]{X\in\PowNE{S}}{\rA}
    \proofstepwidestar[2]{y\in\Phi(X)}{\rA}
    \proofstepwidestar[]{\SemigroupAlg{T}{\MogDiamond}}{%
      \FormulaRefAuto{MogiljanskajaSecondSemigroup}[thm]}
    \proofstepwidestar[]{S\subseteq T}{\FormulaRefAuto{MogiljanskajaSSubsetT}[thm]}
    \proofstepwidestar[]{\Phi\colon\PowNE{S}\to\PowNE{T}}{%
      \FormulaRefAuto{Bijektive Funktion}[def]{\FormulaRefAuto{MogiljanskajaPhiBijective}[thm]}}
    \proofstepwidestar[1]{\Phi(X)\in\PowNE{T}}{%
      \FormulaRefAuto{F\colon A\to B\dsep x\in A\vdash F(x)\in B}[thm]{5,1}}
    \proofstepwidestar[1,2]{y\in T}{\FormulaRefAuto{PowerSemigroupCarrierElement}[thm]{6,2}}
    \proofstepwidestar[1]{\Phi(X)\cap L=X\cap L}{%
      \rAEa{\rAEb{\FormulaRefAuto{MogiljanskajaPhiInvariants}[thm]{1}}}}
    \proofstepwidestar[1]{X\subseteq L\leftrightarrow\Phi(X)\subseteq L}{%
      \rAEb{\rAEb{\FormulaRefAuto{MogiljanskajaPhiInvariants}[thm]{1}}}}
    \proofstepwidestar[]{y\in L\lor y\notin L}{\FormulaRefAuto{P\lor\neg P}[thm]}
    \proofstepwidestar[11]{y\in L}{\rA}
    \proofstepwidestar[2,11]{y\in\Phi(X)\cap L}{%
      \FormulaRefAuto{x\in A\cap B\eqvdash x\in A\land x\in B}[thm]{\rAI{2,11}}}
    \proofstepwidestar[1,2,11]{y\in X\cap L}{\rIE{8,12}}
    \proofstepwidestar[1,2,11]{y\in X}{%
      \rAEa{\FormulaRefAuto{x\in A\cap B\eqvdash x\in A\land x\in B}[thm]{13}}}
    \proofstepwidestar[1,2]{y\mathrel{\equiv^{\mathrm{tr}}_{T,\MogDiamond}}y}{%
      \FormulaRefAuto{SemigroupTranslationIndistinguishableReflexive}[thm]{3,7}}
    \proofstepwidestar[1,2,11]{\exists x\in X\,
      \bigl(y\mathrel{\equiv^{\mathrm{tr}}_{T,\MogDiamond}}x\bigr)}{\rEI{\rAI{14,15}}}
    \proofstepwidestar[17]{y\notin L}{\rA}
    \proofstepwidestar[2,17]{\Phi(X)\not\subseteq L}{%
      \FormulaRefAuto{OutsideWitnessBlocksSubset}[thm]{2,17}}
    \proofstepwidestar[1,2,17]{X\not\subseteq L}{%
      \FormulaRefAuto{P\leftrightarrow Q,\neg Q\vdash\neg P}[thm]{9,18}}
    \proofstepwidestar[1,2,17]{\exists x\in X\,(x\notin L)}{%
      \FormulaRefAuto{NonSubsetWitness}[thm]{19}}
    \proofstepwidestar[21]{x\in X\land x\notin L}{\rA}
    \proofstepwidestar[21]{x\in X}{\rAEa{21}}
    \proofstepwidestar[21]{x\notin L}{\rAEb{21}}
    \proofstepwidestar[1,21]{x\in S}{\FormulaRefAuto{PowerSemigroupCarrierElement}[thm]{1,22}}
    \proofstepwidestar[1,21]{x\in T}{%
      \FormulaRefAuto{A\subseteq B\dsep x\in A\vdash x\in B}[thm]{4,24}}
    \proofstepwidestar[1,2,17,21]{y\mathrel{\equiv^{\mathrm{tr}}_{T,\MogDiamond}}x}{%
      \FormulaRefAuto{MogiljanskajaOutsideLTranslationIndistinguishable}[thm]{7,25,17,23}}
    \proofstepwidestar[1,2,17,21]{\exists x\in X\,
      \bigl(y\mathrel{\equiv^{\mathrm{tr}}_{T,\MogDiamond}}x\bigr)}{\rEI{\rAI{22,26}}}
    \proofstepwidestar[1,2,17]{\exists x\in X\,
      \bigl(y\mathrel{\equiv^{\mathrm{tr}}_{T,\MogDiamond}}x\bigr)}{\rEE{20,21,27}}
    \proofstepwidestar[1,2]{\exists x\in X\,
      \bigl(y\mathrel{\equiv^{\mathrm{tr}}_{T,\MogDiamond}}x\bigr)}{\rOE{10,11,16,17,28}}
    \proofstepwidestar[1]{\forall y\in\Phi(X)\,\exists x\in X\,
      \bigl(y\mathrel{\equiv^{\mathrm{tr}}_{T,\MogDiamond}}x\bigr)}{\rUI{\rRI{2,29}}}
  \closeproofpartwideindR

  \proofpartwide{Schluss}
    \proofstepwidestar[1]{X\in\PowNE{S}}{\rA}
    \proofstepwidestar[1]{\forall x\in X\,\exists y\in\Phi(X)\,
      \bigl(x\mathrel{\equiv^{\mathrm{tr}}_{T,\MogDiamond}}y\bigr)}{%
      \FormulaRefAuto{MogiljanskajaPhiTranslationSupportForward}[thm]{1}}
    \proofstepwidestar[1]{\forall y\in\Phi(X)\,\exists x\in X\,
      \bigl(y\mathrel{\equiv^{\mathrm{tr}}_{T,\MogDiamond}}x\bigr)}{%
      \FormulaRefAuto{MogiljanskajaPhiTranslationSupportBackward}[thm]{1}}
    \proofstepwidestar[1]{%
      \begin{aligned}[t]
        &\Bigl(\forall x\in X\,\exists y\in\Phi(X)\,
          \bigl(x\mathrel{\equiv^{\mathrm{tr}}_{T,\MogDiamond}}y\bigr)\Bigr)\\[-2pt]
        &\land\Bigl(\forall y\in\Phi(X)\,\exists x\in X\,
          \bigl(y\mathrel{\equiv^{\mathrm{tr}}_{T,\MogDiamond}}x\bigr)\Bigr)
      \end{aligned}}{\rAI{2,3}}
  \closeproofpartwide
\end{tabproofsplitwide}

\FormulaThmDeltaK[Die globale Abbildung wahrt das Translationsprofil
in der Potenzhalbgruppe]{%
  X\in\PowNE{S}\vdash
  X\mathrel{\equiv^{\mathrm{tr}}_{\PowNE{T},\MogDiamond_{\mathcal P}}}\Phi(X)%
}{MogiljanskajaPhiPowerTranslationIndistinguishable}{%
  \DeltaRow{Mengen}{S\dsep T\dsep X}
  \DeltaRow{Funktionen}{\Phi\dsep\MogDiamond\dsep\MogDiamond_{\mathcal P}}
}
\begin{tabproofwide}
  \proofstepwidestar[1]{X\in\PowNE{S}}{\rA}
  \proofstepwidestar[]{\SemigroupAlg{T}{\MogDiamond}}{%
    \FormulaRefAuto{MogiljanskajaSecondSemigroup}[thm]}
  \proofstepwidestar[]{S\subseteq T}{\FormulaRefAuto{MogiljanskajaSSubsetT}[thm]}
  \proofstepwidestar[1]{X\in\PowNE{T}}{\FormulaRefAuto{NonemptyPowersetMonotone}[thm]{3,1}}
  \proofstepwidestar[]{\Phi\colon\PowNE{S}\to\PowNE{T}}{%
    \FormulaRefAuto{Bijektive Funktion}[def]{\FormulaRefAuto{MogiljanskajaPhiBijective}[thm]}}
  \proofstepwidestar[1]{\Phi(X)\in\PowNE{T}}{%
    \FormulaRefAuto{F\colon A\to B\dsep x\in A\vdash F(x)\in B}[thm]{5,1}}
  \proofstepwidestar[1]{\forall x\in X\,\exists y\in\Phi(X)\,
    \bigl(x\mathrel{\equiv^{\mathrm{tr}}_{T,\MogDiamond}}y\bigr)}{%
    \rAEa{\FormulaRefAuto{MogiljanskajaPhiTranslationSupport}[thm]{1}}}
  \proofstepwidestar[1]{\forall y\in\Phi(X)\,\exists x\in X\,
    \bigl(y\mathrel{\equiv^{\mathrm{tr}}_{T,\MogDiamond}}x\bigr)}{%
    \rAEb{\FormulaRefAuto{MogiljanskajaPhiTranslationSupport}[thm]{1}}}
  \proofstepwidestar[1]{%
    X\mathrel{\equiv^{\mathrm{tr}}_{\PowNE{T},\MogDiamond_{\mathcal P}}}\Phi(X)}{%
    \FormulaRefAuto{PowerSemigroupTranslationSupportCriterion}[thm]{2,4,6,7,8}}
\end{tabproofwide}

\FormulaThmDeltaK[Das erste Mengenprodukt ist die Einschr\"ankung des
zweiten]{%
  X,Y\in\PowNE{S}\vdash
  X\MogStar_{\mathcal P}Y=X\MogDiamond_{\mathcal P}Y%
}{MogiljanskajaPowerProductRestriction}{%
  \DeltaRow{Mengen}{S\dsep T\dsep X\dsep Y\dsep z\dsep x\dsep y}
  \DeltaRow{Funktionen}{\MogStar\dsep\MogDiamond
    \dsep\MogStar_{\mathcal P}\dsep\MogDiamond_{\mathcal P}}
}
\begin{tabproofsplitwide}
  \proofpartwideindR[Erste Teilmengenrichtung]{%
    X,Y\in\PowNE{S}\vdash
    X\MogStar_{\mathcal P}Y\subseteq X\MogDiamond_{\mathcal P}Y
  }{%
    X,Y\in\PowNE{S}\vdash
    X\MogStar_{\mathcal P}Y\subseteq X\MogDiamond_{\mathcal P}Y
  }[MogiljanskajaPowerProductRestrictionForward]
    \proofstepwidestar[1]{X\in\PowNE{S}}{\rA}
    \proofstepwidestar[2]{Y\in\PowNE{S}}{\rA}
    \proofstepwidestar[]{\SemigroupAlg{S}{\MogStar}}{%
      \FormulaRefAuto{MogiljanskajaFirstSemigroup}[thm]}
    \proofstepwidestar[]{\SemigroupAlg{T}{\MogDiamond}}{%
      \FormulaRefAuto{MogiljanskajaSecondSemigroup}[thm]}
    \proofstepwidestar[]{S\subseteq T}{\FormulaRefAuto{MogiljanskajaSSubsetT}[thm]}
    \proofstepwidestar[1]{X\in\PowNE{T}}{\FormulaRefAuto{NonemptyPowersetMonotone}[thm]{5,1}}
    \proofstepwidestar[2]{Y\in\PowNE{T}}{\FormulaRefAuto{NonemptyPowersetMonotone}[thm]{5,2}}
    \proofstepwidestar[8]{z\in X\MogStar_{\mathcal P}Y}{\rA}
    \proofstepwidestar[1,2,8]{\exists x\in X\,\exists y\in Y\,(z=x\MogStar y)}{%
      \FormulaRefAuto{PowerSemigroupProductWitnesses}[thm]{3,1,2,8}}
    \proofstepwidestar[10]{x\in X\land\exists y\in Y\,(z=x\MogStar y)}{\rA}
    \proofstepwidestar[10]{x\in X}{\rAEa{10}}
    \proofstepwidestar[10]{\exists y\in Y\,(z=x\MogStar y)}{\rAEb{10}}
    \proofstepwidestar[13]{y\in Y\land z=x\MogStar y}{\rA}
    \proofstepwidestar[13]{y\in Y}{\rAEa{13}}
    \proofstepwidestar[13]{z=x\MogStar y}{\rAEb{13}}
    \proofstepwidestar[1,10]{x\in S}{\FormulaRefAuto{PowerSemigroupCarrierElement}[thm]{1,11}}
    \proofstepwidestar[2,13]{y\in S}{\FormulaRefAuto{PowerSemigroupCarrierElement}[thm]{2,14}}
    \proofstepwidestar[1,2,10,13]{x\MogDiamond y=x\MogStar y}{%
      \FormulaRefAuto{MogiljanskajaDiamondDef}[def]{16,17}}
    \proofstepwidestar[1,2,10,13]{z=x\MogDiamond y}{%
      \rChain{15,\FormulaRefAuto{a=b\vdash b=a}[thm]{18}}}
    \proofstepwidestar[1,2,10,13]{x\MogDiamond y\in X\MogDiamond_{\mathcal P}Y}{%
      \FormulaRefAuto{PowerSemigroupProductContains}[thm]{4,6,7,11,14}}
    \proofstepwidestar[1,2,10,13]{z\in X\MogDiamond_{\mathcal P}Y}{\rIE{19,20}}
    \proofstepwidestar[1,2,10]{z\in X\MogDiamond_{\mathcal P}Y}{\rEE{12,13,21}}
    \proofstepwidestar[1,2,8]{z\in X\MogDiamond_{\mathcal P}Y}{\rEE{9,10,22}}
    \proofstepwidestar[1,2]{\forall z\,(z\in X\MogStar_{\mathcal P}Y\rightarrow
      z\in X\MogDiamond_{\mathcal P}Y)}{\rUI{\rRI{8,23}}}
    \proofstepwidestar[1,2]{X\MogStar_{\mathcal P}Y\subseteq X\MogDiamond_{\mathcal P}Y}{%
      \FormulaRefAuto{A\subseteq B\coloneq\forall x\,(x\in A\rightarrow x\in B)}[def]{24}}
  \closeproofpartwideindR

  \proofpartwideindR[Zweite Teilmengenrichtung]{%
    X,Y\in\PowNE{S}\vdash
    X\MogDiamond_{\mathcal P}Y\subseteq X\MogStar_{\mathcal P}Y
  }{%
    X,Y\in\PowNE{S}\vdash
    X\MogDiamond_{\mathcal P}Y\subseteq X\MogStar_{\mathcal P}Y
  }[MogiljanskajaPowerProductRestrictionBackward]
    \proofstepwidestar[1]{X\in\PowNE{S}}{\rA}
    \proofstepwidestar[2]{Y\in\PowNE{S}}{\rA}
    \proofstepwidestar[]{\SemigroupAlg{S}{\MogStar}}{%
      \FormulaRefAuto{MogiljanskajaFirstSemigroup}[thm]}
    \proofstepwidestar[]{\SemigroupAlg{T}{\MogDiamond}}{%
      \FormulaRefAuto{MogiljanskajaSecondSemigroup}[thm]}
    \proofstepwidestar[]{S\subseteq T}{\FormulaRefAuto{MogiljanskajaSSubsetT}[thm]}
    \proofstepwidestar[1]{X\in\PowNE{T}}{\FormulaRefAuto{NonemptyPowersetMonotone}[thm]{5,1}}
    \proofstepwidestar[2]{Y\in\PowNE{T}}{\FormulaRefAuto{NonemptyPowersetMonotone}[thm]{5,2}}
    \proofstepwidestar[8]{z\in X\MogDiamond_{\mathcal P}Y}{\rA}
    \proofstepwidestar[4,6,7,8]{\exists x\in X\,\exists y\in Y\,(z=x\MogDiamond y)}{%
      \FormulaRefAuto{PowerSemigroupProductWitnesses}[thm]{4,6,7,8}}
    \proofstepwidestar[10]{x\in X\land\exists y\in Y\,(z=x\MogDiamond y)}{\rA}
    \proofstepwidestar[10]{x\in X}{\rAEa{10}}
    \proofstepwidestar[10]{\exists y\in Y\,(z=x\MogDiamond y)}{\rAEb{10}}
    \proofstepwidestar[13]{y\in Y\land z=x\MogDiamond y}{\rA}
    \proofstepwidestar[13]{y\in Y}{\rAEa{13}}
    \proofstepwidestar[13]{z=x\MogDiamond y}{\rAEb{13}}
    \proofstepwidestar[1,10]{x\in S}{\FormulaRefAuto{PowerSemigroupCarrierElement}[thm]{1,11}}
    \proofstepwidestar[2,13]{y\in S}{\FormulaRefAuto{PowerSemigroupCarrierElement}[thm]{2,14}}
    \proofstepwidestar[1,2,10,13]{x\MogDiamond y=x\MogStar y}{%
      \FormulaRefAuto{MogiljanskajaDiamondDef}[def]{16,17}}
    \proofstepwidestar[1,2,10,13]{z=x\MogStar y}{\rChain{15,18}}
    \proofstepwidestar[1,2,10,13]{x\MogStar y\in X\MogStar_{\mathcal P}Y}{%
      \FormulaRefAuto{PowerSemigroupProductContains}[thm]{3,1,2,11,14}}
    \proofstepwidestar[1,2,10,13]{z\in X\MogStar_{\mathcal P}Y}{\rIE{19,20}}
    \proofstepwidestar[1,2,10]{z\in X\MogStar_{\mathcal P}Y}{\rEE{12,13,21}}
    \proofstepwidestar[1,2,8]{z\in X\MogStar_{\mathcal P}Y}{\rEE{9,10,22}}
    \proofstepwidestar[1,2]{\forall z\,(z\in X\MogDiamond_{\mathcal P}Y\rightarrow
      z\in X\MogStar_{\mathcal P}Y)}{\rUI{\rRI{8,23}}}
    \proofstepwidestar[1,2]{X\MogDiamond_{\mathcal P}Y\subseteq X\MogStar_{\mathcal P}Y}{%
      \FormulaRefAuto{A\subseteq B\coloneq\forall x\,(x\in A\rightarrow x\in B)}[def]{24}}
  \closeproofpartwideindR

  \proofpartwide{Schluss}
    \proofstepwidestar[1]{X,Y\in\PowNE{S}}{\rA}
    \proofstepwidestar[1]{X\MogStar_{\mathcal P}Y\subseteq X\MogDiamond_{\mathcal P}Y}{%
      \FormulaRefAuto{MogiljanskajaPowerProductRestrictionForward}[thm]{1}}
    \proofstepwidestar[1]{X\MogDiamond_{\mathcal P}Y\subseteq X\MogStar_{\mathcal P}Y}{%
      \FormulaRefAuto{MogiljanskajaPowerProductRestrictionBackward}[thm]{1}}
    \proofstepwidestar[1]{X\MogStar_{\mathcal P}Y=X\MogDiamond_{\mathcal P}Y}{%
      \FormulaRefAuto{A\subseteq B\dsep B\subseteq A\vdash A=B}[thm]{2,3}}
  \closeproofpartwide
\end{tabproofsplitwide}

\FormulaThmDeltaK[Der Produktschichtanteil ist genau das Rechteck]{%
  X,Y\in\PowNE{S}\vdash(X\MogStar_{\mathcal P}Y)\cap D=R(X,Y)%
}{MogiljanskajaStarPowerProductDPart}{%
  \DeltaRow{Mengen}{S\dsep D\dsep L\dsep\mathbf0\dsep
    X\dsep Y\dsep x\dsep y\dsep z\dsep R(X,Y)}
  \DeltaRow{Funktionen}{\MogStar\dsep\MogStar_{\mathcal P}}
}
\begin{tabproofsplitwide}
  \proofpartwideindR[Produktzeugen]{%
    \begin{aligned}[t]
      &X,Y\in\PowNE{S}\dsep z\in X\MogStar_{\mathcal P}Y\\[-2pt]
      &\vdash\exists x\in X\,\exists y\in Y\,(z=x\MogStar y)
    \end{aligned}
  }{%
    \begin{aligned}[t]
      &X,Y\in\PowNE{S}\dsep z\in X\MogStar_{\mathcal P}Y\\[-2pt]
      &\vdash\exists x\in X\,\exists y\in Y\,(z=x\MogStar y)
    \end{aligned}
  }[MogiljanskajaStarPowerProductDPartForwardProductWitnesses]
    \proofstepwidestar[1]{X\in\PowNE{S}}{\rA}
    \proofstepwidestar[2]{Y\in\PowNE{S}}{\rA}
    \proofstepwidestar[3]{z\in X\MogStar_{\mathcal P}Y}{\rA}
    \proofstepwidestar[]{\SemigroupAlg{S}{\MogStar}}{%
      \FormulaRefAuto{MogiljanskajaFirstSemigroup}[thm]}
    \proofstepwidestarbreak[1,2,3]{%
      \exists x\in X\,\exists y\in Y\,(z=x\MogStar y)}{%
      \FormulaRefAuto{PowerSemigroupProductWitnesses}[thm]{4,1,2,3}}
  \closeproofpartwideindR

  \proofpartwideindR[Ausschluss des Nullfalls und Faktoren in \(L\)]{%
    \begin{aligned}[t]
      &x\in S\dsep y\in S\dsep z=x\MogStar y\dsep z\in D\\[-2pt]
      &\vdash x\in L\land y\in L
    \end{aligned}
  }{%
    \begin{aligned}[t]
      &x\in S\dsep y\in S\dsep z=x\MogStar y\dsep z\in D\\[-2pt]
      &\vdash x\in L\land y\in L
    \end{aligned}
  }[MogiljanskajaStarPowerProductDPartForwardFactorsInL]
    \proofstepwidestar[1]{x\in S}{\rA}
    \proofstepwidestar[2]{y\in S}{\rA}
    \proofstepwidestar[3]{z=x\MogStar y}{\rA}
    \proofstepwidestar[4]{z\in D}{\rA}
    \proofstepwidestar[]{(x\in L\land y\in L)\lor(x\notin L\lor y\notin L)}{%
      \FormulaRefAuto{(P\land Q)\lor(\neg P\lor\neg Q)}[thm]}
    \proofstepwidestar[6]{x\in L\land y\in L}{\rA}
    \proofstepwidestar[7]{x\notin L\lor y\notin L}{\rA}
    \proofstepwidestar[1,2,7]{x\MogStar y=\mathbf0}{%
      \FormulaRefAuto{MogiljanskajaStarDef}[def]{1,2,7}}
    \proofstepwidestar[1,2,3,7]{z=\mathbf0}{\rChain{3,8}}
    \proofstepwidestar[1,2,3,4,7]{\mathbf0\in D}{\rIE{9,4}}
    \proofstepwidestar[]{\mathbf0\notin D}{%
      \FormulaRefAuto{MogiljanskajaZeroNotInD}[thm]{%
        \FormulaRefAuto{MogiljanskajaZeroOutsideLayers}[thm]}}
    \proofstepwidestar[1,2,3,4,7]{\bot}{\rBI{11,10}}
    \proofstepwidestar[1,2,3,4,7]{x\in L\land y\in L}{%
      \FormulaRefAuto{P\dsep\neg P\vdash Q}[thm]{10,11}}
    \proofstepwidestar[1,2,3,4]{x\in L\land y\in L}{%
      \rOE{5,6,6,7,13}}
  \closeproofpartwideindR

  \proofpartwideindR[Indexzeugen und Rechteckmitgliedschaft]{%
    \begin{aligned}[t]
      &x\in X\dsep y\in Y\dsep z=x\MogStar y\dsep z\in D\dsep
        x\in L\land y\in L\\[-2pt]
      &\vdash z\in R(X,Y)
    \end{aligned}
  }{%
    \begin{aligned}[t]
      &x\in X\dsep y\in Y\dsep z=x\MogStar y\dsep z\in D\dsep
        x\in L\land y\in L\\[-2pt]
      &\vdash z\in R(X,Y)
    \end{aligned}
  }[MogiljanskajaStarPowerProductDPartForwardRectangleMembership]
    \proofstepwidestar[1]{x\in X}{\rA}
    \proofstepwidestar[2]{y\in Y}{\rA}
    \proofstepwidestar[3]{z=x\MogStar y}{\rA}
    \proofstepwidestar[4]{z\in D}{\rA}
    \proofstepwidestar[5]{x\in L\land y\in L}{\rA}
    \proofstepwidestar[5]{\exists i,j\in\mathbb N\,(x=a_i\land y=a_j)}{%
      \FormulaRefAuto{MogiljanskajaLayerLMembership}[thm]{5}}
    \proofstepwidestar[7]{i,j\in\mathbb N\land x=a_i\land y=a_j}{\rA}
    \proofstepwidestar[7]{i,j\in\mathbb N}{\rAEa{7}}
    \proofstepwidestar[7]{x=a_i}{\rAEa{\rAEb{7}}}
    \proofstepwidestar[7]{y=a_j}{\rAEb{\rAEb{7}}}
    \proofstepwidestar[7]{x\MogStar y=b_{ij}}{%
      \FormulaRefAuto{MogiljanskajaStarDef}[def]{7}}
    \proofstepwidestar[3,7]{z=b_{ij}}{\rChain{3,11}}
    \proofstepwidestar[1,7]{a_i\in X}{\rIE{9,1}}
    \proofstepwidestar[5,7]{a_i\in L}{\rIE{9,\rAEa{5}}}
    \proofstepwidestarbreak[1,5,7]{a_i\in X\cap L}{%
      \FormulaRefAuto{x\in A\cap B\eqvdash x\in A\land x\in B}[thm]{%
        \rAI{13,14}}}
    \proofstepwidestar[2,7]{a_j\in Y}{\rIE{10,2}}
    \proofstepwidestar[5,7]{a_j\in L}{\rIE{10,\rAEb{5}}}
    \proofstepwidestarbreak[2,5,7]{a_j\in Y\cap L}{%
      \FormulaRefAuto{x\in A\cap B\eqvdash x\in A\land x\in B}[thm]{%
        \rAI{16,17}}}
    \proofstepwidestarbreak[1,2,3,5,7]{%
      z=b_{ij}\land a_i\in X\cap L\land a_j\in Y\cap L}{%
      \rAI{12,\rAI{15,18}}}
    \proofstepwidestarbreak[1,2,3,5,7]{%
      \exists i,j\in\mathbb N\,
      (z=b_{ij}\land a_i\in X\cap L\land a_j\in Y\cap L)}{%
      \rEI{\rAI{8,19}}}
    \proofstepwidestarbreak[1,2,3,4,5,7]{z\in R(X,Y)}{%
      \rIE{\FormulaRefAuto{MogiljanskajaProductPartsDef}[def],
        \FormulaRefAuto{x\in\{u\in A\mid P(u)\}\eqvdash x\in A\land P(x)}[thm]{%
          \rAI{4,20}}}}
    \proofstepwidestar[1,2,3,4,5]{z\in R(X,Y)}{\rEE{6,7,21}}
  \closeproofpartwideindR

  \proofpartwideindR[Hinrichtung]{%
    X,Y\in\PowNE{S}\vdash(X\MogStar_{\mathcal P}Y)\cap D\subseteq R(X,Y)
  }{%
    X,Y\in\PowNE{S}\vdash(X\MogStar_{\mathcal P}Y)\cap D\subseteq R(X,Y)
  }[MogiljanskajaStarPowerProductDPartForward]
    \proofstepwidestar[1]{X\in\PowNE{S}}{\rA}
    \proofstepwidestar[2]{Y\in\PowNE{S}}{\rA}
    \proofstepwidestar[3]{z\in(X\MogStar_{\mathcal P}Y)\cap D}{\rA}
    \proofstepwidestarbreak[3]{%
      z\in X\MogStar_{\mathcal P}Y\land z\in D}{%
      \FormulaRefAuto{x\in A\cap B\eqvdash x\in A\land x\in B}[thm]{3}}
    \proofstepwidestar[3]{z\in X\MogStar_{\mathcal P}Y}{\rAEa{4}}
    \proofstepwidestar[3]{z\in D}{\rAEb{4}}
    \proofstepwidestarbreak[1,2,3]{%
      \exists x\in X\,\exists y\in Y\,(z=x\MogStar y)}{%
      \FormulaRefAuto{MogiljanskajaStarPowerProductDPartForwardProductWitnesses}[thm]{%
        1,2,5}}
    \proofstepwidestar[8]{x\in X\land\exists y\in Y\,(z=x\MogStar y)}{\rA}
    \proofstepwidestar[8]{x\in X}{\rAEa{8}}
    \proofstepwidestar[8]{\exists y\in Y\,(z=x\MogStar y)}{\rAEb{8}}
    \proofstepwidestar[11]{y\in Y\land z=x\MogStar y}{\rA}
    \proofstepwidestar[11]{y\in Y}{\rAEa{11}}
    \proofstepwidestar[11]{z=x\MogStar y}{\rAEb{11}}
    \proofstepwidestar[1,8]{x\in S}{%
      \FormulaRefAuto{PowerSemigroupCarrierElement}[thm]{1,9}}
    \proofstepwidestar[2,11]{y\in S}{%
      \FormulaRefAuto{PowerSemigroupCarrierElement}[thm]{2,12}}
    \proofstepwidestarbreak[1,2,3,8,11]{x\in L\land y\in L}{%
      \FormulaRefAuto{MogiljanskajaStarPowerProductDPartForwardFactorsInL}[thm]{%
        14,15,13,6}}
    \proofstepwidestarbreak[1,2,3,8,11]{z\in R(X,Y)}{%
      \FormulaRefAuto{MogiljanskajaStarPowerProductDPartForwardRectangleMembership}[thm]{%
        9,12,13,6,16}}
    \proofstepwidestar[1,2,3,8]{z\in R(X,Y)}{\rEE{10,11,17}}
    \proofstepwidestar[1,2,3]{z\in R(X,Y)}{\rEE{7,8,18}}
    \proofstepwidestarbreak[1,2]{%
      \forall z\,(z\in(X\MogStar_{\mathcal P}Y)\cap D\rightarrow z\in R(X,Y))}{%
      \rUI{\rRI{3,19}}}
    \proofstepwidestar[1,2]{(X\MogStar_{\mathcal P}Y)\cap D\subseteq R(X,Y)}{%
      \FormulaRefAuto{A\subseteq B\coloneq\forall x\,(x\in A\rightarrow x\in B)}[def]{20}}
  \closeproofpartwideindR

  \proofpartwideindR[Rechteckzeugen auspacken]{%
    \begin{aligned}[t]
      z\in R(X,Y)\vdash z\in D\land\exists i,j\in\mathbb N\,
      (&z=b_{ij}\land{}\\[-2pt]
       &a_i\in X\cap L\land a_j\in Y\cap L)
    \end{aligned}
  }{%
    \begin{aligned}[t]
      z\in R(X,Y)\vdash z\in D\land\exists i,j\in\mathbb N\,
      (&z=b_{ij}\land{}\\[-2pt]
       &a_i\in X\cap L\land a_j\in Y\cap L)
    \end{aligned}
  }[MogiljanskajaStarPowerProductDPartBackwardRectangleWitnesses]
    \proofstepwidestar[1]{z\in R(X,Y)}{\rA}
    \proofstepwidestarbreak[1]{%
      z\in D\land\exists i,j\in\mathbb N\,
      (z=b_{ij}\land a_i\in X\cap L\land a_j\in Y\cap L)}{%
      \rIE{\FormulaRefAuto{MogiljanskajaProductPartsDef}[def],
        \FormulaRefAuto{x\in\{u\in A\mid P(u)\}\eqvdash x\in A\land P(x)}[thm]{1}}}
  \closeproofpartwideindR

  \proofpartwideindR[Mengenprodukt und Schnitt]{%
    \begin{aligned}[t]
      &X,Y\in\PowNE{S}\dsep z\in D\dsep i,j\in\mathbb N\dsep
        z=b_{ij}\dsep{}\\[-2pt]
      &a_i\in X\cap L\dsep a_j\in Y\cap L
        \vdash z\in(X\MogStar_{\mathcal P}Y)\cap D
    \end{aligned}
  }{%
    \begin{aligned}[t]
      &X,Y\in\PowNE{S}\dsep z\in D\dsep i,j\in\mathbb N\dsep
        z=b_{ij}\dsep{}\\[-2pt]
      &a_i\in X\cap L\dsep a_j\in Y\cap L
        \vdash z\in(X\MogStar_{\mathcal P}Y)\cap D
    \end{aligned}
  }[MogiljanskajaStarPowerProductDPartBackwardProductIntersection]
    \proofstepwidestar[1]{X\in\PowNE{S}}{\rA}
    \proofstepwidestar[2]{Y\in\PowNE{S}}{\rA}
    \proofstepwidestar[3]{z\in D}{\rA}
    \proofstepwidestar[4]{i,j\in\mathbb N}{\rA}
    \proofstepwidestar[5]{z=b_{ij}}{\rA}
    \proofstepwidestar[6]{a_i\in X\cap L}{\rA}
    \proofstepwidestar[7]{a_j\in Y\cap L}{\rA}
    \proofstepwidestarbreak[6]{a_i\in X}{%
      \rAEa{\FormulaRefAuto{x\in A\cap B\eqvdash x\in A\land x\in B}[thm]{6}}}
    \proofstepwidestarbreak[7]{a_j\in Y}{%
      \rAEa{\FormulaRefAuto{x\in A\cap B\eqvdash x\in A\land x\in B}[thm]{7}}}
    \proofstepwidestar[]{\SemigroupAlg{S}{\MogStar}}{%
      \FormulaRefAuto{MogiljanskajaFirstSemigroup}[thm]}
    \proofstepwidestar[4]{a_i\MogStar a_j=b_{ij}}{%
      \FormulaRefAuto{MogiljanskajaStarDef}[def]{4}}
    \proofstepwidestarbreak[4,5]{z=a_i\MogStar a_j}{%
      \rChain{5,\FormulaRefAuto{a=b\vdash b=a}[thm]{11}}}
    \proofstepwidestarbreak[1,2,6,7]{%
      a_i\MogStar a_j\in X\MogStar_{\mathcal P}Y}{%
      \FormulaRefAuto{PowerSemigroupProductContains}[thm]{10,1,2,8,9}}
    \proofstepwidestar[1,2,4,5,6,7]{z\in X\MogStar_{\mathcal P}Y}{%
      \rIE{12,13}}
    \proofstepwidestarbreak[1,2,3,4,5,6,7]{%
      z\in X\MogStar_{\mathcal P}Y\land z\in D}{\rAI{14,3}}
    \proofstepwidestarbreak[1,2,3,4,5,6,7]{%
      z\in(X\MogStar_{\mathcal P}Y)\cap D}{%
      \FormulaRefAuto{x\in A\cap B\eqvdash x\in A\land x\in B}[thm]{15}}
  \closeproofpartwideindR

  \proofpartwideindR[R\"uckrichtung]{%
    X,Y\in\PowNE{S}\vdash R(X,Y)\subseteq(X\MogStar_{\mathcal P}Y)\cap D
  }{%
    X,Y\in\PowNE{S}\vdash R(X,Y)\subseteq(X\MogStar_{\mathcal P}Y)\cap D
  }[MogiljanskajaStarPowerProductDPartBackward]
    \proofstepwidestar[1]{X\in\PowNE{S}}{\rA}
    \proofstepwidestar[2]{Y\in\PowNE{S}}{\rA}
    \proofstepwidestar[3]{z\in R(X,Y)}{\rA}
    \proofstepwidestarbreak[3]{%
      z\in D\land\exists i,j\in\mathbb N\,
      (z=b_{ij}\land a_i\in X\cap L\land a_j\in Y\cap L)}{%
      \FormulaRefAuto{MogiljanskajaStarPowerProductDPartBackwardRectangleWitnesses}[thm]{3}}
    \proofstepwidestar[3]{z\in D}{\rAEa{4}}
    \proofstepwidestarbreak[3]{%
      \exists i,j\in\mathbb N\,
      (z=b_{ij}\land a_i\in X\cap L\land a_j\in Y\cap L)}{\rAEb{4}}
    \proofstepwidestarbreak[7]{%
      i,j\in\mathbb N\land z=b_{ij}\land
      a_i\in X\cap L\land a_j\in Y\cap L}{\rA}
    \proofstepwidestar[7]{i,j\in\mathbb N}{\rAEa{7}}
    \proofstepwidestar[7]{z=b_{ij}}{\rAEa{\rAEb{7}}}
    \proofstepwidestar[7]{a_i\in X\cap L}{\rAEa{\rAEb{\rAEb{7}}}}
    \proofstepwidestar[7]{a_j\in Y\cap L}{\rAEb{\rAEb{\rAEb{7}}}}
    \proofstepwidestarbreak[1,2,3,7]{z\in(X\MogStar_{\mathcal P}Y)\cap D}{%
      \FormulaRefAuto{MogiljanskajaStarPowerProductDPartBackwardProductIntersection}[thm]{%
        1,2,5,8,9,10,11}}
    \proofstepwidestar[1,2,3]{z\in(X\MogStar_{\mathcal P}Y)\cap D}{%
      \rEE{6,7,12}}
    \proofstepwidestarbreak[1,2]{%
      \forall z\,(z\in R(X,Y)\rightarrow
      z\in(X\MogStar_{\mathcal P}Y)\cap D)}{\rUI{\rRI{3,13}}}
    \proofstepwidestar[1,2]{R(X,Y)\subseteq(X\MogStar_{\mathcal P}Y)\cap D}{%
      \FormulaRefAuto{A\subseteq B\coloneq\forall x\,(x\in A\rightarrow x\in B)}[def]{14}}
  \closeproofpartwideindR

  \proofpartwide{Schluss}
    \proofstepwidestar[1]{X,Y\in\PowNE{S}}{\rA}
    \proofstepwidestar[1]{(X\MogStar_{\mathcal P}Y)\cap D\subseteq R(X,Y)}{%
      \FormulaRefAuto{MogiljanskajaStarPowerProductDPartForward}[thm]{1}}
    \proofstepwidestar[1]{R(X,Y)\subseteq(X\MogStar_{\mathcal P}Y)\cap D}{%
      \FormulaRefAuto{MogiljanskajaStarPowerProductDPartBackward}[thm]{1}}
    \proofstepwidestar[1]{(X\MogStar_{\mathcal P}Y)\cap D=R(X,Y)}{%
      \FormulaRefAuto{A\subseteq B\dsep B\subseteq A\vdash A=B}[thm]{2,3}}
  \closeproofpartwide
\end{tabproofsplitwide}

\FormulaThmDeltaK[Kein Produktrechteck liegt in der Reservefamilie]{%
  X,Y\in\powerset(T)\vdash R(X,Y)\notin\mathcal C%
}{MogiljanskajaRectangleOutsideReserve}{%
  \DeltaRow{Mengen}{T\dsep D\dsep D'\dsep P\dsep V\dsep L\dsep
    \mathcal C\dsep X\dsep Y\dsep c_0\dsep c_1\dsep b_{01}\dsep R(X,Y)}
}
\begin{tabproofwide}
  \proofstepwidestar[1]{X\in\powerset(T)}{\rA}
  \proofstepwidestar[2]{Y\in\powerset(T)}{\rA}
  \proofstepwidestar[3]{R(X,Y)\in\mathcal C}{\rA}
  \proofstepwidestarbreak[3]{P\subseteq R(X,Y)\land R(X,Y)\subseteq P\cup V}{%
    \FormulaRefAuto{CoreFamilyMembership}[thm]{3,
      \FormulaRefAuto{MogiljanskajaReserveFamiliesDef}[def]}}
  \proofstepwidestarbreak[]{c_0,c_1\in P}{%
    \FormulaRefAuto{x\in\{a,b\}\eqvdash(x=a\lor x=b)}[thm]{%
      \FormulaRefAuto{MogiljanskajaDPrimeDef}[def]}}
  \proofstepwidestarbreak[3]{c_0,c_1\in R(X,Y)}{%
    \FormulaRefAuto{A\subseteq B\dsep x\in A\vdash x\in B}[thm]{\rAEa{4},5}}
  \proofstepwidestarbreak[1,2,3]{%
    a_0\in X\cap L\land a_1\in Y\cap L}{%
    \FormulaRefAuto{MogiljanskajaProductPartsDef}[def]{6,
      \FormulaRefAuto{MogiljanskajaDPrimeDef}[def],
      \FormulaRefAuto{MogiljanskajaBIndexUnique}[thm]}}
  \proofstepwidestarbreak[]{b_{01}\in D}{%
    \FormulaRefAuto{MogiljanskajaLayerDMembership}[thm]{%
      \FormulaRefAuto{PeanoZero}[ax],
      \FormulaRefAuto{PeanoOneInN}[thm]}}
  \proofstepwidestarbreak[]{0,1\in\mathbb N}{%
    \rAI{\FormulaRefAuto{PeanoZero}[ax],
      \FormulaRefAuto{PeanoOneInN}[thm]}}
  \proofstepwidestarbreak[1,2,3]{%
    \begin{aligned}[t]
      &b_{01}=b_{01}\land a_0\in X\cap L\land{}\\[-2pt]
      &a_1\in Y\cap L
    \end{aligned}}{%
    \rAI{\rII,7}}
  \proofstepwidestarbreak[1,2,3]{%
    \begin{aligned}[t]
      &\exists i,j\in\mathbb N\,\bigl(b_{01}=b_{ij}\land{}\\[-2pt]
      &\qquad a_i\in X\cap L\land a_j\in Y\cap L\bigr)
    \end{aligned}}{%
    \rEI{\rAI{9,10}}}
  \proofstepwidestarbreak[1,2,3]{%
    \begin{aligned}[t]
      &b_{01}\in D\land\exists i,j\in\mathbb N\,\bigl(b_{01}=b_{ij}\land{}\\[-2pt]
      &\qquad a_i\in X\cap L\land a_j\in Y\cap L\bigr)
    \end{aligned}}{%
    \rAI{8,11}}
  \proofstepwidestarbreak[1,2,3]{b_{01}\in R(X,Y)}{%
    \rIE{\FormulaRefAuto{MogiljanskajaProductPartsDef}[def],
      \FormulaRefAuto{x\in\{u\in A\mid P(u)\}\eqvdash
        x\in A\land P(x)}[thm]{12}}}
  \proofstepwidestarbreak[]{P\cup V=D'}{%
    \FormulaRefAuto{MogiljanskajaDPrimeDef}[def]}
  \proofstepwidestarbreak[3]{R(X,Y)\subseteq D'}{\rIE{14,\rAEb{4}}}
  \proofstepwidestarbreak[1,2,3]{b_{01}\in D'}{%
    \FormulaRefAuto{A\subseteq B\dsep x\in A\vdash x\in B}[thm]{15,13}}
  \proofstepwidestarbreak[]{b_{01}\notin D'}{%
    \rAEb{\rAEb{\rAEb{\rAEb{\FormulaRefAuto{MogiljanskajaReserveMarkerFacts}[thm]}}}}}
  \proofstepwidestar[1,2,3]{\bot}{\rBI{17,16}}
  \proofstepwidestar[1,2]{R(X,Y)\notin\mathcal C}{\rCI{3,18}}
\end{tabproofwide}

\FormulaThmDeltaK[Die globale Abbildung fixiert jedes erste Mengenprodukt]{%
  X,Y\in\PowNE{S}\vdash
  \Phi(X\MogStar_{\mathcal P}Y)=X\MogStar_{\mathcal P}Y%
}{MogiljanskajaPhiFixesStarProduct}{%
  \DeltaRow{Mengen}{A_*\dsep D\dsep S\dsep T\dsep X\dsep Y\dsep R(X,Y)}
  \DeltaRow{Funktionen}{\eta\dsep\Phi\dsep\MogStar_{\mathcal P}}
}
\begin{tabproofwide}
  \proofstepwidestar[1]{X\in\PowNE{S}}{\rA}
  \proofstepwidestar[2]{Y\in\PowNE{S}}{\rA}
  \proofstepwidestar[1]{X\in\powerset(S)}{%
    \FormulaRefAuto{A\setminus B\subseteq A}[thm]{1}}
  \proofstepwidestar[1]{X\subseteq S}{%
    \FormulaRefAuto{x\in\powerset(A)\eqvdash x\subseteq A}[thm]{3}}
  \proofstepwidestar[2]{Y\in\powerset(S)}{%
    \FormulaRefAuto{A\setminus B\subseteq A}[thm]{2}}
  \proofstepwidestar[2]{Y\subseteq S}{%
    \FormulaRefAuto{x\in\powerset(A)\eqvdash x\subseteq A}[thm]{5}}
  \proofstepwidestar[]{S\subseteq T}{%
    \FormulaRefAuto{MogiljanskajaSSubsetT}[thm]}
  \proofstepwidestar[1]{X\subseteq T}{%
    \FormulaRefAuto{A\subseteq B\dsep B\subseteq C\vdash A\subseteq C}[thm]{4,7}}
  \proofstepwidestar[2]{Y\subseteq T}{%
    \FormulaRefAuto{A\subseteq B\dsep B\subseteq C\vdash A\subseteq C}[thm]{6,7}}
  \proofstepwidestar[1]{X\in\powerset(T)}{%
    \FormulaRefAuto{x\in\powerset(A)\eqvdash x\subseteq A}[thm]{8}}
  \proofstepwidestar[2]{Y\in\powerset(T)}{%
    \FormulaRefAuto{x\in\powerset(A)\eqvdash x\subseteq A}[thm]{9}}
  \MogProofStepWideStarBreak[1,2]{R(X,Y)\notin\mathcal C}{%
    \FormulaRefAuto{MogiljanskajaRectangleOutsideReserve}[thm]{10,11}}
  \proofstepwidestar[1,2]{R(X,Y)\subseteq D}{%
    \FormulaRefAuto{MogiljanskajaProductPartsDef}[def]}
  \proofstepwidestar[1,2]{R(X,Y)\in\powerset(D)}{%
    \FormulaRefAuto{x\in\powerset(A)\eqvdash x\subseteq A}[thm]{13}}
  \proofstepwidestar[1,2]{\eta(R(X,Y))=R(X,Y)}{%
    \FormulaRefAuto{MogiljanskajaEtaFixesOutsideReserve}[thm]{14,12}}
  \proofstepwidestar[1,2]{(X\MogStar_{\mathcal P}Y)\cap D=R(X,Y)}{%
    \FormulaRefAuto{MogiljanskajaStarPowerProductDPart}[thm]{1,2}}
  \proofstepwidestar[1,2]{\eta((X\MogStar_{\mathcal P}Y)\cap D)
    =(X\MogStar_{\mathcal P}Y)\cap D}{\rIE{16,15}}
  \proofstepwidestar[]{\SemigroupAlg{S}{\MogStar}}{%
    \FormulaRefAuto{MogiljanskajaFirstSemigroup}[thm]}
  \proofstepwidestar[1,2]{X\MogStar_{\mathcal P}Y\in\PowNE{S}}{%
    \FormulaRefAuto{PowerSemigroupProductClosure}[thm]{18,1,2}}
  \proofstepwidestar[1,2]{X\MogStar_{\mathcal P}Y\in\powerset(S)}{%
    \FormulaRefAuto{A\setminus B\subseteq A}[thm]{19}}
  \proofstepwidestar[]{S=A_*\cup D}{%
    \rAEa{\rAEb{\rAEb{\FormulaRefAuto{MogiljanskajaLayerCarrierFacts}[thm]}}}}
  \MogProofStepWideStarBreak[1,2]{%
    X\MogStar_{\mathcal P}Y\in\powerset(A_*\cup D)}{\rIE{21,20}}
  \proofstepwidestar[]{\eta\colon\powerset(D)\bij\powerset(D\cup\{k\})}{%
    \rAEa{\FormulaRefAuto{MogiljanskajaEtaBijective}[thm]}}
  \proofstepwidestar[]{\eta\colon\powerset(D)\to\powerset(D\cup\{k\})}{%
    \FormulaRefAuto{Bijektive Funktion}[def]{23}}
  \MogProofStepWideStarBreak[1,2]{%
    \LayerPowerMap{A_*}{D}{D\cup\{k\}}{\eta}(X\MogStar_{\mathcal P}Y)
      =X\MogStar_{\mathcal P}Y}{%
    \FormulaRefAuto{LayerPowerMapFixesStableUnion}[thm]{24,22,17}}
  \proofstepwidestar[]{\PowNE{S}\subseteq\powerset(S)}{%
    \FormulaRefAuto{A\setminus B\subseteq A}[thm]}
  \proofstepwidestar[]{\PowNE{S}\subseteq\powerset(A_*\cup D)}{%
    \rIE{21,26}}
  \MogProofStepWideStarBreak[1,2]{\Phi(X\MogStar_{\mathcal P}Y)=
    \LayerPowerMap{A_*}{D}{D\cup\{k\}}{\eta}(X\MogStar_{\mathcal P}Y)}{%
    \FormulaRefAuto{C\subseteq A\dsep x\in C\vdash F\restriction_C(x)=F(x)}[thm]{%
      \FormulaRefAuto{MogiljanskajaPhiDef}[def],27,19}}
  \proofstepwidestar[1,2]{\Phi(X\MogStar_{\mathcal P}Y)=X\MogStar_{\mathcal P}Y}{%
    \rChain{28,25}}
\end{tabproofwide}

\FormulaThmDeltaK[Die konstruierten Potenzhalbgruppen sind isomorph]{%
  \SemigroupIso{\Phi}{\PowNE{S},\MogStar_{\mathcal P}}
    {\PowNE{T},\MogDiamond_{\mathcal P}}%
}{MogiljanskajaPowerSemigroupIsomorphism}{%
  \DeltaRow{Mengen}{S\dsep T\dsep X\dsep Y}
  \DeltaRow{Funktionen}{\Phi\dsep\MogStar\dsep\MogDiamond
    \dsep\MogStar_{\mathcal P}\dsep\MogDiamond_{\mathcal P}}
}
\begin{tabproofwide}
  \proofstepwidestar[]{\SemigroupAlg{S}{\MogStar}}{%
    \FormulaRefAuto{MogiljanskajaFirstSemigroup}[thm]}
  \proofstepwidestar[]{\SemigroupAlg{T}{\MogDiamond}}{%
    \FormulaRefAuto{MogiljanskajaSecondSemigroup}[thm]}
  \proofstepwidestar[]{\SemigroupAlg{\PowNE{S}}{\MogStar_{\mathcal P}}}{%
    \FormulaRefAuto{NonemptyPowerSemigroup}[thm]{1}}
  \proofstepwidestar[]{\SemigroupAlg{\PowNE{T}}{\MogDiamond_{\mathcal P}}}{%
    \FormulaRefAuto{NonemptyPowerSemigroup}[thm]{2}}
  \proofstepwidestar[5]{X\in\PowNE{S}}{\rA}
  \proofstepwidestar[5]{X\in\PowNE{T}}{%
    \FormulaRefAuto{NonemptyPowersetMonotone}[thm]{%
      \FormulaRefAuto{MogiljanskajaSSubsetT}[thm],5}}
  \proofstepwidestar[]{\PowNE{S}\subseteq\PowNE{T}}{%
    \FormulaRefAuto{A\subseteq B\coloneq\forall x\,(x\in A\rightarrow x\in B)}[def]{%
      \rUI{\rRI{5,6}}}}
  \proofstepwidestar[]{\Phi\colon\PowNE{S}\bij\PowNE{T}}{%
    \FormulaRefAuto{MogiljanskajaPhiBijective}[thm]}
  \proofstepwidestar[9]{X\in\PowNE{S}}{\rA}
  \proofstepwidestar[10]{Y\in\PowNE{S}}{\rA}
  \proofstepwidestar[9,10]{X\MogStar_{\mathcal P}Y=X\MogDiamond_{\mathcal P}Y}{%
    \FormulaRefAuto{MogiljanskajaPowerProductRestriction}[thm]{9,10}}
  \proofstepwidestar[]{\forall X,Y\in\PowNE{S}\,
    (X\MogStar_{\mathcal P}Y=X\MogDiamond_{\mathcal P}Y)}{%
    \rUI{\rRI{9,\rUI{\rRI{10,11}}}}}
  \proofstepwidestar[13]{X\in\PowNE{S}}{\rA}
  \proofstepwidestar[13]{X\mathrel{\equiv^{\mathrm{tr}}_{\PowNE{T},
    \MogDiamond_{\mathcal P}}}\Phi(X)}{%
    \FormulaRefAuto{MogiljanskajaPhiPowerTranslationIndistinguishable}[thm]{13}}
  \proofstepwidestar[]{\forall X\in\PowNE{S}\,
    \bigl(X\mathrel{\equiv^{\mathrm{tr}}_{\PowNE{T},
      \MogDiamond_{\mathcal P}}}\Phi(X)\bigr)}{\rUI{\rRI{13,14}}}
  \proofstepwidestar[16]{X\in\PowNE{S}}{\rA}
  \proofstepwidestar[17]{Y\in\PowNE{S}}{\rA}
  \proofstepwidestar[16,17]{\Phi(X\MogStar_{\mathcal P}Y)
    =X\MogStar_{\mathcal P}Y}{%
    \FormulaRefAuto{MogiljanskajaPhiFixesStarProduct}[thm]{16,17}}
  \proofstepwidestar[]{\forall X,Y\in\PowNE{S}\,
    \Phi(X\MogStar_{\mathcal P}Y)=X\MogStar_{\mathcal P}Y}{%
    \rUI{\rRI{16,\rUI{\rRI{17,18}}}}}
  \proofstepwidestar[]{\SemigroupIso{\Phi}{\PowNE{S},\MogStar_{\mathcal P}}
    {\PowNE{T},\MogDiamond_{\mathcal P}}}{%
    \FormulaRefAuto{SemigroupFixedProductsIsomorphismCriterion}[thm]{%
      3,4,7,8,12,15,19}}
\end{tabproofwide}

\FormulaThmDeltaK[Das Mogiljanskaja-Paar ist ein unendliches
Gegenbeispiel]{%
  \begin{aligned}[t]
    &\SemigroupAlg{S}{\MogStar}\land
      \SemigroupAlg{T}{\MogDiamond}\land{}\\[-2pt]
    &\neg\Finite{S}\land\neg\Finite{T}\land{}\\[-2pt]
    &\neg\exists f\,\SemigroupIso{f}{S,\MogStar}{T,\MogDiamond}\land{}\\[-2pt]
    &{}
      \exists F\,\SemigroupIso{F}
        {\PowNE{S},\MogStar_{\mathcal P}}
        {\PowNE{T},\MogDiamond_{\mathcal P}}
  \end{aligned}
}{MogiljanskajaPairCounterexample}{%
  \DeltaRow{Mengen}{S\dsep T}
  \DeltaRow{Funktionen}{f\dsep F\dsep\Phi\dsep
    \MogStar\dsep\MogDiamond\dsep
    \MogStar_{\mathcal P}\dsep\MogDiamond_{\mathcal P}}
}
\begin{tabproofwide}
  \proofstepwidestar[]{\SemigroupAlg{S}{\MogStar}}{%
    \FormulaRefAuto{MogiljanskajaFirstSemigroup}[thm]}
  \proofstepwidestar[]{\SemigroupAlg{T}{\MogDiamond}}{%
    \FormulaRefAuto{MogiljanskajaSecondSemigroup}[thm]}
  \proofstepwidestar[]{\neg\Finite{S}}{%
    \FormulaRefAuto{MogiljanskajaCarrierSInfinite}[thm]}
  \proofstepwidestar[]{\neg\Finite{T}}{%
    \rAEb{\FormulaRefAuto{MogiljanskajaCarrierTInfinite}[thm]}}
  \proofstepwidestar[]{\neg\exists f\,
    \SemigroupIso{f}{S,\MogStar}{T,\MogDiamond}}{%
    \FormulaRefAuto{MogiljanskajaNotIsomorphic}[thm]}
  \proofstepwidestarbreak[]{\SemigroupIso{\Phi}
    {\PowNE{S},\MogStar_{\mathcal P}}
    {\PowNE{T},\MogDiamond_{\mathcal P}}}{%
    \FormulaRefAuto{MogiljanskajaPowerSemigroupIsomorphism}[thm]}
  \proofstepwidestarbreak[]{\exists F\,\SemigroupIso{F}
    {\PowNE{S},\MogStar_{\mathcal P}}
    {\PowNE{T},\MogDiamond_{\mathcal P}}}{\rEI{6}}
  \proofstepwidestarbreak[]{%
    \begin{aligned}[t]
      &\SemigroupAlg{S}{\MogStar}\land
        \SemigroupAlg{T}{\MogDiamond}\land{}\\[-2pt]
      &\neg\Finite{S}\land\neg\Finite{T}\land{}\\[-2pt]
      &\neg\exists f\,\SemigroupIso{f}{S,\MogStar}{T,\MogDiamond}\land{}\\[-2pt]
      &{}
        \exists F\,\SemigroupIso{F}
          {\PowNE{S},\MogStar_{\mathcal P}}
          {\PowNE{T},\MogDiamond_{\mathcal P}}
    \end{aligned}}{%
    \rAI{1,\rAI{2,\rAI{3,\rAI{4,\rAI{5,7}}}}}}
\end{tabproofwide}

Damit ist das Mogiljanskaja-Paar ein Gegenbeispiel zur Umkehrung von
\FormulaRefAuto{SemigroupIsoImpliesPowerSemigroupIso}[thm].
Die Umkehrung ist also ohne Endlichkeitsvoraussetzung falsch.

\end{document}
